% !TEX root = print.tex

\part{Sobolev and Bessel potential spaces}
\label{part:sobolev}

This part treats the Sobolev spaces $\Sob{m}{p}$ defined by weak derivatives and Mathlib's
Bessel potential spaces $H^{s,p}$ on $\R^d$. \Cref{ch:sobolev-rd} identifies Mathlib's $H^s$,
$s \geq 0$, with $\TL{s}{2}{2}$ (and, for $s > 0$, with $\Besov{s}{2}{2}$), proves Rellich's
theorem, and, from the vector-valued singular integrals of \Cref{ch:singular-integrals}, proves
$\Sob{m}{p} = \TL{m}{p}{2}$ for $1 < p < \infty$ and the Sobolev embeddings;
\Cref{ch:sobolev-torus} defines weak derivatives in $L_p(\T^d)$ through Fourier coefficients and
proves $\Sob{m}{p}(\T^d) = \TL{m}{p}{2}(\T^d)$ for $1 < p < \infty$ with the Sobolev
embeddings on the torus;
\Cref{ch:bessel-potential-spaces} proves, also from the singular integrals,
$H^{s,p} = \TL{s}{p}{2}$ for tempered distributions, every real $s$ and $1 < p < \infty$, and
Calder\'on's theorem $\Sob{m}{p} = H^{m,p}$;
\Cref{ch:fourier-multipliers} treats the symbols of order $\sigma$, their Fourier multipliers
on $\TL{s}{p}{q}$ and $H^{s,p}$, and Mihlin's multiplier theorem for symbols of order $0$.

\chapter{Sobolev spaces on \texorpdfstring{$\R^d$}{R\^{}d}}
\label{ch:sobolev-rd}

This chapter works on a finite-dimensional real inner product space $V$ of dimension $n$, such
as $\R^d$, with Lebesgue measure, and defines Sobolev spaces by weak derivatives. A weak
derivative is an identity between distributions on an open set $\Omega \subseteq V$, in Mathlib's
space $\mathcal D'(\Omega)$ of distributions tested against real smooth compactly supported
functions (\Cref{sec:schwartz-cutoff}). On the whole space $V$ and for $L_p$ functions it is the
same as an identity between tempered distributions, the setting of Mathlib's Fourier transform
and Bessel potential spaces (\Cref{sec:weak-derivatives}).
\Cref{sec:weak-derivative-change-of-variables} shows that Lipschitz functions are weakly
differentiable, bounds products with cutoffs in $L_p$ and proves the chain rule under Lipschitz
maps that preserve Lebesgue measure. For $p = 2$,
\Cref{sec:sobolev-two-littlewood-paley} identifies Mathlib's Bessel potential spaces $H^s(V)$,
$s \geq 0$, with the Triebel--Lizorkin spaces $\TL{s}{2}{2}$ (and, for $s > 0$, with the Besov
spaces $\Besov{s}{2}{2}$) and proves $\Sob{1}{2}(V) = H^1(V) = \Besov{1}{2}{2}$.
\Cref{sec:rellich-compactness} proves Rellich's theorem: the restriction from $H^1(V)$ to
$L_2(\Omega)$ is compact for every bounded $\Omega \subseteq V$. \Cref{sec:sobolev-scale}
builds, for every real $s$, the inclusions $H^s(V) \subseteq H^t(V)$, $t \leq s$, the $L_2$
functions of $H^s(V)$, $s \geq 0$, and the derivatives $H^{s+1}(V) \to H^s(V)$ with the exact
recursion of the norms, and gives, for $s > 0$, $H^s(V) = \Besov{s}{2}{2}$ by differences with
explicit constants. \Cref{sec:leibniz-rule}
proves the Leibniz rule for the weak derivatives of every order of a product with a smooth
compactly supported function. \Cref{sec:derivative-kernels}
writes the Littlewood--Paley blocks as sums of derivatives of Schwartz kernels,
\Cref{sec:divergence-kernels} writes the difference of two dilates of a smooth function as a
divergence of smooth kernels (used for the extension theorems of \Cref{sec:sobolev-extension}),
and \Cref{sec:sobolev-triebel-lizorkin} proves
$\Sob{m}{p}(V) = \TL{m}{p}{2}$ for $1 < p < \infty$,
with named constants and the Besov sandwich, from the strong type of vector-valued singular
integrals (\Cref{ch:singular-integrals}). \Cref{sec:sobolev-embedding} proves the Nikolskii
inequality for the blocks and from it the embeddings of Besov spaces into $L_r(V)$ and into the
bounded continuous functions, and the Sobolev embeddings $\Sob{m}{p}(V) \hookrightarrow L_r(V)$
for $m > n(1/p - 1/r)$ and $\Sob{m}{p}(V) \hookrightarrow C_b(V)$ for $m > n/p$, all with named
constants. These named constants are formulas in norms and Schwartz seminorms of fixed kernels,
proved finite. \Cref{sec:critical-sobolev-embedding} proves the critical case
$1/r = 1/p - m/n$, $1 < p < \infty$, and Jawerth's embedding
$\TL{s_0}{p_0}{\infty} \hookrightarrow \TL{s_1}{p_1}{q}$, following Johnsen and Sickel.
\Cref{sec:gagliardo-nirenberg-sobolev} carries Mathlib's Gagliardo--Nirenberg--Sobolev
inequality to weakly differentiable functions, which gives the case $p = 1$:
$\Sob{1}{p}(V) \hookrightarrow L_r(V)$ for $1/r = 1/p - 1/n$ and every $1 \leq p$, and
$\Sob{m}{1}(V) \hookrightarrow L_r(V)$ for $1/r = 1 - m/n$.
\Cref{sec:sobolev-interpolation} identifies the real interpolation spaces of
$L_p(V)$ and $\Sob{m}{p}(V)$, and of two Sobolev spaces, with Besov spaces. The Bessel
potential spaces $H^{s,p}(V)$ for $1 < p < \infty$ follow in
\Cref{ch:bessel-potential-spaces}, the Sobolev spaces on open sets and on bounded Lipschitz
domains in \Cref{ch:sobolev-domains}.

\section{Weak derivatives on open sets}
\label{sec:weak-derivatives}

Throughout this section $\Omega \subseteq V$ is open. Mathlib identifies a locally integrable
function $f$ on $\Omega$ with the distribution $\varphi \mapsto \int \varphi f$ and a function
$f \in L_p(V)$, $1 \leq p \leq \infty$, with the tempered distribution $u \mapsto \int u f$.

\begin{definition}[Weak derivatives and Sobolev spaces]
\label{def:weak-line-derivative}
\lean{Rd.HasWeakLineDeriv, Rd.MemWeakSobolev, Rd.MemSobolev}
\leanok
For every $v \in V^k$, $g$ is the weak derivative $\partial^v f$ of $f$ on $\Omega$ if $f$ and $g$
are locally integrable on $\Omega$ and $\partial^v f = g$ as distributions on $\Omega$. For every
$m \in \N$ and every $1 \leq p \leq \infty$, $f$ lies in the Sobolev space $\Sob{m}{p}(\Omega)$ if
$f \in L_p(\Omega)$ and for every $k \leq m$ and every $v \in V^k$ the weak derivative
$\partial^v f$ exists on $\Omega$ and lies in $L_p(\Omega)$. For every $f \in L_p(V)$,
$f \in \Sob{m}{p}(V)$ means $f \in \Sob{m}{p}(\Omega)$ with $\Omega = V$.
\end{definition}

Every function in $L_p(\Omega)$ with $p \geq 1$ is locally integrable on $\Omega$
(\Cref{lem:weak-derivative-basic}), so the weak derivatives in the definition of
$\Sob{m}{p}(\Omega)$ need no further hypothesis. The definition quantifies over all directions;
a basis of $V$ enters only in proofs.

\begin{lemma}[Weak derivatives by integration by parts]
\label{lem:weak-derivative-integral}
\lean{Rd.HasWeakLineDeriv.integral_eq, Rd.hasWeakLineDeriv_iff_integral}
\leanok
\uses{def:weak-line-derivative, lem:distribution-transpose}
For all locally integrable $f$, $g$ on $\Omega$ and every $v \in V^k$, $g$ is the weak
derivative $\partial^v f$ on $\Omega$ if and only if for every $\varphi \in \mathcal D(\Omega)$,
\[
  \int_V \varphi \, g = (-1)^k \int_V (\partial^v \varphi) \, f .
\]
\end{lemma}

\begin{proof}
\leanok
Both distributions are determined by their values on test functions, the distribution of a
locally integrable function is $\varphi \mapsto \int \varphi f$ (Mathlib's
\texttt{Distribution.\allowbreak ofFun\_\allowbreak apply}), and
$\partial^v f(\varphi) = (-1)^k f(\partial^v \varphi)$ by \Cref{lem:distribution-transpose}.
\end{proof}

\begin{lemma}[Properties of weak derivatives]
\label{lem:weak-derivative-basic}
\lean{Rd.HasWeakLineDeriv.unique, Rd.hasWeakLineDeriv_fin_zero, Rd.HasWeakLineDeriv.add,
  Rd.HasWeakLineDeriv.neg, Rd.HasWeakLineDeriv.sub, Rd.HasWeakLineDeriv.smul,
  Rd.HasWeakLineDeriv.congr_ae, Rd.HasWeakLineDeriv.mono, Rd.hasWeakLineDeriv_zero,
  Rd.HasWeakLineDeriv.ae_eq_zero_of_ae_eq_zero,
  Rd.locallyIntegrableOn_of_memLp, Rd.memWeakSobolev_zero_iff, Rd.MemWeakSobolev.mono_order,
  Rd.MemWeakSobolev.mono_set, Rd.zero_memWeakSobolev, Rd.MemWeakSobolev.add,
  Rd.MemWeakSobolev.smul, Rd.MemWeakSobolev.congr_ae, Rd.zero_memSobolev, Rd.MemSobolev.add,
  Rd.MemSobolev.smul, Rd.HasWeakLineDeriv.sum, Rd.memSobolev_iff_exists_hasWeakLineDeriv,
  Rd.MemSobolev.sum, Rd.HasWeakLineDeriv.eq_of_top, Rd.memSobolev_one_of_forall_hasWeakLineDeriv,
  Rd.HasWeakLineDeriv.locallyIntegrable, Rd.HasWeakLineDeriv.locallyIntegrable_deriv,
  Rd.HasWeakLineDeriv.of_cons}
\leanok
\uses{def:weak-line-derivative, lem:weak-derivative-integral, lem:directional-derivatives-commute}
For every $v \in V^k$, all functions $f, f', g, g'$ and every $c \in \C$:
\begin{enumerate}
  \item if $g$ and $g'$ are weak derivatives $\partial^v f$ on $\Omega$, then $g = g'$ almost
    everywhere on $\Omega$; for $\Omega = V$, every $r \in [0, \infty]$ and
    $g, g' \in L_r(V)$, $g = g'$ in $L_r(V)$;
  \item for $k = 0$, every locally integrable $f$ on $\Omega$ is its own weak derivative;
  \item if $g = \partial^v f$ and $g' = \partial^v f'$ weakly on $\Omega$, then
    $g + g' = \partial^v (f + f')$, $-g = \partial^v(-f)$, $g - g' = \partial^v (f - f')$ and
    $c g = \partial^v (c f)$ weakly on $\Omega$; for a finite set $s$ and functions $F_j$,
    $G_j$ with $G_j = \partial^v F_j$ weakly on $\Omega$ for every $j \in s$,
    $\sum_{j \in s} G_j = \partial^v \sum_{j \in s} F_j$ weakly on $\Omega$;
  \item if $g = \partial^v f$ weakly on $\Omega$, $f = f'$ and $g = g'$ almost everywhere on
    $\Omega$, then $g' = \partial^v f'$ weakly on $\Omega$;
  \item if $g = \partial^v f$ weakly on $\Omega$, then $g = \partial^v f$ weakly on every open
    $\Omega' \subseteq \Omega$;
  \item $0 = \partial^v 0$ weakly on $\Omega$;
  \item if $g = \partial^v f$ weakly on $\Omega$, $O \subseteq \Omega$ is open and $f = 0$
    almost everywhere on $O$, then $g = 0$ almost everywhere on $O$;
  \item if $g = \partial^v f$ weakly on $V$, then $f$ and $g$ are locally integrable on $V$;
  \item for every $w \in V$ and all $G$, $H$ with $G = \partial^v f$ and
    $H = \partial^{(w, v)} f$ weakly on $\Omega$, $H = \partial_w G$ weakly on $\Omega$.
\end{enumerate}
For every $1 \leq p \leq \infty$, every $f \in L_p(\Omega)$ is locally integrable on $\Omega$
and $\Sob{0}{p}(\Omega) = L_p(\Omega)$. For all $m, l \in \N$ with
$l \leq m$, every $1 \leq p \leq \infty$ and every open $\Omega' \subseteq \Omega$,
$\Sob{m}{p}(\Omega) \subseteq \Sob{l}{p}(\Omega)$, every $f \in \Sob{m}{p}(\Omega)$ lies in
$\Sob{m}{p}(\Omega')$, and $\Sob{m}{p}(\Omega)$ is a complex vector space. For every
$f \in \Sob{m}{p}(\Omega)$ and every $f'$ with $f' = f$ almost everywhere on $\Omega$,
$f' \in \Sob{m}{p}(\Omega)$. For $\Omega = V$, $\Sob{m}{p}(V)$ is a complex linear subspace of
$L_p(V)$, and $f \in L_p(V)$ lies in $\Sob{m}{p}(V)$ if and only if for every $k \leq m$ and
every $v \in V^k$ some $g \in L_p(V)$ is the weak derivative $\partial^v f$ on $V$; for $m = 1$
it suffices that for every $v \in V$ some $g \in L_p(V)$ is the weak derivative $\partial_v f$
on $V$.
\end{lemma}

\begin{proof}
\leanok
(1) A locally integrable function is determined almost everywhere on $\Omega$ by its
distribution (Mathlib's \texttt{Distribution.\allowbreak ofFun\_\allowbreak injective}). (2)
$\partial^v$ is the identity for $k = 0$. (3) The distribution of a function and $\partial^v$ are
additive, and finite sums follow by induction; for complex $c$ use
\Cref{lem:weak-derivative-integral}, since $c$ leaves both integrals. (4) The distribution of a
function depends on it only almost everywhere on $\Omega$.
(5) A test function $\varphi$ on $\Omega'$ is one on $\Omega$, with the same derivatives
(\Cref{lem:directional-derivatives-commute}), so the identity of
\Cref{lem:weak-derivative-integral} on $\Omega$ gives it on $\Omega'$. (6) The distribution of
$0$ is $0$, and so are its derivatives. (7) By (5) and (4), $g = \partial^v 0$ weakly on $O$,
so $g = 0$ almost everywhere on $O$ by (6) and (1). (8) is the definition for $\Omega = V$.
(9) For distributions $\partial^{(w, v)} = \partial_w \partial^v$, so
$\partial_w G = \partial_w \partial^v f = \partial^{(w,v)} f = H$ as distributions on $\Omega$.
A function in $L_p(\Omega)$, $p \geq 1$, is
integrable on every compact subset of $\Omega$, which has finite measure. The statements on
$\Sob{m}{p}$ follow, with $L_p(\Omega') \supseteq L_p(\Omega)|_{\Omega'}$ and (4); in $L_p(V)$,
the zero class, and the class of a sum or a multiple, agree almost everywhere with $0$ and with
the sum or the multiple of the representatives, and a weak derivative in $L_p(V)$ may be
replaced by its class (4). For $m = 1$ the order $k = 0$ is (2).
\end{proof}

A function of $\Sob{1}{p}(\Omega)$ together with its weak derivatives along a basis is a single
element of a product of $L_p$ spaces, a jet. Linear operators on $\Sob{1}{p}$, such as the
extension operator of \Cref{thm:reflection-extension-operator}, act on jets.

\begin{definition}[Sobolev jets]
\label{def:sobolev-jets}
\lean{Rd.weakSobolevJet, Rd.mem_weakSobolevJet, Rd.sobolevJet, Rd.mem_sobolevJet}
\leanok
\uses{def:weak-line-derivative, lem:weak-derivative-basic}
Let $(e_i)_{i \in I}$ be the orthonormal basis of the dyadic grid of $V$, $\Omega \subseteq V$
open and $1 \leq p \leq \infty$. The \emph{jets} of $\Sob{1}{p}(\Omega)$ are the pairs
$(f, (g_i)_{i \in I})$ of elements of $L_p(\Omega)$ such that $g_i = \partial_{e_i} f$ weakly on
$\Omega$ for every $i \in I$; for $\Omega = V$ these are the jets of $\Sob{1}{p}(V)$. Both are
complex linear subspaces of $L_p \times L_p^I$, with the norm
$\max(\|f\|_p, \max_i \|g_i\|_p)$.
\end{definition}

\begin{lemma}[Properties of Sobolev jets]
\label{lem:sobolev-jets}
\lean{Rd.norm_le_norm_fst_add_sum, Rd.norm_fst_add_sum_le, Rd.toReal_enorm_add_sum,
  Rd.weakSobolevJet.snd_eq_of_fst_eq, Rd.weakSobolevJet.ext_of_fst_eq,
  Rd.exists_mem_weakSobolevJet_of_memWeakSobolev, Rd.sobolevJet.snd_eq_of_fst_eq,
  Rd.sobolevJet.ext_of_fst_eq, Rd.sobolevJet.lift, Rd.sobolevJet.lift_apply_fst,
  Rd.sobolevJet.lift_apply_snd, Rd.sobolevJet.restrictL, Rd.sobolevJet.restrictL_apply_fst,
  Rd.sobolevJet.restrictL_apply_snd}
\leanok
\uses{def:sobolev-jets, lem:weak-derivative-basic}
For every open $\Omega \subseteq V$ and every $1 \leq p \leq \infty$:
\begin{enumerate}
  \item for every jet $u = (f, (g_i)_i)$,
    $\|u\| \leq \|f\|_p + \sum_i \|g_i\|_p \leq (n + 1) \|u\|$;
  \item two jets of $\Sob{1}{p}(\Omega)$ with the same value $f$ are equal, and so are two jets of
    $\Sob{1}{p}(V)$;
  \item for every $f \in \Sob{1}{p}(\Omega)$ some jet of $\Sob{1}{p}(\Omega)$ has a value equal to
    $f$ almost everywhere on $\Omega$;
  \item for every complex vector space $M$ and every linear map $L \colon M \to L_p(V)$ such that
    for every $x \in M$ and every $i$ some $G \in L_p(V)$ is the weak derivative
    $\partial_{e_i} (L x)$ on $V$, the map $x \mapsto (L x, (\partial_{e_i} (L x))_i)$ is a
    linear map into the jets of $\Sob{1}{p}(V)$;
  \item the restriction $(F, (G_i)_i) \mapsto (F|_\Omega, (G_i|_\Omega)_i)$ is a bounded linear
    operator from the jets of $\Sob{1}{p}(V)$ to the jets of $\Sob{1}{p}(\Omega)$.
\end{enumerate}
\end{lemma}

\begin{proof}
\leanok
(1) Product and finite product norms are maxima. (2) Weak derivatives are unique almost
everywhere on $\Omega$, and in $L_p(V)$ for $\Omega = V$ (\Cref{lem:weak-derivative-basic}).
(3) Take the classes of $f$ and of its weak derivatives $\partial_{e_i} f \in L_p(\Omega)$.
(4) By (2) the map is additive and homogeneous, since $L$ is. (5) Restriction is bounded on each
component, and weak derivatives restrict to $\Omega$ (\Cref{lem:weak-derivative-basic}).
\end{proof}

The converse of (3) holds on every open set: the value of every jet of $\Sob{1}{p}(\Omega)$
lies in $\Sob{1}{p}(\Omega)$ (\Cref{lem:weak-derivative-direction-linear} (3)); on a bounded
Lipschitz domain it also follows from the extension operator
(\Cref{thm:reflection-extension-operator}).

\begin{lemma}[Complex test functions]
\label{lem:weak-derivative-complex-test}
\lean{Rd.exists_testSchwartz_eq_of_tsupport_subset, Rd.HasWeakLineDeriv.integral_mul_eq}
\leanok
\uses{lem:weak-derivative-integral, def:test-schwartz, lem:directional-derivatives-commute}
Every smooth compactly supported $w \colon V \to \C$ with support in $\Omega$ is
$w = a + i b$ for real test functions $a, b \in \mathcal D(\Omega)$. For every $v \in V^k$, all
$f$, $g$ with $g = \partial^v f$ weakly on $\Omega$ and every such $w$,
\[
  \int_V w \, g = (-1)^k \int_V (\partial^v w) \, f .
\]
\end{lemma}

\begin{proof}
\leanok
Take $a = \operatorname{Re} w$ and $b = \operatorname{Im} w$, whose supports lie in that of $w$.
Then $\partial^v w = \partial^v a + i \, \partial^v b$
(\Cref{lem:directional-derivatives-commute}), and \Cref{lem:weak-derivative-integral} for $a$ and
$b$ gives the identity.
\end{proof}

\begin{lemma}[The product rule for test functions]
\label{lem:test-function-product}
\lean{Rd.testFunctionMul, Rd.testFunctionMul_apply, Rd.lineDerivOp_testFunctionMul_apply}
\leanok
\uses{lem:directional-derivatives-commute}
For every test function $\varphi \in \mathcal D(\Omega)$ and every smooth real function $\chi$
on $V$, $\varphi \chi \in \mathcal D(\Omega)$ and
$\partial_w (\varphi \chi) = (\partial_w \varphi) \chi + \varphi \, \partial_w \chi$ for every
$w \in V$.
\end{lemma}

\begin{proof}
\leanok
The closed support of $\varphi \chi$ lies in that of $\varphi$, and the product rule for
derivatives gives $\partial_w (\varphi \chi)$, since $\partial_w \psi = D\psi(\cdot) \, w$ for
test functions $\psi$ (\Cref{lem:directional-derivatives-commute}).
\end{proof}

\begin{lemma}[The product rule with a cutoff]
\label{lem:weak-derivative-product-cutoff}
\lean{Rd.testFunctionMulOn, Rd.testFunctionMulOn_apply, Rd.lineDerivOp_testFunctionMulOn_apply,
  Rd.locallyIntegrableOn_mul_of_tsupport_subset, Rd.HasWeakLineDeriv.mul_of_tsupport_subset}
\leanok
\uses{lem:weak-derivative-integral, lem:test-function-product,
  lem:directional-derivatives-commute}
For all open $O, B \subseteq V$ and every smooth real function $\chi$ on $V$ with closed support
in $B$:
\begin{enumerate}
  \item for every $\varphi \in \mathcal D(O)$, $\varphi \chi \in \mathcal D(O \cap B)$ and
    $\partial_w (\varphi \chi) = (\partial_w \varphi) \chi + \varphi \, \partial_w \chi$ for every
    $w \in V$;
  \item for every continuous $c \colon V \to \C$ with closed support in $B$ and every locally
    integrable $u$ on $O \cap B$, $c u$ is locally integrable on $O$;
  \item for every $w \in V$ and all $f$, $g$ with $g = \partial_w f$ weakly on $O \cap B$,
    $\chi g + (\partial_w \chi) f = \partial_w (\chi f)$ weakly on $O$.
\end{enumerate}
\end{lemma}

\begin{proof}
\leanok
(1) The closed support of $\varphi \chi$ lies in those of $\varphi$ and of $\chi$, hence in
$O \cap B$, and as a function $\varphi \chi$ is the test function of
\Cref{lem:test-function-product} on $O$, with the same derivative. (2) Near a point of $B$, $c u$
is the product of a continuous function with a locally integrable function on the open set
$O \cap B$; near a point of $O$ outside $B$, which lies outside the closed support of $c$, $c u$
vanishes. (3) The closed support of $\partial_w \chi$ lies in $B$ as well, so by (2) the functions
$\chi f$, $\chi g$ and $(\partial_w \chi) f$ are locally integrable on $O$. For
$\varphi \in \mathcal D(O)$,
\[
  \int \varphi \bigl(\chi g + (\partial_w \chi) f\bigr) = -\int \partial_w(\varphi \chi) \, f
  + \int \varphi \, (\partial_w \chi) f = -\int (\partial_w \varphi) \, \chi f
\]
by \Cref{lem:weak-derivative-integral} on $O \cap B$ for the test function
$\varphi \chi \in \mathcal D(O \cap B)$ of (1), and by the formula for its derivative in (1).
\end{proof}

\begin{lemma}[The product rule for weak derivatives]
\label{lem:weak-derivative-product}
\lean{Rd.HasWeakLineDeriv.mul_contDiff}
\leanok
\uses{lem:weak-derivative-product-cutoff}
For every $w \in V$, all $f$, $g$ with $g = \partial_w f$ weakly on $\Omega$ and every smooth
real $\chi$, $\chi g + (\partial_w \chi) f = \partial_w (\chi f)$ weakly on $\Omega$.
\end{lemma}

\begin{proof}
\leanok
\Cref{lem:weak-derivative-product-cutoff} (3) with $\Omega$ for $O$ and $V$ for $B$, since
$\Omega \cap V = \Omega$.
\end{proof}

\begin{lemma}[Pairing with a test function]
\label{lem:test-function-pairing}
\lean{Rd.eLpNorm_testFunction_smul_le, Rd.eventually_memLp_of_tendsto,
  Rd.tendsto_integral_testFunction_smul}
\leanok
For every $1 \leq p \leq \infty$, every $\varphi \in \mathcal D(\Omega)$ with support $K$ and
every $h$ that is almost everywhere strongly measurable on $\Omega$,
\[
  \|\varphi h\|_{L_1(\Omega)} \leq \|\varphi\|_\infty \, |K|^{1 - 1/p} \, \|h\|_{L_p(\Omega)} .
\]
For every filter, every $h \in L_p(\Omega)$ and all $h_i$ with
$\|h_i - h\|_{L_p(\Omega)} \to 0$, eventually $h_i \in L_p(\Omega)$ and
$\int_V \varphi h_i \to \int_V \varphi h$.
\end{lemma}

\begin{proof}
\leanok
$\varphi h$ vanishes off the compact set $K \subseteq \Omega$, so its $L_1(\Omega)$ norm is that on
$K$, which is at most $|K|^{1 - 1/p}$ times its $L_p(K)$ norm (H\"older's inequality on the set
$K$ of finite measure), and $\|\varphi h\|_{L_p(K)} \leq \|\varphi\|_\infty \|h\|_{L_p(\Omega)}$.
Eventually $\|h_i - h\|_{L_p(\Omega)} < 1$, and then $h_i = (h_i - h) + h \in L_p(\Omega)$. The
integrals over $V$ are integrals over $\Omega$, and the bound applied to $h_i - h$ gives the
convergence in $L_1(\Omega)$ of $\varphi h_i$ to $\varphi h$.
\end{proof}

\begin{theorem}[Weak derivatives pass to limits]
\label{thm:weak-derivative-limits}
\lean{Rd.HasWeakLineDeriv.of_tendsto}
\leanok
\uses{lem:weak-derivative-integral, lem:test-function-pairing}
For all $1 \leq p, r \leq \infty$, every $v \in V^k$, every nontrivial filter, every
$f \in L_p(\Omega)$, every $g \in L_r(\Omega)$ and all $f_i$, $g_i$ with $g_i = \partial^v f_i$
weakly on $\Omega$ for every $i$, $\|f_i - f\|_{L_p(\Omega)} \to 0$ and
$\|g_i - g\|_{L_r(\Omega)} \to 0$, $g$ is the weak derivative $\partial^v f$ on $\Omega$.
\end{theorem}

\begin{proof}
\leanok
For every $\varphi \in \mathcal D(\Omega)$,
$\int \varphi g_i = (-1)^k \int (\partial^v \varphi) f_i$ by
\Cref{lem:weak-derivative-integral}; both sides converge by \Cref{lem:test-function-pairing},
to $\int \varphi g$ and to $(-1)^k \int (\partial^v \varphi) f$, and limits in $\C$ are unique.
\end{proof}

The bridge to tempered distributions rests on \Cref{thm:test-functions-determine}. For
$1 \leq p \leq \infty$ and $f \in L_p(V)$ let $T_f$ be $f$ as a distribution on $V$ and
$S_f$ be $f$ as a tempered distribution.

\begin{lemma}[$L_p$ functions as distributions on $V$]
\label{lem:lp-distribution-whole-space}
\lean{Rd.volume_restrict_top, Rd.locallyIntegrableOn_top_Lp, Rd.ofFun_top_apply,
  Rd.iteratedLineDerivOp_ofFun_top_apply}
\leanok
\uses{lem:distribution-transpose, def:test-schwartz, lem:directional-derivatives-commute,
  lem:weak-derivative-basic}
Lebesgue measure restricted to $\Omega = V$ is Lebesgue measure, so $L_p(\Omega) = L_p(V)$. For
every $1 \leq p \leq \infty$, every $f \in L_p(V)$, every $\varphi \in \mathcal D(V)$ and
every $v \in V^k$, $f$ is locally integrable on $V$, and
\[
  T_f(\varphi) = S_f(\varphi), \qquad \partial^v T_f(\varphi) = \partial^v S_f(\varphi),
\]
with $\varphi$ as a Schwartz function on the right.
\end{lemma}

\begin{proof}
\leanok
Both sides of the first identity are $\int \varphi f$. For the second, both sides are
$(-1)^k \int (\partial^v \varphi) f$ by \Cref{lem:distribution-transpose}, since the Schwartz
function $\partial^v \varphi$ is the iterated derivative of the Schwartz function $\varphi$
(\Cref{lem:directional-derivatives-commute}).
\end{proof}

\begin{theorem}[Weak derivatives are derivatives of tempered distributions]
\label{thm:weak-derivative-tempered}
\lean{Rd.hasWeakLineDeriv_top_iff}
\leanok
\uses{lem:lp-distribution-whole-space, thm:test-functions-determine, def:weak-line-derivative}
For all $1 \leq p, r \leq \infty$, every $f \in L_p(V)$, every $g \in L_r(V)$ and every
$v \in V^k$, $g$ is the weak derivative $\partial^v f$ on $V$ if and only if $\partial^v f = g$
as tempered distributions.
\end{theorem}

\begin{proof}
\leanok
By \Cref{lem:lp-distribution-whole-space} the distribution $\partial^v T_f - T_g$ on $V$ and the
tempered distribution $\partial^v S_f - S_g$ agree on test functions. If the tempered distribution
is zero, so is the distribution. If the distribution is zero, the tempered distribution vanishes
on test functions and is zero by \Cref{thm:test-functions-determine}.
\end{proof}

\begin{corollary}[Sobolev spaces by tempered distributions]
\label{cor:sobolev-tempered}
\lean{Rd.memSobolev_iff_tempered}
\leanok
\uses{thm:weak-derivative-tempered, lem:weak-derivative-basic}
For every $1 \leq p \leq \infty$, every $m \in \N$ and every $f \in L_p(V)$, $f \in \Sob{m}{p}(V)$
if and only if for every $k \leq m$ and every $v \in V^k$ there is $g \in L_p(V)$ with
$\partial^v f = g$ as tempered distributions.
\end{corollary}

\begin{proof}
\leanok
\Cref{thm:weak-derivative-tempered}, with a function in $L_p(V)$ replaced by its class, which
changes it only on a null set (\Cref{lem:weak-derivative-basic}).
\end{proof}

\begin{corollary}[Schwartz functions]
\label{cor:weak-derivative-schwartz}
\lean{Rd.iteratedLineDerivOp_toLp_eq, Rd.hasWeakLineDeriv_schwartz, Rd.memSobolev_toLp}
\leanok
\uses{lem:distribution-transpose, thm:weak-derivative-tempered, cor:sobolev-tempered,
  lem:weak-derivative-basic}
For every Schwartz function $\eta$, every $v \in V^k$ and every open $\Omega \subseteq V$, the
classical derivative $\partial^v \eta$ is the weak derivative of $\eta$ on $\Omega$. For every
$1 \leq p \leq \infty$ and every $m \in \N$, $\eta \in \Sob{m}{p}(V)$.
\end{corollary}

\begin{proof}
\leanok
As tempered distributions $\partial^v \eta$ is the Schwartz function $\partial^v \eta$
(\Cref{lem:distribution-transpose}); \Cref{thm:weak-derivative-tempered} with $p = r = 2$ gives
the weak derivative on $V$, and \Cref{lem:weak-derivative-basic} (5) on $\Omega$. The Sobolev
statement is \Cref{cor:sobolev-tempered} with $g = \partial^v \eta$.
\end{proof}

\section{Changes of variables and cutoffs}
\label{sec:weak-derivative-change-of-variables}

Throughout this section $(e_i)_{i \in I}$ is the orthonormal basis of the dyadic grid of $V$
(\Cref{def:dyadic-cubes}). For a function $f$ on $V$, $D_v f(x)$ is the line derivative
$\lim_{t \to 0} (f(x + t v) - f(x))/t$ where it exists and $0$ elsewhere (Mathlib's
\texttt{lineDeriv}), and $\partial_v f$ is a weak derivative; for a map $\Phi \colon V \to V$,
$D\Phi(x)$ is the derivative where it exists and $0$ elsewhere (Mathlib's \texttt{fderiv}). For
Lipschitz $f$ and $\Phi$ both exist almost everywhere by Rademacher's theorem. The section
proves that Lipschitz functions are weakly differentiable, that weak derivatives are linear in
the direction, $L_p$ bounds for the products with a cutoff $\chi$ supported in $B$, an
approximation by mollification on open sets in $L_1$, and the chain rule for Lipschitz maps that
preserve Lebesgue measure. These are the tools of the extension by reflection across Lipschitz
graphs \cite[Section~4.4]{EvansGariepy}, \cite[Chapter~9]{Brezis2011}.

\paragraph{Lipschitz functions.}

\begin{lemma}[Line derivatives of Lipschitz functions integrate to zero]
\label{lem:integral-line-derivative-zero}
\lean{Rd.integral_lineDeriv_eq_zero, Rd.integrable_lineDeriv, Rd.ae_re_lineDeriv_eq,
  Rd.ae_im_lineDeriv_eq, Rd.integral_lineDeriv_eq_zero_complex}
\leanok
For every finite-dimensional real normed space $F$ with a Haar measure, every Lipschitz function
$f \colon F \to \R$ with compact support and every $v \in F$, $\int D_v f = 0$. For every
real Banach space $G$, every Lipschitz function $f \colon F \to G$ with compact support and every
$v \in F$, $D_v f$ is integrable. For every Lipschitz function $f \colon F \to \C$ and every
$v \in F$, $\operatorname{Re} D_v f = D_v \operatorname{Re} f$ and
$\operatorname{Im} D_v f = D_v \operatorname{Im} f$ almost everywhere; if $f$ has
compact support, then $\int D_v f = 0$.
\end{lemma}

\begin{proof}
\leanok
Integration by parts for Lipschitz functions with the constant function $1$, whose line
derivatives vanish (Mathlib's \texttt{LipschitzWith.\allowbreak integral\_\allowbreak
lineDeriv\_\allowbreak mul\_\allowbreak eq}). The line derivative $D_v f$ is measurable (Mathlib's
\texttt{aestronglyMeasurable\_\allowbreak lineDeriv}), bounded by $C \|v\|$ for the Lipschitz
constant $C$, and zero outside the closed support of $f$. For complex $f$, $f$ is
differentiable almost everywhere by Rademacher's theorem, so the real and imaginary parts of
$D_v f$ are almost everywhere the line derivatives of the Lipschitz functions
$\operatorname{Re} f$ and $\operatorname{Im} f$.
\end{proof}

\begin{lemma}[Lipschitz functions are weakly differentiable]
\label{lem:lipschitz-weak-derivative}
\lean{Rd.integral_testFunction_mul_lineDeriv, Rd.hasWeakLineDeriv_of_lipschitzWith}
\leanok
\uses{lem:integral-line-derivative-zero, lem:weak-derivative-integral,
  lem:directional-derivatives-commute}
For every open $\Omega \subseteq V$, every $\varphi \in \mathcal D(\Omega)$, every Lipschitz
$u \colon V \to \R$ and every $v \in V$,
\[
  \int_V \varphi \, D_v u = -\int_V (\partial_v \varphi) \, u .
\]
For every Lipschitz $f \colon V \to \C$, every open $\Omega \subseteq V$ and every $v \in V$,
$D_v f$ is the weak derivative of $f$ on $\Omega$ in the direction $v$.
\end{lemma}

\begin{proof}
\leanok
The test function $\varphi$ is smooth with compact support, hence Lipschitz (Mathlib's
\texttt{ContDiff.\allowbreak lipschitzWith\_\allowbreak of\_\allowbreak hasCompactSupport}), and
Mathlib's integration by parts for Lipschitz functions gives
$\int (D_v u) \varphi = \int (D_{-v} \varphi) u$, with
$D_{-v} \varphi = -\partial_v \varphi$ (\Cref{lem:directional-derivatives-commute}). For
complex $f$, $f$ is continuous and $D_v f$ is measurable and bounded, so both are locally
integrable. The real and imaginary parts of $\int \varphi \, D_v f$ are
$\int \varphi \, D_v \operatorname{Re} f$ and
$\int \varphi \, D_v \operatorname{Im} f$ (\Cref{lem:integral-line-derivative-zero}), and
the first part for the Lipschitz functions $\operatorname{Re} f$ and $\operatorname{Im} f$ gives
the identity of \Cref{lem:weak-derivative-integral}.
\end{proof}

\paragraph{Directions.}

\begin{lemma}[Weak derivatives are linear in the direction]
\label{lem:weak-derivative-direction-linear}
\lean{Rd.locallyIntegrableOn_sum_smul, Rd.HasWeakLineDeriv.sum_smul_dir,
  Rd.locallyIntegrableOn_of_forall_gridBasis, Rd.hasWeakLineDeriv_of_forall_gridBasis,
  Rd.memSobolev_one_of_forall_gridBasis, Rd.memWeakSobolev_one_of_forall_gridBasis}
\leanok
\uses{lem:weak-derivative-integral, lem:weak-derivative-basic,
  lem:directional-derivatives-commute, def:dyadic-cubes}
For every open $\Omega \subseteq V$:
\begin{enumerate}
  \item for every finite set $s$, all locally integrable $g_i$ on $\Omega$ and all $c_i$ in $\R$
    or $\C$ ($i \in s$), $\sum_{i \in s} c_i g_i$ is locally integrable on $\Omega$;
  \item for every finite set $s$, all $v_i \in V$ and $c_i \in \R$ ($i \in s$), every locally
    integrable $f$ on $\Omega$ and all $g_i$ with $g_i = \partial_{v_i} f$ weakly on $\Omega$ for
    every $i \in s$, $\sum_{i \in s} c_i g_i = \partial_v f$ weakly on $\Omega$ for
    $v = \sum_{i \in s} c_i v_i$;
  \item for every $f$ and all $g_i$ with $g_i = \partial_{e_i} f$ weakly on $\Omega$ for every
    $i \in I$, $f$ is locally integrable on $\Omega$, and for every $v \in V$,
    $\sum_{i \in I} \langle v, e_i \rangle g_i = \partial_v f$ weakly on $\Omega$;
  \item for every $1 \leq p \leq \infty$, every $f \in L_p(V)$ and all $g_i \in L_p(V)$ with
    $g_i = \partial_{e_i} f$ weakly on $V$ for every $i \in I$, $f \in \Sob{1}{p}(V)$;
  \item for every $1 \leq p \leq \infty$, every $f \in L_p(\Omega)$ and all $g_i \in L_p(\Omega)$
    with $g_i = \partial_{e_i} f$ weakly on $\Omega$ for every $i \in I$,
    $f \in \Sob{1}{p}(\Omega)$.
\end{enumerate}
\end{lemma}

\begin{proof}
\leanok
(1) Induction on $s$. (2) For every $\varphi \in \mathcal D(\Omega)$,
$\partial_v \varphi = \sum_i c_i \, \partial_{v_i} \varphi$, since
$\partial_w \varphi = D\varphi(\cdot) \, w$ is linear in $w$
(\Cref{lem:directional-derivatives-commute}); integrate against $f$ and apply
\Cref{lem:weak-derivative-integral} for every $i$. (3) For $n \geq 1$, $f$ is locally integrable
as a function with a weak derivative; for $n = 0$, $V$ is a point and every function on it is
continuous. Then (2) with $v = \sum_i \langle v, e_i \rangle e_i$. (4) By (3) the class of
$\sum_i \langle v, e_i \rangle g_i$ in $L_p(V)$ is the weak derivative $\partial_v f$ for every
$v$, and \Cref{lem:weak-derivative-basic} for $m = 1$ applies. (5) By (3),
$\sum_i \langle v, e_i \rangle g_i \in L_p(\Omega)$ is the weak derivative $\partial_v f$ on
$\Omega$ for every $v$, and $f$ is its own derivative of order $0$
(\Cref{lem:weak-derivative-basic}).
\end{proof}

\paragraph{Cutoffs.}

\begin{lemma}[Bounded multipliers on $L_p$]
\label{lem:bounded-multiplier-lp}
\lean{MeasureTheory.eLpNorm_mul_le_of_forall_norm_le}
\leanok
For every measure space $(X, \mu)$, every $0 \leq p \leq \infty$, every $C \geq 0$, every almost
everywhere strongly measurable $c \colon X \to \C$ with $|c| \leq C$ and every $u \colon X \to \C$,
$\|c u\|_{L_p(\mu)} \leq C \, \|u\|_{L_p(\mu)}$.
\end{lemma}

\begin{proof}
\leanok
For almost everywhere strongly measurable $u$ this is the pointwise bound $|c u| \leq C |u|$.
For the other $u$, $\|u\|_{L_p(\mu)} = \infty$ in Mathlib, and for $C = 0$ the product $c u$ is
$0$.
\end{proof}

The product rule with a cutoff (\Cref{lem:weak-derivative-product-cutoff}) carries a weak
derivative from $O \cap B$ to $O$; the next lemma bounds the $L_p$ norms of the products.

\begin{lemma}[Cutoffs in $L_p$]
\label{lem:weak-derivative-cutoff}
\lean{Rd.eLpNorm_restrict_eq_inf_of_support_subset, Rd.eLpNorm_mul_add_fderiv_mul_le,
  Rd.eLpNorm_mul_le_of_tsupport_subset}
\leanok
\uses{lem:bounded-multiplier-lp}
For all open $O, B \subseteq V$:
\begin{enumerate}
  \item for every $1 \leq p \leq \infty$, every $w \in V$, every smooth real function $\chi$ on
    $V$ with closed support in $B$, all $C, C' \geq 0$ with $|\chi| \leq C$ and
    $|\partial_w \chi| \leq C'$ and all $f$, $g$,
    \[
      \|\chi g + (\partial_w \chi) f\|_{L_p(O)} \leq C \, \|g\|_{L_p(O \cap B)}
        + C' \, \|f\|_{L_p(O \cap B)} ;
    \]
  \item for every $0 \leq p \leq \infty$, every $h$ that vanishes outside $B$ has
    $\|h\|_{L_p(O)} = \|h\|_{L_p(O \cap B)}$, and for every continuous real $\chi$ with closed
    support in $B$ and $|\chi| \leq C$ and every $f$,
    $\|\chi f\|_{L_p(O)} \leq C \, \|f\|_{L_p(O \cap B)}$.
\end{enumerate}
\end{lemma}

\begin{proof}
\leanok
(1) The function vanishes outside $B$, so its norm in $L_p(O)$ is its norm in $L_p(O \cap B)$ by
(2); Minkowski's inequality and \Cref{lem:bounded-multiplier-lp} give the estimate. (2)
$h = \mathbf 1_B h$, and the $L_p(O)$ norm of $\mathbf 1_B h$ is the $L_p(O \cap B)$ norm of $h$;
$\chi f$ vanishes outside $B$, and \Cref{lem:bounded-multiplier-lp} bounds it.
\end{proof}

\begin{lemma}[Smooth cutoffs]
\label{lem:smooth-cutoff-compact}
\lean{Rd.exists_contDiff_eventuallyEq_one, Rd.exists_norm_fderiv_gridBasis_le,
  Rd.norm_le_one_of_forall_mem_Icc}
\leanok
For every compact $K \subseteq V$ and every open $U \supseteq K$ there is a smooth
$\chi \colon V \to [0, 1]$ with compact closed support in $U$ that equals $1$ on a neighbourhood
of $K$. For every smooth real $\chi$ with compact support there are $d_i \geq 0$ with
$|\partial_{e_i} \chi| \leq d_i$ for every $i \in I$, and every real function with values in
$[0, 1]$ has absolute value at most $1$.
\end{lemma}

\begin{proof}
\leanok
Choose a compact $L$ with $K \subseteq \operatorname{int} L$ and $L \subseteq U$, and an open
$W \supseteq L$ with compact closure in $U$ (Mathlib's \texttt{exists\_\allowbreak
compact\_\allowbreak between} and \texttt{exists\_\allowbreak open\_\allowbreak between\_\allowbreak
and\_\allowbreak isCompact\_\allowbreak closure}). Mathlib's
\texttt{exists\_\allowbreak contDiff\_\allowbreak support\_\allowbreak eq\_\allowbreak
eq\_\allowbreak one\_\allowbreak iff} gives a smooth $\chi$ with values in $[0, 1]$, support $W$
and $\chi = 1$ on $L$. For the bounds, $\partial_{e_i} \chi$ is continuous with compact support.
\end{proof}

\begin{lemma}[The cutoff bound on the whole space]
\label{lem:cutoff-bound-whole-space}
\lean{Rd.eLpNorm_mul_add_sum_le_of_tsupport_subset}
\leanok
\uses{lem:weak-derivative-cutoff, lem:smooth-cutoff-compact}
For every $1 \leq p \leq \infty$, every open $O \subseteq V$, all $d_i \geq 0$, every smooth real
$\chi$ with values in $[0, 1]$, closed support in $O$ and $|\partial_{e_i} \chi| \leq d_i$ for
every $i \in I$, and all $f$, $g_i \colon V \to \C$,
\begin{multline*}
  \|\chi f\|_{L_p(V)} + \sum_{i \in I} \|\chi g_i + (\partial_{e_i} \chi) f\|_{L_p(V)} \\
  \leq \Bigl(1 + \sum_{i \in I} d_i\Bigr) \Bigl(\|f\|_{L_p(O)} + \sum_{i \in I}
  \|g_i\|_{L_p(O)}\Bigr) .
\end{multline*}
\end{lemma}

\begin{proof}
\leanok
Here $|\chi| \leq 1$ (\Cref{lem:smooth-cutoff-compact}), so \Cref{lem:weak-derivative-cutoff} (2)
and (1) with $V$ for $O$ and $O$ for $B$ give $\|\chi f\|_{L_p(V)} \leq \|f\|_{L_p(O)}$ and
$\|\chi g_i + (\partial_{e_i} \chi) f\|_{L_p(V)} \leq \|g_i\|_{L_p(O)} + d_i \|f\|_{L_p(O)}$.
The sum of these bounds is at most the right side.
\end{proof}

\paragraph{Mollification on open sets.}

\begin{lemma}[The derivative of a mollification]
\label{lem:mollification-weak-derivative}
\lean{Rd.coe_lineDerivOp_toSchwartzMap, Rd.fderiv_conv_apply_eq_conv}
\leanok
\uses{lem:conv-derivatives, lem:weak-derivative-convolution-local, lem:weak-derivative-basic}
For every smooth $K \colon V \to \C$ with compact support and every $v \in V$, $K$ is a Schwartz
function with $\partial_v K = DK(\cdot) \, v$, and for all $h$, $g$ with $g = \partial_v h$
weakly on $V$ and every $x \in V$,
\[
  D(K * h)(x) \, v = (K * g)(x) .
\]
\end{lemma}

\begin{proof}
\leanok
A smooth function with compact support is a Schwartz function, and its directional derivative
as a Schwartz function is $DK(\cdot) v$. The function $h$ is locally integrable
(\Cref{lem:weak-derivative-basic}). By \Cref{lem:conv-derivatives},
$D(K * h)(x) v = ((\partial_v K) * h)(x)$, which is $(K * g)(x)$ by
\Cref{lem:weak-derivative-convolution-local} with $\Omega = V$.
\end{proof}

Let $\eta$ be the standard bump of \Cref{thm:smooth-atomic-decomposition}, a smooth nonnegative
function with integral $1$ supported in the unit ball, and $\eta_c(x) = c^n \eta(c x)$.

\begin{lemma}[Mollification on open sets]
\label{lem:mollification-open-set}
\lean{MeasureTheory.LocallyIntegrable.integrable_of_hasCompactSupport,
  Rd.tendsto_eLpNorm_mollify_sub_restrict, Rd.HasWeakLineDeriv.cutoff_top,
  Rd.exists_tendsto_contDiff_of_hasWeakLineDeriv}
\leanok
\uses{lem:weak-derivative-product-cutoff, lem:smooth-cutoff-compact,
  lem:mollification-weak-derivative, thm:approx-identity, lem:dilate, lem:conv-derivatives,
  thm:smooth-atomic-decomposition,
  lem:weak-derivative-basic}
\begin{enumerate}
  \item For every topological space $X$ with a measure and every locally integrable $u$ on $X$
    with compact support, $u$ is integrable. For every
    integrable $h$, every $K \subseteq V$ and every $h'$ with $h = h'$ almost everywhere on $K$,
    $\|\eta_{k+1} * h - h'\|_{L_1(K)} \to 0$ as $k \to \infty$.
  \item For every open $O \subseteq V$, every $w \in V$, all $f$, $g$ with $g = \partial_w f$
    weakly on $O$ and every smooth real $\chi$ with compact closed support in $O$, the functions
    $\chi f$ and $\chi g + (\partial_w \chi) f$ are integrable on $V$, and
    $\chi g + (\partial_w \chi) f = \partial_w (\chi f)$ weakly on $V$.
  \item For every open $O \subseteq V$, every compact $K \subseteq O$, every family
    $(w_i)_{i \in J}$ in $V$, every locally integrable $f$ on $O$ and all $g_i$ with
    $g_i = \partial_{w_i} f$ weakly on $O$ for every $i \in J$, there are smooth $u_k$ with
    compact support ($k \in \N$) such that $\|u_k - f\|_{L_1(K)} \to 0$ and
    $\|D u_k(\cdot) \, w_i - g_i\|_{L_1(K)} \to 0$ for every $i \in J$.
\end{enumerate}
\end{lemma}

\begin{proof}
\leanok
(1) Such a function is integrable on its compact closed support and zero outside it. The
mollifications converge to $h$ in $L_1(V)$ (\Cref{thm:approx-identity}), and on $K$, $h = h'$
almost everywhere. (2) \Cref{lem:weak-derivative-product-cutoff} (3) with $V$ for $O$ and $O$
for $B$; both
functions are locally integrable on $V$ (\Cref{lem:weak-derivative-basic}) with compact support,
so (1) applies. (3) Take $\chi$ of \Cref{lem:smooth-cutoff-compact} for $K \subseteq O$ and
$u_k = \eta_{k+1} * (\chi f)$, which is smooth with compact support (\Cref{lem:conv-derivatives},
\Cref{lem:dilate}). By \Cref{lem:mollification-weak-derivative} and (2),
$D u_k(\cdot) \, w_i = \eta_{k+1} * (\chi g_i + (\partial_{w_i} \chi) f)$. On $K$, $\chi = 1$ and
$D\chi = 0$, so $\chi f = f$ and $\chi g_i + (\partial_{w_i} \chi) f = g_i$ there, and (1) gives
both limits.
\end{proof}

\paragraph{The chain rule.}

\begin{lemma}[Compositions with measure-preserving maps]
\label{lem:composition-measure-preserving}
\lean{Rd.aestronglyMeasurable_comp_of_mapsTo, Rd.eLpNorm_comp_le_of_mapsTo}
\leanok
For all measure spaces $(X, \mu)$ and $(Y, \nu)$, every measurable $\Phi \colon X \to Y$ with
$\Phi_* \mu = \nu$, every $U \subseteq X$ and every measurable $O \subseteq Y$ with
$\Phi(U) \subseteq O$, every normed group $E$ and every $h \colon Y \to E$: if $h$ is almost
everywhere strongly measurable on $O$, then $h \circ \Phi$ is almost everywhere strongly
measurable on $U$; and for every $0 \leq p \leq \infty$,
$\|h \circ \Phi\|_{L_p(U)} \leq \|h\|_{L_p(O)}$.
\end{lemma}

\begin{proof}
\leanok
The map $\Phi$ carries $\mu$ restricted to $\Phi^{-1}(O)$ to $\nu$ restricted to $O$
(Mathlib's \texttt{MeasurePreserving.\allowbreak restrict\_\allowbreak preimage}), and
$U \subseteq \Phi^{-1}(O)$. Hence $h \circ \Phi$ is almost everywhere strongly measurable on
$\Phi^{-1}(O)$, so on $U$, and
$\|h \circ \Phi\|_{L_p(U)} \leq \|h \circ \Phi\|_{L_p(\Phi^{-1}(O))} = \|h\|_{L_p(O)}$. In
Mathlib the $L_p$ norm of a function that is not almost everywhere strongly measurable is
$\infty$, so the bound holds for every $h$.
\end{proof}

\begin{lemma}[The chain rule for smooth functions]
\label{lem:chain-rule-smooth}
\lean{Rd.lineDeriv_comp_eq_sum, Rd.hasWeakLineDeriv_comp_of_contDiff,
  Rd.eLpNorm_sum_inner_fderiv_mul_comp_le}
\leanok
\uses{lem:lipschitz-weak-derivative, lem:composition-measure-preserving, def:dyadic-cubes,
  lem:weak-derivative-basic, lem:bounded-multiplier-lp}
\begin{enumerate}
  \item For all $u \colon V \to \C$, $\Phi \colon V \to V$ and $x \in V$ such that $u$ is
    differentiable at $\Phi(x)$ and $\Phi$ at $x$, and for every $v \in V$,
    $D_v (u \circ \Phi)(x) = \sum_{i \in I} \langle D\Phi(x) v, e_i \rangle \,
    Du(\Phi(x)) \, e_i$.
  \item For every smooth $u \colon V \to \C$ with compact support, every Lipschitz
    $\Phi \colon V \to V$, every open $\Omega \subseteq V$ and every $v \in V$, the function
    $x \mapsto \sum_{i \in I} \langle D\Phi(x) v, e_i \rangle \, Du(\Phi(x)) \, e_i$ is the weak
    derivative of $u \circ \Phi$ on $\Omega$ in the direction $v$.
  \item For every $1 \leq p \leq \infty$, every $K \geq 0$, every $K$-Lipschitz
    $\Phi \colon V \to V$ that preserves Lebesgue measure, every $U \subseteq V$ and every
    measurable $O \subseteq V$ with $\Phi(U) \subseteq O$, every $v \in V$ and all $h_i$
    ($i \in I$),
    \[
      \Bigl\| \sum_{i \in I} \langle D\Phi(\cdot) v, e_i \rangle \, h_i \circ \Phi
      \Bigr\|_{L_p(U)} \leq K \|v\| \sum_{i \in I} \|h_i\|_{L_p(O)} .
    \]
\end{enumerate}
\end{lemma}

\begin{proof}
\leanok
(1) The chain rule and $D\Phi(x) v = \sum_i \langle e_i, D\Phi(x) v \rangle e_i$. (2) The smooth
$u$ with compact support is Lipschitz, hence so is $u \circ \Phi$, and
\Cref{lem:lipschitz-weak-derivative} shows that $D_v (u \circ \Phi)$ is its weak derivative. By
Rademacher's theorem (Mathlib's \texttt{LipschitzWith.\allowbreak ae\_\allowbreak
differentiableAt}) $\Phi$ is differentiable almost everywhere, where (1) applies, and weak
derivatives only see functions almost everywhere (\Cref{lem:weak-derivative-basic}).
(3) $\|D\Phi(x)\| \leq K$ for every $x$ (Mathlib's \texttt{norm\_\allowbreak fderiv\_\allowbreak
le\_\allowbreak of\_\allowbreak lipschitz}), so
$|\langle D\Phi(x) v, e_i \rangle| \leq K \|v\|$, and the coefficients are measurable (Mathlib's
\texttt{measurable\_\allowbreak fderiv\_\allowbreak apply\_\allowbreak const}). Minkowski's
inequality, \Cref{lem:bounded-multiplier-lp} and \Cref{lem:composition-measure-preserving} give
the bound.
\end{proof}

\begin{lemma}[Weak derivatives from compact pieces]
\label{lem:weak-derivative-local-to-global}
\lean{Rd.hasWeakLineDeriv_of_forall_isCompact}
\leanok
\uses{lem:weak-derivative-integral, lem:directional-derivatives-commute}
For every open $U \subseteq V$, every $v \in V^k$ and all $F$, $G$ such that every compact
$K \subseteq U$ lies in an open $U'$ with $G = \partial^v F$ weakly on $U'$,
$G = \partial^v F$ weakly on $U$.
\end{lemma}

\begin{proof}
\leanok
For $x \in U$ the compact set $\{x\}$ lies in such a $U'$, a neighbourhood of $x$, so $F$ and $G$
are locally integrable on $U$. Every $\varphi \in \mathcal D(U)$ is a test function on an open
$U'$ containing its closed support, with the same derivatives
(\Cref{lem:directional-derivatives-commute}), and \Cref{lem:weak-derivative-integral} on $U'$
gives the identity on $U$.
\end{proof}

\begin{theorem}[The chain rule for Lipschitz maps]
\label{thm:weak-derivative-chain-rule}
\lean{Rd.HasWeakLineDeriv.comp_lipschitzWith_of_mapsTo_isCompact,
  Rd.HasWeakLineDeriv.comp_lipschitzWith}
\leanok
\uses{lem:mollification-open-set, lem:chain-rule-smooth, lem:composition-measure-preserving,
  lem:weak-derivative-local-to-global, lem:weak-derivative-direction-linear,
  thm:weak-derivative-limits}
\begin{enumerate}
  \item For all open $O, U \subseteq V$, every compact $L \subseteq O$, every locally integrable
    $f$ on $O$, all $g_i$ with $g_i = \partial_{e_i} f$ weakly on $O$ for every $i \in I$, every
    $K \geq 0$, every $K$-Lipschitz $\Phi \colon V \to V$ that preserves Lebesgue measure with
    $\Phi(U) \subseteq L$, and every $v \in V$,
    $\partial_v (f \circ \Phi) = \sum_{i \in I} \langle D\Phi(\cdot) v, e_i \rangle \,
    g_i \circ \Phi$ weakly on $U$.
  \item For all open $O, U \subseteq V$, every $f$ and all $g_i$ with $g_i = \partial_{e_i} f$
    weakly on $O$ for every $i \in I$, every $K \geq 0$, every $K$-Lipschitz
    $\Phi \colon V \to V$ that preserves Lebesgue measure and maps $U$ into $O$, and every
    $v \in V$,
    \[
      \partial_v (f \circ \Phi)(x) = \sum_{i \in I} \langle D\Phi(x) v, e_i \rangle \,
      g_i(\Phi(x))
    \]
    weakly on $U$.
\end{enumerate}
\end{theorem}

\begin{proof}
\leanok
(1) Let $u_k$ be the smooth functions of \Cref{lem:mollification-open-set} for $L$ and the
directions $e_i$, and $G$ the right side. By \Cref{lem:chain-rule-smooth} (2),
$G_k(x) = \sum_i \langle D\Phi(x) v, e_i \rangle \, D u_k(\Phi(x)) \, e_i$ is the weak derivative
$\partial_v (u_k \circ \Phi)$ on $U$. By \Cref{lem:composition-measure-preserving},
$\|u_k \circ \Phi - f \circ \Phi\|_{L_1(U)} \leq \|u_k - f\|_{L_1(L)} \to 0$, and by
\Cref{lem:chain-rule-smooth} (3),
$\|G_k - G\|_{L_1(U)} \leq K \|v\| \sum_i \|D u_k(\cdot) \, e_i - g_i\|_{L_1(L)} \to 0$. The
same bounds give $f \circ \Phi, G \in L_1(U)$, since $f$ and the $g_i$ are integrable on the
compact $L \subseteq O$. \Cref{thm:weak-derivative-limits} with $p = r = 1$ gives
$G = \partial_v (f \circ \Phi)$ weakly on $U$.

(2) The function $f$ is locally integrable on $O$ (\Cref{lem:weak-derivative-direction-linear}).
By \Cref{lem:weak-derivative-local-to-global} it suffices to treat a compact $K' \subseteq U$. Its
image $\Phi(K')$ is a compact subset of $O$; choose a compact $L \subseteq O$ with
$\Phi(K') \subseteq \operatorname{int} L$. The open set $U' = \Phi^{-1}(\operatorname{int} L)$
contains $K'$, and $\Phi(U') \subseteq L$, so (1) applies on $U'$.
\end{proof}

Measure preservation alone carries the convergence in $L_1(L)$ to the convergence in $L_1(U)$,
so every Lipschitz map that preserves Lebesgue measure qualifies.

\section{Mathlib's Sobolev spaces \texorpdfstring{$H^s$}{H\^{}s} by Littlewood--Paley blocks}
\label{sec:sobolev-two-littlewood-paley}

Mathlib defines the Bessel potential $J_s u = \mathcal F^{-1}(\langle \xi \rangle^s \widehat u)$ of
a tempered distribution $u$, with $\langle \xi \rangle^s = (1 + |\xi|^2)^{s/2}$, and the Sobolev
space $H^s(V)$ of the tempered distributions $u$ with $J_s u \in L_2(V)$
(\texttt{TemperedDistribution.\allowbreak besselPotential},
\texttt{TemperedDistribution.\allowbreak MemSobolev}; the bundled space is
\texttt{BesselPotentialSpace}). This section identifies $H^s(V)$, $s \geq 0$, with the
Triebel--Lizorkin space $\TL{s}{2}{2}$ of \Cref{def:triebel-lizorkin} and, for $s > 0$, with the
Besov space $\Besov{s}{2}{2}$, with explicit constants, and proves $\Sob{1}{2}(V) = H^1(V)$. The
Fourier transform $\widehat F$ of $F \in L_2(V)$ is Plancherel's, Mathlib's linear isometry of
$L_2(V)$, and Mathlib identifies it with the Fourier transform of the tempered distribution of
$F$ (\texttt{Lp.\allowbreak fourier\_\allowbreak toTemperedDistribution\_\allowbreak eq}).

\begin{lemma}[Plancherel's theorem]
\label{lem:plancherel-l2}
\lean{Rd.two_ne_top, Rd.enorm_fourier_Lp_two, Rd.enorm_sq_eq_lintegral,
  Rd.lintegral_enorm_sq_eq_enorm_toLp}
\leanok
In $[0, \infty]$, $2 \neq \infty$. For every $F \in L_2(V)$, $\|\widehat F\|_2 = \|F\|_2$. On
every measure space, $\|F\|_2^2 = \int |F|^2$ for every $F \in L_2$, and
$\int |f|^2 = \|[f]\|_2^2$ for every measurable function $f$ with $\int |f|^2 < \infty$ and its
class $[f] \in L_2$.
\end{lemma}

\begin{proof}
\leanok
The Fourier transform on $L_2(V)$ is a linear isometry in Mathlib
(\texttt{Lp.\allowbreak norm\_\allowbreak fourier\_\allowbreak eq}), in particular continuous
(\texttt{FourierTransform.\allowbreak continuous\_\allowbreak fourier}). The third identity is
the definition of the $L_2$ norm, and the last one is the third for $F = [f]$, since $[f] = f$
almost everywhere.
\end{proof}

\begin{lemma}[$L_p$ functions as tempered distributions]
\label{lem:lp-distribution-determined}
\lean{Rd.eq_of_toTemperedDistribution_eq}
\leanok
For every $1 \leq p \leq \infty$, all $f, g \in L_p(V)$ with the same tempered distribution are
equal.
\end{lemma}

\begin{proof}
\leanok
Mathlib's \texttt{Lp.\allowbreak ker\_\allowbreak toTemperedDistributionCLM\_\allowbreak
eq\_\allowbreak bot}.
\end{proof}

\begin{lemma}[Almost everywhere equal classes are the same distribution]
\label{lem:ae-equal-classes-distribution}
\lean{Rd.toTemperedDistribution_eq_of_ae_eq}
\leanok
For all $1 \leq p, r \leq \infty$, every $f \in L_p(V)$ and every $g \in L_r(V)$ with $g = f$
almost everywhere, $g = f$ as tempered distributions.
\end{lemma}

\begin{proof}
\leanok
For every Schwartz function $u$, $\int_V u g = \int_V u f$, since the integrands agree almost
everywhere.
\end{proof}

\begin{lemma}[The norm of Mathlib's $H^{s,p}$]
\label{lem:bessel-space-enorm}
\lean{BesselPotentialSpace.enorm_toLp_eq}
\leanok
For every $s \in \R$, every $1 \leq p \leq \infty$ and every $U \in H^{s,p}(V)$,
$\|J_s U\|_p = \|U\|$ in $[0, \infty]$.
\end{lemma}

\begin{proof}
\leanok
Mathlib's \texttt{BesselPotentialSpace.\allowbreak norm\_\allowbreak toLp\_\allowbreak eq}.
\end{proof}

\begin{theorem}[The convolution theorem on $L_2$]
\label{thm:convolution-theorem-l2}
\lean{Rd.fourier_convLp_toLp, Rd.fourier_convLp, Rd.enorm_convLp_sq, Rd.fourier_lpBlock,
  Rd.enorm_lpBlock_sq}
\leanok
\uses{lem:fourier-integrable, lem:fourier-conv, lem:plancherel-l2, lem:conv-op, def:lp-blocks,
  lem:lp-kernels-basic, lem:lp-distribution-determined}
For every integrable $K$ on $V$ and every $f \in L_2(V)$,
$\widehat{K * f} = \widehat K \, \widehat f$ almost everywhere and
\[
  \|K * f\|_2^2 = \int_V |\widehat K(\xi)|^2 \, |\widehat f(\xi)|^2 \, d\xi .
\]
In particular, for every $f \in L_2(V)$ and every $j \in \N$,
$\widehat{\Delta_j f} = \widehat{\psi_j} \, \widehat f$ almost everywhere and
\[
  \|\Delta_j f\|_2^2 = \int_V |\widehat{\psi_j}(\xi)|^2 \, |\widehat f(\xi)|^2 \, d\xi .
\]
\end{theorem}

\begin{proof}
\leanok
Let first $f = \varphi$ be a Schwartz function. Then $K * \varphi$ is integrable, and for every
Schwartz function $u$, by \Cref{lem:fourier-integrable} and
$\widehat{K * \varphi} = \widehat K \, \widehat \varphi$ for integrable functions
(\Cref{lem:fourier-conv}),
\[
  \int \widehat u \, (K * \varphi) = \int u \, \widehat{K * \varphi}
  = \int u \, \widehat K \, \widehat \varphi .
\]
The left side is the Fourier transform of the tempered distribution of $K * \varphi$ applied to
$u$, which is the tempered distribution of the $L_2$ Fourier transform of $K * \varphi$ applied
to $u$. So the tempered distributions of the $L_2$ functions $\widehat{K * \varphi}$ and
$\widehat K \, \widehat \varphi$ agree, and an $L_2$ function is determined by its tempered
distribution (\Cref{lem:lp-distribution-determined}). Both sides of the identity are continuous
in $f \in L_2(V)$: the left by
Young's inequality (\Cref{lem:conv-op}) and \Cref{lem:plancherel-l2}, the right because
$\widehat K \in L_\infty$ (\Cref{lem:fourier-integrable}). Schwartz functions are dense in
$L_2(V)$, so the identity holds for all $f$. The norm identity follows with
\Cref{lem:plancherel-l2}, and the statements on the blocks are the case $K = \psi_j$.
\end{proof}

\begin{definition}[The Bessel weight]
\label{def:bessel-weight}
\lean{Rd.besselWeight}
\leanok
For $s \in \R$ and $\xi \in V$, $\langle \xi \rangle^s = (1 + |\xi|^2)^{s/2}$.
\end{definition}

The Bessel weight is the symbol of Mathlib's Bessel potential $J_s$:
$\widehat{J_s u} = \langle \cdot \rangle^s \widehat u$.

\begin{lemma}[Elementary properties of the weights]
\label{lem:bessel-weight-basic}
\lean{Rd.besselWeight_nonneg, Rd.besselWeight_neg, Rd.hasTemperateGrowth_besselWeight,
  Rd.enorm_besselWeight, Rd.ofReal_besselWeight_sq, Rd.norm_le_besselWeight_one,
  ENNReal.dyadicWeight_sq, Rd.one_add_norm_sq_le_of_norm_lt}
\leanok
\uses{def:bessel-weight, def:lq-norm}
For every $s \in \R$, the Bessel weight $\xi \mapsto \langle \xi \rangle^s$ is nonnegative, even
and smooth with polynomially bounded derivatives, $(\langle \xi \rangle^s)^2 = (1 + |\xi|^2)^s$ and
$|\xi| \leq \langle \xi \rangle^1$ for every $\xi \in V$, and $(2^{js})^2 = (4^j)^s$ for every
$j \in \N$. For every $j \in \N$ and every $\xi \in V$ with $|\xi| < 2^{j+1}$,
$1 + |\xi|^2 \leq 5 \cdot 4^j$.
\end{lemma}

\begin{proof}
\leanok
The weight depends on $|\xi|$ only, and the smoothness and the growth bounds of
$(1 + |\xi|^2)^{s/2}$ are in Mathlib. Further, $|\xi|^2 < 4^{j+1} = 4 \cdot 4^j$ and
$1 \leq 4^j$.
\end{proof}

\begin{definition}[The Bessel potential on Schwartz functions]
\label{def:bessel-schwartz}
\lean{Rd.besselSchwartz}
\leanok
\uses{def:bessel-weight, lem:bessel-weight-basic}
For $t \in \R$ and a Schwartz function $u$,
$J_t u = \mathcal F^{-1}(\langle \cdot \rangle^t \widehat u)$, a continuous linear map of the
Schwartz space.
\end{definition}

\begin{lemma}[$J_t$ on Schwartz functions and on distributions]
\label{lem:bessel-schwartz-distribution}
\lean{Rd.besselPotential_schwartz, Rd.besselSchwartz_apply,
  Rd.besselSchwartz_besselSchwartz_neg}
\leanok
\uses{def:bessel-schwartz}
For every $t \in \R$ and every Schwartz function $u$, the tempered distribution of $J_t u$ is
Mathlib's $J_t$ of the tempered distribution of $u$,
$J_t u(x) = \mathcal F^{-1}(\langle \cdot \rangle^t \widehat u)(x)$ for every $x \in V$, with the
Fourier integrals of functions, and $J_t J_{-t} u = u$.
\end{lemma}

\begin{proof}
\leanok
Mathlib's
\texttt{TemperedDistribution.\allowbreak fourierMultiplierCLM\_\allowbreak
toTemperedDistributionCLM\_\allowbreak eq}; the composition of two Fourier multipliers is the
multiplier of the product of their symbols, here
$\langle \xi \rangle^t \langle \xi \rangle^{-t} = 1$.
\end{proof}

For a function $m$ on $V$ that is smooth with polynomially bounded derivatives,
$m(D) u = \mathcal F^{-1}(m \, \widehat u)$ is Mathlib's Fourier multiplier on Schwartz functions
and on tempered distributions (\texttt{SchwartzMap.\allowbreak fourier\allowbreak MultiplierCLM},
\texttt{TemperedDistribution.\allowbreak fourier\allowbreak MultiplierCLM}); $J_t$ is the case
$m = \langle \cdot \rangle^t$. Mathlib lets $m(D)$ act on a tempered distribution $T$ by
$(m(D) T)(u) = T\bigl(\mathcal F(m \, \mathcal F^{-1} u)\bigr)$. The symbol $m$ is even if
$m(-\xi) = m(\xi)$ for every $\xi \in V$.

\begin{lemma}[Fourier multipliers multiply the Fourier transform]
\label{lem:fourier-multiplier-fourier}
\lean{Rd.fourier_fourierMultiplierCLM_apply}
\leanok
Let $m$ be a function on $V$, smooth with polynomially bounded derivatives. For every Schwartz
function $\varphi$ and every $\xi \in V$,
$\widehat{m(D) \varphi}(\xi) = m(\xi) \, \widehat\varphi(\xi)$.
\end{lemma}

\begin{proof}
\leanok
Fourier inversion on Schwartz functions.
\end{proof}

\begin{lemma}[Even Fourier multipliers act by the transpose]
\label{lem:even-multiplier-transpose}
\lean{Rd.fourierMultiplierCLM_apply_eq_fourier, Rd.fourierMultiplierCLM_apply_apply_of_even}
\leanok
\uses{lem:bessel-schwartz-distribution}
Let $m$ be an even function on $V$, smooth with polynomially bounded derivatives. For every
Schwartz function $u$ and every $x \in V$,
$m(D) u(x) = \mathcal F(m \, \mathcal F^{-1} u)(x)$, and for every tempered distribution $T$,
$(m(D) T)(u) = T(m(D) u)$.
\end{lemma}

\begin{proof}
\leanok
$\mathcal F^{-1} h(x) = \widehat{h}(-x)$ and $\mathcal F^{-1} u(\xi) = \widehat u(-\xi)$ (Mathlib's
\texttt{Real.\allowbreak fourierInv\_\allowbreak eq\_\allowbreak fourier\_\allowbreak
comp\_\allowbreak neg}), so
$m(D) u(x) = \mathcal F(m(-\cdot) \, \widehat u(-\cdot))(x) = \mathcal F(m \, \mathcal F^{-1} u)(x)$
since $m$ is even. The second identity is Mathlib's definition of $m(D) T$ with the first.
\end{proof}

\begin{lemma}[The transpose of the Bessel potential]
\label{lem:bessel-potential-transpose}
\lean{Rd.besselSchwartz_apply_eq_fourier, Rd.besselPotential_apply_apply,
  Rd.integral_besselSchwartz_mul}
\leanok
\uses{def:bessel-schwartz, lem:bessel-weight-basic, lem:bessel-schwartz-distribution,
  lem:even-multiplier-transpose}
For every $t \in \R$, every Schwartz function $u$ and every $x \in V$,
$J_t u(x) = \mathcal F(\langle \cdot \rangle^t \, \mathcal F^{-1} u)(x)$. For every $t \in \R$,
every tempered distribution $T$ and all Schwartz functions $\varphi, u$,
\[
  (J_t T)(u) = T(J_t u) \qquad \text{and} \qquad
  \int_V J_t \varphi \, u = \int_V \varphi \, J_t u .
\]
\end{lemma}

\begin{proof}
\leanok
The Bessel weight is even (\Cref{lem:bessel-weight-basic}), so \Cref{lem:even-multiplier-transpose}
gives the first two identities. For the tempered distribution of $\varphi$ the second gives
$\int_V J_t u \, \varphi = (J_t \varphi)(u) = \int_V u \, J_t \varphi$
(\Cref{lem:bessel-schwartz-distribution}).
\end{proof}

\begin{lemma}[Two elementary inequalities]
\label{lem:elementary-square-inequalities}
\lean{Rd.one_le_two_mul_enorm_sq_add, Rd.le_rpow_half_mul_of_sq_le}
\leanok
For all $x, y \in \C$ with $x + y = 1$, $1 \leq 2(|x|^2 + |y|^2)$. For all
$a, b, c \in [0, \infty]$ with $a^2 \leq c \, b^2$, $a \leq c^{1/2} \, b$.
\end{lemma}

\begin{proof}
\leanok
$1 \leq |x| + |y|$ by the triangle inequality, and $(|x| + |y|)^2 \leq 2(|x|^2 + |y|^2)$
(Mathlib's \texttt{add\_\allowbreak sq\_\allowbreak le}); and $(c^{1/2} b)^2 = c \, b^2$.
\end{proof}

\begin{proposition}[The Bessel weight and the Littlewood--Paley weights]
\label{prop:bessel-lp-weights}
\lean{Rd.dyadicWeight_sq_le_of_fourier_psi_ne_zero, Rd.ofReal_besselWeight_sq_le_of_norm_lt,
  Rd.tsum_dyadicWeight_fourier_psi_sq_le, Rd.ofReal_besselWeight_sq_le_tsum}
\leanok
\uses{lem:bessel-weight-basic, lem:elementary-square-inequalities, lem:fourier-psi-window}
For every $s \geq 0$, every $j \in \N$ and every $\xi \in V$: if $\widehat{\psi_j}(\xi) \neq 0$,
then $2^{2js} \leq 4^s \langle \xi \rangle^{2s}$; if $|\xi| < 2^{j+1}$, then
$\langle \xi \rangle^{2s} \leq 5^s \, 2^{2js}$; and
\[
  \sum_{j \in \N} 2^{2js} |\widehat{\psi_j}(\xi)|^2 \leq 2 \cdot 4^s \langle \xi \rangle^{2s},
  \qquad
  \langle \xi \rangle^{2s} \leq 2 \cdot 5^s \sum_{j \in \N} 2^{2js} |\widehat{\psi_j}(\xi)|^2 .
\]
\end{proposition}

\begin{proof}
\leanok
The first two bounds are $(4^j)^s \leq (4(1 + |\xi|^2))^s$ and
$(1 + |\xi|^2)^s \leq (5 \cdot 4^j)^s$ (\Cref{lem:fourier-psi-window} and
\Cref{lem:bessel-weight-basic}), using $s \geq 0$. Let $a$ be as in \Cref{lem:fourier-psi-window}.
The sum has at most the two terms $j = a, a + 1$, each at most $4^s \langle \xi \rangle^{2s}$, since
$|\widehat{\psi_j}(\xi)| \leq 1$. Both $j = a$ and $j = a + 1$ satisfy $|\xi| < 2^{j+1}$, and
$1 \leq 2 (|\widehat{\psi_a}(\xi)|^2 + |\widehat{\psi_{a+1}}(\xi)|^2)$ by
\Cref{lem:elementary-square-inequalities}, so
$\langle \xi \rangle^{2s} \leq 2 \sum_{j \in \{a, a+1\}} \langle \xi \rangle^{2s}
|\widehat{\psi_j}(\xi)|^2 \leq 2 \cdot 5^s \sum_{j \in \{a, a+1\}} 2^{2js}
|\widehat{\psi_j}(\xi)|^2$.
\end{proof}

\begin{lemma}[Temperate weights and $L_p$ functions]
\label{lem:temperate-weight-lp}
\lean{Rd.smulLeftCLM_toTemperedDistribution_eq, Rd.ae_eq_of_smulLeftCLM_eq}
\leanok
Let $w$ be a smooth function on $V$ with polynomially bounded derivatives, and
$1 \leq p, q \leq \infty$. For every $g \in L_p(V)$ with $w g \in L_q(V)$, the product of $w$
with the tempered distribution of $g$ is the tempered distribution of $w g$. For every
$G \in L_p(V)$ and every $H \in L_q(V)$ such that $w G = H$ as tempered distributions,
$w G = H$ almost everywhere.
\end{lemma}

\begin{proof}
\leanok
Both distributions in the first statement send a Schwartz function $u$ to $\int u \, w g$. In
the second, $w G$ and $H$ are locally integrable, and testing with the Schwartz functions
$\varphi \in C^\infty_c(V)$ gives $\int \varphi \, w G = \int \varphi \, H$; such functions
are equal almost everywhere (Mathlib's \texttt{ae\_\allowbreak eq\_\allowbreak of\_\allowbreak
integral\_\allowbreak contDiff\_\allowbreak smul\_\allowbreak eq}).
\end{proof}

\begin{lemma}[Bessel potentials of $L_2$ functions]
\label{lem:bessel-potential-l2}
\lean{Rd.fourier_ae_eq_of_besselPotential_eq, Rd.enorm_sq_eq_lintegral_of_besselPotential_eq,
  Rd.memLp_besselWeight_mul_fourier, Rd.memSobolev_two_of_lintegral_lt_top,
  Rd.exists_toTemperedDistribution_eq_of_besselPotential_eq}
\leanok
\uses{def:bessel-weight, lem:bessel-weight-basic, lem:temperate-weight-lp, lem:plancherel-l2}
Let $s \in \R$. For all $f, h \in L_2(V)$ with $J_s f = h$, $\widehat h = \langle \cdot \rangle^s
\widehat f$ almost everywhere and
\[
  \|h\|_2^2 = \int_V \langle \xi \rangle^{2s} |\widehat f(\xi)|^2 \, d\xi .
\]
For every $f \in L_2(V)$ with $\int \langle \xi \rangle^{2s} |\widehat f(\xi)|^2 \, d\xi < \infty$,
$\langle \cdot \rangle^s \widehat f \in L_2(V)$ and $f \in H^s(V)$. If $s \geq 0$, then every
tempered distribution $u$ with $J_s u \in L_2(V)$ is an $L_2(V)$ function.
\end{lemma}

\begin{proof}
\leanok
The Fourier transform of $J_s f$ is $\langle \cdot \rangle^s \widehat f$ as tempered
distributions, so the first identity is \Cref{lem:temperate-weight-lp}, and the second follows by
\Cref{lem:plancherel-l2}. In the second statement, $\langle \cdot \rangle^s \widehat f$ is the
Fourier transform of $J_s f$ by \Cref{lem:temperate-weight-lp}, so $J_s f$ is the inverse Fourier
transform of an $L_2$ function (Mathlib's \texttt{TemperedDistribution.\allowbreak
memSobolev\_\allowbreak iff\_\allowbreak exists\_\allowbreak smulLeftCLM\_\allowbreak fourier}).
For the third, $u = J_{-s}(J_s u)$, and $J_{-s}$ maps $L_2(V) = H^0(V)$ into $L_2(V)$ for
$-s \leq 0$ (Mathlib's \texttt{TemperedDistribution.\allowbreak MemSobolev.\allowbreak mono}).
\end{proof}

\begin{lemma}[The Triebel--Lizorkin norm of $\TL{s}{2}{2}$ on the frequency side]
\label{lem:tl-two-fourier}
\lean{Rd.triebelLizorkinNorm_two_sq, Rd.triebelLizorkinNorm_two_sq_le,
  Rd.lintegral_besselWeight_sq_le}
\leanok
\uses{def:triebel-lizorkin, thm:tl-diagonal, thm:convolution-theorem-l2,
  prop:bessel-lp-weights}
For every $s \in \R$ and every $f \in L_2(V)$,
\[
  \|f\|_{F^s_{2,2}}^2 = \int_V \Bigl( \sum_{j \in \N} 2^{2js} |\widehat{\psi_j}(\xi)|^2 \Bigr)
  |\widehat f(\xi)|^2 \, d\xi ,
\]
and if $s \geq 0$, then
$\|f\|_{F^s_{2,2}}^2 \leq 2 \cdot 4^s \int_V \langle \xi \rangle^{2s} |\widehat f(\xi)|^2 \, d\xi$
and $\int_V \langle \xi \rangle^{2s} |\widehat f(\xi)|^2 \, d\xi \leq 2 \cdot 5^s
\|f\|_{F^s_{2,2}}^2$.
\end{lemma}

\begin{proof}
\leanok
By \Cref{thm:tl-diagonal}, $\|f\|_{F^s_{2,2}}^2 = \sum_j 2^{2js} \|\Delta_j f\|_2^2$; insert
\Cref{thm:convolution-theorem-l2} and interchange sum and integral (Tonelli). The bounds follow
from \Cref{prop:bessel-lp-weights} under the integral.
\end{proof}

\begin{theorem}[$H^s = \TL{s}{2}{2}$]
\label{thm:sobolev-two-triebel-lizorkin}
\lean{Rd.triebelLizorkinNorm_le_of_besselPotential_eq,
  Rd.enorm_besselPotential_le_triebelLizorkinNorm, Rd.memSobolev_two_iff_memTriebelLizorkin}
\leanok
\uses{lem:bessel-potential-l2, lem:tl-two-fourier, lem:elementary-square-inequalities}
For every $s \geq 0$ and all $f, h \in L_2(V)$ with $J_s f = h$,
\[
  \|f\|_{F^s_{2,2}} \leq (2 \cdot 4^s)^{1/2} \, \|h\|_2
  \quad \text{and} \quad
  \|h\|_2 \leq (2 \cdot 5^s)^{1/2} \, \|f\|_{F^s_{2,2}} .
\]
For every $s \geq 0$, a tempered distribution $u$ lies in $H^s(V)$ if and only if $u$ is a
function $f \in L_2(V)$ with $f \in \TL{s}{2}{2}$.
\end{theorem}

\begin{proof}
\leanok
By \Cref{lem:bessel-potential-l2}, $\|h\|_2^2 = \int \langle \xi \rangle^{2s} |\widehat f|^2$,
and \Cref{lem:tl-two-fourier} gives both bounds after taking square roots
(\Cref{lem:elementary-square-inequalities}). If $u \in H^s(V)$, then $u = f \in L_2(V)$ by
\Cref{lem:bessel-potential-l2}, and the first bound shows $f \in \TL{s}{2}{2}$. Conversely, for
$f \in \TL{s}{2}{2}$, $\int \langle \xi \rangle^{2s} |\widehat f|^2 < \infty$ by
\Cref{lem:tl-two-fourier}, so $f \in H^s(V)$ by \Cref{lem:bessel-potential-l2}.
\end{proof}

\begin{corollary}[$H^s = \Besov{s}{2}{2}$]
\label{cor:sobolev-two-besov}
\lean{Rd.memSobolev_two_iff_exists_memApprox, Rd.memSobolev_two_iff_littlewoodPaleyNorm_lt_top,
  Rd.exists_eq_toDistr_two}
\leanok
\uses{thm:sobolev-two-triebel-lizorkin, thm:tl-diagonal, def:band-space, def:lp-blocks,
  lem:bessel-potential-l2, lem:bessel-space-enorm}
For every $s > 0$, a tempered distribution $u$ lies in $H^s(V)$ if and only if $u$ is a function
$f \in L_2(V)$ with $f \in \Besov{s}{2}{2}$, the approximation space $A^s_2$ of the
band-limited scheme. For every $s \geq 0$, $u \in H^s(V)$ if and only if $u$ is a function
$f \in L_2(V)$ with $\|f\|^{\mathrm{LP}}_{s,2} < \infty$, and every element $F$ of Mathlib's
bundled space $H^s(V)$ is a function $f \in \TL{s}{2}{2}$ with
$\|f\|_{F^s_{2,2}} \leq (2 \cdot 4^s)^{1/2} \|F\|_{H^s}$.
\end{corollary}

\begin{proof}
\leanok
\Cref{thm:sobolev-two-triebel-lizorkin} and \Cref{thm:tl-diagonal} with $p = 2$. An element $F$
of the bundled space carries $J_s F$ as an $L_2$ function of norm $\|F\|_{H^s}$
(\Cref{lem:bessel-space-enorm}).
\end{proof}

\begin{lemma}[Derivatives on the frequency side]
\label{lem:derivative-fourier-l2}
\lean{Rd.fourier_ae_eq_of_lineDerivOp_eq, Rd.norm_two_pi_I_inner_le, Rd.lintegral_inner_sq_eq,
  Rd.ofReal_besselWeight_one_sq, Rd.lintegral_besselWeight_one_sq_eq}
\leanok
\uses{lem:temperate-weight-lp, lem:plancherel-l2, lem:bessel-weight-basic}
For every $v \in V$ and all $f, g \in L_2(V)$ with $\partial_v f = g$ as tempered distributions,
$\widehat g(\xi) = 2 \pi i \langle \xi, v \rangle \widehat f(\xi)$ for almost every $\xi$ and
\[
  (2\pi)^2 \int_V \langle \xi, v \rangle^2 |\widehat f(\xi)|^2 \, d\xi = \|g\|_2^2 ;
\]
moreover $|2 \pi i \langle \xi, v \rangle| \leq 2 \pi |v| \langle \xi \rangle^1$ for every
$\xi \in V$. For every orthonormal basis $(e_i)_{i \in I}$ of $V$,
$(\langle \xi \rangle^1)^2 = 1 + \sum_i \langle \xi, e_i \rangle^2$ for every $\xi \in V$, and
for all $f \in L_2(V)$ and $g_i \in L_2(V)$ with $\partial_{e_i} f = g_i$ as tempered
distributions for every $i \in I$,
\[
  (2\pi)^2 \int_V (\langle \xi \rangle^1)^2 |\widehat f(\xi)|^2 \, d\xi
  = (2\pi)^2 \|f\|_2^2 + \sum_i \|g_i\|_2^2 .
\]
\end{lemma}

\begin{proof}
\leanok
The Fourier transform of $\partial_v f$ is $2 \pi i \langle \cdot, v \rangle \widehat f$ as
tempered distributions (Mathlib's \texttt{TemperedDistribution.\allowbreak fourier\_\allowbreak
lineDerivOp\_\allowbreak eq}), so the first identity is \Cref{lem:temperate-weight-lp} and the
second follows by \Cref{lem:plancherel-l2}. The bound is the Cauchy--Schwarz inequality with
$|\xi| \leq \langle \xi \rangle^1$. Parseval's identity
$|\xi|^2 = \sum_i \langle \xi, e_i \rangle^2$ (Mathlib's \texttt{OrthonormalBasis.\allowbreak
sum\_\allowbreak sq\_\allowbreak inner\_\allowbreak left}) gives the expansion of
$(\langle \xi \rangle^1)^2$; integrate it against $|\widehat f|^2$ and use the first part for
each $e_i$.
\end{proof}

\begin{theorem}[$\Sob{1}{2} = H^1$]
\label{thm:sobolev-one-two}
\lean{Rd.memSobolev_one_two_of_lineDerivOp_eq, Rd.memSobolev_one_iff,
  Rd.enorm_le_of_hasWeakLineDeriv, Rd.enorm_besselPotential_one_sq_eq,
  Rd.sq_norm_besselPotential_one_eq}
\leanok
\uses{def:weak-line-derivative, cor:sobolev-tempered, thm:weak-derivative-tempered,
  lem:derivative-fourier-l2, lem:bessel-potential-l2, def:dyadic-cells}
For every orthonormal basis $(e_i)_{i \in I}$ of $V$ and all $f, g_i \in L_2(V)$ such that
$\partial_{e_i} f = g_i$ as tempered distributions for every $i \in I$, $f \in H^1(V)$. For every
$f \in L_2(V)$, $f \in \Sob{1}{2}(V)$ if and only if $f \in H^1(V)$. For every
$v \in V$ and all $f, g, h \in L_2(V)$ such that $g$ is the weak derivative $\partial_v f$ on $V$
and $J_1 f = h$, $\|g\|_2 \leq 2 \pi |v| \, \|h\|_2$. For every orthonormal basis
$(e_i)_{i \in I}$ of $V$ and all $f, h, g_i \in L_2(V)$ such that $g_i$ is the weak derivative
$\partial_{e_i} f$ on $V$ for every $i \in I$ and $J_1 f = h$,
\[
  (2\pi)^2 \|h\|_2^2 = (2\pi)^2 \|f\|_2^2 + \sum_i \|g_i\|_2^2 .
\]
\end{theorem}

\begin{proof}
\leanok
Weak derivatives on $V$ of $L_2$ functions are derivatives of tempered distributions
(\Cref{thm:weak-derivative-tempered} and \Cref{cor:sobolev-tempered}). For derivatives
$g_i = \partial_{e_i} f \in L_2(V)$ along an orthonormal basis $(e_i)_i$,
$(2\pi)^2 \int (\langle \xi \rangle^1)^2 |\widehat f|^2 = (2\pi)^2 \|f\|_2^2 +
\sum_i \|g_i\|_2^2$ is finite (\Cref{lem:derivative-fourier-l2}), so $f \in H^1(V)$ by
\Cref{lem:bessel-potential-l2}; this applies to $f \in \Sob{1}{2}(V)$ with the orthonormal basis
of \Cref{def:dyadic-cells}. If $f \in H^1(V)$, then $\partial_v f \in
H^0(V) = L_2(V)$ for every $v \in V$ (Mathlib's \texttt{TemperedDistribution.\allowbreak
MemSobolev.\allowbreak lineDerivOp}), and the derivative of order $0$ is $f$ itself. The bound
compares $\widehat g = 2 \pi i \langle \cdot, v \rangle \widehat f$ with
$\widehat h = \langle \cdot \rangle^1 \widehat f$ pointwise
(\Cref{lem:derivative-fourier-l2} and \Cref{lem:bessel-potential-l2}) and uses
\Cref{lem:plancherel-l2}; the identity is \Cref{lem:derivative-fourier-l2} with
$\|h\|_2^2 = \int (\langle \xi \rangle^1)^2 |\widehat f|^2$.
\end{proof}

\begin{corollary}[The $H^1$ norm by a weak gradient]
\label{cor:weak-gradient-sobolev-norm}
\lean{Rd.enorm_besselPotential_one_sq_le, Rd.sq_norm_besselPotential_one_le}
\leanok
\uses{thm:sobolev-one-two}
For every orthonormal basis $(e_i)_{i \in I}$ of $V$ and all $f, h, g_i \in L_2(V)$ such that
$g_i$ is the weak derivative $\partial_{e_i} f$ on $V$ for every $i \in I$ and $J_1 f = h$,
\[
  \|h\|_2^2 \leq \|f\|_2^2 + \sum_i \|g_i\|_2^2 .
\]
\end{corollary}

\begin{proof}
\leanok
The identity $(2\pi)^2 \|h\|_2^2 = (2\pi)^2 \|f\|_2^2 + \sum_i \|g_i\|_2^2$ of
\Cref{thm:sobolev-one-two} and $(2\pi)^2 \geq 1$.
\end{proof}

\begin{theorem}[$\Sob{1}{2} = \Besov{1}{2}{2}$]
\label{thm:sobolev-one-two-besov}
\lean{Rd.memSobolev_one_iff_memApprox, Rd.memSobolev_one_iff_modulus}
\leanok
\uses{thm:sobolev-one-two, cor:sobolev-two-besov, thm:rd-besov, lem:lp-distribution-determined}
For every $F \in L_2(V)$, $F \in \Sob{1}{2}(V)$ if and only if $F$ lies in the approximation space
$A^1_2$ of the band-limited scheme. For every integer $r \geq 2$ and every $F \in L_2(V)$,
$F \in \Sob{1}{2}(V)$ if and only if
$\|F\|_2 + \|(2^k \omega_r(F, 2^{-k})_2)_k\|_{\ell_2} < \infty$.
\end{theorem}

\begin{proof}
\leanok
$\Sob{1}{2}(V) = H^1(V)$ by \Cref{thm:sobolev-one-two}, and $H^1(V) = \Besov{1}{2}{2}$ by
\Cref{cor:sobolev-two-besov}; an $L_2$ function is determined by its tempered distribution
(\Cref{lem:lp-distribution-determined}). The second statement is \Cref{thm:rd-besov} with
$p = q = 2$ and
$s = 1 < r$.
\end{proof}

\section{Rellich compactness}
\label{sec:rellich-compactness}

Throughout this section $\Omega \subseteq V$ is an arbitrary bounded set, $L_2(\Omega)$ is the
$L_2$ space of Lebesgue measure restricted to $\Omega$, and $F \mapsto F|_\Omega$ is the
restriction from $L_2(V)$ to $L_2(\Omega)$, Mathlib's operator of norm at most $1$
(\texttt{LpToLpRestrictCLM}). An element $F$ of Mathlib's bundled Sobolev space
$H^1(V) = H^{1,2}(V, \C)$ (\texttt{BesselPotentialSpace}) is a tempered distribution $u$ together
with the $L_2$ function $J_1 u$, and $\|F\|_{H^1} = \|J_1 u\|_2$. For $K \in \R$, $B_K$ is the
closed ball of radius $K$ around $0$ in $V$. The proof of Rellich's theorem truncates in
frequency: the part of $F$ with frequencies in $B_K$ is a bounded Lipschitz function on $V$, so
the truncated operators are compact by the Arzel\`a--Ascoli theorem on the compact closure of
$\Omega$, and they converge to the restriction in operator norm as $K \to \infty$.

\begin{theorem}[The inverse Fourier transform on $L_1 \cap L_2$]
\label{thm:fourier-inverse-l1-l2}
\lean{Rd.fourierInv_Lp_ae_eq_of_integrable}
\leanok
\uses{lem:fourier-integrable}
For every integrable $h \in L_2(V)$, the $L_2$ inverse Fourier transform of $h$ agrees almost
everywhere with the function
$x \mapsto \int e^{2 \pi i \langle \xi, x \rangle} h(\xi) \, d\xi$.
\end{theorem}

\begin{proof}
\leanok
Both functions are locally integrable, the second because it is continuous
(\Cref{lem:fourier-integrable}). Two locally integrable functions with the same integrals
against all smooth compactly supported $u$ agree almost everywhere (Mathlib's
\texttt{ae\_\allowbreak eq\_\allowbreak of\_\allowbreak integral\_\allowbreak contDiff\_\allowbreak
smul\_\allowbreak eq}), and such $u$ are Schwartz functions. The $L_2$ inverse Fourier transform
agrees with the inverse Fourier transform of tempered distributions (Mathlib's
\texttt{Lp.\allowbreak fourierInv\_\allowbreak toTemperedDistribution\_\allowbreak eq}), so
$\int u \, \mathcal F^{-1} h = h(\mathcal F^{-1} u) = \int (\mathcal F^{-1} u) \, h$, which is
$\int u \, \mathcal F^{-1} h$ for the pointwise inverse transform by
\Cref{lem:fourier-integrable}.
\end{proof}

\begin{lemma}[The Bessel weight of order $-1$]
\label{lem:bessel-weight-inverse}
\lean{Rd.besselWeight_neg_one_eq_inv, Rd.besselWeight_neg_one_le_one,
  Rd.besselWeight_neg_one_le_inv, Rd.norm_besselWeight_neg_one_le_one,
  Rd.memLp_top_besselWeight_neg_one, Rd.besselInvMul, Rd.besselInvMul_apply_ae,
  Rd.norm_besselInvMul_le_one}
\leanok
\uses{def:bessel-weight, lem:bessel-weight-basic}
For every $\xi \in V$, $\langle \xi \rangle^{-1} = (\langle \xi \rangle^1)^{-1} \leq 1$, and for
every $K > 0$ and every $\xi$ with $|\xi| \geq K$, $\langle \xi \rangle^{-1} \leq K^{-1}$. The
multiplication $g \mapsto \langle \cdot \rangle^{-1} g$ is a bounded linear operator on $L_2(V)$
of norm at most $1$.
\end{lemma}

\begin{proof}
\leanok
$\langle \xi \rangle^1 \geq \max(1, |\xi|)$ (\Cref{lem:bessel-weight-basic}). The weight
$\langle \cdot \rangle^{-1}$ is continuous and bounded by $1$, hence lies in $L_\infty(V)$, and
multiplication by an $L_\infty$ function is bounded on $L_2$ (Mathlib's Hölder multiplication
\texttt{ContinuousLinearMap.\allowbreak holderL}), with norm at most the bound $1$ of the weight.
\end{proof}

\begin{definition}[$H^1$ functions as $L_2$ functions]
\label{def:sobolev-one-to-l2}
\lean{Rd.toL2, Rd.toL2_apply}
\leanok
\uses{lem:bessel-weight-inverse, lem:plancherel-l2}
For $F \in H^1(V)$, $\iota F = \mathcal F^{-1}(\langle \cdot \rangle^{-1} \widehat{J_1 F})$,
with the Fourier transform of $L_2(V)$; $\iota : H^1(V) \to L_2(V)$ is a bounded linear operator.
\end{definition}

\begin{lemma}[$\iota F$ represents $F$]
\label{lem:sobolev-one-to-l2}
\lean{Rd.fourier_toL2, Rd.fourier_toL2_ae_eq, Rd.toTemperedDistribution_toL2,
  Rd.toL2_eq_of_toDistr_eq, Rd.toL2_injective, Rd.norm_toL2_le}
\leanok
\uses{def:sobolev-one-to-l2, lem:temperate-weight-lp, lem:bessel-weight-inverse,
  lem:plancherel-l2, lem:lp-distribution-determined}
For every $F \in H^1(V)$, given by the tempered distribution $u$, the Fourier transform of
$\iota F$ is $\langle \cdot \rangle^{-1} \widehat{J_1 u}$, the tempered distribution of $\iota F$
is $u$, and $\|\iota F\|_2 \leq \|F\|_{H^1}$. For every $F \in H^1(V)$ and every $f \in L_2(V)$
whose tempered distribution is that of $F$, $\iota F = f$. The map $\iota$ is injective.
\end{lemma}

\begin{proof}
\leanok
$u = J_{-1}(J_1 u)$, and the Fourier transform of $J_{-1} g$ is $\langle \cdot \rangle^{-1}
\widehat g$; for $g = J_1 u \in L_2(V)$ this is the tempered distribution of the $L_2$ function
$\langle \cdot \rangle^{-1} \widehat g$ (\Cref{lem:temperate-weight-lp}), whose inverse Fourier
transform is $\iota F$. An $L_2$ function is determined by its tempered distribution
(\Cref{lem:lp-distribution-determined}), and $F$ by $u$, so $\iota$ is injective. By
\Cref{lem:plancherel-l2} and
\Cref{lem:bessel-weight-inverse},
$\|\iota F\|_2 = \|\langle \cdot \rangle^{-1} \widehat{J_1 u}\|_2 \leq \|\widehat{J_1 u}\|_2
= \|J_1 u\|_2$.
\end{proof}

\begin{definition}[Truncation in frequency]
\label{def:frequency-truncation}
\lean{Rd.truncL2, Rd.truncL2_apply}
\leanok
\uses{def:sobolev-one-to-l2, lem:indicator-lp}
For $K \in \R$ and $F \in H^1(V)$, $T_K F = \mathcal F^{-1}(1_{B_K} \widehat{\iota F})$; $T_K$
is a bounded linear operator from $H^1(V)$ to $L_2(V)$.
\end{definition}

\begin{proposition}[The truncation error]
\label{prop:frequency-truncation-error}
\lean{Rd.norm_toL2_sub_truncL2_apply_le, Rd.norm_toL2_sub_truncL2_le, Rd.tendsto_truncL2}
\leanok
\uses{def:frequency-truncation, lem:sobolev-one-to-l2, lem:bessel-weight-inverse,
  lem:plancherel-l2}
For every $K > 0$ and every $F \in H^1(V)$, $\|\iota F - T_K F\|_2 \leq K^{-1} \|F\|_{H^1}$. For
every $K > 0$, $\|\iota - T_K\| \leq K^{-1}$, and $T_K \to \iota$ in operator norm as
$K \to \infty$.
\end{proposition}

\begin{proof}
\leanok
By \Cref{lem:plancherel-l2} and \Cref{lem:sobolev-one-to-l2}, $\|\iota F - T_K F\|_2$ is the
$L_2$ norm of $1_{V \setminus B_K} \langle \cdot \rangle^{-1} \widehat{J_1 F}$, and
$\langle \xi \rangle^{-1} \leq K^{-1}$ for $|\xi| > K$ (\Cref{lem:bessel-weight-inverse}).
\end{proof}

\begin{definition}[The square root of the volume of a ball]
\label{def:ball-volume-sqrt}
\lean{Rd.ballVolumeSqrt, Rd.ballVolumeSqrt_nonneg}
\leanok
For $K \in \R$, $C_K = |B_K|^{1/2} \geq 0$, the square root of the Lebesgue measure of $B_K$.
\end{definition}

\begin{lemma}[Cauchy--Schwarz on a ball]
\label{lem:cauchy-schwarz-ball}
\lean{Rd.eLpNorm_one_indicator_closedBall_le, Rd.integrable_indicator_closedBall,
  Rd.integral_norm_indicator_closedBall_le}
\leanok
\uses{def:ball-volume-sqrt}
For every $h \in L_2(V)$ and every $K \in \R$, $1_{B_K} h$ is integrable and
$\int |1_{B_K} h| \leq C_K \|h\|_2$.
\end{lemma}

\begin{proof}
\leanok
Hölder's inequality on the finite measure space $B_K$ (Mathlib's
\texttt{eLpNorm\_\allowbreak le\_\allowbreak eLpNorm\_\allowbreak mul\_\allowbreak rpow\_\allowbreak
measure\_\allowbreak univ}).
\end{proof}

\begin{lemma}[The Fourier character is Lipschitz]
\label{lem:fourier-character-lipschitz}
\lean{Rd.norm_fourierChar_sub_le}
\leanok
For all $a, b \in \R$, $|e^{2 \pi i a} - e^{2 \pi i b}| \leq 2 \pi |a - b|$.
\end{lemma}

\begin{proof}
\leanok
$e^{2 \pi i a} - e^{2 \pi i b} = e^{2 \pi i b}(e^{2 \pi i (a - b)} - 1)$ and
$|e^{i t} - 1| \leq |t|$ (Mathlib's \texttt{Real.\allowbreak norm\_\allowbreak
exp\_\allowbreak I\_\allowbreak mul\_\allowbreak ofReal\_\allowbreak sub\_\allowbreak
one\_\allowbreak le}).
\end{proof}

\begin{theorem}[Band-limited functions are bounded and Lipschitz]
\label{thm:band-limited-lipschitz}
\lean{Rd.norm_fourierInv_indicator_le, Rd.norm_fourierInv_sub_le,
  Rd.lipschitzWith_fourierInv_indicator}
\leanok
\uses{lem:fourier-integrable, lem:cauchy-schwarz-ball,
  lem:fourier-character-lipschitz}
For every integrable $g$ on $V$, every $K \in \R$ such that $g(\xi) = 0$ for almost every $\xi$
with $|\xi| > K$, and all $x, y \in V$,
\[
  |\mathcal F^{-1} g(x) - \mathcal F^{-1} g(y)| \leq 2 \pi K |x - y| \int |g| .
\]
For every $h \in L_2(V)$, every $K \in \R$ and all $x, y \in V$,
$|\mathcal F^{-1}(1_{B_K} h)(x)| \leq C_K \|h\|_2$ and
\[
  |\mathcal F^{-1}(1_{B_K} h)(x) - \mathcal F^{-1}(1_{B_K} h)(y)|
  \leq 2 \pi \max(K, 0) \, C_K \|h\|_2 \, |x - y| .
\]
\end{theorem}

\begin{proof}
\leanok
The difference is $\int (e^{2 \pi i \langle \xi, x \rangle} - e^{2 \pi i \langle \xi, y \rangle})
g(\xi) \, d\xi$, and for $|\xi| \leq K$, \Cref{lem:fourier-character-lipschitz} and the
Cauchy--Schwarz inequality give
$|e^{2 \pi i \langle \xi, x \rangle} - e^{2 \pi i \langle \xi, y \rangle}| \leq
2 \pi |\langle \xi, x - y \rangle| \leq 2 \pi K |x - y|$. For $g = 1_{B_K} h$, which vanishes for
$|\xi| > \max(K, 0)$, combine this and \Cref{lem:fourier-integrable} with
\Cref{lem:cauchy-schwarz-ball}.
\end{proof}

\begin{lemma}[Arzel\`a--Ascoli for Lipschitz functions]
\label{lem:lipschitz-totally-bounded}
\lean{Rd.totallyBounded_setOf_lipschitzWith}
\leanok
For every compact metric space $X$, every $C \in \R$ and every $L \geq 0$, the set of bounded
continuous functions $g : X \to \C$ with $|g| \leq C$ that are Lipschitz with constant $L$ is
totally bounded in the uniform norm.
\end{lemma}

\begin{proof}
\leanok
The set is equicontinuous, since all its members have the Lipschitz constant $L$, and its
functions take values in the compact disc of radius $C$; the Arzel\`a--Ascoli theorem (Mathlib's
\texttt{BoundedContinuousFunction.\allowbreak arzela\_\allowbreak ascoli}) makes its closure
compact.
\end{proof}

\begin{definition}[Continuous functions on the closure as $L_2(\Omega)$ functions]
\label{def:closure-extension}
\lean{Rd.extendClosure, Rd.closureToLp}
\leanok
For a bounded continuous function $g$ on $\overline \Omega$, $\tilde g$ is the function on $V$
equal to $g$ on $\overline \Omega$ and to $0$ elsewhere, and $\Phi g \in L_2(\Omega)$ is its
class.
\end{definition}

\begin{lemma}[The map $\Phi$ is Lipschitz]
\label{lem:closure-extension}
\lean{Rd.extendClosure_of_mem, Rd.measurable_extendClosure, Rd.norm_extendClosure_le,
  Rd.norm_extendClosure_sub_le, Rd.memLp_extendClosure, Rd.dist_closureToLp_le,
  Rd.lipschitzWith_closureToLp}
\leanok
\uses{def:closure-extension}
For every bounded continuous $g$ on $\overline \Omega$, $\tilde g$ is measurable with
$|\tilde g| \leq \|g\|_\infty$ and lies in $L_2(\Omega)$. For all bounded continuous $g, g'$ on
$\overline \Omega$, $|\tilde g - \tilde g'| \leq \|g - g'\|_\infty$ and
$\|\Phi g - \Phi g'\|_{L_2(\Omega)} \leq |\Omega|^{1/2} \|g - g'\|_\infty$.
\end{lemma}

\begin{proof}
\leanok
$\overline \Omega$ is closed, hence measurable, and $\Omega$ has finite measure since it is
bounded. The $L_2(\Omega)$ norm of a function bounded by $c$ is at most $|\Omega|^{1/2} c$.
\end{proof}

\begin{lemma}[The truncations on $\Omega$]
\label{lem:truncation-restriction-image}
\lean{Rd.truncL2_ae_eq_fourierInv_indicator, Rd.restr_truncL2_mem_closureToLp_image}
\leanok
\uses{def:frequency-truncation, lem:sobolev-one-to-l2, thm:fourier-inverse-l1-l2,
  lem:cauchy-schwarz-ball, thm:band-limited-lipschitz, lem:closure-extension}
For every $K \in \R$ and every $F \in H^1(V)$, $T_K F$ agrees almost everywhere with the function
$G(x) = \int e^{2 \pi i \langle \xi, x \rangle} 1_{B_K}(\xi) \widehat{\iota F}(\xi) \, d\xi$.
For every $K \in \R$ and every $F \in H^1(V)$ with $\|F\|_{H^1} \leq 1$, there is a bounded
continuous function $g$ on $\overline \Omega$ with $|g| \leq C_K$, Lipschitz with constant
$2 \pi \max(K, 0) \, C_K$, such that $(T_K F)|_\Omega = \Phi g$.
\end{lemma}

\begin{proof}
\leanok
Let $h = \widehat{\iota F}$. Then $1_{B_K} h$ is integrable (\Cref{lem:cauchy-schwarz-ball}),
so $T_K F = \mathcal F^{-1}(1_{B_K} h) = G$ almost everywhere (\Cref{thm:fourier-inverse-l1-l2}).
For $\|F\|_{H^1} \leq 1$, $\|h\|_2 = \|\iota F\|_2 \leq 1$ by \Cref{lem:sobolev-one-to-l2}; let
$g$ be the restriction of $G$ to $\overline \Omega$. It has the bound and the Lipschitz constant
by \Cref{thm:band-limited-lipschitz}, and $G = \tilde g$ on $\overline \Omega \supseteq \Omega$.
\end{proof}

\begin{proposition}[The truncations are compact]
\label{prop:truncation-compact}
\lean{Rd.isCompactOperator_restr_truncL2}
\leanok
\uses{lem:truncation-restriction-image, lem:lipschitz-totally-bounded, lem:closure-extension}
For every $K \in \R$, $F \mapsto (T_K F)|_\Omega$ is a compact operator from $H^1(V)$ to
$L_2(\Omega)$.
\end{proposition}

\begin{proof}
\leanok
$\overline \Omega$ is compact. By \Cref{lem:truncation-restriction-image}, the image of the
closed unit ball lies in the image under $\Phi$ of the set of \Cref{lem:lipschitz-totally-bounded}
with $X = \overline \Omega$, $C = C_K$ and $L = 2 \pi \max(K, 0) \, C_K$, which is totally
bounded; $\Phi$ is Lipschitz (\Cref{lem:closure-extension}), so the image of the unit ball is
totally bounded, and its closure is compact since $L_2(\Omega)$ is complete.
\end{proof}

\begin{theorem}[Rellich's theorem]
\label{thm:rellich-compactness}
\lean{Rd.isCompactOperator_restr_toL2, Rd.exists_subseq_tendsto_restr}
\leanok
\uses{prop:frequency-truncation-error, prop:truncation-compact}
For every bounded set $\Omega \subseteq V$, the restriction $F \mapsto (\iota F)|_\Omega$ from
$H^1(V)$ to $L_2(\Omega)$ is a compact operator. For every bounded set $\Omega \subseteq V$, every
$M \in \R$ and every sequence $(F_n)_n$ in $H^1(V)$ with $\|F_n\|_{H^1} \leq M$ for every $n$,
there are a strictly increasing $\varphi : \N \to \N$ and $u \in L_2(\Omega)$ with
$(\iota F_{\varphi(k)})|_\Omega \to u$ in $L_2(\Omega)$.
\end{theorem}

\begin{proof}
\leanok
The operators $F \mapsto (T_K F)|_\Omega$ are compact (\Cref{prop:truncation-compact}) and
converge in operator norm to $F \mapsto (\iota F)|_\Omega$ as $K \to \infty$, since the
restriction has norm at most $1$ (\Cref{prop:frequency-truncation-error}). A limit of compact
operators in operator norm is compact (Mathlib's \texttt{isCompactOperator\_\allowbreak
of\_\allowbreak tendsto}). The image of the closed ball of radius $M$ has compact closure, so the
sequence $((\iota F_n)|_\Omega)_n$ has a convergent subsequence.
\end{proof}

By \Cref{lem:sobolev-one-to-l2}, $\iota F$ is $F$ itself, read as an $L_2$ function. The theorem
holds for every bounded set $\Omega$. On a bounded Lipschitz domain the restriction space
$H^1(\Omega)$ of \Cref{def:sobolev-lipschitz-domain} consists of the restrictions of the
functions $\iota F$, since $\Sob{1}{2}(V) = H^1(V)$ (\Cref{thm:sobolev-one-two}) and every
$f \in H^1(V)$ is $\iota F$ for an $F \in H^1(V)$ (\Cref{lem:sobolev-one-to-l2}).

\begin{corollary}[Rellich's theorem by weak gradients]
\label{cor:rellich-weak-gradient}
\lean{Rd.exists_subseq_tendsto_restr_of_hasWeakLineDeriv}
\leanok
\uses{thm:rellich-compactness, thm:weak-derivative-tempered, thm:sobolev-one-two,
  cor:weak-gradient-sobolev-norm, lem:sobolev-one-to-l2}
Let $\Omega \subseteq V$ be bounded, $(e_i)_{i \in I}$ an orthonormal basis of $V$ and
$M \in \R$. For every sequence $(f_n)_n$ in $L_2(V)$ and all $g_{n,i} \in L_2(V)$ such that,
for every $n \in \N$, $g_{n,i}$ is the weak derivative $\partial_{e_i} f_n$ on $V$ for every
$i \in I$ and $\|f_n\|_2^2 + \sum_i \|g_{n,i}\|_2^2 \leq M^2$, there are a strictly increasing
$\varphi : \N \to \N$ and $u \in L_2(\Omega)$ with $f_{\varphi(k)}|_\Omega \to u$ in
$L_2(\Omega)$.
\end{corollary}

\begin{proof}
\leanok
Weak derivatives on $V$ are derivatives of tempered distributions
(\Cref{thm:weak-derivative-tempered}), so each $f_n$ lies in $H^1(V)$
(\Cref{thm:sobolev-one-two}) and $f_n = \iota F_n$ for an $F_n \in H^1(V)$
(\Cref{lem:sobolev-one-to-l2}) with $\|F_n\|_{H^1}^2 \leq M^2$
(\Cref{cor:weak-gradient-sobolev-norm}). Apply \Cref{thm:rellich-compactness} with the bound
$|M|$.
\end{proof}

\section{The scale of Sobolev spaces}
\label{sec:sobolev-scale}

This section builds the maps between Mathlib's Sobolev spaces $H^s(V) = H^{s,2}(V, \C)$ for
every real $s$: the inclusions $\iota_{t,s} : H^s(V) \to H^t(V)$ for $t \leq s$, the $L_2$
functions $\iota_s U$ of the elements $U$ of $H^s(V)$ for $s \geq 0$, and the derivatives
$\partial_v : H^{s+1}(V) \to H^s(V)$, with the exact recursion of the norms. Each map is a
Fourier multiplier between the $L_2$ functions $J_s U$, so it is bounded by the bound of its
symbol. It closes with $H^s(V) = \Besov{s}{2}{2}$ by differences, with explicit constants in
both directions, for every $s > 0$. As in \Cref{sec:rellich-compactness}, an element $U$ of
$H^s(V)$ is a tempered distribution $u$ together with the $L_2$ function $J_s u$, written
$J_s U$, and $\|U\|_{H^s} = \|J_s U\|_2$.

\begin{lemma}[Fourier multipliers of $L_2$ functions]
\label{lem:fourier-multipliers-l2}
\lean{Rd.fourierMultiplierCLM_toTemperedDistribution_eq,
  Rd.enorm_sq_eq_lintegral_of_fourier_ae_eq, Rd.lineDerivOp_besselPotential}
\leanok
\uses{lem:temperate-weight-lp, lem:plancherel-l2, def:bessel-weight}
For all $g, h \in L_2(V)$ and every function $w : V \to \C$ with $\widehat h = w \widehat g$
almost everywhere,
\[
  \|h\|_2^2 = \int_V |w(\xi)|^2 |\widehat g(\xi)|^2 \, d\xi .
\]
If moreover $w$ is smooth with polynomially bounded derivatives, then the Fourier multiplier
$w(D) = \mathcal F^{-1} w \mathcal F$ of tempered distributions maps $g$ to $h$. For every
$t \in \R$, every $v \in V$ and every tempered distribution $u$,
$\partial_v J_t u = J_t \partial_v u$.
\end{lemma}

\begin{proof}
\leanok
The identity is \Cref{lem:plancherel-l2} for $\widehat h$. The Fourier transform of the
tempered distribution $g$ is the $L_2$ function $\widehat g$, the product of $w$ with it is the
tempered distribution of the $L_2$ function $w \widehat g = \widehat h$
(\Cref{lem:temperate-weight-lp}), and its inverse Fourier transform is $h$. Both
$\partial_v = 2 \pi i \langle \cdot, v \rangle(D)$ and $J_t = \langle \cdot \rangle^t(D)$ are
Fourier multipliers (Mathlib's \texttt{TemperedDistribution.\allowbreak lineDeriv\_\allowbreak
eq\_\allowbreak fourierMultiplierCLM}), and the multipliers with two symbols compose to the
multiplier with their product (\texttt{TemperedDistribution.\allowbreak
fourierMultiplierCLM\_\allowbreak fourierMultiplierCLM\_\allowbreak apply}), which does not
depend on the order.
\end{proof}

\begin{lemma}[Bessel weights of nonpositive order as multipliers]
\label{lem:bessel-weight-multiplier}
\lean{Rd.besselWeight_le_one, Rd.norm_besselWeight_le_one, Rd.continuous_ofReal_besselWeight,
  Rd.besselWeightMul, Rd.besselWeightMul_apply_ae, Rd.norm_besselWeightMul_le_one,
  Rd.besselWeightMul_neg_one}
\leanok
\uses{def:bessel-weight, lem:bessel-weight-basic, lem:product-modulus, lem:bessel-weight-inverse}
For every $t \leq 0$ and every $\xi \in V$, $\langle \xi \rangle^t \leq 1$. For every $t \leq 0$
the multiplication $M_t g = \langle \cdot \rangle^t g$ is a bounded linear operator on $L_2(V)$
of norm at most $1$, and $M_{-1}$ is the multiplication by $\langle \cdot \rangle^{-1}$ of
\Cref{lem:bessel-weight-inverse}.
\end{lemma}

\begin{proof}
\leanok
$1 + |\xi|^2 \geq 1$ and $t/2 \leq 0$. The weight is smooth (\Cref{lem:bessel-weight-basic}),
and multiplication by a smooth bounded function is a bounded linear operator on $L_2$
(\Cref{lem:product-modulus} with $r = 0$), of norm at most the bound $1$ of the function. Both
operators in the last statement multiply by $\langle \cdot \rangle^{-1}$ almost everywhere.
\end{proof}

\begin{definition}[The inclusions]
\label{def:sobolev-inclusion}
\lean{Rd.sobolevInclusion, Rd.toLp_sobolevInclusion}
\leanok
\uses{lem:bessel-weight-multiplier, lem:plancherel-l2}
For $t \leq s$ and $U \in H^s(V)$, $\iota_{t,s} U \in H^t(V)$ is the element with the $L_2$
function
\[
  J_t \iota_{t,s} U = \mathcal F^{-1} \bigl(M_{t-s} \widehat{J_s U}\bigr) ,
\]
that is, $\iota_{t,s} U = J_{-t} \mathcal F^{-1}(\langle \cdot \rangle^{t-s} \widehat{J_s U})$;
$\iota_{t,s} : H^s(V) \to H^t(V)$ is a bounded linear operator.
\end{definition}

\begin{lemma}[The inclusions keep the distribution]
\label{lem:sobolev-inclusion}
\lean{Rd.fourier_toLp_sobolevInclusion_ae, Rd.toDistr_sobolevInclusion,
  Rd.norm_sobolevInclusion_apply_le, Rd.sobolevInclusion_injective,
  Rd.sobolevInclusion_sobolevInclusion}
\leanok
\uses{def:sobolev-inclusion, lem:fourier-multipliers-l2, lem:bessel-weight-multiplier,
  lem:plancherel-l2}
For all $t \leq s$ and $U \in H^s(V)$, the Fourier transform of $J_t \iota_{t,s} U$ is
$\langle \cdot \rangle^{t-s} \widehat{J_s U}$ almost everywhere, $\iota_{t,s} U = U$ as tempered
distributions, and $\|\iota_{t,s} U\|_{H^t} \leq \|U\|_{H^s}$. For all $t \leq s$ the operator
$\iota_{t,s}$ is injective, and for all $u \leq t \leq s$,
$\iota_{u,t} \circ \iota_{t,s} = \iota_{u,s}$.
\end{lemma}

\begin{proof}
\leanok
The first statement is \Cref{def:sobolev-inclusion} and $\mathcal F \mathcal F^{-1} = 1$. By
\Cref{lem:fourier-multipliers-l2}, $J_t \iota_{t,s} U = J_{t-s} J_s U$ as tempered
distributions, so $\iota_{t,s} U = J_{-t} J_{t-s} J_s U = J_0 U = U$ (Mathlib's
\texttt{TemperedDistribution.\allowbreak besselPotential\_\allowbreak besselPotential\_\allowbreak
apply} and \texttt{besselPotential\_\allowbreak zero}). By \Cref{lem:plancherel-l2}
and \Cref{lem:bessel-weight-multiplier},
$\|\iota_{t,s} U\|_{H^t} = \|M_{t-s} \widehat{J_s U}\|_2 \leq \|\widehat{J_s U}\|_2 =
\|U\|_{H^s}$. An element of $H^t(V)$ is determined by its tempered distribution
(Mathlib's \texttt{BesselPotentialSpace.\allowbreak ext}), which gives the last two statements.
\end{proof}

\begin{definition}[$L_2$ functions of $H^s(V)$]
\label{def:sobolev-to-l2}
\lean{Rd.sobolevToL2, Rd.sobolevToL2_apply}
\leanok
\uses{def:sobolev-inclusion}
For $s \geq 0$ and $U \in H^s(V)$, $\iota_s U = J_0 \iota_{0,s} U \in L_2(V)$ is the $L_2$
function of $\iota_{0,s} U \in H^0(V)$; $\iota_s : H^s(V) \to L_2(V)$ is a bounded linear
operator.
\end{definition}

\begin{lemma}[$\iota_s U$ represents $U$]
\label{lem:sobolev-to-l2}
\lean{Rd.toTemperedDistribution_sobolevToL2, Rd.besselPotential_sobolevToL2,
  Rd.sobolevToL2_eq_of_toDistr_eq, Rd.sobolevToL2_injective, Rd.norm_sobolevToL2_apply_le,
  Rd.sobolevToL2_one, Rd.sobolevToL2_sobolevInclusion}
\leanok
\uses{def:sobolev-to-l2, lem:sobolev-inclusion, lem:lp-distribution-determined,
  def:sobolev-one-to-l2, lem:sobolev-one-to-l2}
Let $s \geq 0$. For every $U \in H^s(V)$, the tempered distribution of $\iota_s U$ is $U$,
$J_s \iota_s U = J_s U$ and $\|\iota_s U\|_2 \leq \|U\|_{H^s}$. For every $U \in H^s(V)$ and
every $f \in L_2(V)$ whose tempered distribution is $U$, $\iota_s U = f$. The operator
$\iota_s$ is injective, $\iota_1$ is the operator $\iota$ of \Cref{def:sobolev-one-to-l2}, and
for every $0 \leq t \leq s$ and every $U \in H^s(V)$, $\iota_t \iota_{t,s} U = \iota_s U$.
\end{lemma}

\begin{proof}
\leanok
$J_0$ is the identity, so the tempered distribution of $\iota_s U$ is $\iota_{0,s} U = U$
(\Cref{lem:sobolev-inclusion}), and $\|\iota_s U\|_2 = \|\iota_{0,s} U\|_{H^0} \leq
\|U\|_{H^s}$. An $L_2$ function is determined by its tempered distribution
(\Cref{lem:lp-distribution-determined}), and an element of $H^s(V)$ by its distribution; this
gives the remaining statements, with \Cref{lem:sobolev-one-to-l2} for $\iota_1 = \iota$.
\end{proof}

\begin{lemma}[The symbol of $J_{-1} \partial_v$]
\label{lem:line-derivative-symbol}
\lean{Rd.besselLineDerivSymbol, Rd.hasTemperateGrowth_besselLineDerivSymbol,
  Rd.continuous_besselLineDerivSymbol, Rd.norm_besselLineDerivSymbol,
  Rd.norm_besselLineDerivSymbol_le, Rd.besselPotential_neg_one_lineDerivOp}
\leanok
\uses{def:bessel-weight, lem:bessel-weight-basic, lem:bessel-weight-inverse}
For $v \in V$ let $m_v(\xi) = 2 \pi i \langle \xi, v \rangle \langle \xi \rangle^{-1}$. Then
$m_v$ is smooth with polynomially bounded derivatives, and for every $\xi \in V$,
\[
  |m_v(\xi)| = 2 \pi |\langle \xi, v \rangle| \langle \xi \rangle^{-1} \leq 2 \pi |v| .
\]
For every tempered distribution $u$, $J_{-1} \partial_v u = m_v(D) u$.
\end{lemma}

\begin{proof}
\leanok
$|\langle \xi, v \rangle| \leq |\xi| |v|$ by the Cauchy--Schwarz inequality, and
$|\xi| \leq \langle \xi \rangle^1 = (\langle \xi \rangle^{-1})^{-1}$
(\Cref{lem:bessel-weight-basic}, \Cref{lem:bessel-weight-inverse}). The last statement is
Mathlib's \texttt{TemperedDistribution.\allowbreak besselPotential\_\allowbreak neg\_\allowbreak
one\_\allowbreak lineDerivOp\_\allowbreak eq} with the factor $2 \pi i$ moved into the symbol.
\end{proof}

\begin{definition}[Derivatives in the scale]
\label{def:sobolev-line-derivative}
\lean{Rd.sobolevLineDeriv, Rd.toLp_sobolevLineDeriv}
\leanok
\uses{lem:line-derivative-symbol, lem:product-modulus, lem:plancherel-l2}
For $s \in \R$, $v \in V$ and $U \in H^{s+1}(V)$, $\partial_v U \in H^s(V)$ is the element with
the $L_2$ function
\[
  J_s \partial_v U = \mathcal F^{-1} \bigl(m_v \widehat{J_{s+1} U}\bigr) ;
\]
$\partial_v : H^{s+1}(V) \to H^s(V)$ is a bounded linear operator.
\end{definition}

\begin{lemma}[Derivatives of tempered distributions]
\label{lem:sobolev-line-derivative}
\lean{Rd.fourier_toLp_sobolevLineDeriv_ae, Rd.toDistr_sobolevLineDeriv,
  Rd.norm_sobolevLineDeriv_apply_le}
\leanok
\uses{def:sobolev-line-derivative, lem:fourier-multipliers-l2, lem:line-derivative-symbol,
  lem:product-modulus, lem:plancherel-l2}
For all $s \in \R$, $v \in V$ and $U \in H^{s+1}(V)$, the Fourier transform of $J_s \partial_v U$
is $m_v \widehat{J_{s+1} U}$ almost everywhere, $\partial_v U$ is the derivative of the tempered
distribution $U$ in the direction $v$, and
$\|\partial_v U\|_{H^s} \leq 2 \pi |v| \, \|U\|_{H^{s+1}}$.
\end{lemma}

\begin{proof}
\leanok
By \Cref{lem:fourier-multipliers-l2} and \Cref{lem:line-derivative-symbol},
$J_s \partial_v U = m_v(D) J_{s+1} U = J_{-1} \partial_v J_{s+1} U$ as tempered distributions,
so $\partial_v U = J_{-s} J_{-1} \partial_v J_{s+1} U = J_{-s} J_{-1} J_{s+1} \partial_v U$, since
$\partial_v$ commutes with $J_{s+1}$ (\Cref{lem:fourier-multipliers-l2}), and
$J_{-s} J_{-1} J_{s+1} = J_0 = 1$. The bound is \Cref{lem:plancherel-l2} with
$\|m_v F\|_2 \leq 2 \pi |v| \, \|F\|_2$ for $F \in L_2(V)$ (\Cref{lem:line-derivative-symbol},
\Cref{lem:product-modulus}).
\end{proof}

\begin{theorem}[The norm recursion]
\label{thm:sobolev-norm-recursion}
\lean{Rd.besselWeight_neg_one_sq_mul, Rd.sq_besselWeight_add_sum_sq_norm_besselLineDerivSymbol,
  Rd.enorm_sq_sobolevInclusion_add_one, Rd.enorm_sq_sobolevLineDeriv,
  Rd.enorm_sq_eq_add_sum_sobolevLineDeriv, Rd.sq_norm_eq_add_sum_sobolevLineDeriv}
\leanok
\uses{lem:sobolev-inclusion, lem:sobolev-line-derivative, lem:fourier-multipliers-l2,
  lem:line-derivative-symbol, lem:plancherel-l2}
For every orthonormal basis $(e_i)_{i \in I}$ of $V$, every $s \in \R$ and every
$U \in H^{s+1}(V)$,
\[
  (2\pi)^2 \|U\|_{H^{s+1}}^2 = (2\pi)^2 \|\iota_{s,s+1} U\|_{H^s}^2
  + \sum_{i \in I} \|\partial_{e_i} U\|_{H^s}^2 .
\]
\end{theorem}

\begin{proof}
\leanok
Let $G = \widehat{J_{s+1} U}$. By \Cref{lem:plancherel-l2}, $\|U\|_{H^{s+1}}^2 = \int |G|^2$; by
\Cref{lem:sobolev-inclusion}, \Cref{lem:sobolev-line-derivative} and
\Cref{lem:fourier-multipliers-l2},
$\|\iota_{s,s+1} U\|_{H^s}^2 = \int (\langle \xi \rangle^{-1})^2 |G(\xi)|^2 \, d\xi$ and
$\|\partial_{e_i} U\|_{H^s}^2 = \int |m_{e_i}(\xi)|^2 |G(\xi)|^2 \, d\xi$. Parseval's identity
$\sum_i \langle \xi, e_i \rangle^2 = |\xi|^2$ (Mathlib's \texttt{OrthonormalBasis.\allowbreak
sum\_\allowbreak sq\_\allowbreak inner\_\allowbreak left}) and
$(\langle \xi \rangle^{-1})^2 (1 + |\xi|^2) = 1$ give
$(2\pi)^2 (\langle \xi \rangle^{-1})^2 + \sum_i |m_{e_i}(\xi)|^2 = (2\pi)^2$ for every $\xi$
(\Cref{lem:line-derivative-symbol}); multiply by $|G(\xi)|^2$ and integrate.
\end{proof}

At $s = 0$ the recursion is the identity of \Cref{thm:sobolev-one-two}: $\iota_{0,1} U$ is the
$L_2$ function $\iota U$, and $\partial_{e_i} U$ the weak derivative.

\begin{theorem}[$H^s(V)$ and $\Besov{s}{2}{2}$ by differences, with explicit constants]
\label{thm:sobolev-besov-constants}
\lean{Rd.sobolevBesovConst, Rd.sobolevBesovConst_ne_top, Rd.modNorm_sobolevToL2_le,
  Rd.sobolevBesovInvConst, Rd.sobolevBesovInvConst_ne_top, Rd.enorm_le_modNorm_sobolevToL2}
\leanok
\uses{thm:modulus-coincidence, thm:rd-besov, thm:tl-besov-sandwich, lem:tl-sandwich-constants,
  thm:sobolev-two-triebel-lizorkin, lem:sobolev-to-l2, lem:bessel-space-enorm}
Let $s > 0$ and let $r$ be an integer with $s < r$. Let $K_{s,r}$ be the constant of the inverse
direction of \Cref{thm:modulus-coincidence} for the band-limited scheme with the moduli of order
$r$ (\Cref{thm:rd-besov}) at smoothness $s$, and
$C_{s,r} = K_{s,r} \cdot 2 (1 - 2^{-s})^{-1} \cdot (2 \cdot 4^s)^{1/2}$. Then
$C_{s,r} < \infty$, and for every $U \in H^s(V)$,
\[
  \|\iota_s U\|_2 + \|(2^{ks} \omega_r(\iota_s U, 2^{-k})_2)_k\|_{\ell_2} \leq C_{s,r} \|U\| .
\]
Let $s \geq 0$ and $r \in \N$, let $D_{s,r}$ be the constant of the direct direction of
\Cref{thm:modulus-coincidence} for the band-limited scheme with the moduli of order $r$ at
smoothness $s$, $\Phi$ the kernel of \Cref{thm:tl-besov-sandwich}, and
$C'_{s,r} = (2 \cdot 5^s)^{1/2} \cdot 8 \cdot 8^s \, \|\Phi\|_1 D_{s,r}$. Then
$C'_{s,r} < \infty$, and for every $U \in H^s(V)$,
\[
  \|U\| \leq C'_{s,r} \bigl(\|\iota_s U\|_2 +
  \|(2^{ks} \omega_r(\iota_s U, 2^{-k})_2)_k\|_{\ell_2}\bigr) .
\]
\end{theorem}

\begin{proof}
\leanok
$J_s \iota_s U$ is the $L_2$ function of $U$ (\Cref{lem:sobolev-to-l2}), of norm $\|U\|$
(\Cref{lem:bessel-space-enorm}). For the first bound, the left side is at most $K_{s,r}$ times
the approximation norm of $\iota_s U$ in $A^s_2$ of the band-limited scheme
(\Cref{thm:modulus-coincidence}), which is at most
$2 (1 - 2^{-s})^{-1} \|\iota_s U\|_{F^s_{2,2}}$ (\Cref{thm:tl-besov-sandwich}), and
$\|\iota_s U\|_{F^s_{2,2}} \leq (2 \cdot 4^s)^{1/2} \|J_s \iota_s U\|_2$
(\Cref{thm:sobolev-two-triebel-lizorkin}). For the second, $\|J_s \iota_s U\|_2 \leq
(2 \cdot 5^s)^{1/2} \|\iota_s U\|_{F^s_{2,2}}$ (\Cref{thm:sobolev-two-triebel-lizorkin}), which
is at most $8 \cdot 8^s \, \|\Phi\|_1$ times the approximation norm of $\iota_s U$ in $A^s_2$
(\Cref{thm:tl-besov-sandwich}), which is at most $D_{s,r}$ times the right side without
$C'_{s,r}$ (\Cref{thm:modulus-coincidence}).
\end{proof}

The second bound is informative for $s < r$: for $r \leq s$ its right side is infinite unless
$U = 0$. At $s = 1$ and $r = 2$ the constants $C_{1,2}$ and $C'_{1,2}$ are the constants $C_1$
and $C_2$ of \Cref{thm:sobolev-one-besov-constant} and \Cref{thm:besov-sobolev-one-constant},
which state the two bounds for the operator $\iota = \iota_1$.

\begin{lemma}[$\Besov{s}{2}{2}$ functions in $H^s(V)$]
\label{lem:sobolev-from-differences}
\lean{Rd.exists_sobolevToL2_eq_of_modNorm_lt_top}
\leanok
\uses{thm:rd-besov, cor:sobolev-two-besov, lem:sobolev-to-l2}
Let $s > 0$ and let $r$ be an integer with $s < r$. For every $F \in L_2(V)$ with
$\|F\|_2 + \|(2^{ks} \omega_r(F, 2^{-k})_2)_k\|_{\ell_2} < \infty$ there is $U \in H^s(V)$ with
$\iota_s U = F$.
\end{lemma}

\begin{proof}
\leanok
$F$ lies in the approximation space $A^s_2$ of the band-limited scheme (\Cref{thm:rd-besov}),
hence in $H^s(V)$ (\Cref{cor:sobolev-two-besov}); let $U$ be $F$ as an element of the bundled
space. Its tempered distribution is that of $F$, so $\iota_s U = F$ (\Cref{lem:sobolev-to-l2}).
\end{proof}

\begin{lemma}[$\Besov{s}{p}{2} \subseteq \Besov{1/2}{p}{1}$ for $s > 1/2$]
\label{lem:besov-half-one-general}
\lean{Rd.modNorm_half_one_le_of_lt, Rd.modNorm_sobolevToL2_lt_top}
\leanok
\uses{lem:lq-cauchy-schwarz-sequences, lem:lq-reindex, thm:sobolev-besov-constants}
Let $1 \leq p \leq \infty$, $s > 1/2$, $r \in \N$ and
$c_s = \|(2^{-k(s - 1/2)})_k\|_{\ell_2} = (1 - 2^{-2(s - 1/2)})^{-1/2}$. For every
$F \in L_p(V)$,
\[
  \|F\|_p + \sum_k 2^{k/2} \omega_r(F, 2^{-k})_p
  \leq c_s \bigl(\|F\|_p + \|(2^{ks} \omega_r(F, 2^{-k})_p)_k\|_{\ell_2}\bigr) .
\]
For every $1/2 < s < 2$ and every $U \in H^s(V)$,
\[
  \|\iota_s U\|_2 + \|(2^{k((s - 1/2) + 1/2)} \omega_2(\iota_s U, 2^{-k})_2)_k\|_{\ell_2}
  < \infty .
\]
\end{lemma}

\begin{proof}
\leanok
$2^{k/2} \omega_k = 2^{-k(s - 1/2)} (2^{ks} \omega_k)$, so
\Cref{lem:lq-cauchy-schwarz-sequences} bounds the sum by
$c_s \|(2^{ks} \omega_k)_k\|_{\ell_2}$, with the value of $c_s$ from \Cref{lem:lq-reindex};
and $c_s \geq 1$. The second statement is the first bound of \Cref{thm:sobolev-besov-constants}
with $r = 2 > s$, since $(s - 1/2) + 1/2 = s$.
\end{proof}

The finiteness in \Cref{lem:besov-half-one-general} is the hypothesis of the boundary trace
theorems (\Cref{thm:boundary-trace-lp}, \Cref{thm:boundary-trace-bounded}) with smoothness
$s - 1/2$, $p = q = 2$ and order $2$, whose condition $(s - 1/2) + 1/2 < 2$ is $s < 2$. The
finiteness holds for every $0 < s < 2$, and $s > 1/2$ is the range of the trace theorems. At
$s = 1$ it is that of
\Cref{thm:sobolev-one-besov-constant}, and the first bound is \Cref{lem:besov-half-one} with
$c_1 = \sqrt 2$.

\section{The Leibniz rule}
\label{sec:leibniz-rule}

Let $(e_i)_{i \in I}$ be the orthonormal basis of \Cref{def:dyadic-cells}. A \emph{word} of
length $k$ is a tuple $w = (w_0, \dots, w_{k-1}) \in I^k$; it stands for the derivative
$\partial^{e_w} = \partial_{e_{w_0}} \cdots \partial_{e_{w_{k-1}}}$. For directions
$u = (u_0, \dots, u_{k-1}) \in V^k$ and $v \in V$, $(u, v) \in V^{k+1}$ is the tuple $u$ followed
by $v$, so that $\partial^{(u, v)} = \partial^u \partial_v$. This section proves the Leibniz rule
for $\chi f$ with a smooth compactly supported $\chi$, with an explicit constant, by induction on
the order: $\partial_{e_j} (\chi f) = \chi \, \partial_{e_j} f + (\partial_{e_j} \chi) f$
(\Cref{lem:weak-derivative-product-cutoff}). The Leibniz rule is the tool for pasting local
extensions in $\Sob{m}{p}(V)$, as \Cref{lem:product-rule-lp} is for $m = 1$.

\paragraph{Words and directions.}

\begin{lemma}[Composition of weak derivatives]
\label{lem:weak-derivative-snoc}
\lean{Rd.HasWeakLineDeriv.snoc, Rd.HasWeakLineDeriv.of_snoc}
\leanok
\uses{def:weak-line-derivative}
For every open $\Omega \subseteq V$, every $k \in \N$, all $u \in V^k$ and $v \in V$ and all
$f$, $h$, $g$ with $h = \partial_v f$ weakly on $\Omega$: $g = \partial^u h$ weakly on $\Omega$ if
and only if $g = \partial^{(u, v)} f$ weakly on $\Omega$.
\end{lemma}

\begin{proof}
\leanok
For distributions, $\partial^{(u, v)} T = \partial^u (\partial_v T)$ (Mathlib's
\texttt{LineDeriv.\allowbreak iteratedLineDerivOp\_\allowbreak succ\_\allowbreak right}), and
$\partial_v f = h$ as distributions on $\Omega$.
\end{proof}

\begin{lemma}[Derivatives along all directions from the words]
\label{lem:weak-derivative-words}
\lean{Rd.iteratedLineDerivOp_eq_sum_prod_smul, Rd.ofFun_sum_smul,
  Rd.hasWeakLineDeriv_of_gridBasis_words, Rd.memSobolev_of_forall_gridBasis_words}
\leanok
\uses{def:weak-line-derivative, def:dyadic-cells, lem:weak-derivative-basic,
  lem:weak-derivative-direction-linear}
\begin{enumerate}
  \item For every real vector space $E$ with derivatives $\partial_v \colon E \to E$ ($v \in V$)
    that are additive and real homogeneous both in $v$ and in their argument, every finite set
    $J$, all $b_j \in V$ ($j \in J$), every $k \in \N$, all $c_{l,j} \in \R$ ($l < k$, $j \in J$)
    and every $T \in E$,
    \[
      \partial^v T = \sum_{i \in J^k} \Bigl(\prod_{l < k} c_{l, i_l}\Bigr)
        \partial_{b_{i_0}} \cdots \partial_{b_{i_{k-1}}} T
      \qquad \text{for } v_l = \sum_{j \in J} c_{l,j} b_j .
    \]
  \item For every open $\Omega \subseteq V$, every finite set $J$, all $c_j \in \R$ and all
    locally integrable $g_j$ on $\Omega$ ($j \in J$), the distribution of $\sum_j c_j g_j$ on
    $\Omega$ is $\sum_j c_j$ times the distribution of $g_j$.
  \item For every open $\Omega \subseteq V$, every $k \in \N$, every locally integrable $f$ on
    $\Omega$, all $g_w$ with $g_w = \partial^{e_w} f$ weakly on $\Omega$ for every word
    $w \in I^k$ and every $v \in V^k$,
    \[
      \sum_{w \in I^k} \Bigl(\prod_{l < k} \langle v_l, e_{w_l} \rangle\Bigr) g_w = \partial^v f
      \quad \text{weakly on } \Omega .
    \]
  \item For every $1 \leq p \leq \infty$, every $m \in \N$ and every $F \in L_p(V)$ with a weak
    derivative $\partial^{e_w} F \in L_p(V)$ on $V$ for every $k \leq m$ and every word
    $w \in I^k$, $F \in \Sob{m}{p}(V)$.
\end{enumerate}
\end{lemma}

\begin{proof}
\leanok
(1) Induction on $k$: $\partial^v T = \partial_{v_0} \partial^{(v_1, \dots, v_{k-1})} T$, the
induction hypothesis for the tail, linearity of $\partial_{v_0}$ in $v_0$ and in its argument,
and the splitting of the words $i \in J^{k+1}$ into their first letter and the rest.
(2) Induction on $J$. (3) (1) for distributions on $\Omega$ with
$v_l = \sum_{j \in I} \langle v_l, e_j \rangle e_j$, and (2). (4) By (3) the derivative along
$v \in V^k$ is a finite combination of the $\partial^{e_w} F$, hence lies in $L_p(V)$.
\end{proof}

\paragraph{The Leibniz constant.}

\begin{definition}[The Leibniz constant]
\label{def:leibniz-constant}
\lean{Rd.leibnizConst}
\leanok
For every smooth $\chi \colon V \to \R$ and every $k \in \N$,
\[
  C_k(\chi) = 2^k \sup_{x \in V} \sum_{j \leq k} \|D^j \chi(x)\| \in [0, \infty] .
\]
\end{definition}

\begin{lemma}[Properties of the Leibniz constant]
\label{lem:leibniz-constant}
\lean{Rd.leibnizConst_ne_top, Rd.enorm_le_leibnizConst_zero,
  Rd.enorm_iteratedFDeriv_fderiv_apply_le, Rd.leibnizConst_add_le}
\leanok
\uses{def:leibniz-constant}
For every smooth $\chi \colon V \to \R$:
\begin{enumerate}
  \item if $\chi$ has compact support, then $C_k(\chi) < \infty$ for every $k \in \N$;
  \item $|\chi(x)| \leq C_0(\chi)$ for every $x \in V$;
  \item for every $v \in V$ with $\|v\| \leq 1$, every $j \in \N$ and every $x \in V$,
    $\|D^j (\partial_v \chi)(x)\| \leq \|D^{j+1} \chi(x)\|$, and
    $C_k(\chi) + C_k(\partial_v \chi) \leq C_{k+1}(\chi)$ for every $k \in \N$.
\end{enumerate}
\end{lemma}

\begin{proof}
\leanok
(1) Each $D^j \chi$ is continuous with compact support, hence bounded. (2) $D^0 \chi = \chi$.
(3) $D^j (\partial_v \chi)(x)$ is $D^{j+1} \chi(x)$ with $v$ inserted as one argument (Mathlib's
\texttt{norm\_\allowbreak iteratedFDeriv\_\allowbreak clm\_\allowbreak apply\_\allowbreak const}
and \texttt{norm\_\allowbreak iteratedFDeriv\_\allowbreak fderiv}). Hence both
$\sum_{j \leq k} \|D^j \chi(x)\|$ and $\sum_{j \leq k} \|D^j (\partial_v \chi)(x)\|$ are at most
$\sum_{j \leq k+1} \|D^j \chi(x)\|$, and $2^k + 2^k = 2^{k+1}$.
\end{proof}

\paragraph{The Leibniz rule.}

\begin{lemma}[Words with a fixed last letter]
\label{lem:words-snoc}
\lean{Rd.sum_snoc_le, Rd.sum_range_sum_snoc_le, Rd.gridBasis_snoc}
\leanok
For every finite set $J$, every $j \in J$ and every $k \in \N$:
\begin{enumerate}
  \item for every $a \colon J^{k+1} \to [0, \infty]$,
    $\sum_{w \in J^k} a(w, j) \leq \sum_{w \in J^{k+1}} a(w)$;
  \item for all $a_{k'} \colon J^{k'} \to [0, \infty]$ ($k' \in \N$),
    $\sum_{k' \leq k} \sum_{w \in J^{k'}} a_{k'+1}(w, j)
      \leq \sum_{k' \leq k+1} \sum_{w \in J^{k'}} a_{k'}(w)$;
  \item for $J = I$ and every $u \in I^k$, the directions of the word $(u, j)$ are
    $(e_{u_0}, \dots, e_{u_{k-1}}, e_j)$.
\end{enumerate}
\end{lemma}

\begin{proof}
\leanok
(1) The map $w \mapsto (w, j)$ is injective. (2) (1) for every $k'$, and the term $k' = 0$ on the
right is nonnegative. (3) By definition.
\end{proof}

\begin{lemma}[Sums and choices over words]
\label{lem:sums-over-words}
\lean{Rd.le_sum_range_sum_words, Rd.sum_range_sum_words_mono, Rd.sum_range_sum_words_sum_mul,
  Rd.exists_forall_words_of_forall_le}
\leanok
For every finite set $J$ and every $m \in \N$:
\begin{enumerate}
  \item for all $a_{k,w} \in [0, \infty]$ ($k \in \N$, $w \in J^k$),
    $a_{k,w} \leq \sum_{k' \leq m} \sum_{w' \in J^{k'}} a_{k',w'}$ for every $k \leq m$ and
    every $w \in J^k$;
  \item for all $a_{k,w} \in [0, \infty]$ ($k \in \N$, $w \in J^k$),
    $\sum_{k' \leq k} \sum_{w' \in J^{k'}} a_{k',w'} \leq \sum_{k' \leq m}
    \sum_{w' \in J^{k'}} a_{k',w'}$ for every $k \leq m$;
  \item for every finite set $\kappa$ and all $\lambda_{l,k}, c_l \in [0, \infty]$
    ($l \in \kappa$, $k \in \N$),
    $\sum_{k \leq m} \sum_{w \in J^k} \sum_{l \in \kappa} \lambda_{l,k} c_l
      = \sum_{l \in \kappa} \bigl(\sum_{k \leq m} \sum_{w \in J^k} \lambda_{l,k}\bigr) c_l$;
  \item for every nonempty set $A$ and every property $P(k, w, a)$ of $k \in \N$, $w \in J^k$
    and $a \in A$: if for every $k \leq m$ and every $w \in J^k$ some $a \in A$ satisfies
    $P(k, w, a)$, then there is a family $(a_{k,w})_{k \in \N, w \in J^k}$ in $A$ with
    $P(k, w, a_{k,w})$ for every $k \leq m$ and every $w \in J^k$.
\end{enumerate}
\end{lemma}

\begin{proof}
\leanok
(1) and (2) The terms are nonnegative. (3) Exchange the sums. (4) Choose $a_{k,w}$ for $k \leq m$
and any element of $A$ for $k > m$.
\end{proof}

\begin{theorem}[The Leibniz rule with a bound]
\label{thm:leibniz-rule}
\lean{Rd.exists_hasWeakLineDeriv_mul_zero, Rd.HasWeakLineDeriv.mul_snoc,
  Rd.exists_hasWeakLineDeriv_mul_of_tsupport_subset}
\leanok
\uses{lem:weak-derivative-snoc, lem:weak-derivative-product-cutoff, lem:weak-derivative-cutoff,
  lem:leibniz-constant, lem:words-snoc, lem:weak-derivative-basic, def:leibniz-constant}
For all open $O, B \subseteq V$, every $1 \leq p \leq \infty$, every $m \in \N$, every $f$, all
$g_{k,w}$ ($k \leq m$, $w \in I^k$) with $g_{k,w} = \partial^{e_w} f$ weakly on $O \cap B$, every
smooth $\chi \colon V \to \R$ with compact closed support in $B$, every $k \leq m$ and every
$w \in I^k$, there is $G$ with $G = \partial^{e_w} (\chi f)$ weakly on $O$ and
\[
  \|G\|_{L_p(O)} \leq C_k(\chi) \sum_{k' \leq k} \sum_{w' \in I^{k'}}
    \|g_{k',w'}\|_{L_p(O \cap B)} .
\]
\end{theorem}

\begin{proof}
\leanok
Induction on $k$, for all $m$, $f$, $g$ and $\chi$ at once. For $k = 0$, $G = \chi f$, and
$\|\chi f\|_{L_p(O)} \leq C_0(\chi) \|f\|_{L_p(O \cap B)}$
(\Cref{lem:weak-derivative-cutoff} (2), \Cref{lem:leibniz-constant} (2)) with $f = g_{0,()}$
almost everywhere on $O \cap B$. For a word $(u, j)$ of length $k + 1$, let
$f' = g_{1,(j)} = \partial_{e_j} f$. Then $g_{k'+1,(w', j)} = \partial^{e_{w'}} f'$ on $O \cap B$
(\Cref{lem:weak-derivative-snoc}), and the induction hypothesis for $\chi f'$ and for
$(\partial_{e_j} \chi) f$ along $u$ gives $G_1$ and $G_2$. Since
$\partial_{e_j} (\chi f) = \chi f' + (\partial_{e_j} \chi) f$ on $O$
(\Cref{lem:weak-derivative-product-cutoff} (3)), $G = G_1 + G_2$ is the derivative along $(u, j)$
(\Cref{lem:weak-derivative-snoc}). Both bounds are at most the sum over the words of length at most
$k + 1$ (\Cref{lem:words-snoc}), and
$C_k(\chi) + C_k(\partial_{e_j} \chi) \leq C_{k+1}(\chi)$ (\Cref{lem:leibniz-constant} (3)).
\end{proof}

In the next corollary, $\chi F$ is the product in $L_p(V)$ of \Cref{lem:product-modulus}, formed
with the continuity of $\chi$ and a bound $|\chi| \leq C$; both follow from smoothness and compact
support.

\begin{corollary}[Smooth compactly supported multipliers on $\Sob{m}{p}(V)$]
\label{cor:sobolev-multiplier}
\lean{Rd.MemSobolev.exists_hasWeakLineDeriv_mul, Rd.MemSobolev.mulLp_of_hasCompactSupport}
\leanok
\uses{thm:leibniz-rule, lem:weak-derivative-words, lem:leibniz-constant, lem:product-modulus,
  def:weak-line-derivative}
For every $1 \leq p \leq \infty$, every $m \in \N$, every $F \in \Sob{m}{p}(V)$ and every smooth
$\chi \colon V \to \R$ with compact support:
\begin{enumerate}
  \item for every $k \leq m$ and every $w \in I^k$, $\chi F$ has a weak derivative
    $\partial^{e_w} (\chi F) \in L_p(V)$ on $V$;
  \item $\chi F \in \Sob{m}{p}(V)$.
\end{enumerate}
\end{corollary}

\begin{proof}
\leanok
(1) \Cref{thm:leibniz-rule} with $O = B = V$; the bound is finite since $C_k(\chi) < \infty$
(\Cref{lem:leibniz-constant} (1)) and the derivatives of $F$ lie in $L_p(V)$. (2) (1) and
\Cref{lem:weak-derivative-words} (4).
\end{proof}

\begin{lemma}[The product rule in $L_p$]
\label{lem:product-rule-lp}
\lean{Rd.hasWeakLineDeriv_mulLp, Rd.MemSobolev.mulLp}
\leanok
\uses{lem:weak-derivative-product, lem:product-modulus, lem:weak-derivative-basic,
  cor:sobolev-multiplier}
Let $1 \leq p \leq \infty$, $w \in V$, $F, G \in L_p(V)$ with $G = \partial_w F$ weakly on $V$,
and $\chi$ a smooth real function with $\chi$ and $\partial_w \chi$ bounded. Then
$\chi G + (\partial_w \chi) F = \partial_w (\chi F)$ weakly on $V$, with the products in
$L_p(V)$. For every $F \in \Sob{1}{p}(V)$ and every smooth real $\chi$ with compact support,
$\chi F \in \Sob{1}{p}(V)$.
\end{lemma}

\begin{proof}
\leanok
\Cref{lem:weak-derivative-product} for the representatives, and
\Cref{lem:weak-derivative-basic}. The second statement is the case $m = 1$ of
\Cref{cor:sobolev-multiplier}.
\end{proof}

\section{Derivative kernels}
\label{sec:derivative-kernels}

The Sobolev space $\Sob{m}{p}(V)$ is described by weak derivatives, the space $\TL{m}{p}{2}$ by the
Littlewood--Paley blocks. This section provides the kernels that connect the two: the block
$\psi_1$ is a sum of derivatives of order $m$ of Schwartz functions, and after dilation
$2^{jm} \psi_{j+1}$ is a sum of derivatives of the dilated kernels. Let $e_0, \dots, e_{n-1}$ be
the orthonormal basis of \Cref{def:dyadic-cells}. For $m \in \N$ a multi-index is a tuple
$\alpha = (\alpha_0, \dots, \alpha_{m-1})$ with entries in $\{0, \dots, n-1\}$; write
$e_\alpha = (e_{\alpha_0}, \dots, e_{\alpha_{m-1}}) \in V^m$, so that
$\partial^{e_\alpha} = \partial_{e_{\alpha_0}} \cdots \partial_{e_{\alpha_{m-1}}}$, and sums over
$\alpha$ run over all $n^m$ multi-indices. On the Fourier side $\partial^{e_\alpha}$ multiplies by
$\prod_k 2\pi i \langle \xi, e_{\alpha_k} \rangle$, and
$\sum_\alpha \prod_k \langle \xi, e_{\alpha_k} \rangle^2 = \|\xi\|^{2m}$. The kernel $\eta_\alpha$
is chosen so that $\partial^{e_\alpha} \eta_\alpha$ has the Fourier transform
$\prod_k \langle \xi, e_{\alpha_k} \rangle^2 \|\xi\|^{-2m} \widehat{\psi_1}(\xi)$; the division by
$\|\xi\|^{2m}$ keeps $\widehat{\eta_\alpha}$ smooth because $\widehat{\psi_1}$ vanishes near $0$.
For every $w \in V^k$, the rescaled derivatives $2^{-jk} \partial^w \tilde\psi_j$ of the synthesis
kernels (\Cref{def:psi-dual}) form Calder\'on--Zygmund families as well.

\begin{lemma}[Fourier transforms of iterated derivatives]
\label{lem:fourier-iterated-derivative}
\lean{SchwartzMap.fourier_iteratedLineDerivOp_apply}
\leanok
For every complex normed space $E$, every Schwartz function $G \colon V \to E$, every $k \in \N$,
all $w \in V^k$ and every $\xi \in V$,
\[
  \widehat{\partial^w G}(\xi) = \Bigl( \prod_{l < k} 2\pi i \langle \xi, w_l \rangle \Bigr)
  \widehat G(\xi) .
\]
\end{lemma}

\begin{proof}
\leanok
Induction on $k$, with $\widehat{\partial_v G}(\xi) = 2\pi i \langle \xi, v \rangle \widehat G(\xi)$
(Mathlib's \texttt{SchwartzMap.\allowbreak fourier\_\allowbreak lineDerivOp\_\allowbreak eq}).
\end{proof}

\begin{lemma}[Powers of the norm in coordinates]
\label{lem:norm-power-coordinates}
\lean{Rd.norm_pow_eq_sum_prod_inner}
\leanok
\uses{def:dyadic-cells}
For every $\xi \in V$ and every $m \in \N$,
\[
  \|\xi\|^{2m} = \sum_\alpha \prod_{k < m} \langle \xi, e_{\alpha_k} \rangle^2 .
\]
\end{lemma}

\begin{proof}
\leanok
$\|\xi\|^2 = \sum_{i < n} \langle \xi, e_i \rangle^2$ (Mathlib's
\texttt{OrthonormalBasis.\allowbreak sum\_\allowbreak sq\_\allowbreak inner\_\allowbreak left}),
and the $m$-th power of this sum expands into the sum over all multi-indices (Mathlib's
\texttt{Fintype.\allowbreak sum\_\allowbreak pow}).
\end{proof}

\begin{definition}[Derivative symbols]
\label{def:derivative-symbol}
\lean{Rd.derivSymbol}
\leanok
\uses{def:dyadic-cells}
For every $m \in \N$, every multi-index $\alpha$ and every $\xi \in V$,
\[
  \sigma_\alpha(\xi) = \frac{\prod_{k < m} \langle \xi, e_{\alpha_k} \rangle}
    {(2\pi i)^m \|\xi\|^{2m}} ,
\]
with the convention $t / 0 = 0$.
\end{definition}

\begin{lemma}[Properties of the derivative symbols]
\label{lem:derivative-symbol}
\lean{Rd.sum_prod_mul_derivSymbol, Rd.norm_derivSymbol_le, Rd.contDiffOn_derivSymbol}
\leanok
\uses{def:derivative-symbol, lem:norm-power-coordinates}
Let $m \in \N$. For every $\xi \in V \setminus \{0\}$,
\[
  \sum_\alpha \Bigl( \prod_{k < m} 2\pi i \langle \xi, e_{\alpha_k} \rangle \Bigr)
  \sigma_\alpha(\xi) = 1 .
\]
For every multi-index $\alpha$, $|\sigma_\alpha(\xi)| \leq 1$ for every $\xi \in V$ with
$\|\xi\| \geq 1$, and $\sigma_\alpha$ is smooth on $V \setminus \{0\}$.
\end{lemma}

\begin{proof}
\leanok
The term of $\alpha$ in the sum is $\prod_k \langle \xi, e_{\alpha_k} \rangle^2 / \|\xi\|^{2m}$, and
these terms add up to $1$ by \Cref{lem:norm-power-coordinates}. Since
$|\langle \xi, e_i \rangle| \leq \|\xi\|$,
$|\sigma_\alpha(\xi)| \leq (2\pi)^{-m} \|\xi\|^m / \|\xi\|^{2m} \leq 1$ for $\|\xi\| \geq 1$.
Products of linear functions and $\|\xi\|^{2m} = (\|\xi\|^2)^m$ are smooth, and the latter
vanishes only at $0$.
\end{proof}

\begin{lemma}[The block $\psi_1$]
\label{lem:psi-one}
\lean{Rd.psiOneS, Rd.coe_psiOneS, Rd.fourier_psiOneS_apply, Rd.dilate_psi_one,
  Rd.fourier_psi_succ_eq_fourier_psi_one, Rd.one_lt_norm_of_fourier_psi_one_ne_zero}
\leanok
\uses{def:psi-base, lem:psi-base, def:moment-kernel, lem:fourier-dilate,
  lem:lp-kernels-frequency}
The block $\psi_1$ is the Schwartz function $\delta_2 G^{\psi}_1$. For every $j \in \N$ and every
$\xi \in V$, $\psi_{j+1} = (\psi_1)_{2^j}$ and
$\widehat{\psi_{j+1}}(\xi) = \widehat{\psi_1}(2^{-j} \xi)$. For every $\xi \in V$ with
$\widehat{\psi_1}(\xi) \neq 0$, $\|\xi\| > 1$.
\end{lemma}

\begin{proof}
\leanok
$\psi_{j+1} = \delta_{2^{j+1}} G^{\psi}_{j+1}$ (\Cref{lem:psi-base}) and
$G^{\psi}_{j+1} = G^{\psi}_1$ (\Cref{def:psi-base}); with $j = 0$ this is the first claim, and
$\delta_{2^j} \delta_2 = \delta_{2^{j+1}}$ (\Cref{def:moment-kernel}) gives the second. The third
is \Cref{lem:fourier-dilate}, and the last is the case $j = 0$ of
\Cref{lem:lp-kernels-frequency}.
\end{proof}

\begin{lemma}[Derivative kernels]
\label{lem:derivative-kernel}
\lean{Rd.contDiff_derivSymbol_mul_fourier_psi_one,
  Rd.hasCompactSupport_derivSymbol_mul_fourier_psi_one, Rd.derivKernelHat, Rd.derivKernel,
  Rd.fourier_derivKernel}
\leanok
\uses{lem:derivative-symbol, lem:psi-one, lem:lp-kernels-frequency, def:lp-kernels}
For every $m \in \N$ and every multi-index $\alpha$, the function
$\sigma_\alpha \widehat{\psi_1}$ is smooth with compact support, and the derivative kernel
\[
  \eta_\alpha = \mathcal F^{-1} \bigl( \sigma_\alpha \widehat{\psi_1} \bigr)
\]
is a Schwartz function with $\widehat{\eta_\alpha} = \sigma_\alpha \widehat{\psi_1}$.
\end{lemma}

\begin{proof}
\leanok
$\widehat{\psi_1}$ is the Fourier transform of a Schwartz function (\Cref{lem:psi-one}), hence
smooth, and it vanishes on the closed unit ball and outside the ball of radius $4$
(\Cref{lem:lp-kernels-frequency}). So the product vanishes near $0$ and outside the closed ball
of radius $4$, and away from $0$ it is smooth by \Cref{lem:derivative-symbol}. Smooth functions
with compact support are Schwartz functions (Mathlib's
\texttt{HasCompactSupport.\allowbreak toSchwartzMap}), and $\mathcal F$ inverts
$\mathcal F^{-1}$ on Schwartz functions.
\end{proof}

\begin{theorem}[$\psi_1$ as a sum of derivatives]
\label{thm:psi-one-derivatives}
\lean{Rd.sum_iteratedLineDerivOp_derivKernel}
\leanok
\uses{lem:fourier-iterated-derivative, lem:derivative-symbol, lem:derivative-kernel,
  lem:psi-one, lem:lp-kernels-frequency}
For every $m \in \N$,
\[
  \psi_1 = \sum_\alpha \partial^{e_\alpha} \eta_\alpha .
\]
\end{theorem}

\begin{proof}
\leanok
By \Cref{lem:fourier-iterated-derivative} and \Cref{lem:derivative-kernel}, the Fourier transform
of the right side at $\xi$ is
$\bigl( \sum_\alpha \prod_k 2\pi i \langle \xi, e_{\alpha_k} \rangle \, \sigma_\alpha(\xi) \bigr)
\widehat{\psi_1}(\xi)$. For $\xi \neq 0$ the sum is $1$ (\Cref{lem:derivative-symbol}), and
$\widehat{\psi_1}(0) = 0$ (\Cref{lem:lp-kernels-frequency}). The Fourier transform is injective on
Schwartz functions.
\end{proof}

\begin{corollary}[The dilated identity]
\label{cor:psi-derivatives-dilated}
\lean{Rd.two_pow_smul_psi_succ_eq_sum}
\leanok
\uses{thm:psi-one-derivatives, lem:derivative-dilate, lem:psi-one}
For all $m, j \in \N$,
\[
  2^{jm} \psi_{j+1} = \sum_\alpha \partial^{e_\alpha} \bigl( (\eta_\alpha)_{2^j} \bigr) .
\]
\end{corollary}

\begin{proof}
\leanok
By \Cref{lem:derivative-dilate},
$\partial^{e_\alpha} ((\eta_\alpha)_{2^j}) = 2^{jm} (\partial^{e_\alpha} \eta_\alpha)_{2^j}$, and
the sum over $\alpha$ is $2^{jm} (\psi_1)_{2^j} = 2^{jm} \psi_{j+1}$ by
\Cref{thm:psi-one-derivatives}, the linearity of the dilation and \Cref{lem:psi-one}.
\end{proof}

\paragraph{The derivative families.}

\begin{definition}[Derivative families]
\label{def:derivative-family}
\lean{Rd.derivFamily, Rd.derivKernelBound}
\leanok
\uses{lem:derivative-kernel, def:uniform-schwartz}
For every $m \in \N$, every multi-index $\alpha$ and every $j \in \N$ let
$\theta^\alpha_j = (\eta_\alpha)_{2^j}$, and for all $k, l \in \N$ let
$A^\alpha_{k,l} = \sup_x \|x\|^k \|D^l \eta_\alpha(x)\|$.
\end{definition}

\begin{lemma}[Frequencies of the derivative families]
\label{lem:derivative-family-frequency}
\lean{Rd.fourier_derivFamily, Rd.enorm_fourier_derivFamily_le}
\leanok
\uses{def:derivative-family, lem:fourier-dilate, lem:derivative-kernel, lem:psi-one,
  lem:derivative-symbol}
For every $m \in \N$, every multi-index $\alpha$, every $j \in \N$ and every $\xi \in V$,
\[
  \widehat{\theta^\alpha_j}(\xi) = \sigma_\alpha(2^{-j} \xi) \, \widehat{\psi_{j+1}}(\xi)
  \quad \text{and} \quad
  |\widehat{\theta^\alpha_j}(\xi)| \leq |\widehat{\psi_{j+1}}(\xi)| .
\]
\end{lemma}

\begin{proof}
\leanok
The identity follows from \Cref{lem:fourier-dilate}, \Cref{lem:derivative-kernel} and
\Cref{lem:psi-one}. Where $\widehat{\psi_{j+1}}(\xi) = \widehat{\psi_1}(2^{-j} \xi) \neq 0$,
$\|2^{-j} \xi\| > 1$ (\Cref{lem:psi-one}), so $|\sigma_\alpha(2^{-j} \xi)| \leq 1$
(\Cref{lem:derivative-symbol}).
\end{proof}

\begin{theorem}[The derivative families form Calder\'on--Zygmund families]
\label{thm:cz-family-derivative}
\lean{Rd.tsum_fourier_derivFamily_sq_le, Rd.isUniformSchwartz_derivKernel,
  Rd.isCZFamily_derivFamily}
\leanok
\uses{lem:derivative-family-frequency, thm:cz-family-psi, def:derivative-family,
  lem:uniform-schwartz-constants, thm:cz-family-dilation, def:hormander-const}
For every $m \in \N$, every multi-index $\alpha$ and every $\xi \in V$,
$\sum_j |\widehat{\theta^\alpha_j}(\xi)|^2 \leq 2$. The constant family $(\eta_\alpha)_j$ is
uniformly Schwartz with bounds $A^\alpha$, and $(\theta^\alpha_j)_j$ is a Calder\'on--Zygmund
family with overlap $2$ and H\"ormander constant $4 H_{A^\alpha}$.
\end{theorem}

\begin{proof}
\leanok
By \Cref{lem:derivative-family-frequency},
$\sum_j |\widehat{\theta^\alpha_j}(\xi)|^2 \leq \sum_j |\widehat{\psi_{j+1}}(\xi)|^2 \leq 2$
(\Cref{thm:cz-family-psi}). The bounds are the seminorms of $\eta_\alpha$ themselves
(\Cref{lem:uniform-schwartz-constants}), and \Cref{thm:cz-family-dilation} gives the rest.
\end{proof}

\paragraph{The dual derivative families.}

\begin{definition}[Dual derivative families]
\label{def:dual-derivative-family}
\lean{Rd.derivDualBase, Rd.derivDualFamily, Rd.derivDualBound, Rd.derivDualOverlap,
  Rd.derivDualOverlap_ne_top}
\leanok
\uses{def:psi-dual, def:schwartz-dilate}
For every $k \in \N$, all $w \in V^k$ and every $j \in \N$ let
$\tilde G^w_j = \partial^w \tilde G_j$ and $\tilde\theta^w_j = (\tilde G^w_j)_{2^j}$. For all
$a, l \in \N$ let $\tilde A^w_{a,l}$ be the larger of the seminorms
$\sup_x \|x\|^a \|D^l \tilde G^w_0(x)\|$ and $\sup_x \|x\|^a \|D^l \tilde G^w_2(x)\|$, and let
\[
  B_w = 4 \Bigl( (8\pi)^k \prod_{l < k} \|w_l\| \Bigr)^2 ,
\]
which is finite.
\end{definition}

\begin{lemma}[Derivatives of the synthesis kernels]
\label{lem:psi-dual-derivatives}
\lean{Rd.psiDualS, Rd.coe_psiDualS, Rd.iteratedLineDerivOp_psiDualS,
  Rd.isUniformSchwartz_derivDualBase}
\leanok
\uses{def:dual-derivative-family, lem:psi-dual-basic, lem:derivative-dilate,
  lem:uniform-schwartz-constants}
For every $j \in \N$, the synthesis kernel $\tilde\psi_j$ is the Schwartz function
$(\tilde G_j)_{2^j}$. For every $k \in \N$, all $w \in V^k$ and every $j \in \N$,
\[
  \partial^w \tilde\psi_j = 2^{jk} \, \tilde\theta^w_j ,
\]
and $(\tilde G^w_j)_j$ is uniformly Schwartz with bounds $\tilde A^w$.
\end{lemma}

\begin{proof}
\leanok
$\tilde\psi_j = (\tilde G_j)_{2^j}$ (\Cref{lem:psi-dual-basic}), and \Cref{lem:derivative-dilate}
with $c = 2^j$ gives the identity. The family $(\tilde G^w_j)_j$ takes the two values
$\tilde G^w_0$ and $\tilde G^w_2$ (\Cref{lem:uniform-schwartz-constants}).
\end{proof}

\begin{lemma}[Frequencies of the dual derivative families]
\label{lem:dual-derivative-frequency}
\lean{Rd.fourier_derivDualFamily, Rd.enorm_prod_inner_le, Rd.enorm_fourier_derivDualFamily_le}
\leanok
\uses{def:dual-derivative-family, lem:fourier-dilate, lem:fourier-iterated-derivative,
  lem:psi-dual-basic, lem:psi-dual-frequency}
Let $k \in \N$ and $w \in V^k$. For every real $R$ and every $\zeta \in V$ with
$\|\zeta\| \leq R$,
$|\prod_{l < k} 2\pi i \langle \zeta, w_l \rangle| \leq (2\pi R)^k \prod_{l < k} \|w_l\|$. For every
$j \in \N$ and every $\xi \in V$,
\[
  \widehat{\tilde\theta^w_j}(\xi) = \Bigl( \prod_{l < k} 2\pi i \langle 2^{-j} \xi, w_l \rangle
  \Bigr) \widehat{\tilde\psi_j}(\xi)
\]
and
\[
  |\widehat{\tilde\theta^w_j}(\xi)| \leq (8\pi)^k \Bigl( \prod_{l < k} \|w_l\| \Bigr)
  |\widehat{\tilde\psi_j}(\xi)| .
\]
\end{lemma}

\begin{proof}
\leanok
The first bound is the Cauchy--Schwarz inequality
$|\langle \zeta, w_l \rangle| \leq \|\zeta\| \|w_l\|$. The identity follows from
\Cref{lem:fourier-dilate}, \Cref{lem:fourier-iterated-derivative} and
$\tilde\psi_j = (\tilde G_j)_{2^j}$ (\Cref{lem:psi-dual-basic}). Where
$\widehat{\tilde\psi_j}(\xi) \neq 0$, $\|\xi\| < 2^{j+2}$ (\Cref{lem:psi-dual-frequency}), so
$\|2^{-j} \xi\| < 4$ and the first bound with $R = 4$ applies.
\end{proof}

\begin{theorem}[The dual derivative families form Calder\'on--Zygmund families]
\label{thm:cz-family-dual-derivative}
\lean{Rd.tsum_fourier_derivDualFamily_sq_le, Rd.isCZFamily_derivDualFamily}
\leanok
\uses{lem:dual-derivative-frequency, thm:cz-family-psi-dual, lem:psi-dual-derivatives,
  thm:cz-family-dilation, def:hormander-const}
For every $k \in \N$, all $w \in V^k$ and every $\xi \in V$,
$\sum_j |\widehat{\tilde\theta^w_j}(\xi)|^2 \leq B_w$, and $(\tilde\theta^w_j)_j$ is a
Calder\'on--Zygmund family with overlap $B_w$ and H\"ormander constant $4 H_{\tilde A^w}$.
\end{theorem}

\begin{proof}
\leanok
By \Cref{lem:dual-derivative-frequency},
$\sum_j |\widehat{\tilde\theta^w_j}(\xi)|^2 \leq ((8\pi)^k \prod_l \|w_l\|)^2
\sum_j |\widehat{\tilde\psi_j}(\xi)|^2$, and the last sum is at most $4$
(\Cref{thm:cz-family-psi-dual}). \Cref{thm:cz-family-dilation} and
\Cref{lem:psi-dual-derivatives} give the rest.
\end{proof}

On the torus the derivative families act through the periodic lift, and the strong type there
needs their far mass $M_\theta$ of \Cref{def:torus-far-mass}.

\begin{lemma}[The far mass of the derivative families]
\label{lem:torus-far-mass-derivative-families}
\lean{Torus.farMass_derivFamily_le, Torus.farMass_derivDualFamily_le}
\leanok
\uses{lem:torus-far-mass-dilate, def:derivative-family, def:dual-derivative-family,
  thm:cz-family-derivative, lem:psi-dual-derivatives}
Let $V = \R^d$. For every $m \in \N$ and every multi-index $\alpha$,
$M_{\theta^\alpha} \leq 2 T_{A^\alpha,0,1}$, and for every $k \in \N$ and all $w \in V^k$,
$M_{\tilde\theta^w} \leq 2 T_{\tilde A^w,0,1}$.
\end{lemma}

\begin{proof}
\leanok
The families $\theta^\alpha_j = (\eta_\alpha)_{2^j}$ and $\tilde\theta^w_j = (\tilde G^w_j)_{2^j}$
are dilates of uniform Schwartz families with bounds $A^\alpha$ and $\tilde A^w$
(\Cref{thm:cz-family-derivative}, \Cref{lem:psi-dual-derivatives}), and
\Cref{lem:torus-far-mass-dilate} applies.
\end{proof}

\section{Divergence kernels}
\label{sec:divergence-kernels}

This section writes the difference $m_2 - m$ of two dilates of a smooth function $m$ as a
divergence $\sum_i \partial_{e_i} K_i$ of smooth kernels that vanish where $m$ vanishes along the
rays from the origin, with mean zero when the first moments of $m$ vanish. Dilates of a
mean-zero Schwartz function form a Calder\'on--Zygmund family, so the divergence kernels carry
the strong type of \Cref{thm:vector-valued-strong-type}. Throughout, $(e_i)_{i \in I}$ is an
orthonormal basis of $V$ and $\Phi_c(x) = c^n \Phi(c x)$ (\Cref{lem:dilate}).

\paragraph{Parametric integrals.}

\begin{lemma}[Smoothness of parametric integrals]
\label{lem:parametric-integral-smooth}
\lean{intervalIntegral.hasFDerivAt_integral_of_contDiff,
  intervalIntegral.fderiv_integral_apply_of_contDiff,
  intervalIntegral.contDiff_parametric_intervalIntegral_of_contDiff_aux,
  intervalIntegral.contDiff_parametric_intervalIntegral_of_contDiff}
\leanok
Let $H$ be a finite-dimensional real normed space, $E$ a real normed space and $a, b \in \R$.
For every continuously differentiable $F \colon H \times \R \to E$ and all $x_0, v \in H$, the
function $x \mapsto \int_a^b F(x, t) \, dt$ is differentiable at $x_0$ with
\[
  D\Bigl(\int_a^b F(\cdot, t) \, dt\Bigr)(x_0) \, v = \int_a^b D_x F(x_0, t) \, v \, dt .
\]
For every smooth $F \colon H \times \R \to E$ the function $x \mapsto \int_a^b F(x, t) \, dt$
is smooth.
\end{lemma}

\begin{proof}
\leanok
The partial derivative $(x, t) \mapsto D_x F(x, t)$ is continuous, hence bounded on the compact
set $\bar B(x_0, 1) \times [a, b]$, and Mathlib's differentiation under the integral sign
with a dominating bound gives the derivative. For smooth $F$ the derivative is again a
parametric integral, of the smooth function
$D_x F$ with values in the bounded linear maps $H \to E$, so the parametric integral is
$k$ times continuously differentiable for every $k$, by induction on $k$. The induction runs
through the spaces $E$, $L(H, E)$, $L(H, L(H, E))$, \dots, which Lean needs in one universe; a
lift of $E$ to the universe of $H$ provides it.
\end{proof}

Mathlib has the derivative under the integral sign and the continuity of parametric integrals;
\Cref{lem:parametric-integral-smooth} adds their smoothness.

\paragraph{The dilation identity.}

\begin{lemma}[Exponential dilations]
\label{lem:exponential-dilation}
\lean{Rd.contDiff_dilate_exp, Rd.continuous_dilate_exp, Rd.fderiv_dilate_apply,
  Rd.differentiableAt_dilate, Rd.hasDerivAt_dilate_exp}
\leanok
\uses{lem:dilate}
For every continuous $m \colon V \to \C$ the function $(x, s) \mapsto m_{e^s}(x)$ on
$V \times \R$ is continuous, and it is $k$ times continuously differentiable if $m$ is. For every
differentiable $m$ and every $c \in \R$, $m_c$ is differentiable, for all $x, v \in V$,
$D(m_c)(x) \, v = c^n \, Dm(c x)(c v)$, and for every $x \in V$ and every $s \in \R$,
\[
  \frac{d}{ds} \, m_{e^s}(x) = n \, m_{e^s}(x) + D(m_{e^s})(x) \, x .
\]
\end{lemma}

\begin{proof}
\leanok
$m_{e^s}(x) = e^{ns} m(e^s x)$; the product and chain rules give
$\frac{d}{ds} m_{e^s}(x) = n e^{ns} m(e^s x) + e^{ns} Dm(e^s x)(e^s x)$, and the second term is
$D(m_{e^s})(x) \, x$.
\end{proof}

\begin{definition}[Divergence kernels]
\label{def:divergence-kernel}
\lean{Rd.dilationAverage, Rd.divergenceKernel}
\leanok
\uses{lem:exponential-dilation}
For $m \colon V \to \C$ and $w \in V$, the \emph{dilation average} and the \emph{divergence
kernel} in the direction $w$ are
\[
  A_m(x) = \int_0^{\log 2} m_{e^s}(x) \, ds , \qquad K^m_w(x) = \langle w, x \rangle \, A_m(x) .
\]
\end{definition}

The substitution $t = e^s$ gives the form $\int_1^2 t^{-1} m_t(x) \, dt$ of the dilation average.

\begin{lemma}[The dilation average after substitution]
\label{lem:dilation-average-substitution}
\lean{Rd.dilationAverage_eq_integral_inv_mul}
\leanok
\uses{def:divergence-kernel}
For every continuous $m \colon V \to \C$ and every $x \in V$,
$A_m(x) = \int_1^2 t^{-1} m_t(x) \, dt$.
\end{lemma}

\begin{proof}
\leanok
Substitute $t = e^s$, $dt = e^s \, ds$.
\end{proof}

\begin{lemma}[Smoothness of the divergence kernels]
\label{lem:divergence-kernel-smooth}
\lean{Rd.contDiff_dilationAverage, Rd.contDiff_divergenceKernel,
  Rd.differentiable_dilationAverage}
\leanok
\uses{def:divergence-kernel, lem:parametric-integral-smooth, lem:exponential-dilation}
For every smooth $m \colon V \to \C$ and every $w \in V$, $A_m$ and $K^m_w$ are smooth. For every
continuously differentiable $m$, $A_m$ is differentiable.
\end{lemma}

\begin{proof}
\leanok
$(x, s) \mapsto m_{e^s}(x)$ is as smooth as $m$ (\Cref{lem:exponential-dilation}), and
\Cref{lem:parametric-integral-smooth} applies.
\end{proof}

\begin{lemma}[Supports of the divergence kernels]
\label{lem:divergence-kernel-support}
\lean{Rd.dilationAverage_eq_zero, Rd.divergenceKernel_eq_zero,
  Rd.divergenceKernel_eq_zero_of_lt, Rd.hasCompactSupport_divergenceKernel}
\leanok
\uses{def:divergence-kernel}
For every $m \colon V \to \C$, every $w \in V$ and every $x \in V$ with $m(t x) = 0$ for every
$t \in [1, 2]$, $A_m(x) = 0$ and $K^m_w(x) = 0$. For every $R \in \R$ and every $m$ vanishing
outside the closed ball $\bar B(0, R)$, every $K^m_w$ vanishes outside $\bar B(0, R)$ and has
compact support.
\end{lemma}

\begin{proof}
\leanok
For $s \in [0, \log 2]$, $e^s \in [1, 2]$, so the integrand of $A_m(x)$ vanishes. For
$\|x\| > R$ and $t \geq 1$, $\|t x\| \geq \|x\| > R$.
\end{proof}

\begin{lemma}[Supports of dyadic dilates]
\label{lem:dyadic-dilate-support}
\lean{Rd.tsupport_dilate_two_pow_subset}
\leanok
\uses{lem:dilate}
For every closed $S \subseteq V$ with $y \in S$ whenever $c y \in S$ for some $c \geq 1$, every
$\Psi \colon V \to \C$ that vanishes outside $S$ and every $j \in \N$, the dilate $\Psi_{2^j}$ is
supported in $S$.
\end{lemma}

\begin{proof}
\leanok
$\Psi_{2^j}(y) \neq 0$ gives $\Psi(2^j y) \neq 0$, so $2^j y \in S$ and $y \in S$; the support is
the closure of this set, and $S$ is closed.
\end{proof}

\begin{definition}[Divergence kernels as Schwartz functions]
\label{def:divergence-schwartz}
\lean{Rd.divergenceSchwartz, Rd.coe_divergenceSchwartz}
\leanok
\uses{lem:divergence-kernel-smooth, lem:divergence-kernel-support}
For every smooth $m \colon V \to \C$ that vanishes outside a ball and every $w \in V$, the
divergence kernel $K^m_w$, smooth with compact support, is a Schwartz function.
\end{definition}

\begin{lemma}[The divergence of $x A(x)$]
\label{lem:radial-divergence}
\lean{Rd.sum_fderiv_inner_mul}
\leanok
For every orthonormal basis $(b_j)_{j \in J}$ of $V$, every $x \in V$ and every
$A \colon V \to \C$ differentiable at $x$,
\[
  \sum_{j \in J} \partial_{b_j} \bigl(\langle b_j, \cdot \rangle A\bigr)(x)
    = n \, A(x) + DA(x) \, x .
\]
\end{lemma}

\begin{proof}
\leanok
By the product rule,
\[
  \partial_{b_j} \bigl(\langle b_j, \cdot \rangle A\bigr)(x)
    = \langle b_j, b_j \rangle A(x) + \langle b_j, x \rangle \, DA(x) \, b_j ,
\]
and $\sum_j \langle b_j, x \rangle b_j = x$.
\end{proof}

\begin{theorem}[The dilation identity]
\label{thm:dilation-divergence}
\lean{Rd.fderiv_dilationAverage_apply, Rd.continuous_fderiv_dilate_exp_apply,
  Rd.dilate_two_sub_eq, Rd.sum_fderiv_divergenceKernel, Rd.sum_lineDerivOp_divergenceSchwartz,
  Rd.sum_lineDerivOp_schwartzDilate_divergenceSchwartz}
\leanok
\uses{def:divergence-kernel, lem:exponential-dilation, lem:parametric-integral-smooth,
  lem:divergence-kernel-smooth, def:divergence-schwartz, lem:derivative-dilate,
  lem:radial-divergence}
For every continuously differentiable $m \colon V \to \C$ and every $x \in V$,
\[
  m_2(x) - m(x) = n \, A_m(x) + D A_m(x) \, x = \sum_{i \in I} \partial_{e_i} K^m_{e_i}(x) .
\]
If $m$ is smooth and vanishes outside a ball, then for every $c \neq 0$ and every $x \in V$,
\[
  \sum_{i \in I} \partial_{e_i} \bigl(K^m_{e_i}\bigr)_c (x) = c \, (m_2 - m)_c(x) .
\]
\end{theorem}

\begin{proof}
\leanok
By \Cref{lem:parametric-integral-smooth},
$D A_m(x) \, x = \int_0^{\log 2} D(m_{e^s})(x) \, x \, ds$,
so by \Cref{lem:exponential-dilation} $n A_m(x) + D A_m(x) \, x$ is the integral of the
derivative of $s \mapsto m_{e^s}(x)$ over $[0, \log 2]$, which is
$m_2(x) - m_1(x) = m_2(x) - m(x)$ by the fundamental theorem of calculus, and
$\sum_i \partial_{e_i} K^m_{e_i}(x) = n A_m(x) + D A_m(x) \, x$ by \Cref{lem:radial-divergence}.
The second identity follows with $\partial_w (G_c) = c \, (\partial_w G)_c$
(\Cref{lem:derivative-dilate}).
\end{proof}

\begin{lemma}[The local means kernels as dilation differences]
\label{lem:local-means-kernel-dilation}
\lean{Rd.LocalMeansBump.kernel_succ_eq_dilate}
\leanok
\uses{def:local-means, lem:dilate}
For every bump $g$, every $L \in \N$ and every $j \in \N$, the moment kernel $m$ of order $L$ and
the local means kernels $k_j$ of order $L$ satisfy $k_{j+1} = (m_2 - m)_{2^j}$.
\end{lemma}

\begin{proof}
\leanok
$k_{j+1} = m_{2^{j+1}} - m_{2^j}$, and $(m_2)_{2^j} = m_{2^{j+1}}$ (\Cref{lem:dilate}).
\end{proof}

\paragraph{The mean.}

\begin{lemma}[The divergence kernels have mean zero]
\label{lem:divergence-kernel-mean}
\lean{Rd.integrable_inner_mul_dilate_exp, Rd.integral_divergenceKernel}
\leanok
\uses{def:divergence-kernel, lem:first-moments-dilate}
For every continuous $m \colon V \to \C$ that vanishes outside a ball and every $w \in V$ with
$\int_V \langle w, x \rangle \, m(x) \, dx = 0$, $\int_V K^m_w = 0$.
\end{lemma}

\begin{proof}
\leanok
The function $(x, s) \mapsto \langle w, x \rangle m_{e^s}(x)$ is continuous and vanishes for
$s \in [0, \log 2]$ outside the compact set $\bar B(0, R) \times [0, \log 2]$, so it is
integrable on
$V \times (0, \log 2]$, and Fubini's theorem gives
$\int K^m_w = \int_0^{\log 2} \int_V \langle w, x \rangle m_{e^s}(x) \, dx \, ds
= \int_0^{\log 2} e^{-s} \cdot 0 \, ds = 0$ by \Cref{lem:first-moments-dilate}.
\end{proof}

\paragraph{Mean-zero kernels as Calder\'on--Zygmund families.}

\begin{lemma}[Fourier transforms of mean-zero kernels]
\label{lem:mean-zero-fourier}
\lean{Rd.norm_fourier_le_of_integral_eq_zero, Rd.ofReal_mul_enorm_mul_enorm_fourier_le,
  Rd.enorm_fourier_le_mul_min}
\leanok
\uses{lem:fourier-character-lipschitz, lem:fourier-integrable, def:gradient-mass,
  lem:gradient-mass}
For every Schwartz function $G$ and every $\xi \in V$,
$2 \pi \|\xi\| \, |\widehat G(\xi)| \leq C_G$ with the gradient mass $C_G$ of
\Cref{def:gradient-mass}. If moreover $\int G = 0$, then
$|\widehat G(\xi)| \leq 2 \pi \|\xi\| \int_V \|x\| \, |G(x)| \, dx$ and
\[
  |\widehat G(\xi)| \leq C \min\bigl(\|\xi\|, \|\xi\|^{-1}\bigr), \qquad
  C = 2 \pi \int_V \|x\| \, |G(x)| \, dx + \frac{C_G}{2 \pi} .
\]
\end{lemma}

\begin{proof}
\leanok
$\widehat{\partial_\xi G}(\xi) = 2 \pi i \|\xi\|^2 \widehat G(\xi)$ (Mathlib's
\texttt{SchwartzMap.\allowbreak fourier\_\allowbreak lineDerivOp\_\allowbreak eq}), and
$|\widehat{\partial_\xi G}(\xi)| \leq \|\partial_\xi G\|_1 \leq \|\xi\| C_G$
(\Cref{lem:fourier-integrable}, \Cref{lem:gradient-mass}). If $\int G = 0$,
$\widehat G(\xi) = \int (e^{-2 \pi i \langle x, \xi \rangle} - 1) \, G(x) \, dx$, and
$|e^{-2 \pi i a} - 1| \leq 2 \pi |a|$ (\Cref{lem:fourier-character-lipschitz}) with
$|\langle x, \xi \rangle| \leq \|x\| \|\xi\|$.
\end{proof}

\begin{theorem}[Dilates of mean-zero kernels]
\label{thm:mean-zero-cz-family}
\lean{Rd.meanZeroOverlap, Rd.meanZeroOverlap_ne_top, Rd.tsum_enorm_fourier_dilate_sq_le,
  Rd.isCZFamily_dilate_of_integral_eq_zero, Rd.isCZFamily_divergenceSchwartz}
\leanok
\uses{lem:mean-zero-fourier, lem:dyadic-min-sum, lem:fourier-dilate, thm:cz-family-dilation,
  lem:uniform-schwartz-constants, def:cz-family, def:hormander-const,
  lem:divergence-kernel-mean, def:divergence-schwartz}
For every Schwartz function $G$ with $\int G = 0$ and every $\xi \in V$,
\[
  \sum_{j \in \N} \bigl|\widehat{G_{2^j}}(\xi)\bigr|^2 \leq 4 C^2
\]
with the constant $C$ of \Cref{lem:mean-zero-fourier}, and the dilates $(G_{2^j})_{j \in \N}$
form a Calder\'on--Zygmund family with Fourier overlap $4 C^2$ and H\"ormander constant $4 H_A$
for the seminorms $A_{k,l}$ of $G$ (\Cref{def:hormander-const}). In particular, for every smooth
$m$ that vanishes outside a ball and every $w \in V$ with
$\int_V \langle w, x \rangle \, m(x) \, dx = 0$, the dilates
$\bigl((K^m_w)_{2^j}\bigr)_{j \in \N}$ form a Calder\'on--Zygmund family.
\end{theorem}

\begin{proof}
\leanok
$\widehat{G_{2^j}}(\xi) = \widehat G(2^{-j} \xi)$ (\Cref{lem:fourier-dilate}). For $\xi = 0$ every
term vanishes. Otherwise, with $a = \|\xi\|^{-1}$ and $t_j = 2^j a$, \Cref{lem:mean-zero-fourier}
gives $|\widehat G(2^{-j} \xi)|^2 \leq C^2 \min(t_j, t_j^{-1})^2 \leq C^2 \min(t_j, t_j^{-1})$,
and $\sum_j \min(t_j, t_j^{-1}) \leq 4$ (\Cref{lem:dyadic-min-sum}). The constant family $G$ is
uniformly Schwartz (\Cref{lem:uniform-schwartz-constants}), and \Cref{thm:cz-family-dilation}
applies. The divergence kernels have mean zero by \Cref{lem:divergence-kernel-mean}.
\end{proof}

\section{Sobolev spaces as Triebel--Lizorkin spaces}
\label{sec:sobolev-triebel-lizorkin}

Let $n \geq 1$ and $1 < p < \infty$. This section proves $\Sob{m}{p}(V) = \TL{m}{p}{2}$ for every
$m \in \N$, with named constants in both directions, formulas in the Schwartz seminorms of fixed
kernels that are proved finite. For $f \in \Sob{m}{p}(V)$, the blocks
$2^{jm} \psi_{j+1} * f$ are sums of convolutions of the weak derivatives of order $m$ with the
dilated derivative kernels (\Cref{cor:psi-derivatives-dilated}), and the analysis bound of their
Calder\'on--Zygmund families gives $f \in \TL{m}{p}{2}$. Conversely, for $f \in \TL{k}{p}{2}$ and
$w \in V^k$, the derivatives $\sum_{j < J} (\partial^w \tilde\psi_j) * (\psi_j * f)$
of the reproducing sums converge in $L_p(V)$ by the synthesis bound of the dual derivative
families (\Cref{thm:cz-family-dual-derivative}), and their limit is $\partial^w f$. With the
Besov sandwich (\Cref{thm:tl-besov-sandwich}) this places $\Sob{m}{p}(V)$ between
$\Besov{m}{p}{\min(p,2)}$ and $\Besov{m}{p}{\max(p,2)}$ for $m \geq 1$. The identity
$\Sob{m}{p}(V) = \TL{m}{p}{2}$ is classical; see Triebel~\cite{Triebel}. In this section
$e_\alpha$, $\eta_\alpha$, $\theta^\alpha_j$ and $\tilde\theta^w_j$ are those of
\Cref{sec:derivative-kernels}.

\paragraph{From weak derivatives to the blocks.}

\begin{definition}[The constants of the embedding into $\TL{m}{p}{2}$]
\label{def:sobolev-to-tl-const}
\lean{Rd.sobolevToTLConst, Rd.sobolevToTLConst_ne_top}
\leanok
\uses{def:cz-const, def:hormander-const, def:derivative-family}
For $1 < p < \infty$, every $m \in \N$ and every multi-index $\alpha$,
\[
  K_{V,p,\alpha} = S_{V,p}\bigl(2, 4 H_{A^\alpha}\bigr) ,
\]
which is finite.
\end{definition}

\begin{lemma}[The blocks as convolutions with the derivatives]
\label{lem:blocks-as-derivative-convolutions}
\lean{Rd.enorm_two_pow_complex, Rd.two_pow_mul_blockFun_succ_ae,
  Rd.lqNorm_blockFun_succ_le_sum_ae}
\leanok
\uses{cor:psi-derivatives-dilated, thm:weak-derivative-convolution, lem:averaged-differences,
  def:derivative-family, def:triebel-lizorkin, lem:mixed-norm-tools}
For every $l \in \N$, the complex number $2^l$ has absolute value $|2^l| = 2^l$.
Let $1 \leq p < \infty$, $m \in \N$ and $f \in L_p(V)$, and for every multi-index $\alpha$ let
$g_\alpha \in L_p(V)$ be the weak derivative $\partial^{e_\alpha} f$ on $V$. For almost every
$x \in V$ and every $j \in \N$,
\[
  2^{jm} (\psi_{j+1} * f)(x) = \sum_\alpha (\theta^\alpha_j * g_\alpha)(x) ,
\]
and for almost every $x \in V$,
\[
  \bigl\| (2^{(j+1)m} |(\psi_{j+1} * f)(x)|)_j \bigr\|_{\ell_2}
  \leq 2^m \sum_\alpha \bigl\| ((\theta^\alpha_j * g_\alpha)(x))_j \bigr\|_{\ell_2} .
\]
\end{lemma}

\begin{proof}
\leanok
By \Cref{cor:psi-derivatives-dilated}, $2^{jm} \psi_{j+1} * f$ is the convolution of $f$ with
$\sum_\alpha \partial^{e_\alpha} \theta^\alpha_j$, which is
$\sum_\alpha (\partial^{e_\alpha} \theta^\alpha_j) * f$ almost everywhere
(\Cref{lem:averaged-differences}); each term is $\theta^\alpha_j * g_\alpha$ at every point
(\Cref{thm:weak-derivative-convolution}). For the bound, $|2^{jm}| = 2^{jm}$ in $\C$, so the
triangle inequality in $\C$ gives
$2^{jm} |(\psi_{j+1} * f)(x)| \leq \sum_\alpha |(\theta^\alpha_j * g_\alpha)(x)|$, and
\Cref{lem:mixed-norm-tools} (5) gives the claim.
\end{proof}

\begin{theorem}[$\Sob{m}{p} \subseteq \TL{m}{p}{2}$]
\label{thm:sobolev-to-triebel-lizorkin}
\lean{Rd.eLpNorm_lqNorm_blockFun_succ_le, Rd.triebelLizorkinNorm_le_of_hasWeakLineDeriv,
  Rd.memTriebelLizorkin_of_hasWeakLineDeriv, Rd.memTriebelLizorkin_of_memSobolev}
\leanok
\uses{lem:blocks-as-derivative-convolutions, def:sobolev-to-tl-const, thm:cz-family-derivative,
  thm:vector-valued-strong-type, lem:tl-pointwise-blocks, lem:mixed-norm-tools,
  def:weak-line-derivative, lem:weak-derivative-basic}
Let $n \geq 1$, $1 < p < \infty$, $m \in \N$ and $f \in L_p(V)$, and for every multi-index
$\alpha$ let $g_\alpha \in L_p(V)$ be the weak derivative $\partial^{e_\alpha} f$ on $V$. Then
\[
  \bigl\| \|(2^{(j+1)m} |(\psi_{j+1} * f)(x)|)_j\|_{\ell_2} \bigr\|_p
  \leq 2^m \sum_\alpha K_{V,p,\alpha} \|g_\alpha\|_p
\]
and
\[
  \|f\|_{\TL{m}{p}{2}} \leq \|\Phi\|_1 \, \|f\|_p + 2^m \sum_\alpha K_{V,p,\alpha} \,
  \|g_\alpha\|_p .
\]
In particular, $f \in \TL{m}{p}{2}$, and every $f \in \Sob{m}{p}(V)$ lies in $\TL{m}{p}{2}$.
\end{theorem}

\begin{proof}
\leanok
For the first bound, take the $L_p$ norm in \Cref{lem:blocks-as-derivative-convolutions} with
Minkowski's inequality in $L_p$ (\Cref{lem:mixed-norm-tools} (3)), and bound each term by the
strong type of the analysis operator (\Cref{thm:vector-valued-strong-type}) of the
Calder\'on--Zygmund family $(\theta^\alpha_j)_j$ with overlap $2$ and H\"ormander constant
$4 H_{A^\alpha}$ (\Cref{thm:cz-family-derivative}). Splitting off the block $j = 0$
(\Cref{lem:mixed-norm-tools} (4)) and $\|\psi_0 * f\|_p \leq \|\Phi\|_1 \, \|f\|_p$
(\Cref{lem:tl-pointwise-blocks}) give the second bound.
For $f \in \Sob{m}{p}(V)$ the weak derivatives $\partial^{e_\alpha} f$ exist in $L_p(V)$
(\Cref{lem:weak-derivative-basic}), and the right side is finite.
\end{proof}

\paragraph{From the blocks to weak derivatives.}

\begin{definition}[The constants of the embedding into $\Sob{m}{p}$]
\label{def:tl-to-sobolev-const}
\lean{Rd.tlToSobolevConst, Rd.tlToSobolevConst_ne_top}
\leanok
\uses{def:cz-const, def:hormander-const, def:dual-derivative-family}
For $1 < p < \infty$, every $k \in \N$ and all $w \in V^k$,
\[
  C_{V,p,w} = S_{V,p}\bigl(B_w, 4 H_{\tilde A^w}\bigr) ,
\]
which is finite.
\end{definition}

\begin{lemma}[The synthesis bound with scaled coefficients]
\label{lem:scaled-synthesis-bound}
\lean{Rd.IsCZFamily.eLpNorm_sum_conv_smul_le, Rd.IsCZFamily.eLpNorm_sum_conv_smul_blockFun_le}
\leanok
\uses{thm:vector-valued-strong-type, def:cz-const, def:triebel-lizorkin, lem:tl-pointwise-blocks}
Let $n \geq 1$, $1 < p < \infty$, $(\theta_j)_j$ a Calder\'on--Zygmund family with overlap $B$ and
H\"ormander constant $H$, $\sigma \in \R$ and $c_j \in \C$ with $|c_j| = 2^{j\sigma}$. For every
finite set $t \subseteq \N$ and all $g_j \in L_p(V)$,
\[
  \Bigl\| \sum_{j \in t} \theta_j * (c_j g_j) \Bigr\|_p
  \leq S_{V,p}(B, H) \, \Bigl\| x \mapsto \bigl\| \bigl(\mathbf 1_t(j) \, 2^{j\sigma}
  |g_j(x)|\bigr)_j \bigr\|_{\ell_2} \Bigr\|_p .
\]
\end{lemma}

\begin{proof}
\leanok
The functions $\mathbf 1_t(j) \, c_j g_j$ lie in $L_p(V)$ and have the absolute values
$\mathbf 1_t(j) \, 2^{j\sigma} |g_j|$; the synthesis bound of \Cref{thm:vector-valued-strong-type}
applies.
\end{proof}

For $f \in L_p(V)$ the blocks $g_j = \psi_j * f$ lie in $L_p(V)$ (\Cref{lem:tl-pointwise-blocks}),
and the lemma bounds $\sum_{j \in t} \theta_j * (c_j \, \psi_j * f)$ by the weighted blocks of $f$.

\begin{lemma}[Convergent synthesis sums]
\label{lem:synthesis-sums-convergent}
\lean{Rd.memLp_sum_conv_blockFun, Rd.eLpNorm_sum_conv_blockFun_sub_le,
  Rd.cauchySeq_toLp_sum_conv_blockFun, Rd.exists_tendsto_sum_conv_blockFun}
\leanok
\uses{def:triebel-lizorkin, thm:young, lem:weighted-block-tails, lem:tl-pointwise-blocks,
  lem:mixed-norm-tools}
Let $\sigma$ be real, $f \in L_p(V)$, $C \in [0, \infty]$ and $K_j$ Schwartz functions with
\[
  \Bigl\| \sum_{j \in t} K_j * (\psi_j * f) \Bigr\|_p \leq C \, \Bigl\| x \mapsto
  \bigl\| \bigl(\mathbf 1_t(j) \, 2^{j\sigma} |(\psi_j * f)(x)|\bigr)_j \bigr\|_{\ell_2}
  \Bigr\|_p
\]
for every finite set $t \subseteq \N$.
\begin{enumerate}
  \item For all $N \leq N'$,
    \[
      \Bigl\| \sum_{N \leq j < N'} K_j * (\psi_j * f) \Bigr\|_p \leq C \, \Bigl\| x \mapsto
      \bigl\| \bigl(2^{j\sigma} |(\psi_j * f)(x)|\bigr)_{j \geq N} \bigr\|_{\ell_2} \Bigr\|_p ,
    \]
    where the $\ell_2$ norm is taken of the sequence with the terms $j < N$ replaced by $0$.
  \item For $1 \leq p < \infty$, the sums $\sum_{j \in t} K_j * (\psi_j * f)$ lie in $L_p(V)$ for
    every finite set $t \subseteq \N$.
  \item For $1 \leq p < \infty$, $C < \infty$ and $f \in \TL{\sigma}{p}{2}$, the partial sums
    $\sum_{j < J} K_j * (\psi_j * f)$ form a Cauchy sequence in $L_p(V)$, and they converge to
    some $g \in L_p(V)$ with $\|g\|_p \leq C \|f\|_{F^\sigma_{p,2}}$.
\end{enumerate}
\end{lemma}

\begin{proof}
\leanok
(1) The sum is the sum over $t = \{N, \dots, N' - 1\}$, and setting the terms $j < N$ instead of
the terms $j \notin t$ to zero enlarges the mixed norm (\Cref{lem:mixed-norm-tools} (1)).
(2) The blocks $\psi_j * f$ lie in $L_p(V)$ (\Cref{lem:tl-pointwise-blocks}), and Young's
inequality (\Cref{thm:young}) applies. (3) By (1) and \Cref{lem:mixed-norm-tools} (2) for the
functions $2^{j\sigma} |\psi_j * f|$, whose mixed norm $\|f\|_{F^\sigma_{p,2}}$ is finite
(\Cref{def:triebel-lizorkin}), the partial sums form a Cauchy sequence. The limit $g$ exists
since $L_p(V)$ is complete, and each partial sum has norm at most $C \|f\|_{F^\sigma_{p,2}}$
(\Cref{lem:weighted-block-tails}).
\end{proof}

\begin{lemma}[The synthesis bound for derivatives of the synthesis kernels]
\label{lem:derivative-synthesis-bound}
\lean{Rd.eLpNorm_sum_conv_derivDual_le, Rd.derivSynthesisSum, Rd.memLp_derivSynthesisSum,
  Rd.eLpNorm_derivSynthesisSum_sub_le}
\leanok
\uses{def:tl-to-sobolev-const, lem:psi-dual-derivatives, thm:cz-family-dual-derivative,
  lem:blocks-as-derivative-convolutions, lem:dyadic-weights-integer, lem:scaled-synthesis-bound,
  lem:synthesis-sums-convergent}
Let $n \geq 1$, $1 < p < \infty$, $k \in \N$, $w \in V^k$ and $f \in L_p(V)$. For every finite
set $s \subseteq \N$,
\begin{multline*}
  \Bigl\| \sum_{j \in s} (\partial^w \tilde\psi_j) * (\psi_j * f) \Bigr\|_p \\
  \leq C_{V,p,w} \, \bigl\| \|(\mathbf 1_s(j) \, 2^{jk} |(\psi_j * f)(x)|)_j\|_{\ell_2}
  \bigr\|_p .
\end{multline*}
The partial sums $D_J = \sum_{j < J} (\partial^w \tilde\psi_j) * (\psi_j * f)$ lie in $L_p(V)$,
and for all $N \leq N'$,
\[
  \|D_{N'} - D_N\|_p \leq C_{V,p,w} \,
  \bigl\| \|(2^{jk} |(\psi_j * f)(x)|)_{j \geq N}\|_{\ell_2} \bigr\|_p .
\]
\end{lemma}

\begin{proof}
\leanok
$\partial^w \tilde\psi_j = 2^{jk} \tilde\theta^w_j$ (\Cref{lem:psi-dual-derivatives}), so the
sum is $\sum_{j \in s} \tilde\theta^w_j * (c_j \, \psi_j * f)$ with $c_j = 2^{jk}$, and
$|c_j| = 2^{jk}$ in $\C$ (\Cref{lem:blocks-as-derivative-convolutions}), the weight of
smoothness $k$ (\Cref{lem:dyadic-weights-integer}). The synthesis bound with scaled blocks
(\Cref{lem:scaled-synthesis-bound}) for the Calder\'on--Zygmund family $(\tilde\theta^w_j)_j$
(\Cref{thm:cz-family-dual-derivative}) gives the first bound. The partial sums lie in $L_p(V)$,
and the first bound gives the second (\Cref{lem:synthesis-sums-convergent} (2) and (1)).
\end{proof}

\begin{lemma}[The partial sums of the derivative]
\label{lem:derivative-partial-sums}
\lean{Rd.cauchySeq_toLp_derivSynthesisSum, Rd.hasWeakLineDeriv_sum_conv_psiDual}
\leanok
\uses{lem:derivative-synthesis-bound, lem:synthesis-sums-convergent,
  thm:weak-derivative-convolution, lem:weak-derivative-basic, lem:psi-dual-derivatives}
Let $n \geq 1$, $1 < p < \infty$, $k \in \N$, $w \in V^k$ and $f \in L_p(V)$. For every
$J \in \N$, $D_J$ is the weak derivative $\partial^w$ of $\sum_{j < J} \tilde\psi_j * (\psi_j * f)$
on $V$. If $f \in \TL{k}{p}{2}$, then $(D_J)_J$ is a Cauchy sequence in $L_p(V)$.
\end{lemma}

\begin{proof}
\leanok
Each term $(\partial^w \tilde\psi_j) * (\psi_j * f)$ is the weak derivative of
$\tilde\psi_j * (\psi_j * f)$ (\Cref{thm:weak-derivative-convolution}), and weak derivatives
add up (\Cref{lem:weak-derivative-basic}). The first bound of
\Cref{lem:derivative-synthesis-bound} and \Cref{lem:synthesis-sums-convergent} (3) give the
Cauchy property.
\end{proof}

\begin{theorem}[$\TL{m}{p}{2} \subseteq \Sob{m}{p}$]
\label{thm:triebel-lizorkin-to-sobolev}
\lean{Rd.exists_hasWeakLineDeriv_of_memTriebelLizorkin,
  Rd.exists_hasWeakLineDeriv_of_memTriebelLizorkin_of_le, Rd.memSobolev_of_memTriebelLizorkin}
\leanok
\uses{lem:derivative-partial-sums, lem:derivative-synthesis-bound, lem:synthesis-sums-convergent,
  lem:square-function-reproducing, thm:weak-derivative-limits, lem:weak-derivative-basic,
  def:tl-to-sobolev-const, prop:tl-elementary}
Let $n \geq 1$, $1 < p < \infty$, $k \in \N$ and $w \in V^k$. For every $f \in \TL{k}{p}{2}$ the
weak derivative $\partial^w f$ exists on $V$ as a function $g \in L_p(V)$ with
\[
  \|g\|_p \leq C_{V,p,w} \, \|f\|_{\TL{k}{p}{2}} .
\]
For every $m \geq k$ and every $f \in \TL{m}{p}{2}$, the weak derivative $g = \partial^w f$
satisfies $\|g\|_p \leq C_{V,p,w} \, \|f\|_{\TL{m}{p}{2}}$. In particular,
$\TL{m}{p}{2} \subseteq \Sob{m}{p}(V)$ for every $m \in \N$.
\end{theorem}

\begin{proof}
\leanok
By the first bound of \Cref{lem:derivative-synthesis-bound} and
\Cref{lem:synthesis-sums-convergent} (3), $D_J$ converges in $L_p(V)$ to some $g$ with
$\|g\|_p \leq C_{V,p,w} \|f\|_{\TL{k}{p}{2}}$. The reproducing sums tend to $f$ in $L_p(V)$
(\Cref{lem:square-function-reproducing}), and $D_J$ is their weak derivative $\partial^w$
(\Cref{lem:derivative-partial-sums}), so $g = \partial^w f$ weakly on $V$
(\Cref{thm:weak-derivative-limits}). For $m \geq k$, $\TL{m}{p}{2} \subseteq \TL{k}{p}{2}$ with
$\|f\|_{\TL{k}{p}{2}} \leq \|f\|_{\TL{m}{p}{2}}$ (\Cref{prop:tl-elementary}). The last claim is
\Cref{lem:weak-derivative-basic}.
\end{proof}

\paragraph{Consequences.}

\begin{corollary}[$\Sob{m}{p} = \TL{m}{p}{2}$]
\label{cor:sobolev-triebel-lizorkin}
\lean{Rd.memSobolev_iff_memTriebelLizorkin}
\leanok
\uses{thm:sobolev-to-triebel-lizorkin, thm:triebel-lizorkin-to-sobolev}
For every $n \geq 1$, every $1 < p < \infty$, every $m \in \N$ and every $f \in L_p(V)$,
$f \in \Sob{m}{p}(V)$ if and only if $f \in \TL{m}{p}{2}$.
\end{corollary}

\begin{proof}
\leanok
\Cref{thm:sobolev-to-triebel-lizorkin} and \Cref{thm:triebel-lizorkin-to-sobolev}.
\end{proof}

For $m = 0$ both sides are all of $L_p(V)$; the case $m = 0$ of $\TL{m}{p}{2}$ is the square
function theorem (\Cref{cor:square-function-theorem}).

\begin{corollary}[The Besov sandwich for $\Sob{m}{p}$]
\label{cor:sobolev-besov-sandwich}
\lean{Rd.memApprox_of_memSobolev, Rd.memSobolev_of_memApprox,
  Rd.approxNormDyadic_le_of_hasWeakLineDeriv, Rd.exists_hasWeakLineDeriv_of_memApprox}
\leanok
\uses{cor:sobolev-triebel-lizorkin, thm:sobolev-to-triebel-lizorkin,
  thm:triebel-lizorkin-to-sobolev, thm:tl-besov-sandwich, cor:square-function-theorem}
Let $n \geq 1$ and $1 < p < \infty$. For every $m \in \N$,
\[
  \Besov{m}{p}{\min(p,2)} \subseteq \Sob{m}{p}(V) ,
\]
and for every $m \geq 1$,
\[
  \Sob{m}{p}(V) \subseteq \Besov{m}{p}{\max(p,2)} .
\]
For every $m \geq 1$ and every $f \in L_p(V)$ with weak derivatives
$g_\alpha = \partial^{e_\alpha} f \in L_p(V)$ for every multi-index $\alpha$ of length $m$,
\[
  \|f\|^{\mathrm{dy}}_{A^m_{\max(p,2)}} \leq 2 (1 - 2^{-m})^{-1}
  \Bigl( \|\Phi\|_1 \, \|f\|_p + 2^m \sum_\alpha K_{V,p,\alpha} \, \|g_\alpha\|_p \Bigr) .
\]
For every $m \in \N$, every $f \in \Besov{m}{p}{\min(p,2)}$, every $k \leq m$ and all
$w \in V^k$, the weak derivative $\partial^w f$ is a function $g \in L_p(V)$ with
\[
  \|g\|_p \leq C_{V,p,w} \cdot 8 \cdot 8^m \, \|\Phi\|_1 \, \|f\|^{\mathrm{dy}}_{A^m_{\min(p,2)}} .
\]
\end{corollary}

\begin{proof}
\leanok
The inclusion $\Besov{m}{p}{\min(p,2)} \subseteq \TL{m}{p}{2}$
(\Cref{cor:square-function-theorem}) and \Cref{cor:sobolev-triebel-lizorkin} give the first
inclusion, and the Besov sandwich for $s = m > 0$ gives the second. The bounds combine
\Cref{thm:sobolev-to-triebel-lizorkin} and \Cref{thm:triebel-lizorkin-to-sobolev} with the
bounds of \Cref{thm:tl-besov-sandwich}.
\end{proof}

\begin{corollary}[$\Sob{m}{2} = \Besov{m}{2}{2}$]
\label{cor:sobolev-integer-two-besov}
\lean{Rd.memSobolev_two_iff_memApprox}
\leanok
\uses{cor:sobolev-triebel-lizorkin, thm:tl-diagonal}
For every $n \geq 1$, every $m \geq 1$ and every $f \in L_2(V)$, $f \in \Sob{m}{2}(V)$ if and
only if $f \in \Besov{m}{2}{2}$.
\end{corollary}

\begin{proof}
\leanok
\Cref{cor:sobolev-triebel-lizorkin} and $\TL{m}{2}{2} = \Besov{m}{2}{2}$ (\Cref{thm:tl-diagonal}).
\end{proof}

This extends \Cref{thm:sobolev-one-two-besov}, the case $m = 1$ (in every dimension), to every
order $m$ when $n \geq 1$.

\section{Sobolev embeddings}
\label{sec:sobolev-embedding}

Let $1 \leq p < \infty$ and, for $p \leq r \leq \infty$, let
$\alpha(n, p, r) = n (1/p - 1/r)$ be the Nikolskii exponent of \Cref{def:nikolskii-property}
with $D = n$, so that $\alpha(n, p, \infty) = n/p$. This section proves the embeddings of
$\Besov{s}{p}{q}$ into $L_r(V)$ for $s > \alpha(n, p, r)$ and into the bounded continuous
functions for $s > n/p$, directly from the Littlewood--Paley blocks: a block is band-limited, so
its supremum is at most $N_p 2^{jn/p}$ times its $L_p$ norm (\Cref{def:nikolskii-block-const}),
and interpolation gives the Nikolskii inequality
$\|\psi_j * f\|_r \lesssim 2^{j \alpha(n,p,r)} \|\psi_j * f\|_p$. For $s > \alpha$
the series $\sum_j \psi_j * f$ then converges absolutely in $L_r(V)$, and for $s > n/p$ uniformly.
With the bound of $\|f\|_{\TL{m}{p}{2}}$ by the weak derivatives
(\Cref{thm:sobolev-to-triebel-lizorkin}) and $\Sob{m}{p}(V) \subseteq \Besov{m}{p}{\max(p,2)}$
(\Cref{cor:sobolev-besov-sandwich}) this gives the Sobolev embeddings
$\Sob{m}{p}(V) \hookrightarrow L_r(V)$ for $m > n(1/p - 1/r)$ and
$\Sob{m}{p}(V) \hookrightarrow C_b(V)$ for $m > n/p$, with named constants, formulas in norms
and Schwartz seminorms of fixed kernels that are proved finite. These embeddings are classical;
see Nikol'ski\u{\i}~\cite{Nikolskii} and Triebel~\cite{Triebel}.

\paragraph{Continuity of the blocks.}

\begin{lemma}[Convolutions with integrable bounded kernels]
\label{lem:bounded-kernel-conv-continuous}
\lean{Rd.integrable_mul_sub_of_enorm_le, Rd.conv_sub_apply_of_enorm_le,
  Rd.continuous_conv_of_enorm_le, Rd.continuous_conv_schwartz, Rd.coe_schwartzDilate_psiBase,
  Rd.coe_schwartzDilate_bandKernelS, Rd.continuous_blockFun}
\leanok
\uses{def:conv, lem:conv-pointwise-bound, lem:lp-interpolation, lem:schwartz-decay,
  def:schwartz-dilate, lem:psi-base, def:lp-kernels, def:triebel-lizorkin}
Let $1 \leq p \leq \infty$, let $K : V \to \C$ be integrable with $|K(y)| \leq B$ for all
$y \in V$ and some $B < \infty$, and let $f, g \in L_p(V)$. For every $x \in V$ the function
$y \mapsto K(y) f(x - y)$ is integrable, and
\[
  (K * f)(x) - (K * g)(x) = (K * (f - g))(x) .
\]
If $p < \infty$, then $K * f$ is continuous; in particular $G * f$ is continuous for every
$G \in \mathcal S(V)$. For every $j \in \N$, $\psi_j$ and $\Theta_j$ are Schwartz functions: the
Schwartz dilations $(G^{\psi}_j)_{2^j}$ and $\Theta_{2^j}$ of \Cref{def:schwartz-dilate}. For
every $1 \leq p < \infty$, every $f \in L_p(V)$ and every $j \in \N$ the block $\psi_j * f$ is
continuous.
\end{lemma}

\begin{proof}
\leanok
Since $K \in L_1$ and $|K| \leq B$, $K \in L_{p'}$ (\Cref{lem:lp-interpolation}), so the integrand
is integrable by H\"older's inequality, and the identity is the linearity of the integral. For
$p < \infty$ choose continuous functions $g_k$ with compact support and
$\|f - g_k\|_p \leq 1/k$ (Mathlib's density theorem); each $K * g_k$ is continuous.
\Cref{lem:conv-pointwise-bound} gives
$|(K * f)(x) - (K * g_k)(x)| \leq \|K\|_1^{1-1/p} B^{1/p} / k$ for every $x$, so $K * g_k \to K * f$
uniformly and $K * f$ is continuous. A Schwartz function $G$ is integrable and bounded by its
decay constant (\Cref{lem:schwartz-decay}). The kernels $\psi_j$ and $\Theta_j$ are dilations of
Schwartz functions (\Cref{def:schwartz-dilate}, \Cref{lem:psi-base}, \Cref{def:lp-kernels}).
\end{proof}

\paragraph{The Nikolskii inequality for blocks.}

\begin{definition}[The Nikolskii constant of the blocks]
\label{def:nikolskii-block-const}
\lean{Rd.iSup_enorm_bandKernel_lt_top, Rd.nikolskiiBlockConst, Rd.nikolskiiBlockConst_ne_top}
\leanok
\uses{def:lp-kernels, lem:lp-kernels-basic, lem:conv-pointwise-bound}
For $1 \leq p \leq \infty$,
\[
  N_p = \|\Theta\|_1^{1 - 1/p} \, \bigl(\sup |\Theta|\bigr)^{1/p} ,
\]
with $1/\infty = 0$. The supremum $\sup |\Theta|$ and $N_p$ are finite.
\end{definition}

\begin{lemma}[The blocks are band-limited at every point]
\label{lem:blocks-band-limited-pointwise}
\lean{Rd.enorm_Theta_le, Rd.conv_Theta_blockFun}
\leanok
\uses{def:lp-kernels, lem:lp-kernels-basic, lem:schwartz-decay, thm:kernel-identities,
  lem:conv-assoc, lem:bounded-kernel-conv-continuous, lem:tl-pointwise-blocks}
For every $j \in \N$ and every $y \in V$, $|\Theta_j(y)| \leq 2^{jn} \sup |\Theta|$. For every
$1 \leq p < \infty$, every $f \in L_p(V)$, every $j \in \N$ and every $x \in V$,
\[
  (\Theta_j * (\psi_j * f))(x) = (\psi_j * f)(x) .
\]
\end{lemma}

\begin{proof}
\leanok
$\Theta_j(y) = 2^{jn} \Theta(2^j y)$. Almost everywhere,
$\Theta_j * (\psi_j * f) = (\Theta_j * \psi_j) * f = \psi_j * f$ by \Cref{lem:conv-assoc} and
\Cref{thm:kernel-identities}. Both sides are continuous (\Cref{lem:bounded-kernel-conv-continuous},
with $\psi_j * f \in L_p$ by \Cref{lem:tl-pointwise-blocks}), and continuous functions that agree
almost everywhere agree everywhere.
\end{proof}

\begin{theorem}[The Nikolskii inequality for blocks]
\label{thm:nikolskii-blocks}
\lean{Rd.enorm_blockFun_le_nikolskii, Rd.eLpNorm_blockFun_le_nikolskii,
  Rd.memLp_blockFun_of_le}
\leanok
\uses{def:nikolskii-block-const, lem:blocks-band-limited-pointwise, lem:conv-pointwise-bound,
  lem:lp-kernels-basic, lem:lp-interpolation, def:nikolskii-property, def:lq-norm,
  lem:tl-pointwise-blocks}
Let $1 \leq p < \infty$ and $f \in L_p(V)$. For every $j \in \N$ and every $x \in V$,
\[
  |(\psi_j * f)(x)| \leq N_p \, 2^{jn/p} \, \|\psi_j * f\|_p ,
\]
and for every $p \leq r \leq \infty$,
\[
  \|\psi_j * f\|_r \leq N_p^{1 - p/r} \, 2^{j \alpha(n, p, r)} \, \|\psi_j * f\|_p ;
\]
in particular $\psi_j * f \in L_r(V)$.
\end{theorem}

\begin{proof}
\leanok
By \Cref{lem:blocks-band-limited-pointwise}, $\psi_j * f = \Theta_j * (\psi_j * f)$ at every
point, and \Cref{lem:conv-pointwise-bound} with $\|\Theta_j\|_1 = \|\Theta\|_1$
(\Cref{lem:lp-kernels-basic}) and $B = 2^{jn} \sup |\Theta|$ gives the first bound, since
$(2^{jn})^{1/p} = 2^{jn/p}$. For $r = \infty$ the first bound is the second. For $r < \infty$,
interpolation (\Cref{lem:lp-interpolation}) with $M = N_p 2^{jn/p} \|\psi_j * f\|_p$ gives
$\|\psi_j * f\|_r \leq M^{1 - p/r} \|\psi_j * f\|_p^{p/r}$, and
$(n/p)(1 - p/r) = \alpha(n, p, r)$.
\end{proof}

\paragraph{Besov embeddings.}

\begin{lemma}[The series of the blocks converges almost everywhere]
\label{lem:blocks-ae-convergence}
\lean{Rd.ae_hasSum_blockFun}
\leanok
\uses{def:lp-blocks, def:lq-norm, lem:hardy, lem:lp-blocks, def:triebel-lizorkin}
Let $1 \leq p < \infty$, $s > 0$, $1 \leq q \leq \infty$ and $f \in L_p(V)$ with
$\|f\|^{\mathrm{LP}}_{s,q} < \infty$. Then for almost every $x \in V$,
\[
  \sum_{j=0}^\infty (\psi_j * f)(x) = f(x) .
\]
\end{lemma}

\begin{proof}
\leanok
$\sum_j \|\Delta_j f\|_p \leq (1 - 2^{-s})^{-1} \|f\|^{\mathrm{LP}}_{s,q} < \infty$
(\Cref{lem:hardy}), and $\sum_j \Delta_j f = f$ in $L_p$ (\Cref{lem:lp-blocks}). A series that
converges absolutely in $L_p$ converges almost everywhere to its sum (Mathlib's
\texttt{MeasureTheory.\allowbreak Lp.\allowbreak hasSum\_\allowbreak coeFn\_\allowbreak tsum}),
and $\Delta_j f = \psi_j * f$ almost everywhere (\Cref{def:triebel-lizorkin}).
\end{proof}

\begin{definition}[The constant of the Besov embedding]
\label{def:besov-embedding-const}
\lean{Rd.besovEmbeddingConst, Rd.besovEmbeddingConst_ne_top}
\leanok
\uses{def:nikolskii-block-const, def:nikolskii-property}
For $1 \leq p \leq r \leq \infty$ and $s \in \R$,
\[
  K_{V,p,r,s} = N_p^{1 - p/r} \, \bigl(1 - 2^{-(s - \alpha(n, p, r))}\bigr)^{-1} ,
\]
which is finite for $s > \alpha(n, p, r)$.
\end{definition}

\begin{lemma}[The Nikolskii bounds are summable]
\label{lem:nikolskii-blocks-summable}
\lean{Rd.tsum_nikolskii_le, Rd.tsum_eLpNorm_blockFun_le}
\leanok
\uses{def:besov-embedding-const, thm:nikolskii-blocks, lem:hardy, def:lq-norm, def:lp-blocks,
  def:triebel-lizorkin}
Let $1 \leq p \leq r \leq \infty$, $p < \infty$, $s \in \R$, $1 \leq q \leq \infty$ and
$f \in L_p(V)$. Then
\[
  \sum_{j=0}^\infty \|\psi_j * f\|_r
  \leq \sum_{j=0}^\infty N_p^{1 - p/r} \, 2^{j \alpha(n,p,r)} \, \|\psi_j * f\|_p
  \leq K_{V,p,r,s} \, \|f\|^{\mathrm{LP}}_{s,q} .
\]
\end{lemma}

\begin{proof}
\leanok
The first inequality is \Cref{thm:nikolskii-blocks}. For the second, with
$\alpha = \alpha(n, p, r)$, $2^{j(s - \alpha)} 2^{j\alpha} = 2^{js}$ and
$\|\psi_j * f\|_p = \|\Delta_j f\|_p$, so \Cref{lem:hardy} for the sequence
$(2^{j\alpha} \|\psi_j * f\|_p)_j$ and the weight $2^{j(s-\alpha)}$ gives the bound.
\end{proof}

\begin{theorem}[The Besov embedding into $L_r$]
\label{thm:besov-embedding-rd}
\lean{Rd.exists_memLp_ae_eq_of_memApprox}
\leanok
\uses{lem:nikolskii-blocks-summable, thm:nikolskii-blocks, lem:blocks-ae-convergence,
  def:besov-embedding-const, thm:lp-band, def:nikolskii-property}
Let $1 \leq p \leq r \leq \infty$, $p < \infty$, $s > \alpha(n, p, r)$ and $1 \leq q \leq \infty$.
For every $f \in \Besov{s}{p}{q}$ there is $g \in L_r(V)$ with $g = f$ almost everywhere and
\[
  \|g\|_r \leq K_{V,p,r,s} \, \|f\|^{\mathrm{LP}}_{s,q} .
\]
\end{theorem}

\begin{proof}
\leanok
Since $\alpha(n, p, r) \geq 0$, $s > 0$, and $\|f\|^{\mathrm{LP}}_{s,q} < \infty$
(\Cref{thm:lp-band}). The blocks lie in $L_r(V)$ (\Cref{thm:nikolskii-blocks}), and the series of
their $L_r$ norms is at most $K_{V,p,r,s} \|f\|^{\mathrm{LP}}_{s,q} < \infty$
(\Cref{lem:nikolskii-blocks-summable}). So $\sum_j \psi_j * f$ converges absolutely in $L_r(V)$
to some $g$ with $\|g\|_r \leq K_{V,p,r,s} \|f\|^{\mathrm{LP}}_{s,q}$, and almost everywhere to
$g$. By \Cref{lem:blocks-ae-convergence} it converges almost everywhere to $f$, so $g = f$ almost
everywhere.
\end{proof}

\begin{theorem}[The Besov embedding into bounded continuous functions]
\label{thm:besov-embedding-bounded-continuous-rd}
\lean{Rd.exists_boundedContinuous_ae_eq_of_memApprox}
\leanok
\uses{lem:nikolskii-blocks-summable, thm:nikolskii-blocks, lem:bounded-kernel-conv-continuous,
  lem:blocks-ae-convergence, def:besov-embedding-const, thm:lp-band, def:nikolskii-property}
Let $1 \leq p < \infty$, $s > n/p$ and $1 \leq q \leq \infty$. For every $f \in \Besov{s}{p}{q}$
there is a bounded continuous function $a : V \to \C$ with $a = f$ almost everywhere and
\[
  \sup_{x \in V} |a(x)| \leq K_{V,p,\infty,s} \, \|f\|^{\mathrm{LP}}_{s,q} .
\]
\end{theorem}

\begin{proof}
\leanok
As before, $\|f\|^{\mathrm{LP}}_{s,q} < \infty$ (\Cref{thm:lp-band}). The blocks are continuous
(\Cref{lem:bounded-kernel-conv-continuous}) and bounded at every point by
$b_j = N_p 2^{jn/p} \|\psi_j * f\|_p$ (\Cref{thm:nikolskii-blocks}), and
$\sum_j b_j \leq K_{V,p,\infty,s} \|f\|^{\mathrm{LP}}_{s,q} < \infty$
(\Cref{lem:nikolskii-blocks-summable}). So $a(x) = \sum_j (\psi_j * f)(x)$ converges uniformly,
$a$ is continuous with $|a(x)| \leq \sum_j b_j$ for every $x$ (Weierstrass' test, Mathlib's
\texttt{continuous\_\allowbreak tsum}), so $\sup_{x \in V} |a(x)| \leq \sum_j b_j$, and $a = f$
almost everywhere by \Cref{lem:blocks-ae-convergence}.
\end{proof}

\paragraph{Sobolev embeddings.}

\begin{corollary}[Sobolev embeddings]
\label{cor:sobolev-embedding}
\lean{Rd.exists_memLp_ae_eq_of_hasWeakLineDeriv,
  Rd.exists_boundedContinuous_ae_eq_of_hasWeakLineDeriv, Rd.exists_memLp_ae_eq_of_memSobolev,
  Rd.exists_boundedContinuous_ae_eq_of_memSobolev}
\leanok
\uses{thm:sobolev-to-triebel-lizorkin, def:sobolev-to-tl-const, thm:tl-besov-sandwich,
  cor:sobolev-besov-sandwich, thm:besov-embedding-rd, thm:besov-embedding-bounded-continuous-rd,
  def:besov-embedding-const, def:nikolskii-property}
Let $n \geq 1$, $1 < p < \infty$ and $m \in \N$, and let $f \in L_p(V)$ have the weak
derivatives $g_\alpha = \partial^{e_\alpha} f \in L_p(V)$ of \Cref{thm:sobolev-to-triebel-lizorkin}
for every multi-index $\alpha$ of order $m$. Put
\[
  M_f = \|\Phi\|_1 \, \|f\|_p + 2^m \sum_\alpha K_{V,p,\alpha} \, \|g_\alpha\|_p .
\]
\begin{enumerate}
  \item For every $p \leq r \leq \infty$ with $m > n(1/p - 1/r)$ there is $u \in L_r(V)$ with
    $u = f$ almost everywhere and $\|u\|_r \leq K_{V,p,r,m} \, M_f$.
  \item If $m > n/p$, then there is a bounded continuous function $a : V \to \C$ with $a = f$
    almost everywhere and $\sup_{x \in V} |a(x)| \leq K_{V,p,\infty,m} \, M_f$.
\end{enumerate}
In particular, for every $p \leq r \leq \infty$ with $m > n(1/p - 1/r)$, every
$f \in \Sob{m}{p}(V)$ has a representative in $L_r(V)$, and if $m > n/p$, every
$f \in \Sob{m}{p}(V)$ has a bounded continuous representative.
\end{corollary}

\begin{proof}
\leanok
In both cases $m > 0$, so $m \geq 1$. By \Cref{thm:sobolev-to-triebel-lizorkin},
$f \in \TL{m}{p}{2}$ with $\|f\|_{\TL{m}{p}{2}} \leq M_f$, and \Cref{thm:tl-besov-sandwich} gives
$f \in \Besov{m}{p}{\max(p,2)}$ with $\|f\|^{\mathrm{LP}}_{m,\max(p,2)} \leq M_f$.
\Cref{thm:besov-embedding-rd} and \Cref{thm:besov-embedding-bounded-continuous-rd} with $s = m$
and $q = \max(p, 2)$ give the claims. For $f \in \Sob{m}{p}(V)$,
$f \in \Besov{m}{p}{\max(p,2)}$ by \Cref{cor:sobolev-besov-sandwich}.
\end{proof}

For $r < \infty$ the condition $m > n(1/p - 1/r)$ says $1/r > 1/p - m/n$. The critical case
$1/r = 1/p - m/n$ with $m < n/p$, Sobolev's inequality in the inhomogeneous form
$\|f\|_{r} \lesssim \|f\|_{\Sob{m}{p}}$, is \Cref{thm:critical-sobolev-embedding}; the homogeneous
form for $m = 1$ is the Gagliardo--Nirenberg--Sobolev inequality
(\Cref{sec:gagliardo-nirenberg-sobolev}).

\section{The critical Sobolev embedding}
\label{sec:critical-sobolev-embedding}

Let $n \geq 1$, $1 \leq p_0 < p_1 < \infty$, $\theta = p_0/p_1$ and
$\alpha = \alpha(n, p_0, p_1) = n(1/p_0 - 1/p_1)$. The embeddings of
\Cref{sec:sobolev-embedding} need $s > \alpha$; this section proves the critical case
$\TL{\alpha}{p_0}{\infty} \hookrightarrow L_{p_1}(V)$ and with it Sobolev's inequality
$\Sob{m}{p}(V) \hookrightarrow L_r(V)$ for $1 < p < \infty$ and $1/r = 1/p - m/n > 0$. The proof is
the direct one of Johnsen and Sickel~\cite[Proposition~1]{JohnsenSickel2007}. At a point $x$ the
Nikolskii inequality for blocks bounds $\sup_j 2^{j(s_0 - n/p_0)} |(\psi_j * f)(x)|$ by the
Littlewood--Paley quasi-norm of $\Besov{s_0}{p_0}{\infty}$, an interpolation inequality for
weighted suprema of sequences, due to Brezis and Mironescu~\cite{BrezisMironescu2001}
(\Cref{lem:dyadic-interpolation}), combines this bound with
$G(x) = \sup_j 2^{j s_0} |(\psi_j * f)(x)|$, and integration turns $G^\theta$ into
$\|G\|_{p_0}^\theta = \|f\|_{F^{s_0}_{p_0,\infty}}^\theta$. The same argument gives Jawerth's
embedding $\TL{s_0}{p_0}{\infty} \hookrightarrow \TL{s_1}{p_1}{q}$ for
$s_0 - s_1 = \alpha$ and every $q$ (Jawerth~\cite{Jawerth1977}, Triebel~\cite[\S 2.7.1]{Triebel},
Vyb\'{\i}ral~\cite[Theorem~0.1~(ii)]{VybiralJF2008}). The classical proofs of Sobolev's inequality
go through Riesz potentials and the Hardy--Littlewood--Sobolev inequality, for instance by
Hedberg's inequality~\cite{Hedberg1972}.

\paragraph{The pointwise bound and its integral.}

\begin{lemma}[The Nikolskii inequality in supremum form]
\label{lem:nikolskii-sup}
\lean{Rd.iSup_dyadicWeight_enorm_blockFun_le}
\leanok
\uses{thm:nikolskii-blocks, def:nikolskii-block-const, def:lp-blocks, def:triebel-lizorkin}
For every $1 \leq p < \infty$, every $s \in \R$, every $f \in L_p(V)$ and every $x \in V$,
\[
  \sup_{j \in \N} 2^{j(s - n/p)} |(\psi_j * f)(x)| \leq N_p \, \|f\|^{\mathrm{LP}}_{s,\infty} .
\]
\end{lemma}

\begin{proof}
\leanok
By \Cref{thm:nikolskii-blocks},
$2^{j(s - n/p)} |(\psi_j * f)(x)| \leq N_p \, 2^{js} \|\psi_j * f\|_p
\leq N_p \|f\|^{\mathrm{LP}}_{s,\infty}$ for every $j$.
\end{proof}

\begin{lemma}[The pointwise bound]
\label{lem:critical-pointwise-bound}
\lean{ENNReal.toReal_div_toReal_mem_Ioo, nikolskiiExponent_interp_eq,
  Rd.tsum_dyadicWeight_enorm_blockFun_le}
\leanok
\uses{lem:dyadic-interpolation, lem:nikolskii-sup, def:nikolskii-property}
Let $n \geq 1$, $1 \leq p_0 < p_1 < \infty$, $\theta = p_0/p_1$ and $s_0, s_1 \in \R$ with
$s_0 - s_1 = n(1/p_0 - 1/p_1)$. Then $0 < \theta < 1$, and for every $f \in L_{p_0}(V)$ and every
$x \in V$,
\[
  \sum_{j \in \N} 2^{j s_1} |(\psi_j * f)(x)|
  \leq c(n/p_0, \theta) \, G(x)^\theta \,
    \bigl(N_{p_0} \|f\|^{\mathrm{LP}}_{s_0,\infty}\bigr)^{1-\theta} ,
\]
where $G(x) = \sup_{j \in \N} 2^{j s_0} |(\psi_j * f)(x)|$.
\end{lemma}

\begin{proof}
\leanok
Apply \Cref{lem:dyadic-interpolation} with $\sigma_0 = s_0$, $\sigma_1 = s_0 - n/p_0$ and
$a_j = |(\psi_j * f)(x)|$: $\sigma_0 - \sigma_1 = n/p_0 > 0$, and
$\theta s_0 + (1 - \theta)(s_0 - n/p_0) = s_0 - n/p_0 + n/p_1 = s_1$, since $\theta/p_0 = 1/p_1$.
\Cref{lem:nikolskii-sup} bounds $\sup_j 2^{j\sigma_1} a_j$.
\end{proof}

\begin{definition}[The constant of the critical embedding]
\label{def:critical-embedding-const}
\lean{Rd.criticalEmbeddingConst, Rd.criticalEmbeddingConst_ne_top}
\leanok
\uses{lem:dyadic-interpolation, def:nikolskii-block-const}
For $1 \leq p_0 \leq p_1 \leq \infty$ and $\theta = p_0/p_1$,
\[
  K_{V,p_0,p_1} = c(n/p_0, \theta) \, N_{p_0}^{1-\theta} ,
\]
which is finite for $n \geq 1$ and $p_0 < p_1 < \infty$.
\end{definition}

\begin{lemma}[The integrated bound]
\label{lem:critical-integrated-bound}
\lean{ENNReal.mul_ofReal_toReal_div_toReal, Rd.eLpNorm_rpow_lqNorm_blockFun_eq,
  Rd.eLpNorm_tsum_dyadicWeight_enorm_blockFun_le}
\leanok
\uses{lem:critical-pointwise-bound, def:critical-embedding-const, lem:lp-norm-homogeneity,
  thm:tl-besov-sandwich, def:triebel-lizorkin}
Let $n \geq 1$, $1 \leq p_0 < p_1 < \infty$ and $s_0, s_1 \in \R$ with
$s_0 - s_1 = n(1/p_0 - 1/p_1)$. For every $f \in L_{p_0}(V)$,
\[
  \Bigl\| x \mapsto \sum_{j \in \N} 2^{j s_1} |(\psi_j * f)(x)| \Bigr\|_{p_1}
  \leq K_{V,p_0,p_1} \, \|f\|_{F^{s_0}_{p_0,\infty}} .
\]
\end{lemma}

\begin{proof}
\leanok
With $A = \|f\|^{\mathrm{LP}}_{s_0,\infty}$ and $T = \|f\|_{F^{s_0}_{p_0,\infty}} = \|G\|_{p_0}$,
\Cref{lem:critical-pointwise-bound} and \Cref{lem:lp-norm-homogeneity} bound the left side by
$c(n/p_0, \theta) (N_{p_0} A)^{1-\theta} \|G^\theta\|_{p_1}$, and
$\|G^\theta\|_{p_1} = \|G\|_{p_1 \theta}^\theta = T^\theta$ since $p_1\theta = p_0$. Finally
$A \leq T$ (\Cref{thm:tl-besov-sandwich} with $q = \infty$), so
$A^{1-\theta} T^\theta \leq T$.
\end{proof}

\paragraph{The embeddings.}

\begin{theorem}[$\TL{\alpha}{p_0}{q} \hookrightarrow L_{p_1}$]
\label{thm:triebel-lizorkin-into-lp}
\lean{Rd.exists_memLp_ae_eq_of_triebelLizorkinNorm_lt_top,
  Rd.exists_memLp_ae_eq_of_memTriebelLizorkin}
\leanok
\uses{lem:critical-integrated-bound, lem:blocks-ae-convergence, thm:tl-besov-sandwich,
  prop:tl-elementary, def:nikolskii-property}
Let $n \geq 1$, $1 \leq p_0 < p_1 < \infty$ and $\alpha = n(1/p_0 - 1/p_1)$. For every
$1 \leq q \leq \infty$ and every $f \in \TL{\alpha}{p_0}{q}$ there is $g \in L_{p_1}(V)$ with
$g = f$ almost everywhere and
\[
  \|g\|_{p_1} \leq K_{V,p_0,p_1} \, \|f\|_{F^{\alpha}_{p_0,\infty}}
  \leq K_{V,p_0,p_1} \, \|f\|_{F^{\alpha}_{p_0,q}} .
\]
\end{theorem}

\begin{proof}
\leanok
The second inequality is \Cref{prop:tl-elementary}. Since $\alpha > 0$ and
$\|f\|^{\mathrm{LP}}_{\alpha,\infty} \leq \|f\|_{F^\alpha_{p_0,\infty}} < \infty$
(\Cref{thm:tl-besov-sandwich}), $\sum_j (\psi_j * f)(x) = f(x)$ for almost every $x$
(\Cref{lem:blocks-ae-convergence}), so $|f| \leq \sum_j |\psi_j * f|$ almost everywhere.
\Cref{lem:critical-integrated-bound} with $s_0 = \alpha$ and $s_1 = 0$ bounds the $L_{p_1}$ norm
of the right side, so $f \in L_{p_1}(V)$ with the stated bound.
\end{proof}

For $p_0 = 1$ this is the embedding $\TL{n(1 - 1/p_1)}{1}{\infty} \hookrightarrow L_{p_1}(V)$.

\begin{corollary}[$\TL{s}{p_0}{q} \hookrightarrow L_{p_1}$ for $p_0 \leq p_1$]
\label{cor:triebel-lizorkin-into-lp-le}
\lean{Rd.exists_memLp_ae_eq_of_memTriebelLizorkin_of_le}
\leanok
\uses{thm:triebel-lizorkin-into-lp, prop:tl-elementary}
Let $n \geq 1$, $1 \leq p_0 \leq p_1 < \infty$, $1 \leq q \leq \infty$ and
$n(1/p_0 - 1/p_1) \leq s$. Every $f \in \TL{s}{p_0}{q}$ has a representative $g \in L_{p_1}(V)$,
$g = f$ almost everywhere.
\end{corollary}

\begin{proof}
\leanok
For $p_1 = p_0$ take $g = f$. For $p_0 < p_1$, $f \in \TL{\alpha}{p_0}{\infty}$ with
$\alpha = n(1/p_0 - 1/p_1)$ (\Cref{prop:tl-elementary}), and \Cref{thm:triebel-lizorkin-into-lp}
applies.
\end{proof}

\begin{theorem}[Jawerth's embedding]
\label{thm:jawerth-embedding}
\lean{Rd.exists_triebelLizorkinNorm_le_of_nikolskiiExponent}
\leanok
\uses{thm:triebel-lizorkin-into-lp, lem:critical-integrated-bound, prop:tl-elementary,
  lem:tl-pointwise-blocks, lem:lq-basic}
Let $n \geq 1$, $1 \leq p_0 < p_1 < \infty$, $s_1 \geq 0$ and $s_0 - s_1 = n(1/p_0 - 1/p_1)$.
For every $f \in \TL{s_0}{p_0}{\infty}$ there is $g \in L_{p_1}(V)$ with $g = f$ almost
everywhere such that for every $1 \leq q \leq \infty$,
\[
  \|g\|_{F^{s_1}_{p_1,q}} \leq K_{V,p_0,p_1} \, \|f\|_{F^{s_0}_{p_0,\infty}} .
\]
\end{theorem}

\begin{proof}
\leanok
Since $s_0 - s_1 \leq s_0$, $f \in \TL{s_0 - s_1}{p_0}{\infty}$ (\Cref{prop:tl-elementary}), and
\Cref{thm:triebel-lizorkin-into-lp} gives $g$. The blocks of $g$ are those of $f$
(\Cref{lem:tl-pointwise-blocks}), and
$\|\cdot\|_{\ell_q} \leq \|\cdot\|_{\ell_1}$ (\Cref{lem:lq-basic}) bounds
$\|g\|_{F^{s_1}_{p_1,q}}$ by the left side of \Cref{lem:critical-integrated-bound}.
\end{proof}

The companion embeddings with Besov spaces on one side,
$\TL{s_0}{p_0}{q} \hookrightarrow \Besov{s_1}{p_1}{p_0}$ (Jawerth~\cite{Jawerth1977}) and
$\Besov{s_0}{p_0}{p_1} \hookrightarrow \TL{s_1}{p_1}{q}$ (Franke~\cite{Franke1986}), rest on
decreasing rearrangements; Vyb\'{\i}ral~\cite{VybiralJF2008} gives a short proof of both.

\begin{theorem}[The critical Sobolev embedding]
\label{thm:critical-sobolev-embedding}
\lean{Rd.exists_lp_zero_ae_eq, Rd.exists_memLp_ae_eq_of_hasWeakLineDeriv_of_eq,
  Rd.exists_memLp_ae_eq_of_memSobolev_of_le, Rd.exists_memLp_ae_eq_of_memSobolev_of_eq}
\leanok
\uses{thm:triebel-lizorkin-into-lp, thm:sobolev-to-triebel-lizorkin, def:sobolev-to-tl-const,
  def:critical-embedding-const, prop:tl-elementary, def:nikolskii-property,
  cor:triebel-lizorkin-into-lp-le}
Let $n \geq 1$ and $1 < p < \infty$.
\begin{enumerate}
  \item For every $m \in \N$ with $m \geq 1$, every $r < \infty$ with $m = n(1/p - 1/r)$, that is
    $1/r = 1/p - m/n$, and every $f \in L_p(V)$ with weak derivatives
    $g_\alpha = \partial^{e_\alpha} f \in L_p(V)$ for every multi-index $\alpha$ of order $m$, as
    in \Cref{thm:sobolev-to-triebel-lizorkin}, there is $u \in L_r(V)$ with $u = f$ almost
    everywhere and
    \[
      \|u\|_r \leq K_{V,p,r} \Bigl( \|\Phi\|_1 \, \|f\|_p
        + 2^m \sum_\alpha K_{V,p,\alpha} \, \|g_\alpha\|_p \Bigr) .
    \]
  \item For every $m \in \N$, every $p \leq r < \infty$ with $n(1/p - 1/r) \leq m$ and every
    $f \in \Sob{m}{p}(V)$ there is $g \in L_r(V)$ with $g = f$ almost everywhere. In particular,
    this holds for every $r < \infty$ with $n(1/p - 1/r) = m$.
\end{enumerate}
\end{theorem}

\begin{proof}
\leanok
(1) Since $n(1/p - 1/r) = m > 0$, $p < r$. By \Cref{prop:tl-elementary} and
\Cref{thm:sobolev-to-triebel-lizorkin}, $\|f\|_{F^m_{p,\infty}} \leq \|f\|_{F^m_{p,2}}$ is at
most the bracket, and \Cref{thm:triebel-lizorkin-into-lp} with $p_0 = p$, $p_1 = r$ gives $u$.

(2) For $r = p$ take $g = f$. For $p < r$,
$f \in \TL{m}{p}{2} \subseteq \TL{\alpha(n, p, r)}{p}{\infty}$
(\Cref{thm:sobolev-to-triebel-lizorkin}, \Cref{prop:tl-elementary}), and
\Cref{thm:triebel-lizorkin-into-lp} applies. If $n(1/p - 1/r) = m \geq 0$, then $p \leq r$
(\Cref{def:nikolskii-property}).
\end{proof}

In Lean, $1/0 = 0$, so the exponent $r = 0$ satisfies $n(1/p - 1/r) = n/p$; there every
$f \in L_p(V)$ is a class of $L_0(V)$ of norm $0$. For $p = 1$ the space $\Sob{m}{1}(V)$ is larger
than $\TL{m}{1}{2}$, and the argument above starts from $\TL{m}{p}{2}$;
\Cref{sec:gagliardo-nirenberg-sobolev} treats $p = 1$.

\section{The Gagliardo--Nirenberg--Sobolev inequality}
\label{sec:gagliardo-nirenberg-sobolev}

Let $n \geq 1$, let $e_0, \dots, e_{n-1}$ be the orthonormal basis of \Cref{def:dyadic-cells},
$1 \leq p < \infty$ and $1/p' = 1/p - 1/n$. The case $p = 1$ of the Gagliardo--Nirenberg--Sobolev
inequality $\|f\|_{p'} \leq C \sum_k \|\partial_{e_k} f\|_p$ is due to
Gagliardo~\cite{Gagliardo1958} and, independently, Nirenberg~\cite{Nirenberg1959}; applied to
$|u|^\gamma$ it gives $1 < p < n$. Mathlib proves the inequality for compactly supported $C^1$
functions: for every $1 \leq p < \infty$ there is a constant $C_{V,p} < \infty$
(\texttt{MeasureTheory.\allowbreak eLpNormLESNormFDerivOfEqInnerConst}) such that
$\|u\|_{p'} \leq C_{V,p} \|Du\|_p$ for every $p' \geq 0$ with $1/p' = 1/p - 1/n$, every real
inner product space $H$ and every $C^1$ function $u \colon V \to H$ with compact support
(\texttt{MeasureTheory.\allowbreak eLpNorm\_\allowbreak le\_\allowbreak eLpNorm\_\allowbreak
fderiv\_\allowbreak of\_\allowbreak eq\_\allowbreak inner}), by the grid-lines argument of
Gagliardo and Nirenberg for $p = 1$ and by applying it to $|u|^\gamma$ for $p > 1$. This section
carries the inequality to every $f \in L_p(V)$ with weak derivatives $\partial_{e_k} f$ by lower
semicontinuity: for $h$ with compact support the mollifications $\eta_c * h$ satisfy the
inequality and converge to $h$ in measure, which passes the bound to $h$; the cut-offs
$\chi_R f$ have compact support and converge to $f$ pointwise, and Fatou's lemma passes their
bounds to $f$. The result is the Sobolev embedding $\Sob{1}{p}(V) \hookrightarrow L_r(V)$ for
$1/r = 1/p - 1/n$ and every $1 \leq p$, $p = 1$ included, where
\Cref{sec:critical-sobolev-embedding} needs $1 < p$, and with it
$\Sob{m}{1}(V) \hookrightarrow L_r(V)$ for $1/r = 1 - m/n$. In Lean $1/0 = 0$, so the exponent
$p' = 0$ satisfies $1/p' = 1/p - 1/n$ for $p = n$; there the left side is the norm of exponent
$0$, which is $0$.

\paragraph{Fatou's lemma and bases.}

\begin{lemma}[Fatou's lemma with bounds]
\label{lem:fatou-bounded-norms}
\lean{MeasureTheory.eLpNorm_le_of_ae_tendsto_of_tendsto}
\leanok
For every measure space, every $0 \leq p \leq \infty$, all almost everywhere strongly measurable
$u_i$ ($i \in \N$), every $f$ with $u_i \to f$ almost everywhere, and all
$B_i, b \in [0, \infty]$ with $\|u_i\|_p \leq B_i$ for every $i$ and $B_i \to b$,
$\|f\|_p \leq b$.
\end{lemma}

\begin{proof}
\leanok
As an almost everywhere limit, $f$ is almost everywhere strongly measurable (Mathlib's
\texttt{aestronglyMeasurable\_\allowbreak of\_\allowbreak tendsto\_\allowbreak ae}). By Fatou's
lemma (Mathlib's \texttt{MeasureTheory.\allowbreak Lp.\allowbreak
eLpNorm\_\allowbreak lim\_\allowbreak le\_\allowbreak liminf\_\allowbreak eLpNorm}),
$\|f\|_p \leq \liminf_i \|u_i\|_p \leq \liminf_i B_i = b$.
\end{proof}

Mathlib's \texttt{MeasureTheory.\allowbreak Lp.\allowbreak eLpNorm\_\allowbreak le\_\allowbreak
of\_\allowbreak ae\_\allowbreak tendsto} is the case of a constant bound $B_i = b$.

\begin{lemma}[Derivatives through a basis]
\label{lem:derivative-basis-bound}
\lean{ContinuousLinearMap.opNorm_le_sum_norm_apply_orthonormalBasis,
  Rd.norm_le_sum_norm_apply_gridBasis, Rd.eLpNorm_fderiv_le_sum}
\leanok
\uses{def:dyadic-cells}
For every real inner product space $E$ with an orthonormal basis $(b_i)_{i \in I}$, $I$ finite,
every real normed space $F$ and every continuous linear $L \colon E \to F$,
$\|L\| \leq \sum_i \|L b_i\|$; in particular $\|L\| \leq \sum_k \|L e_k\|$ for every continuous
linear $L \colon V \to F$. For every $C^1$ function $u \colon V \to \C$ and every
$1 \leq p \leq \infty$, $\|Du\|_p \leq \sum_k \|D_{e_k} u\|_p$.
\end{lemma}

\begin{proof}
\leanok
$x = \sum_i \langle b_i, x \rangle b_i$ and $|\langle b_i, x \rangle| \leq \|x\|$ give
$\|L x\| \leq \|x\| \sum_i \|L b_i\|$. Apply this to $L = Du(x)$ and use the triangle
inequality in $L_p$.
\end{proof}

\paragraph{The inequality.}

\begin{lemma}[The inequality for a mollification]
\label{lem:gns-mollification}
\lean{Rd.eLpNorm_conv_le_of_hasWeakLineDeriv}
\leanok
\uses{lem:mollification-weak-derivative, lem:derivative-basis-bound, lem:conv-derivatives,
  thm:young, lem:weak-derivative-basic}
Let $n \geq 1$. For every $1 \leq p < \infty$, every $p' \geq 0$ with $1/p' = 1/p - 1/n$, every
smooth $K \colon V \to \C$ with compact support, and all $h$ with compact support and $g_k$ with
$g_k = \partial_{e_k} h$ weakly on $V$ for every $k$,
\[
  \|K * h\|_{p'} \leq C_{V,p} \, \|K\|_1 \sum_k \|g_k\|_p .
\]
\end{lemma}

\begin{proof}
\leanok
The function $h$ is locally integrable (\Cref{lem:weak-derivative-basic}), so $K * h$ is smooth
with compact support (\Cref{lem:conv-derivatives}). By Mathlib's inequality,
\Cref{lem:derivative-basis-bound} and \Cref{lem:mollification-weak-derivative},
$\|K * h\|_{p'} \leq C_{V,p} \sum_k \|K * g_k\|_p$, and Young's inequality (\Cref{thm:young})
gives $\|K * g_k\|_p \leq \|K\|_1 \|g_k\|_p$.
\end{proof}

\begin{lemma}[The inequality with compact support]
\label{lem:gns-compact-support}
\lean{Rd.eLpNorm_le_of_hasCompactSupport_of_hasWeakLineDeriv}
\leanok
\uses{lem:gns-mollification, thm:smooth-atomic-decomposition, lem:dilate, thm:approx-identity,
  lem:weak-derivative-basic, lem:conv-derivatives, lem:dual-kernel-derivatives}
Let $n \geq 1$. For every $1 \leq p < \infty$, every $p' \geq 0$ with $1/p' = 1/p - 1/n$, and
all $h$ with compact support and $g_k$ with $g_k = \partial_{e_k} h$ weakly on $V$ for every $k$,
\[
  \|h\|_{p'} \leq C_{V,p} \sum_k \|g_k\|_p .
\]
\end{lemma}

\begin{proof}
\leanok
The function $h$ is locally integrable (\Cref{lem:weak-derivative-basic}) with compact support,
hence integrable. Let $\eta$ be the standard bump of \Cref{thm:smooth-atomic-decomposition},
with $\|\eta\|_1 = 1$, and $c_i = i + 1$. By \Cref{lem:gns-mollification} and
\Cref{lem:dilate}, $\|\eta_{c_i} * h\|_{p'} \leq C_{V,p} \sum_k \|g_k\|_p$ for every $i$. The
functions $\eta_{c_i} * h$ are smooth (\Cref{lem:dual-kernel-derivatives} and
\Cref{lem:conv-derivatives}) and converge to $h$ in $L_1(V)$ (\Cref{thm:approx-identity}), hence
in measure. A bound of the $L_{p'}$ norms along a sequence that converges in measure holds for
its limit (Mathlib's \texttt{MeasureTheory.\allowbreak eLpNorm\_\allowbreak le\_\allowbreak
of\_\allowbreak tendstoInMeasure}).
\end{proof}

\begin{definition}[The cut-off]
\label{def:real-cutoff}
\lean{Rd.cutoffBump, Rd.realCutoff, Rd.ofReal_realCutoff}
\leanok
\uses{def:lp-kernels}
Let $b = \chi_{1,2} \colon V \to [0, 1]$ be the smooth bump of \Cref{def:lp-kernels}, equal to
$1$ on $\bar B(0, 1)$ and to $0$ outside $B(0, 2)$. For every $R > 0$, $\chi_R(x) = b(x / R)$.
\end{definition}

In Lean $\chi_R(x) = b(R^{-1} x)$ for every real $R$, with $0^{-1} = 0$.

\begin{lemma}[Properties of the cut-off]
\label{lem:real-cutoff}
\lean{Rd.contDiff_realCutoff, Rd.hasCompactSupport_realCutoff, Rd.realCutoff_eq_one,
  Rd.abs_realCutoff_le_one, Rd.abs_fderiv_cutoffBump_le, Rd.abs_fderiv_realCutoff_le,
  Rd.eLpNorm_realCutoff_mul_deriv_le}
\leanok
\uses{def:real-cutoff, def:lp-kernels, def:dyadic-cells}
Let $M = p_{0,1}(\chi_{1,2}) = \sup_y \|D\chi_{1,2}(y)\|$ be the Schwartz seminorm of
\Cref{def:lp-kernels}. For every $R > 0$, $\chi_R$ is smooth with compact support and
$|\chi_R| \leq 1$, and $\chi_R = 1$ on $\bar B(0, R)$. For all $y, w \in V$,
$|Db(y) w| \leq M \|w\|$, and for every $R > 0$ and all $x, v \in V$,
\[
  |D\chi_R(x) v| \leq M R^{-1} \|v\| .
\]
For every $R > 0$, every $M' \in \R$ with $|D\chi_R(x) v| \leq M' R^{-1} \|v\|$ for all
$x, v \in V$, every $1 \leq p \leq \infty$, all almost everywhere strongly measurable $f$, $g$
and every $k$, $\|\chi_R g + (\partial_{e_k} \chi_R) f\|_p \leq \|g\|_p + M' R^{-1} \|f\|_p$.
\end{lemma}

\begin{proof}
\leanok
$|Db(y) w| = |D\chi_{1,2}(y) w| \leq \|D\chi_{1,2}(y)\| \|w\| \leq M \|w\|$, and
$D\chi_R(x) v = Db(x/R)(v/R)$. The last bound follows from $|\chi_R| \leq 1$, $\|e_k\| = 1$ and
the triangle inequality in $L_p$.
\end{proof}

\begin{theorem}[The Gagliardo--Nirenberg--Sobolev inequality on $\Sob{1}{p}$]
\label{thm:gns-weak}
\lean{Rd.eLpNorm_le_of_hasWeakLineDeriv}
\leanok
\uses{lem:gns-compact-support, lem:real-cutoff, lem:weak-derivative-product,
  lem:fatou-bounded-norms, lem:weak-derivative-basic, def:schwartz-cutoff}
Let $n \geq 1$. For every $1 \leq p < \infty$, every $p' \geq 0$ with $1/p' = 1/p - 1/n$, every
$f \in L_p(V)$ and all $g_k$ with $g_k = \partial_{e_k} f$ weakly on $V$ for every $k$,
\[
  \|f\|_{p'} \leq C_{V,p} \sum_k \|g_k\|_p .
\]
\end{theorem}

\begin{proof}
\leanok
Let $R_i = i + 1$ and $M$ as in \Cref{lem:real-cutoff}. The function $\chi_{R_i} f$ has compact
support (\Cref{def:schwartz-cutoff}) and weak derivatives
$\chi_{R_i} g_k + (\partial_{e_k} \chi_{R_i}) f$ (\Cref{lem:weak-derivative-product}), of $L_p$
norm at most $\|g_k\|_p + M R_i^{-1} \|f\|_p$ (\Cref{lem:real-cutoff}), since $g_k$ is locally
integrable (\Cref{lem:weak-derivative-basic}). By \Cref{lem:gns-compact-support},
$\|\chi_{R_i} f\|_{p'} \leq C_{V,p} \sum_k (\|g_k\|_p + M R_i^{-1} \|f\|_p)$, which tends to
$C_{V,p} \sum_k \|g_k\|_p$ since $\|f\|_p < \infty$. At every $x$, $\chi_{R_i}(x) f(x) = f(x)$
as soon as $R_i \geq \|x\|$, and \Cref{lem:fatou-bounded-norms} gives the bound.
\end{proof}

\paragraph{Sobolev embeddings.}

Inside the exponent $n(1/p - 1/r)$ of \Cref{def:nikolskii-property}, Lean reads
$1/0 = 1/\infty = 0$, and the statements below follow this convention.

\begin{lemma}[Exponents]
\label{lem:gns-exponents}
\lean{toReal_inv_eq_of_nikolskiiExponent_eq_one, nikolskiiExponent_eq_one_of_inv_add_inv,
  one_lt_of_nikolskiiExponent_one_eq, nikolskiiExponent_conjExponent_of_eq}
\leanok
\uses{def:nikolskii-property}
For every $D \geq 0$, every $n \in \N$ and all $p, r \in [0, \infty]$, with $1/0 = 1/\infty = 0$
in $D(1/p - 1/r)$ as in Lean:
\begin{enumerate}
  \item if $D(1/p - 1/r) = 1$, then $D > 0$, $p < \infty$ and $1/r = 1/p - 1/D$;
  \item if $1 \leq p$ and $1/r + 1/n = 1/p$ in $[0, \infty]$, where $1/0 = \infty$, then
    $p < \infty$ and $n(1/p - 1/r) = 1$;
  \item for every $m \geq 2$ with $n(1 - 1/r) = m$, $n > 1$, and the conjugate exponent
    $q = n/(n - 1)$ of $n$ satisfies $1 < q < \infty$, $n(1 - 1/q) = 1$ and
    $n(1/q - 1/r) = m - 1$.
\end{enumerate}
\end{lemma}

\begin{proof}
\leanok
(1) For $D = 0$ the left side is $0$, and for $p = \infty$ it is $-D/r \leq 0$. (2) Since
$1/p \leq 1$, both $1/r$ and $1/n$ are finite, so $n \geq 1$ and $r > 0$, and $p = \infty$ would
force $1/n = 0$. (3) $n - n/r = m \geq 2$ and $1/r \geq 0$ give $n > 1$, so $1 < q < \infty$.
Hölder conjugacy (Mathlib's \texttt{ENNReal.\allowbreak conjExponent}) gives $1/q + 1/n = 1$,
hence $n(1 - 1/q) = 1$ by (2) with $p = 1$, and $n(1/q - 1/r) = n(1 - 1/r) - 1 = m - 1$.
\end{proof}

\begin{theorem}[$\Sob{1}{p} \hookrightarrow L_r$ for $1/r = 1/p - 1/n$]
\label{thm:sobolev-one-embedding}
\lean{Rd.exists_memLp_ae_eq_of_hasWeakLineDeriv_one, Rd.exists_memLp_ae_eq_of_memSobolev_one,
  Rd.exists_memLp_ae_eq_of_memSobolev_one_of_inv_add_inv}
\leanok
\uses{thm:gns-weak, lem:gns-exponents, def:weak-line-derivative, lem:weak-derivative-basic}
\begin{enumerate}
  \item For every $1 \leq p \leq \infty$, every $r < \infty$ with $n(1/p - 1/r) = 1$, every
    $f \in L_p(V)$ and all $g_k \in L_p(V)$ with $g_k = \partial_{e_k} f$ weakly on $V$ for
    every $k$, there is $u \in L_r(V)$ with $u = f$ almost everywhere and
    $\|u\|_r \leq C_{V,p} \sum_k \|g_k\|_p$.
  \item For every $1 \leq p \leq \infty$, every $r < \infty$ with $n(1/p - 1/r) = 1$ and every
    $f \in \Sob{1}{p}(V)$ there are $g_k \in L_p(V)$ with $g_k = \partial_{e_k} f$ weakly on $V$
    for every $k$ and $u \in L_r(V)$ with $u = f$ almost everywhere and
    $\|u\|_r \leq C_{V,p} \sum_k \|g_k\|_p$.
  \item For every $1 \leq p \leq \infty$, every $r < \infty$ with $1/r + 1/n = 1/p$ in
    $[0, \infty]$, where $1/0 = \infty$, and every $f \in \Sob{1}{p}(V)$ the conclusion of (2)
    holds.
\end{enumerate}
\end{theorem}

\begin{proof}
\leanok
(1) By \Cref{lem:gns-exponents}, $n \geq 1$, $p < \infty$ and $1/r = 1/p - 1/n$, so
\Cref{thm:gns-weak} with $p' = r$ bounds $\|f\|_r$, and $f \in L_r(V)$. (2) The weak derivatives
$\partial_{e_k} f \in L_p(V)$ exist by \Cref{def:weak-line-derivative} and
\Cref{lem:weak-derivative-basic}, and (1) applies. (3) By \Cref{lem:gns-exponents},
$n(1/p - 1/r) = 1$, and (2) applies.
\end{proof}

The condition $n(1/p - 1/r) = 1$ forces $p < \infty$ and $1/r = 1/p - 1/n$
(\Cref{lem:gns-exponents}). For $p = 1$ the space $\Sob{1}{1}(V)$ is larger than $\TL{1}{1}{2}$,
so this case lies outside \Cref{thm:critical-sobolev-embedding}.

\begin{lemma}[The derivatives of $\Sob{m}{1}$ of lower order]
\label{lem:sobolev-l1-derivatives}
\lean{Rd.exists_hasWeakLineDeriv_memLp_of_memSobolev_of_lt}
\leanok
\uses{thm:sobolev-one-embedding, lem:weak-derivative-basic}
For every $m \in \N$, every $f \in \Sob{m}{1}(V)$, every $q < \infty$ with $n(1 - 1/q) = 1$,
every $k < m$ and every $v \in V^k$, there is $u \in L_q(V)$ with $u = \partial^v f$ weakly on
$V$.
\end{lemma}

\begin{proof}
\leanok
Let $G = \partial^v f \in L_1(V)$ and $H_j = \partial^{(e_j, v)} f \in L_1(V)$, which exist since
$k + 1 \leq m$. By \Cref{lem:weak-derivative-basic}, $H_j = \partial_{e_j} G$ weakly, and
\Cref{thm:sobolev-one-embedding} with $p = 1$ gives $u \in L_q(V)$ with $u = G$ almost
everywhere.
\end{proof}

\begin{theorem}[$\Sob{m}{1} \hookrightarrow L_r$ for $1/r = 1 - m/n$]
\label{thm:sobolev-l1-embedding}
\lean{Rd.exists_memLp_ae_eq_of_memSobolev_one_of_eq}
\leanok
\uses{thm:sobolev-one-embedding, lem:sobolev-l1-derivatives, lem:gns-exponents,
  thm:critical-sobolev-embedding, lem:weak-derivative-basic}
For every $m \geq 1$, every $r < \infty$ with $n(1 - 1/r) = m$, that is $1/r = 1 - m/n$, and
every $f \in \Sob{m}{1}(V)$ there is $g \in L_r(V)$ with $g = f$ almost everywhere.
\end{theorem}

\begin{proof}
\leanok
For $m = 1$ this is \Cref{thm:sobolev-one-embedding}. For $m \geq 2$ let $q = n/(n - 1)$
(\Cref{lem:gns-exponents}). By \Cref{lem:sobolev-l1-derivatives} with $k = 0$, $f$ agrees
almost everywhere with some $f' \in L_q(V)$, and every weak derivative of $f$ of order at most
$m - 1$ lies in $L_q(V)$, so $f' \in \Sob{m-1}{q}(V)$. Since $1 < q < \infty$ and
$n(1/q - 1/r) = m - 1$, \Cref{thm:critical-sobolev-embedding} gives $g \in L_r(V)$ with $g = f'$
almost everywhere.
\end{proof}

\section{Sobolev spaces and real interpolation}
\label{sec:sobolev-interpolation}

Let $n \geq 1$ and $1 < p < \infty$. This section gives $\Sob{m}{p}(V)$ a norm, the grid norm,
and places it between the Besov quasi-norms of $\Besov{m}{p}{\min(p,2)}$ and
$\Besov{m}{p}{\max(p,2)}$ with named constants (\Cref{thm:sobolev-grid-norm-sandwich}, from
\Cref{cor:sobolev-besov-sandwich}). The sandwich principle
(\Cref{cor:approximation-space-sandwich}) then identifies the real interpolation spaces of
$L_p(V)$ and $\Sob{m}{p}(V)$, and of two Sobolev spaces, with Besov spaces; this is
$(H^{s_0}_p, H^{s_1}_p)_{\theta,q} = \Besov{s}{p}{q}$ of~\cite[Theorem~6.4.5(4)]{BerghLofstrom}
for integer orders, with $H^m_p = \Sob{m}{p}$ from~\cite[Theorem~6.2.3]{BerghLofstrom}. As in
\Cref{sec:sobolev-triebel-lizorkin}, $e_\alpha$ is the tuple
$(e_{\alpha_0}, \dots, e_{\alpha_{m-1}})$ of basis vectors of a multi-index $\alpha$ of length
$m$ with entries below $n$, and $K_{V,p,\alpha}$ and $C_{V,p,w}$ are the constants of
\Cref{def:sobolev-to-tl-const} and \Cref{def:tl-to-sobolev-const}.

\paragraph{The grid norm.}

\begin{definition}[The grid norm of $\Sob{m}{p}$]
\label{def:sobolev-grid-norm}
\lean{Rd.sobolevGridNorm, Rd.sobolevGridNorm_le, Rd.sobolevGridNorm_eq}
\leanok
\uses{def:weak-line-derivative, lem:weak-derivative-basic}
For every $m \in \N$ and every $f \in L_p(V)$,
\[
  \|f\|^{\mathrm{gr}}_{m} = \|f\|_p + \inf \sum_\alpha \|g_\alpha\|_p ,
\]
the infimum over all families $(g_\alpha)_\alpha$ in $L_p(V)$, indexed by the multi-indices
$\alpha$ of length $m$, such that $g_\alpha$ is the weak derivative $\partial^{e_\alpha} f$ on
$V$ for every $\alpha$; the infimum of the empty set is $\infty$. In particular
$\|f\|^{\mathrm{gr}}_{m} \leq \|f\|_p + \sum_\alpha \|g_\alpha\|_p$ for every such family.
\end{definition}

If the family exists, it is unique (\Cref{lem:weak-derivative-basic}) and the grid norm equals
$\|f\|_p + \sum_\alpha \|\partial^{e_\alpha} f\|_p$.

\begin{lemma}[The grid norm is finite exactly on $\Sob{m}{p}$]
\label{lem:sobolev-grid-norm-finite}
\lean{Rd.exists_hasWeakLineDeriv_of_sobolevGridNorm_lt_top,
  Rd.sobolevGridNorm_lt_top_of_memSobolev, Rd.sobolevGridNorm_lt_top_iff}
\leanok
\uses{def:sobolev-grid-norm, lem:weak-derivative-basic, thm:sobolev-to-triebel-lizorkin,
  cor:sobolev-triebel-lizorkin}
Let $n \geq 1$ and $1 < p < \infty$. For every $m \in \N$ and every $f \in L_p(V)$,
$\|f\|^{\mathrm{gr}}_{m} < \infty$ if and only if $f \in \Sob{m}{p}(V)$.
\end{lemma}

\begin{proof}
\leanok
If the grid norm is finite, the infimum is taken over a nonempty set, so the weak derivatives
$\partial^{e_\alpha} f \in L_p(V)$ exist for all $\alpha$, and
\Cref{thm:sobolev-to-triebel-lizorkin}
and \Cref{cor:sobolev-triebel-lizorkin} give $f \in \TL{m}{p}{2} = \Sob{m}{p}(V)$. Conversely,
for $f \in \Sob{m}{p}(V)$ the weak derivatives along all $e_\alpha$ lie in $L_p(V)$
(\Cref{lem:weak-derivative-basic}), and the finite sum of their norms bounds the infimum.
\end{proof}

\paragraph{The grid norm between two Besov quasi-norms.}

\begin{definition}[The constants of the grid norm]
\label{def:sobolev-grid-consts}
\lean{Rd.approxSobolevGridConst, Rd.approxSobolevGridConst_ne_top,
  Rd.approxSobolevGridConst_ne_zero, Rd.sobolevGridApproxConst,
  Rd.sobolevGridApproxConst_ne_top, Rd.sobolevGridApproxConst_ne_zero}
\leanok
\uses{def:sobolev-to-tl-const, def:tl-to-sobolev-const, lem:lp-kernels-basic}
Let $1 < p < \infty$. For every $m \in \N$,
\[
  K'_{V,p,m} = 2 (1 - 2^{-m})^{-1} \Bigl( \|\Phi\|_1 + 2^m \sum_\alpha K_{V,p,\alpha} \Bigr)
\]
and
\[
  K''_{V,p,m} = 1 + \Bigl( \sum_\alpha C_{V,p,e_\alpha} \Bigr) \cdot 8 \cdot 8^m \, \|\Phi\|_1 ,
\]
the sums over the multi-indices $\alpha$ of length $m$. Both are nonzero, $K''_{V,p,m}$ is
finite, and $K'_{V,p,m}$ is finite for $m \geq 1$.
\end{definition}

\begin{theorem}[The grid norm between two Besov quasi-norms]
\label{thm:sobolev-grid-norm-sandwich}
\lean{Rd.approxNormDyadic_le_of_hasWeakLineDeriv_sum, Rd.approxNormDyadic_le_sobolevGridNorm,
  Rd.sobolevGridNorm_le_approxNormDyadic}
\leanok
\uses{def:sobolev-grid-norm, def:sobolev-grid-consts, cor:sobolev-besov-sandwich,
  lem:sum-of-products-le-product-of-sums, def:band-space, def:dyadic-norm}
Let $n \geq 1$ and $1 < p < \infty$. For every $m \geq 1$ and every $f \in L_p(V)$,
\[
  \|f\|^{\mathrm{dy}}_{A^m_{\max(p,2)}} \leq K'_{V,p,m} \, \|f\|^{\mathrm{gr}}_{m} ,
\]
and for every $m \in \N$ and every $f \in L_p(V)$,
\[
  \|f\|^{\mathrm{gr}}_{m} \leq K''_{V,p,m} \, \|f\|^{\mathrm{dy}}_{A^m_{\min(p,2)}} ,
\]
with the dyadic quasi-norms of the band-limited scheme.
\end{theorem}

\begin{proof}
\leanok
For one family $(g_\alpha)$ of weak derivatives, \Cref{cor:sobolev-besov-sandwich} bounds
$\|f\|^{\mathrm{dy}}_{A^m_{\max(p,2)}}$ by
$2 (1 - 2^{-m})^{-1} (\|\Phi\|_1 \|f\|_p + 2^m \sum_\alpha K_{V,p,\alpha} \|g_\alpha\|_p)$, and
\Cref{lem:sum-of-products-le-product-of-sums} bounds the bracket by
$(\|\Phi\|_1 + 2^m \sum_\alpha K_{V,p,\alpha})(\|f\|_p + \sum_\alpha \|g_\alpha\|_p)$. The
infimum over the families gives the first bound, since $K'_{V,p,m}$ is finite and nonzero. For
the second, if $f \in \Besov{m}{p}{\min(p,2)}$, \Cref{cor:sobolev-besov-sandwich} gives weak
derivatives $g_\alpha = \partial^{e_\alpha} f$ with
$\|g_\alpha\|_p \leq C_{V,p,e_\alpha} \cdot 8 \cdot 8^m \|\Phi\|_1 \,
\|f\|^{\mathrm{dy}}_{A^m_{\min(p,2)}}$,
and $\|f\|_p \leq \|f\|^{\mathrm{dy}}_{A^m_{\min(p,2)}}$; otherwise the right side is $\infty$.
\end{proof}

\paragraph{Real interpolation.}

\begin{definition}[Sobolev couples]
\label{def:sobolev-couples}
\lean{Rd.lpSobolevCouple, Rd.lpSobolevCouple_norm0, Rd.lpSobolevCouple_norm1,
  Rd.sobolevCouple, Rd.sobolevCouple_norm0, Rd.sobolevCouple_norm1}
\leanok
\uses{def:couple, def:sobolev-grid-norm}
For all $m, m_0, m_1 \in \N$, the couple $(L_p, \Sob{m}{p})$ on $L_p(V)$ is
$(\|\cdot\|_p, \|\cdot\|^{\mathrm{gr}}_{m})$, and the couple $(\Sob{m_0}{p}, \Sob{m_1}{p})$ is
$(\|\cdot\|^{\mathrm{gr}}_{m_0}, \|\cdot\|^{\mathrm{gr}}_{m_1})$.
\end{definition}

\begin{theorem}[Sobolev spaces interpolate to Besov spaces]
\label{thm:sobolev-interpolate}
\lean{Rd.memKMethod_lpSobolevCouple_iff, Rd.memKMethod_sobolevCouple_iff}
\leanok
\uses{def:sobolev-couples, thm:sobolev-grid-norm-sandwich, cor:approximation-space-sandwich,
  def:couple-k-method}
Let $n \geq 1$, $1 < p < \infty$, $1 \leq q \leq \infty$ and $f \in L_p(V)$.
\begin{enumerate}
  \item For every $m \geq 1$ and every $0 < \theta < 1$,
    $f \in (L_p, \Sob{m}{p})_{\theta, q}$ if and only if $f \in \Besov{\theta m}{p}{q}$.
  \item For all integers $1 \leq m_0 < m_1$ and every $0 < \eta < 1$,
    $f \in (\Sob{m_0}{p}, \Sob{m_1}{p})_{\eta, q}$ if and only if
    $f \in \Besov{(1 - \eta) m_0 + \eta m_1}{p}{q}$.
\end{enumerate}
\end{theorem}

\begin{proof}
\leanok
By \Cref{thm:sobolev-grid-norm-sandwich} the grid norm of order $m \geq 1$ lies between the
dyadic quasi-norms of $A^m_{\min(p,2)}$ and $A^m_{\max(p,2)}$ with finite constants, and
$1 \leq \min(p, 2) \leq \max(p, 2)$. Part (1) of \Cref{cor:approximation-space-sandwich} with
$r = m$ gives (1), and part (2) with $s_i = m_i$ gives (2).
\end{proof}

\begin{corollary}[Operators bounded on $L_p$ and on $\Sob{m}{p}$]
\label{cor:sobolev-interpolation-operators}
\lean{Rd.kMethodNorm_map_le_of_bounded_lp_sobolev, Rd.memApprox_map_of_bounded_lp_sobolev}
\leanok
\uses{thm:sobolev-interpolate, thm:interpolation-property}
Let $n \geq 1$, $1 < p < \infty$, $m \geq 1$, $0 < \theta < 1$ and $1 \leq q \leq \infty$, and
let $T : L_p(V) \to L_p(V)$ be
additive with $\|T f\|_p \leq M_0 \|f\|_p$ and
$\|T f\|^{\mathrm{gr}}_{m} \leq M_1 \|f\|^{\mathrm{gr}}_{m}$ for every $f \in L_p(V)$, with
$M_0, M_1 \in [0, \infty)$. For every $f \in \Besov{\theta m}{p}{q}$,
$T f \in \Besov{\theta m}{p}{q}$. If $M_0, M_1 > 0$, then for every $f \in L_p(V)$, in the
$K$-norm of $(L_p, \Sob{m}{p})$,
\[
  \|T f\|_{K, \theta, q} \leq M_0^{1 - \theta} M_1^{\theta} \, \|f\|_{K, \theta, q} .
\]
\end{corollary}

\begin{proof}
\leanok
\Cref{thm:interpolation-property} for the couple $(L_p, \Sob{m}{p})$ and
\Cref{thm:sobolev-interpolate}.
\end{proof}

\chapter{Sobolev spaces on the torus}
\label{ch:sobolev-torus}

This chapter defines weak derivatives in $L_p(\T^d)$ through Fourier coefficients, as
\Cref{def:weak-derivative} does for multi-indices. The rule
$\hat g(m) = \sigma_v(m) \hat f(m)$ for $g = \partial^v f$ is the distributional identity
$\int g \, \varphi = (-1)^k \int f \, \partial^v \varphi$ tested with the characters
$\varphi = e_{-m}$, whose span is dense in the smooth functions on $\T^d$.
\Cref{sec:torus-lp-fourier-coefficients} shows that an element of
$L_p(\T^d)$, $1 \leq p \leq \infty$, is determined by its Fourier coefficients, and computes the
coefficients of convolutions with kernels on $\R^d$ (\Cref{def:torus-econv}).
\Cref{sec:torus-lp-weak-derivatives} defines the derivatives $\partial^v$ along directions
$v \in (\R^d)^k$, proves their calculus, moves them through convolutions with Schwartz kernels,
and describes $\Sob{m}{p}(\T^d)$ by directional derivatives.
\Cref{sec:torus-periodic-distributions} identifies these weak derivatives with the distributional
derivatives of the periodic lifts on $\R^d$, and classical derivatives with weak derivatives.
\Cref{sec:torus-sobolev-triebel-lizorkin} proves $\Sob{m}{p}(\T^d) = \TL{m}{p}{2}(\T^d)$ for
$1 < p < \infty$, and \Cref{sec:torus-sobolev-embedding} derives the Sobolev embeddings on the
torus.

\section{Fourier coefficients in \texorpdfstring{$L_p$}{L\_p} on the torus}
\label{sec:torus-lp-fourier-coefficients}

On the probability space $\T^d$, $|\hat f(m)| \leq \|f\|_1 \leq \|f\|_p$ for every
$1 \leq p \leq \infty$, so every Fourier coefficient is a bounded linear functional on
$L_p(\T^d)$. Since the characters span a dense subspace of $L_p(\T^d)$ for $p < \infty$, identities
between such functionals are checked on the characters.

\begin{lemma}[Fourier coefficients in $L_p$]
\label{lem:torus-lp-coefficients}
\lean{Torus.mFourierCoeff_congr_ae, Torus.norm_mFourier_apply_le,
  Torus.enorm_mFourierCoeff_le_eLpNorm_one, Torus.integrable_mFourier_smul,
  Torus.mFourierCoeff_add, Torus.mFourierCoeff_smul, Torus.enorm_mFourierCoeff_le,
  Torus.norm_mFourierCoeff_le, Torus.mFourierCoeff_sub, Torus.mFourierCoeff_neg,
  Torus.mFourierCoeff_zero, Torus.mFourierCoeff_sum, Torus.mFourierCoeff_mFourier,
  Torus.mFourierCoeff_mFourierLp, Torus.mFourierCoeff_inclusionOfExponentLe}
\leanok
\uses{def:trig-poly, lem:inclusion}
Let $1 \leq p \leq \infty$. For all $f, g \in L_p(\T^d)$, every $c \in \C$ and all
$k, m \in \Z^d$:
\begin{enumerate}
  \item for all functions $f_1, f_2$ on $\T^d$ with $f_1 = f_2$ almost everywhere,
    $\hat f_1(m) = \hat f_2(m)$; for every function $f_1$ on $\T^d$,
    $|\hat f_1(m)| \leq \|f_1\|_1$; and $|e_k(t)| \leq 1$ for every $t \in \T^d$;
  \item $|\hat f(m)| \leq \|f\|_1 \leq \|f\|_p$;
  \item $\widehat{f + g}(m) = \hat f(m) + \hat g(m)$, $\widehat{c f}(m) = c \hat f(m)$,
    $\widehat{f - g}(m) = \hat f(m) - \hat g(m)$, $\widehat{-f}(m) = -\hat f(m)$,
    $\hat 0(m) = 0$, and the coefficients of a finite sum are the sums of the coefficients;
  \item $\hat e_k(m) = 1$ if $k = m$, and $\hat e_k(m) = 0$ otherwise.
\end{enumerate}
For every $r$ with $p \leq r \leq \infty$ and every $f \in L_r(\T^d)$, the class of $f$ in
$L_p(\T^d)$ has the Fourier coefficients of $f$.
\end{lemma}

\begin{proof}
\leanok
(1) The coefficient is an integral, and $|e_k| = 1$. (2) $|e_{-m} f| \leq |f|$, and
$\|f\|_1 \leq \|f\|_p$ on a probability space (Mathlib's
\texttt{eLpNorm\_\allowbreak le\_\allowbreak eLpNorm\_\allowbreak of\_\allowbreak
exponent\_\allowbreak le}).
(3) The integrands $e_{-m} f$ and $e_{-m} g$ are integrable by (2); the remaining identities
follow from linearity. (4) The characters are
orthonormal in $L_2(\T^d)$ (Mathlib's \texttt{mFourierBasis}). The inclusion
$L_r \hookrightarrow L_p$ keeps the function almost everywhere (\Cref{lem:inclusion}).
\end{proof}

\begin{definition}[The coefficient functional]
\label{def:torus-coefficient-functional}
\lean{Torus.coeffCLM, Torus.coeffCLM_apply, Torus.norm_coeffCLM_le}
\leanok
\uses{lem:torus-lp-coefficients}
For $1 \leq p \leq \infty$ and $m \in \Z^d$, the coefficient functional is the bounded linear
map $L_p(\T^d) \to \C$, $f \mapsto \hat f(m)$, of norm at most $1$.
\end{definition}

\begin{theorem}[Fourier coefficients of convolutions]
\label{thm:torus-coefficients-convolution}
\lean{Torus.mFourierCoeff_econvLp, Torus.mFourierCoeff_econv}
\leanok
\uses{def:torus-coefficient-functional, lem:torus-econv-young, thm:torus-econv-fourier,
  lem:torus-econv-identities}
For every $1 \leq p < \infty$, every $\Phi \in L_1(\R^d)$, every $f \in L_p(\T^d)$ and every
$m \in \Z^d$,
\[
  \widehat{\Phi * f}(m) = \mathcal{F}\Phi(m) \, \hat f(m) .
\]
The same holds for every measurable integrable function $f$ on $\T^d$ and the convolution
integral $\Phi * f$ of \Cref{def:torus-econv}.
\end{theorem}

\begin{proof}
\leanok
Both sides are bounded linear functionals of $f$ (\Cref{lem:torus-econv-young},
\Cref{def:torus-coefficient-functional}). On a character $e_k$,
$\Phi * e_k = \mathcal{F}\Phi(k) \, e_k$ (\Cref{thm:torus-econv-fourier}), so both sides are
$\mathcal{F}\Phi(m)$ for $k = m$ and $0$ otherwise. Bounded linear maps that agree on the
characters are equal (\Cref{lem:torus-econv-identities}). For an integrable $f$ take $p = 1$:
the convolution integral represents the operator (\Cref{lem:torus-econv-young}).
\end{proof}

\begin{theorem}[Uniqueness of Fourier coefficients]
\label{thm:torus-coefficients-uniqueness}
\lean{Torus.eq_sum_mFourierCoeff_of_mem_trigPolyLp, Torus.eq_zero_of_mem_trigPolyLp,
  Torus.eq_zero_of_mFourierCoeff_eq_zero_one, Torus.eq_zero_of_mFourierCoeff_eq_zero,
  Torus.eq_of_mFourierCoeff_eq}
\leanok
\uses{lem:torus-lp-coefficients, thm:torus-econv-fourier, thm:torus-coefficients-convolution,
  lem:torus-lp-blocks-trig, lem:torus-lp-blocks, lem:inclusion}
Let $1 \leq p \leq \infty$. For every $n \in \N$, every trigonometric polynomial
$a \in \mathcal{T}_n$ in $L_p(\T^d)$ satisfies
\[
  a = \sum_{\deg m < n} \hat a(m) \, e_m .
\]
For all $f, g \in L_p(\T^d)$ with $\hat f(m) = \hat g(m)$ for every $m \in \Z^d$, $f = g$.
\end{theorem}

\begin{proof}
\leanok
Both sides of the identity are linear in $a$ and agree on the characters $e_k$ with $\deg k < n$
(\Cref{lem:torus-lp-coefficients}), which span $\mathcal{T}_n$ (\Cref{thm:torus-econv-fourier}).
For the second claim, $h = f - g$ has vanishing coefficients. Let first $p = 1$. Every block
$\Delta_j h$ is a trigonometric polynomial (\Cref{lem:torus-lp-blocks-trig}) with the
coefficients $\mathcal{F}\psi_j(m) \, \hat h(m) = 0$ (\Cref{thm:torus-coefficients-convolution}),
hence $\Delta_j h = 0$ by the identity, and $h = \lim_J \sum_{j \leq J} \Delta_j h = 0$
(\Cref{lem:torus-lp-blocks}). For general $p$, the class of $h$ in $L_1(\T^d)$ has the same
coefficients (\Cref{lem:torus-lp-coefficients}), so it is $0$, and the inclusion
$L_p \hookrightarrow L_1$ is injective (\Cref{lem:inclusion}).
\end{proof}

\section{Weak derivatives in \texorpdfstring{$L_p$}{L\_p} on the torus}
\label{sec:torus-lp-weak-derivatives}

For directions $v = (v_0, \dots, v_{k-1}) \in (\R^d)^k$ the derivative
$\partial^v = \partial_{v_0} \cdots \partial_{v_{k-1}}$ of
$x \mapsto e^{2\pi i \langle m, x \rangle}$ is
$\prod_{l < k} 2\pi i \langle m, v_l \rangle \, e^{2\pi i \langle m, x \rangle}$. Directions
make the derivatives independent of a basis of $\R^d$; the coordinate directions give back the
multi-indices of \Cref{def:weak-derivative}.

\begin{definition}[Symbols of directional derivatives]
\label{def:torus-line-symbol}
\lean{Torus.lineSymbol, Torus.coordIndexEquiv, Torus.coordDirs, Torus.coordCount}
\leanok
\uses{def:weak-derivative, lem:torus-proj}
For $k \in \N$, $v \in (\R^d)^k$ and $m \in \Z^d$, regarded as a point of $\R^d$, the symbol of
$\partial^v$ is
\[
  \sigma_v(m) = \prod_{l < k} 2\pi i \langle m, v_l \rangle .
\]
For $\alpha \in \N^d$, the coordinate directions $v^\alpha \in (\R^d)^{|\alpha|}$ list $\alpha_i$
copies of the unit vector $e_i$ for every $i \in d$, in a fixed order. For every map
$c \colon \{0, \dots, k-1\} \to d$, the multiplicities $\alpha^c \in \N^d$ are
$\alpha^c_i = \#\{l : c_l = i\}$.
\end{definition}

\begin{lemma}[Symbols]
\label{lem:torus-line-symbol}
\lean{Torus.lineSymbol_fin_zero, Torus.lineSymbol_neg, Torus.norm_lineSymbol_le,
  Torus.lineSymbol_append, Torus.lineSymbol_coordDirs, Torus.sum_coordCount,
  Torus.prod_eq_derivSymbol_coordCount, Torus.lineSymbol_eq_sum}
\leanok
\uses{def:torus-line-symbol, lem:torus-proj, lem:dual-derivative-frequency}
For all $k, j \in \N$, $v \in (\R^d)^k$, $w \in (\R^d)^j$ and $m \in \Z^d$:
\begin{enumerate}
  \item $\sigma_v(m) = 1$ for $k = 0$, and $\sigma_v(-m) = (-1)^k \sigma_v(m)$;
  \item $|\sigma_v(m)| \leq \prod_{l < k} 2\pi \, |m| \, |v_l|$;
  \item $\sigma_{(v, w)}(m) = \sigma_v(m) \, \sigma_w(m)$ for the concatenation
    $(v, w) \in (\R^d)^{k + j}$;
  \item $\sigma_{v^\alpha}(m) = (2\pi i m)^\alpha$ for every $\alpha \in \N^d$;
  \item for every $c \colon \{0, \dots, k-1\} \to d$, $|\alpha^c| = k$ and
    $\prod_{l < k} 2\pi i m_{c_l} = (2\pi i m)^{\alpha^c}$, and
    \[
      \sigma_v(m) = \sum_{c} \Bigl( \prod_{l < k} (v_l)_{c_l} \Bigr) (2\pi i m)^{\alpha^c} ,
    \]
    the sum over all maps $c \colon \{0, \dots, k-1\} \to d$.
\end{enumerate}
\end{lemma}

\begin{proof}
\leanok
(1) $\langle -m, v_l \rangle = -\langle m, v_l \rangle$ (\Cref{lem:torus-proj}). (2) The
Cauchy--Schwarz inequality, as in \Cref{lem:dual-derivative-frequency}. (3) The product over $k + j$
factors splits. (4) $\langle m, e_i \rangle = m_i$ (\Cref{lem:torus-proj}), and the product over
the slots of $v^\alpha$ groups by coordinates. (5) Group the factors by the value of $c_l$
(Mathlib's \texttt{Finset.\allowbreak prod\_\allowbreak fiberwise'}), and expand
$\prod_l \sum_i 2\pi i \, m_i (v_l)_i$ into a sum over the maps $c$ (Mathlib's
\texttt{Finset.\allowbreak prod\_\allowbreak univ\_\allowbreak sum}).
\end{proof}

\begin{definition}[Weak derivatives along directions on the torus]
\label{def:torus-weak-line-derivative}
\lean{Torus.HasWeakLineDeriv}
\leanok
\uses{def:torus-line-symbol}
For every $1 \leq p \leq \infty$, $k \in \N$ and $v \in (\R^d)^k$, an element $g \in L_p(\T^d)$
is the weak derivative $\partial^v f$ of $f \in L_p(\T^d)$ if $\hat g(m) = \sigma_v(m) \hat f(m)$
for every $m \in \Z^d$.
\end{definition}

Since $\partial^v e_{-m} = \sigma_v(-m) e_{-m} = (-1)^k \sigma_v(m) e_{-m}$
(\Cref{lem:torus-line-symbol}), the definition is the identity
$\int g \, e_{-m} = (-1)^k \int f \, \partial^v e_{-m}$ for every $m \in \Z^d$.

\begin{lemma}[Properties of weak derivatives on the torus]
\label{lem:torus-weak-line-derivative-basic}
\lean{Torus.hasWeakLineDeriv_fin_zero, Torus.hasWeakDeriv_iff_hasWeakLineDeriv,
  Torus.HasWeakLineDeriv.comp, Torus.HasWeakLineDeriv.unique, Torus.HasWeakDeriv.unique,
  Torus.HasWeakLineDeriv.add, Torus.HasWeakLineDeriv.smul, Torus.HasWeakLineDeriv.neg,
  Torus.HasWeakLineDeriv.sub, Torus.HasWeakLineDeriv.sum, Torus.HasWeakDeriv.add,
  Torus.HasWeakDeriv.smul, Torus.HasWeakLineDeriv.of_tendsto,
  Torus.hasWeakLineDeriv_mFourierLp, Torus.HasWeakLineDeriv.of_inclusion,
  Torus.HasWeakDeriv.of_inclusion}
\leanok
\uses{def:torus-weak-line-derivative, def:weak-derivative, lem:torus-line-symbol,
  lem:torus-lp-coefficients, def:torus-coefficient-functional,
  thm:torus-coefficients-uniqueness}
Let $1 \leq p \leq \infty$, $k \in \N$, $v \in (\R^d)^k$, $f, f', g, g' \in L_p(\T^d)$ and
$c \in \C$.
\begin{enumerate}
  \item For $k = 0$, $f = \partial^v f$. For every $\alpha \in \N^d$, $g = \partial^\alpha f$ if
    and only if $g = \partial^{v^\alpha} f$.
  \item For every $j \in \N$, $w \in (\R^d)^j$ and $h \in L_p(\T^d)$ with $g = \partial^v f$ and
    $h = \partial^w g$, $h = \partial^{(w, v)} f$.
  \item If $g$ and $g'$ are weak derivatives $\partial^v f$, then $g = g'$; the
    same holds for $\partial^\alpha$, $\alpha \in \N^d$.
  \item If $g = \partial^v f$ and $g' = \partial^v f'$, then $g + g' = \partial^v (f + f')$,
    $-g = \partial^v (-f)$, $g - g' = \partial^v (f - f')$ and $c g = \partial^v (c f)$; for a
    finite set $s$ and $F_i, G_i \in L_p(\T^d)$ with $G_i = \partial^v F_i$ for every $i \in s$,
    $\sum_{i \in s} G_i = \partial^v \sum_{i \in s} F_i$. Sums and complex multiples of weak
    derivatives $\partial^\alpha$, $\alpha \in \N^d$, are weak derivatives $\partial^\alpha$ in
    the same way.
  \item For every filter $\ell \neq \bot$ on a set $I$ and all $f_i, g_i \in L_p(\T^d)$,
    $i \in I$, with $f_i \to f$ and $g_i \to g$ in $L_p(\T^d)$ along $\ell$ and
    $g_i = \partial^v f_i$ for all $i$ in a set of $\ell$, $g = \partial^v f$.
  \item For every $n \in \Z^d$, $\sigma_v(n) \, e_n = \partial^v e_n$.
  \item For every $r$ with $p \leq r \leq \infty$ and all $f, g \in L_r(\T^d)$ with
    $g = \partial^v f$, the classes of $f$ and $g$ in $L_p(\T^d)$ satisfy $g = \partial^v f$;
    the same holds for $\partial^\alpha$, $\alpha \in \N^d$.
\end{enumerate}
\end{lemma}

\begin{proof}
\leanok
Every claim is an identity between Fourier coefficients. (1) and (2) are (1), (4) and (3) of
\Cref{lem:torus-line-symbol}. (3) Both derivatives have the coefficients $\sigma_v(m) \hat f(m)$,
and \Cref{thm:torus-coefficients-uniqueness} applies. (4) The coefficients are linear
(\Cref{lem:torus-lp-coefficients}, \Cref{def:torus-coefficient-functional}). (5) Every
coefficient is a continuous functional (\Cref{def:torus-coefficient-functional}), so
$\hat g(m) = \lim_i \hat g_i(m) = \lim_i \sigma_v(m) \hat f_i(m) = \sigma_v(m) \hat f(m)$, and
limits in $\C$ are unique. (6) $\hat e_n(m) = 1$ for $m = n$ and $0$ otherwise
(\Cref{lem:torus-lp-coefficients}). (7) The inclusion keeps the coefficients
(\Cref{lem:torus-lp-coefficients}).
\end{proof}

\begin{lemma}[Weak derivatives and convolution on the torus]
\label{lem:torus-weak-derivative-convolution}
\lean{Torus.econvLp_iteratedLineDerivOp, Torus.hasWeakLineDeriv_econvLp,
  Torus.HasWeakLineDeriv.econvLp}
\leanok
\uses{def:torus-weak-line-derivative, thm:torus-coefficients-convolution,
  thm:torus-coefficients-uniqueness, lem:fourier-iterated-derivative}
Let $1 \leq p < \infty$, $k \in \N$, $v \in (\R^d)^k$ and let $\eta$ be a complex Schwartz
function on $\R^d$. For every $f \in L_p(\T^d)$, $(\partial^v \eta) * f$ is the weak derivative
$\partial^v (\eta * f)$. For all $f, g \in L_p(\T^d)$ with $g = \partial^v f$,
\[
  (\partial^v \eta) * f = \eta * g ,
\]
and for every integrable kernel $\Phi$ on $\R^d$, $\Phi * g$ is the weak derivative
$\partial^v (\Phi * f)$.
\end{lemma}

\begin{proof}
\leanok
By \Cref{thm:torus-coefficients-convolution} and \Cref{lem:fourier-iterated-derivative}, the
coefficient of $(\partial^v \eta) * f$ at $m$ is
$\mathcal{F}(\partial^v \eta)(m) \, \hat f(m) = \sigma_v(m) \, \mathcal{F}\eta(m) \, \hat f(m)$.
This is $\sigma_v(m)$ times the coefficient of $\eta * f$, and for $g = \partial^v f$ it is
$\mathcal{F}\eta(m) \, \hat g(m)$, the coefficient of $\eta * g$. Apply
\Cref{thm:torus-coefficients-uniqueness}. The coefficient of $\Phi * g$ at $m$ is
$\mathcal{F}\Phi(m) \, \sigma_v(m) \hat f(m)$, which is $\sigma_v(m)$ times the coefficient of
$\Phi * f$.
\end{proof}

\begin{theorem}[Sobolev spaces on the torus by directional derivatives]
\label{thm:torus-sobolev-directions}
\lean{Torus.hasWeakLineDeriv_sum_coordCount, Torus.memSobolev_iff_exists_hasWeakLineDeriv,
  Torus.memSobolev_zero, Torus.MemSobolev.add, Torus.MemSobolev.smul,
  Torus.MemSobolev.mono_order, Torus.MemSobolev.of_inclusion}
\leanok
\uses{def:weak-derivative, def:torus-weak-line-derivative, lem:torus-line-symbol,
  lem:torus-weak-line-derivative-basic, lem:torus-lp-coefficients}
Let $1 \leq p \leq \infty$ and $m \in \N$. For every $f \in L_p(\T^d)$, $f \in \Sob{m}{p}(\T^d)$
if and only if for every $k \leq m$ and every $v \in (\R^d)^k$ some $g \in L_p(\T^d)$ is the
weak derivative $\partial^v f$. For all $k \in \N$, $v \in (\R^d)^k$ and
$g_c \in L_p(\T^d)$ with $g_c = \partial^{\alpha^c} f$ for every
$c \colon \{0, \dots, k-1\} \to d$,
\[
  \sum_c \Bigl( \prod_{l < k} (v_l)_{c_l} \Bigr) g_c = \partial^v f .
\]
Moreover $\Sob{0}{p}(\T^d) = L_p(\T^d)$, sums and complex
multiples of elements of $\Sob{m}{p}(\T^d)$ lie in $\Sob{m}{p}(\T^d)$,
$\Sob{m}{p}(\T^d) \subseteq \Sob{l}{p}(\T^d)$ for every $l \leq m$, and for every $r$ with
$p \leq r \leq \infty$ the class in $L_p(\T^d)$ of every $f \in \Sob{m}{r}(\T^d)$ lies in
$\Sob{m}{p}(\T^d)$.
\end{theorem}

\begin{proof}
\leanok
Let $f \in \Sob{m}{p}(\T^d)$, $k \leq m$ and $v \in (\R^d)^k$. For every map
$c \colon \{0, \dots, k-1\} \to d$ the multiplicities $\alpha^c$ have order $k$, so
$g_c = \partial^{\alpha^c} f$ exists, and $\sum_c (\prod_{l < k} (v_l)_{c_l}) \, g_c$ is
$\partial^v f$ by the expansion (5) of \Cref{lem:torus-line-symbol} and the linearity of the
coefficients (\Cref{lem:torus-lp-coefficients}). Conversely, $\partial^\alpha f$ is
$\partial^{v^\alpha} f$ (\Cref{lem:torus-weak-line-derivative-basic}). The only multi-index of
order $0$ is $\alpha = 0$, with $\partial^0 f = f$; the other claims follow from (4) and (7) of
\Cref{lem:torus-weak-line-derivative-basic}.
\end{proof}

\section{Periodic distributions}
\label{sec:torus-periodic-distributions}

A function $h$ on $\T^d$ lifts to the periodic function $h \circ \pi$ on $\R^d$, where
$\pi \colon \R^d \to \T^d$ is the quotient map (\Cref{lem:torus-proj}). Periodic distributions,
the distributions on $\R^d$ invariant under the translations by $\Z^d$, are the setting of the
function spaces on the torus in Schmeisser and Triebel~\cite{SchmeisserTriebel}; the lifts of
integrable functions on $\T^d$ are periodic distributions. Mathlib has distributions on open
subsets of $\R^d$; this section shows that the weak derivatives of
\Cref{def:torus-weak-line-derivative} are the weak derivatives of the lifts in the sense of
\Cref{def:weak-line-derivative}, for every $1 \leq p \leq \infty$, and that classical derivatives
of continuously differentiable lifts are weak derivatives. In this section a test function is a
real smooth function with compact support on $\R^d$.

\begin{lemma}[Classical derivatives are weak derivatives]
\label{lem:weak-derivative-classical}
\lean{Rd.fderiv_testSchwartz_apply, Rd.hasWeakLineDeriv_of_contDiff, Rd.hasWeakLineDeriv_const}
\leanok
\uses{def:weak-line-derivative, lem:weak-derivative-integral, def:test-schwartz,
  lem:directional-derivatives-commute}
Let $V$ be a finite-dimensional real inner product space and $\Omega \subseteq V$ open. For every
continuously differentiable $F \colon V \to \C$ and every $w \in V$, the function
$x \mapsto DF(x) w$ is the weak derivative $\partial_w F$ on $\Omega$. For every $w \in V$ and
every $c \in \C$, the weak derivative $\partial_w$ of the constant function $c$ on $\Omega$ is
$0$. For every real test
function $\varphi$ on $\Omega$ and all $x, w \in V$, the derivative of $\varphi$, regarded as a
complex function, at $x$ in the direction $w$ is $(\partial_w \varphi)(x)$.
\end{lemma}

\begin{proof}
\leanok
The second claim is \Cref{lem:directional-derivatives-commute} for the complex Schwartz function
$\varphi$ of \Cref{def:test-schwartz}. Let $\varphi$ be a test function on $\Omega$. Since
$\varphi$ has compact support and $F$ and $DF(\cdot) w$ are continuous, the products
$\varphi \, DF(\cdot) w$, $(\partial_w \varphi) F$ and $\varphi F$ are integrable, and integration
by parts (Mathlib's \texttt{integral\_\allowbreak mul\_\allowbreak fderiv\_\allowbreak
eq\_\allowbreak neg\_\allowbreak fderiv\_\allowbreak mul\_\allowbreak of\_\allowbreak integrable})
gives $\int \varphi \, DF(\cdot) w = -\int (\partial_w \varphi) F$. This is the identity of
\Cref{lem:weak-derivative-integral} for $k = 1$. A constant function has the classical
derivative $0$.
\end{proof}

\begin{definition}[The periodic pairing]
\label{def:torus-periodic-pairing}
\lean{Torus.periodicPairing}
\leanok
\uses{lem:torus-proj}
For every $1 \leq p \leq \infty$ and every test function $\varphi$ on $\R^d$, the periodic pairing
of $\varphi$ is the map $L_p(\T^d) \to \C$,
\[
  h \mapsto \langle \varphi, h \rangle = \int_{\R^d} \varphi(x) \, h(\pi x) \, dx .
\]
\end{definition}

The pairing evaluates the lift $h \circ \pi$, as a distribution on $\R^d$, at $\varphi$.

\begin{lemma}[The periodic pairing]
\label{lem:torus-periodic-pairing}
\lean{Torus.locallyIntegrable_comp_proj_Lp, Torus.integrable_mul_comp_proj,
  Torus.enorm_integral_mul_comp_proj_le, Torus.exists_norm_integral_mul_comp_proj_le,
  Torus.integral_mul_comp_proj_add, Torus.integral_mul_comp_proj_smul,
  Torus.periodicPairing_apply, Torus.periodicPairing_mFourierLp,
  Torus.periodicPairing_iteratedLineDerivOp_mFourierLp}
\leanok
\uses{def:torus-periodic-pairing, lem:period-box-counting, lem:period-box-measure,
  lem:torus-proj, lem:translation-fourier, lem:torus-line-symbol,
  lem:fourier-iterated-derivative, lem:directional-derivatives-commute, def:test-schwartz}
Let $1 \leq p \leq \infty$ and let $\varphi$ be a test function on $\R^d$.
\begin{enumerate}
  \item For every $h \in L_p(\T^d)$, the lift $h \circ \pi$ is locally integrable on $\R^d$, and
    $x \mapsto \varphi(x) \, h(\pi x)$ is integrable. For all $M, R \in \R$ with
    $|\varphi| \leq M$ and $\operatorname{supp} \varphi \subseteq B(0, R)$, the open ball, and
    every $h \in L_p(\T^d)$,
    \[
      |\langle \varphi, h \rangle| \leq M \, (2R + 2)^d \, \|h\|_1
        \leq M \, (2R + 2)^d \, \|h\|_p .
    \]
    The periodic pairing of $\varphi$ is a bounded linear functional on $L_p(\T^d)$.
  \item For every $m \in \Z^d$, $\langle \varphi, e_m \rangle = \mathcal{F}\varphi(-m)$.
  \item For all $k \in \N$, $v \in (\R^d)^k$ and $m \in \Z^d$,
    \[
      \langle \partial^v \varphi, e_m \rangle = (-1)^k \sigma_v(m) \, \langle \varphi, e_m \rangle .
    \]
\end{enumerate}
\end{lemma}

\begin{proof}
\leanok
(1) Local integrability is \Cref{lem:period-box-counting}, since $h$ is integrable on the
probability space $\T^d$; a test function times a locally integrable function is integrable. The
integrand vanishes outside $B(0, R)$, and
$\int_{B(0,R)} |h(\pi x)| \, dx \leq (2R + 2)^d \|h\|_1$ (\Cref{lem:period-box-counting}), and
$\|h\|_1 \leq \|h\|_p$ on the probability space $\T^d$. Linearity holds because the lifts of
almost-everywhere identities on $\T^d$ hold almost everywhere on $\R^d$
(\Cref{lem:period-box-measure}). (2) Since $e_m(t) = e_{-m}(-t)$ for every $t \in \T^d$
(\Cref{lem:translation-fourier}), $\langle \varphi, e_m \rangle$ is the integral
$\int \varphi(y) \, e_{-m}(-\pi y) \, dy$ of \Cref{lem:torus-proj} at $-m$. (3) By (2),
\Cref{lem:directional-derivatives-commute} and \Cref{lem:fourier-iterated-derivative},
$\langle \partial^v \varphi, e_m \rangle = \mathcal{F}(\partial^v \varphi)(-m)
= \sigma_v(-m) \, \mathcal{F}\varphi(-m)$, and $\sigma_v(-m) = (-1)^k \sigma_v(m)$
(\Cref{lem:torus-line-symbol}).
\end{proof}

\begin{lemma}[From the torus to the lifts]
\label{lem:torus-weak-derivative-to-lift}
\lean{Torus.periodicPairing_eq_of_mem_trigPolyLp, Torus.periodicPairing_lpBlock,
  Torus.HasWeakLineDeriv.periodicPairing_eq, Torus.HasWeakLineDeriv.comp_proj}
\leanok
\uses{lem:torus-periodic-pairing, def:torus-weak-line-derivative, def:weak-line-derivative,
  lem:weak-derivative-integral, thm:torus-coefficients-uniqueness, lem:torus-lp-blocks,
  lem:torus-lp-blocks-trig, lem:torus-weak-derivative-convolution,
  lem:torus-weak-line-derivative-basic, lem:period-box-measure}
Let $1 \leq p \leq \infty$, $k \in \N$, $v \in (\R^d)^k$ and let $\varphi$ be a test function on
$\R^d$.
\begin{enumerate}
  \item For every $n \in \N$ and all trigonometric polynomials $a, b \in L_p(\T^d)$ of degree
    $< n$ with $b = \partial^v a$,
    $\langle \varphi, b \rangle = (-1)^k \langle \partial^v \varphi, a \rangle$.
  \item Let $p < \infty$. For all $f, g \in L_p(\T^d)$ with $g = \partial^v f$ and every
    $j \in \N$, $\langle \varphi, \Delta_j g \rangle = (-1)^k \langle \partial^v \varphi,
    \Delta_j f \rangle$ and $\langle \varphi, g \rangle = (-1)^k \langle \partial^v \varphi, f
    \rangle$.
  \item For all $f, g \in L_p(\T^d)$ with $g = \partial^v f$, the lift $g \circ \pi$ is the weak
    derivative $\partial^v (f \circ \pi)$ on $\R^d$.
\end{enumerate}
\end{lemma}

\begin{proof}
\leanok
(1) Write $a = \sum_m \hat a(m) e_m$ and $b = \sum_m \hat b(m) e_m$, sums over the degrees
$< n$ (\Cref{thm:torus-coefficients-uniqueness}). By (3) of \Cref{lem:torus-periodic-pairing},
\[
  \hat b(m) \langle \varphi, e_m \rangle = \sigma_v(m) \hat a(m) \langle \varphi, e_m \rangle
  = (-1)^k \hat a(m) \langle \partial^v \varphi, e_m \rangle ,
\]
since $((-1)^k)^2 = 1$; sum over $m$. (2) The block $\Delta_j g = \psi_j * g$ is the weak
derivative $\partial^v \Delta_j f$ (\Cref{lem:torus-weak-derivative-convolution}), and both blocks
are trigonometric polynomials of degree $< 2^{j+1}$ (\Cref{lem:torus-lp-blocks-trig}), so (1)
applies. The partial sums $\sum_{j < J} \Delta_j$ converge to $f$ and $g$ in $L_p(\T^d)$
(\Cref{lem:torus-lp-blocks}), and both pairings are continuous
(\Cref{lem:torus-periodic-pairing}). (3) The classes of $f$ and $g$ in $L_1(\T^d)$ satisfy
$g = \partial^v f$ (\Cref{lem:torus-weak-line-derivative-basic}), and their lifts agree with
those of $f$ and $g$ almost everywhere (\Cref{lem:period-box-measure}). So (2) for $p = 1$ gives
$\int \varphi \, (g \circ \pi) = (-1)^k \int (\partial^v \varphi) (f \circ \pi)$ for every test
function $\varphi$; both lifts are locally integrable (\Cref{lem:torus-periodic-pairing}), and
\Cref{lem:weak-derivative-integral} applies.
\end{proof}

\begin{definition}[Modulated bumps]
\label{def:torus-modulated-bump}
\lean{Torus.modulatedBump}
\leanok
\uses{lem:torus-proj, thm:smooth-atomic-decomposition}
Let $\chi$ be the normalised smooth bump of radius $1$ around $0$ of
\Cref{thm:smooth-atomic-decomposition}: a smooth function on $\R^d$ with values in $[0, \infty)$,
support in the ball $B(0, 1)$ and $\int \chi = 1$. For every $m \in \Z^d$, regarded as a point
of $\R^d$, the modulated bump is
\[
  \eta_m(x) = \chi(x) \, e^{2\pi i \langle x, m \rangle} \qquad (x \in \R^d) .
\]
\end{definition}

\begin{lemma}[From the lifts to the torus]
\label{lem:torus-weak-derivative-from-lift}
\lean{Torus.econv_eq_of_comp_proj, Torus.modulatedBump_apply,
  Torus.hasCompactSupport_modulatedBump, Torus.fourier_modulatedBump,
  Torus.HasWeakLineDeriv.of_comp_proj}
\leanok
\uses{def:torus-modulated-bump, def:weak-line-derivative, lem:weak-derivative-convolution-local,
  lem:torus-lift-conv, lem:torus-proj, def:torus-econv, thm:torus-coefficients-convolution,
  lem:fourier-iterated-derivative, def:torus-weak-line-derivative}
Let $k \in \N$ and $v \in (\R^d)^k$.
\begin{enumerate}
  \item For all functions $f, g$ on $\T^d$ such that $g \circ \pi$ is the weak derivative
    $\partial^v (f \circ \pi)$ on $\R^d$, and every complex Schwartz function $\eta$ on $\R^d$
    with compact support, $\eta * g = (\partial^v \eta) * f$ at every point of $\T^d$.
  \item For every $m \in \Z^d$, $\eta_m$ is a complex Schwartz function with compact support and
    $\mathcal{F}\eta_m(m) = 1$.
  \item Let $1 \leq p \leq \infty$. For all $f, g \in L_p(\T^d)$ such that $g \circ \pi$ is the
    weak derivative $\partial^v (f \circ \pi)$ on $\R^d$, $g = \partial^v f$.
\end{enumerate}
\end{lemma}

\begin{proof}
\leanok
(1) Every point of $\T^d$ is $\pi x$ for some $x \in \R^d$ (\Cref{lem:torus-proj}), and at $\pi x$
both sides are convolutions on $\R^d$ with the lifts (\Cref{lem:torus-lift-conv}). These agree by
\Cref{lem:weak-derivative-convolution-local}, with $\Omega = \R^d$. (2) $\eta_m$ is smooth with
the compact support of $\chi$, and
$\mathcal{F}\eta_m(m) = \int e^{-2\pi i \langle x, m \rangle} \chi(x)
e^{2\pi i \langle x, m \rangle} \, dx = \int \chi = 1$. (3) Let $m \in \Z^d$. By (1) with
$\eta = \eta_m$, the coefficients at $m$ of $\eta_m * g$ and $(\partial^v \eta_m) * f$ agree. By
\Cref{thm:torus-coefficients-convolution}, \Cref{lem:fourier-iterated-derivative} and (2) they are
$\mathcal{F}\eta_m(m) \, \hat g(m) = \hat g(m)$ and
$\mathcal{F}(\partial^v \eta_m)(m) \, \hat f(m) = \sigma_v(m) \hat f(m)$.
\end{proof}

\begin{theorem}[Weak derivatives on the torus as distributions of the lifts]
\label{thm:torus-weak-derivative-lift}
\lean{Torus.hasWeakLineDeriv_iff_comp_proj, Torus.hasWeakDeriv_iff_comp_proj,
  Torus.memSobolev_iff_comp_proj}
\leanok
\uses{lem:torus-weak-derivative-to-lift, lem:torus-weak-derivative-from-lift,
  def:torus-weak-line-derivative, def:weak-derivative, def:weak-line-derivative,
  thm:torus-sobolev-directions, lem:torus-weak-line-derivative-basic}
Let $1 \leq p \leq \infty$, $k \in \N$ and $v \in (\R^d)^k$. For all $f, g \in L_p(\T^d)$,
$g = \partial^v f$ if and only if $g \circ \pi$ is the weak derivative $\partial^v (f \circ \pi)$
on $\R^d$, that is,
\[
  \int_{\R^d} \varphi(x) \, g(\pi x) \, dx
    = (-1)^k \int_{\R^d} (\partial^v \varphi)(x) \, f(\pi x) \, dx
\]
for every test function $\varphi$ on $\R^d$. For every $\alpha \in \N^d$ and all
$f, g \in L_p(\T^d)$, $g = \partial^\alpha f$ if and only if $g \circ \pi$ is the weak derivative
$\partial^{v^\alpha} (f \circ \pi)$ on $\R^d$. For every $m \in \N$ and every $f \in L_p(\T^d)$,
$f \in \Sob{m}{p}(\T^d)$ if and only if for every $k \leq m$ and every $v \in (\R^d)^k$ some
$g \in L_p(\T^d)$ has the lift $g \circ \pi = \partial^v (f \circ \pi)$, weakly on $\R^d$.
\end{theorem}

\begin{proof}
\leanok
The first claim is (3) of \Cref{lem:torus-weak-derivative-to-lift} and (3) of
\Cref{lem:torus-weak-derivative-from-lift}; the integral form is
\Cref{lem:weak-derivative-integral}. The second follows with $\partial^\alpha =
\partial^{v^\alpha}$ (\Cref{lem:torus-weak-line-derivative-basic}), and the third with
\Cref{thm:torus-sobolev-directions}.
\end{proof}

So $\Sob{m}{p}(\T^d)$ consists of the $f \in L_p(\T^d)$ whose periodic lifts have all
distributional derivatives of order at most $m$ in the lifts of $L_p(\T^d)$.

\begin{theorem}[Classical derivatives on the torus]
\label{thm:torus-classical-derivative}
\lean{Torus.hasWeakLineDeriv_of_contDiff, Torus.hasWeakLineDeriv_of_contDiff_chain,
  Fin.eq_prod_mul_of_succ_eq}
\leanok
\uses{lem:weak-derivative-classical, lem:torus-weak-derivative-from-lift,
  def:torus-weak-line-derivative, def:torus-line-symbol, lem:weak-derivative-basic}
Let $1 \leq p \leq \infty$.
\begin{enumerate}
  \item For every continuously differentiable $F \colon \R^d \to \C$, every $w \in \R^d$ and all
    $f, g \in L_p(\T^d)$ with $f \circ \pi = F$ and $g \circ \pi = DF(\cdot) w$ almost
    everywhere, $g = \partial_w f$.
  \item For every $k \in \N$, every $v \in (\R^d)^k$, all functions
    $F_0, \dots, F_k \colon \R^d \to \C$ such that $F_l$ is continuously differentiable and
    $F_{l+1} = DF_l(\cdot) v_l$ for every $l < k$, and all $f_0, \dots, f_k \in L_p(\T^d)$ with
    $f_l \circ \pi = F_l$ almost everywhere for every $l \leq k$, $f_k = \partial^v f_0$.
\end{enumerate}
\end{theorem}

\begin{proof}
\leanok
(1) $DF(\cdot) w$ is the weak derivative $\partial_w F$ on $\R^d$
(\Cref{lem:weak-derivative-classical}), so also of the lifts, which agree with $F$ and
$DF(\cdot) w$ almost everywhere (\Cref{lem:weak-derivative-basic}). Apply (3) of
\Cref{lem:torus-weak-derivative-from-lift}. (2) By (1),
$\hat f_{l+1}(m) = 2\pi i \langle m, v_l \rangle \hat f_l(m)$ for every $l < k$ and $m \in \Z^d$.
In every commutative monoid, $a_{l+1} = c_l a_l$ for every $l < k$ gives
$a_k = (\prod_{l < k} c_l) \, a_0$ by induction on $k$; so
$\hat f_k(m) = \sigma_v(m) \hat f_0(m)$.
\end{proof}

\section{Sobolev spaces as Triebel--Lizorkin spaces on the torus}
\label{sec:torus-sobolev-triebel-lizorkin}

Let $\#d \geq 1$ and $1 < p < \infty$. This section proves
$\Sob{m}{p}(\T^d) = \TL{m}{p}{2}(\T^d)$ for every $m \in \N$, with named constants in both
directions, by the argument of \Cref{sec:sobolev-triebel-lizorkin}: the strong type on $\R^d$ is
replaced by the strong type of vector-valued singular integrals on the torus
(\Cref{thm:torus-vector-valued-strong-type}), and the weak derivatives on $V$ by those of
\Cref{def:torus-weak-line-derivative}. The identity $\Sob{m}{p}(\T^d) = \TL{m}{p}{2}(\T^d)$ is
classical; see Triebel~\cite{Triebel}. In this section $V = \R^d$, and $e_\alpha$, $\eta_\alpha$,
$\theta^\alpha_j$, $A^\alpha$, $\tilde\theta^w_j$, $\tilde A^w$ and $B_w$ are those of
\Cref{sec:derivative-kernels}; the basis $e$ is the orthonormal basis of $\R^d$ of
\Cref{def:dyadic-cells}, and directional derivatives (\Cref{thm:torus-sobolev-directions}) relate
it to the coordinate basis.

\paragraph{From weak derivatives to the blocks.}

\begin{lemma}[The blocks on the torus as convolutions with the derivatives]
\label{lem:torus-blocks-as-derivative-convolutions}
\lean{Torus.fourier_two_pow_smul_psi_succ, Torus.two_pow_smul_lpBlock_succ,
  Torus.two_pow_mul_blockFun_succ_ae, Torus.lqNorm_blockFun_succ_le_sum_ae}
\leanok
\uses{cor:psi-derivatives-dilated, lem:fourier-iterated-derivative,
  thm:torus-coefficients-convolution, thm:torus-coefficients-uniqueness,
  lem:torus-lp-coefficients, def:torus-weak-line-derivative, def:torus-lp-blocks,
  def:torus-triebel-lizorkin, lem:blocks-as-derivative-convolutions, lem:dyadic-weights-integer,
  lem:lq-basic, lem:countable-minkowski}
For all $m, j \in \N$ and every $\xi \in \R^d$,
\[
  2^{jm} \, \mathcal{F}\psi_{j+1}(\xi)
  = \sum_\alpha \Bigl( \prod_{l < m} 2\pi i \langle \xi, e_{\alpha_l} \rangle \Bigr)
    \mathcal{F}\theta^\alpha_j(\xi) .
\]
Let $1 \leq p < \infty$, $m \in \N$ and $f \in L_p(\T^d)$, and for every multi-index $\alpha$
let $g_\alpha \in L_p(\T^d)$ be the weak derivative $\partial^{e_\alpha} f$. For every $j \in \N$,
\[
  2^{jm} \Delta_{j+1} f = \sum_\alpha \theta^\alpha_j * g_\alpha
\]
in $L_p(\T^d)$; for almost every $t \in \T^d$ and every $j \in \N$,
$2^{jm} (\psi_{j+1} * f)(t) = \sum_\alpha (\theta^\alpha_j * g_\alpha)(t)$; and for almost every
$t \in \T^d$,
\[
  \bigl\| (2^{(j+1)m} |(\psi_{j+1} * f)(t)|)_j \bigr\|_{\ell_2}
  \leq 2^m \sum_\alpha \bigl\| ((\theta^\alpha_j * g_\alpha)(t))_j \bigr\|_{\ell_2} .
\]
\end{lemma}

\begin{proof}
\leanok
The Fourier identity is the Fourier transform of \Cref{cor:psi-derivatives-dilated}: the
transform is linear, and $\mathcal{F}(\partial^{e_\alpha} G)(\xi)
= \prod_l 2\pi i \langle \xi, e_{\alpha_l} \rangle \, \mathcal{F}G(\xi)$ for Schwartz functions
$G$ (\Cref{lem:fourier-iterated-derivative}). At $\xi = m' \in \Z^d$ the product is the symbol
$\sigma_{e_\alpha}(m')$, so both sides of the identity in $L_p(\T^d)$ have the Fourier
coefficients $2^{jm} \mathcal{F}\psi_{j+1}(m') \hat f(m')$
(\Cref{thm:torus-coefficients-convolution}, \Cref{def:torus-weak-line-derivative}), and
\Cref{thm:torus-coefficients-uniqueness} applies. The pointwise identity follows since the
convolution integrals represent the operators (\Cref{def:torus-lp-blocks},
\Cref{def:torus-triebel-lizorkin}), and the bound follows from it as in
\Cref{lem:blocks-as-derivative-convolutions}.
\end{proof}

\begin{definition}[The constants of the embedding into $\TL{m}{p}{2}(\T^d)$]
\label{def:torus-sobolev-to-tl-const}
\lean{Torus.sobolevToTLConst, Torus.sobolevToTLConst_ne_top}
\leanok
\uses{def:periodic-cz-const, def:hormander-const, def:derivative-family,
  lem:torus-far-mass-derivative-families}
For $1 < p < \infty$, every $m \in \N$ and every multi-index $\alpha$,
\[
  K^{\T}_{p,\alpha} = S^{\T}_{p}\bigl(2, 4 H_{A^\alpha}, 2 T_{A^\alpha,0,1}\bigr) ,
\]
which is finite.
\end{definition}

\begin{theorem}[$\Sob{m}{p}(\T^d) \subseteq \TL{m}{p}{2}(\T^d)$]
\label{thm:torus-sobolev-to-triebel-lizorkin}
\lean{Torus.eLpNorm_lqNorm_blockFun_succ_le,
  Torus.triebelLizorkinNorm_le_of_hasWeakLineDeriv, Torus.memTriebelLizorkin_of_hasWeakLineDeriv,
  Torus.memTriebelLizorkin_of_memSobolev}
\leanok
\uses{lem:torus-blocks-as-derivative-convolutions, def:torus-sobolev-to-tl-const,
  thm:cz-family-derivative, lem:torus-far-mass-derivative-families,
  thm:torus-vector-valued-strong-type, prop:torus-tl-elementary, lem:lq-basic,
  lem:lp-norm-nonnegative, lem:lp-norm-homogeneity, thm:torus-sobolev-directions}
Let $\#d \geq 1$, $1 < p < \infty$, $m \in \N$ and
$f \in L_p(\T^d)$, and for every multi-index $\alpha$ let $g_\alpha \in L_p(\T^d)$ be the weak
derivative $\partial^{e_\alpha} f$. Then
\[
  \bigl\| \|(2^{(j+1)m} |(\psi_{j+1} * f)(t)|)_j\|_{\ell_2} \bigr\|_p
  \leq 2^m \sum_\alpha K^{\T}_{p,\alpha} \|g_\alpha\|_p
\]
and
\[
  \|f\|_{\TL{m}{p}{2}(\T^d)} \leq \|\Phi\|_1 \, \|f\|_p + 2^m \sum_\alpha K^{\T}_{p,\alpha} \,
  \|g_\alpha\|_p .
\]
In particular, $f \in \TL{m}{p}{2}(\T^d)$, and every $f \in \Sob{m}{p}(\T^d)$ lies in
$\TL{m}{p}{2}(\T^d)$.
\end{theorem}

\begin{proof}
\leanok
For the first bound, take the $L_p$ norm in
\Cref{lem:torus-blocks-as-derivative-convolutions}, use the homogeneity of the $L_p$ norm
(\Cref{lem:lp-norm-homogeneity}) and Minkowski's inequality in $L_p$, and bound each term by the
strong type of the analysis operator on the torus (\Cref{thm:torus-vector-valued-strong-type}) of
the Calder\'on--Zygmund family $(\theta^\alpha_j)_j$ with overlap $2$ and H\"ormander constant
$4 H_{A^\alpha}$ (\Cref{thm:cz-family-derivative}) and far mass at most $2 T_{A^\alpha,0,1}$
(\Cref{lem:torus-far-mass-derivative-families}). Splitting off the block $j = 0$ from the $\ell_2$
norm (\Cref{lem:lq-basic}), the triangle inequality in $L_p$ and
$\|\psi_0 * f\|_p \leq \|\Phi\|_1 \, \|f\|_p$ (\Cref{prop:torus-tl-elementary}) give the second
bound. For $f \in \Sob{m}{p}(\T^d)$ the weak derivatives $\partial^{e_\alpha} f$ exist in
$L_p(\T^d)$ (\Cref{thm:torus-sobolev-directions}), and the right side is finite.
\end{proof}

\paragraph{From the blocks to weak derivatives.}

\begin{definition}[The constants of the embedding into $\Sob{m}{p}(\T^d)$]
\label{def:torus-tl-to-sobolev-const}
\lean{Torus.tlToSobolevConst, Torus.tlToSobolevConst_ne_top}
\leanok
\uses{def:periodic-cz-const, def:hormander-const, def:dual-derivative-family,
  lem:torus-far-mass-derivative-families}
For $1 < p < \infty$, every $k \in \N$ and all $w \in (\R^d)^k$,
\[
  C^{\T}_{p,w} = S^{\T}_{p}\bigl(B_w, 4 H_{\tilde A^w}, 2 T_{\tilde A^w,0,1}\bigr) ,
\]
which is finite.
\end{definition}

\begin{lemma}[The synthesis bound for derivatives of the synthesis kernels on the torus]
\label{lem:torus-derivative-synthesis-bound}
\lean{Torus.enorm_sum_econvLp_derivDual_le, Torus.derivSynthesisSum,
  Torus.enorm_derivSynthesisSum_sub_le}
\leanok
\uses{def:torus-tl-to-sobolev-const, lem:torus-econv-identities, lem:psi-dual-derivatives,
  thm:cz-family-dual-derivative, lem:torus-far-mass-derivative-families,
  thm:torus-vector-valued-strong-type, def:torus-lp-blocks, def:torus-triebel-lizorkin,
  lem:blocks-as-derivative-convolutions, lem:dyadic-weights-integer, lem:mixed-norm-tails}
Let $\#d \geq 1$, $1 < p < \infty$, $k \in \N$, $w \in (\R^d)^k$ and $f \in L_p(\T^d)$. For
every finite set $s \subseteq \N$,
\begin{multline*}
  \Bigl\| \sum_{j \in s} (\partial^w \tilde\psi_j) * \Delta_j f \Bigr\|_p \\
  \leq C^{\T}_{p,w} \, \bigl\| \|(\mathbf 1_s(j) \, 2^{jk} |(\psi_j * f)(t)|)_j\|_{\ell_2}
  \bigr\|_p .
\end{multline*}
Let $D_J = \sum_{j < J} (\partial^w \tilde\psi_j) * \Delta_j f \in L_p(\T^d)$. For all
$N \leq N'$,
\[
  \|D_{N'} - D_N\|_p \leq C^{\T}_{p,w} \,
  \bigl\| \|(2^{jk} |(\psi_j * f)(t)|)_{j \geq N}\|_{\ell_2} \bigr\|_p .
\]
\end{lemma}

\begin{proof}
\leanok
Since $\partial^w \tilde\psi_j = 2^{jk} \tilde\theta^w_j$ (\Cref{lem:psi-dual-derivatives})
and constants move from the kernel to the function (\Cref{lem:torus-econv-identities}), the sum
is almost everywhere $\sum_{j \in s} \tilde\theta^w_j * h_j$ with
$h_j = \mathbf 1_s(j) 2^{jk} \Delta_j f$ (\Cref{def:torus-lp-blocks}), and
$|h_j| = \mathbf 1_s(j) 2^{jk} |\psi_j * f|$ almost everywhere (\Cref{def:torus-triebel-lizorkin},
\Cref{lem:blocks-as-derivative-convolutions}, \Cref{lem:dyadic-weights-integer}). The strong
type of the synthesis operators on the torus (\Cref{thm:torus-vector-valued-strong-type}) of the
Calder\'on--Zygmund family $(\tilde\theta^w_j)_j$ (\Cref{thm:cz-family-dual-derivative}) with far
mass at most $2 T_{\tilde A^w,0,1}$ (\Cref{lem:torus-far-mass-derivative-families}) gives the
first bound.
$D_{N'} - D_N$ is the sum over $s = \{N, \dots, N' - 1\}$, whose sequence is at most the tail
from $N$; the measurability is \Cref{lem:mixed-norm-tails}.
\end{proof}

\begin{lemma}[The partial sums of the derivative on the torus]
\label{lem:torus-derivative-partial-sums}
\lean{Torus.hasWeakLineDeriv_sum_lpBlock, Torus.cauchySeq_derivSynthesisSum}
\leanok
\uses{lem:torus-derivative-synthesis-bound, lem:torus-weighted-block-tails,
  lem:torus-weak-derivative-convolution, lem:torus-weak-line-derivative-basic,
  lem:torus-square-function-reproducing, lem:psi-dual-derivatives, lem:torus-econv-identities}
Let $1 \leq p < \infty$, $k \in \N$, $w \in (\R^d)^k$ and $f \in L_p(\T^d)$. For every $J \in \N$,
$D_J$ is the weak derivative $\partial^w$ of $\sum_{j < J} \Delta_j f$. If $\#d \geq 1$,
$1 < p < \infty$ and $f \in \TL{k}{p}{2}(\T^d)$, then $(D_J)_J$ is a Cauchy sequence in $L_p(\T^d)$.
\end{lemma}

\begin{proof}
\leanok
Each term $(\partial^w \tilde\psi_j) * \Delta_j f$ is the weak derivative of
$\tilde\psi_j * \Delta_j f$ (\Cref{lem:torus-weak-derivative-convolution}, with the Schwartz
function $\tilde\psi_j$ of \Cref{lem:psi-dual-derivatives}; equal kernels give the same
operator, see after \Cref{lem:torus-econv-identities}), which is $\Delta_j f$
(\Cref{lem:torus-square-function-reproducing}), and weak derivatives add up
(\Cref{lem:torus-weak-line-derivative-basic}). By \Cref{lem:torus-derivative-synthesis-bound},
$\|D_{N'} - D_N\|_p$ is at most $C^{\T}_{p,w}$ times the tail from $N$, which is finite and tends
to zero (\Cref{lem:torus-weighted-block-tails}).
\end{proof}

\begin{theorem}[$\TL{m}{p}{2}(\T^d) \subseteq \Sob{m}{p}(\T^d)$]
\label{thm:torus-triebel-lizorkin-to-sobolev}
\lean{Torus.exists_hasWeakLineDeriv_of_memTriebelLizorkin,
  Torus.exists_hasWeakLineDeriv_of_memTriebelLizorkin_of_le,
  Torus.memSobolev_of_memTriebelLizorkin}
\leanok
\uses{lem:torus-derivative-partial-sums, lem:torus-derivative-synthesis-bound,
  lem:torus-weighted-block-tails, lem:torus-lp-blocks, lem:torus-weak-line-derivative-basic,
  def:torus-tl-to-sobolev-const, prop:torus-tl-elementary, thm:torus-sobolev-directions}
Let $\#d \geq 1$, $1 < p < \infty$, $k \in \N$ and $w \in (\R^d)^k$. For every
$f \in \TL{k}{p}{2}(\T^d)$ the weak derivative $\partial^w f$ exists as some $g \in L_p(\T^d)$
with
\[
  \|g\|_p \leq C^{\T}_{p,w} \, \|f\|_{\TL{k}{p}{2}(\T^d)} .
\]
For every $m \geq k$ and every $f \in \TL{m}{p}{2}(\T^d)$, the weak derivative $g = \partial^w f$
satisfies $\|g\|_p \leq C^{\T}_{p,w} \, \|f\|_{\TL{m}{p}{2}(\T^d)}$. In particular,
$\TL{m}{p}{2}(\T^d) \subseteq \Sob{m}{p}(\T^d)$ for every $m \in \N$.
\end{theorem}

\begin{proof}
\leanok
Let $g$ be the limit of the Cauchy sequence $(D_J)_J$ in $L_p(\T^d)$
(\Cref{lem:torus-derivative-partial-sums}). The partial sums $\sum_{j < J} \Delta_j f$ tend to
$f$ in $L_p(\T^d)$ (\Cref{lem:torus-lp-blocks}), and $D_J$ is their weak derivative
$\partial^w$, so $g = \partial^w f$ (\Cref{lem:torus-weak-line-derivative-basic}). By
\Cref{lem:torus-derivative-synthesis-bound} with $s = \{0, \dots, J - 1\}$ and
\Cref{lem:torus-weighted-block-tails},
$\|D_J\|_p \leq C^{\T}_{p,w} \|f\|_{\TL{k}{p}{2}(\T^d)}$ for every $J$, and the bound passes to
the limit. For $m \geq k$,
$\|f\|_{\TL{k}{p}{2}(\T^d)} \leq \|f\|_{\TL{m}{p}{2}(\T^d)}$ (\Cref{prop:torus-tl-elementary}).
The last claim is \Cref{thm:torus-sobolev-directions}.
\end{proof}

\paragraph{Consequences.}

\begin{corollary}[$\Sob{m}{p}(\T^d) = \TL{m}{p}{2}(\T^d)$]
\label{cor:torus-sobolev-triebel-lizorkin}
\lean{Torus.memSobolev_iff_memTriebelLizorkin}
\leanok
\uses{thm:torus-sobolev-to-triebel-lizorkin, thm:torus-triebel-lizorkin-to-sobolev}
For every $\#d \geq 1$, every $1 < p < \infty$, every $m \in \N$ and every $f \in L_p(\T^d)$,
$f \in \Sob{m}{p}(\T^d)$ if and only if $f \in \TL{m}{p}{2}(\T^d)$.
\end{corollary}

\begin{proof}
\leanok
\Cref{thm:torus-sobolev-to-triebel-lizorkin} and \Cref{thm:torus-triebel-lizorkin-to-sobolev}.
\end{proof}

For $p = 2$ and $m \geq 1$, together with $\TL{m}{2}{2}(\T^d) = \Besov{m}{2}{2}(\T^d)$
(\Cref{thm:torus-tl-besov-sandwich}), this gives $\Sob{m}{2}(\T^d) = \Besov{m}{2}{2}(\T^d)$ once
more, which \Cref{thm:weak-sobolev} proves through the Fourier coefficients.

\begin{corollary}[The Besov sandwich for $\Sob{m}{p}(\T^d)$]
\label{cor:torus-sobolev-besov-sandwich}
\lean{Torus.memBesov_of_memSobolev, Torus.memSobolev_of_memBesov,
  Torus.approxNormDyadic_le_of_hasWeakLineDeriv, Torus.exists_hasWeakLineDeriv_of_memBesov}
\leanok
\uses{cor:torus-sobolev-triebel-lizorkin, thm:torus-sobolev-to-triebel-lizorkin,
  thm:torus-triebel-lizorkin-to-sobolev, thm:torus-tl-besov-sandwich,
  cor:torus-square-function-theorem}
Let $\#d \geq 1$, $1 < p < \infty$ and let $A$ be the scheme of trigonometric polynomials in
$L_p(\T^d)$. For every $m \in \N$,
\[
  \Besov{m}{p}{\min(p,2)}(\T^d) \subseteq \Sob{m}{p}(\T^d) ,
\]
and for every $m \geq 1$,
\[
  \Sob{m}{p}(\T^d) \subseteq \Besov{m}{p}{\max(p,2)}(\T^d) .
\]
For every $m \geq 1$ and every $f \in L_p(\T^d)$ with weak derivatives
$g_\alpha = \partial^{e_\alpha} f \in L_p(\T^d)$ for every multi-index $\alpha$ of length $m$,
\[
  \|f\|^{\mathrm{dy}}_{A^m_{\max(p,2)}} \leq 2 \cdot 2^m (1 - 2^{-m})^{-1}
  \Bigl( \|\Phi\|_1 \, \|f\|_p + 2^m \sum_\alpha K^{\T}_{p,\alpha} \, \|g_\alpha\|_p \Bigr) .
\]
For every $m \in \N$, every $f \in \Besov{m}{p}{\min(p,2)}(\T^d)$, every $k \leq m$ and all
$w \in (\R^d)^k$, the weak derivative $\partial^w f$ is some $g \in L_p(\T^d)$ with
\[
  \|g\|_p \leq C^{\T}_{p,w} \cdot 2 (\#d + 1) \, 2^{(\#d + 1) m} \, \|\Phi\|_1 \,
  \|f\|^{\mathrm{dy}}_{A^m_{\min(p,2)}} .
\]
\end{corollary}

\begin{proof}
\leanok
$\Besov{m}{p}{\min(p,2)}(\T^d) \subseteq \TL{m}{p}{2}(\T^d)$ for every $m \geq 0$
(\Cref{cor:torus-square-function-theorem}), and \Cref{cor:torus-sobolev-triebel-lizorkin} gives
the first inclusion. The Besov sandwich for $s = m > 0$ (\Cref{thm:torus-tl-besov-sandwich})
gives the second. The bounds combine \Cref{thm:torus-sobolev-to-triebel-lizorkin} and
\Cref{thm:torus-triebel-lizorkin-to-sobolev} with the bounds of
\Cref{thm:torus-tl-besov-sandwich}.
\end{proof}

\begin{corollary}[$\Sob{m}{p}(\T^d)$ by local means]
\label{cor:torus-sobolev-local-means}
\lean{Torus.memSobolev_iff_tlLocalMeans}
\leanok
\uses{cor:torus-sobolev-triebel-lizorkin, thm:torus-tl-local-means, def:torus-tl-local-means}
Let $\#d \geq 1$, $1 < p < \infty$, let $g$ be a bump on $\R^d$ and let $m, L \in \N$ with
$1 \leq m < L$. For every $f \in L_p(\T^d)$, $f \in \Sob{m}{p}(\T^d)$ if and only if
$\|f\|^{g,L}_{F^m_{p,2}(\T^d)} < \infty$.
\end{corollary}

\begin{proof}
\leanok
\Cref{cor:torus-sobolev-triebel-lizorkin} and \Cref{thm:torus-tl-local-means} with $q = 2$ and
$s = m$.
\end{proof}

\section{Sobolev embeddings on the torus}
\label{sec:torus-sobolev-embedding}

Through $\Sob{m}{p}(\T^d) \subseteq \Besov{m}{p}{\max(p,2)}(\T^d)$
(\Cref{cor:torus-sobolev-besov-sandwich}), the Besov embeddings of the torus
(\Cref{thm:torus-embeddings}) give the Sobolev embeddings on the torus. Here $\iota_{r,p}$ is the
inclusion $L_r(\T^d) \hookrightarrow L_p(\T^d)$ for $p \leq r$.

\begin{theorem}[Sobolev embeddings on the torus]
\label{thm:torus-sobolev-embedding}
\lean{Torus.memSobolev_exists_inclusion_eq, Torus.memSobolev_exists_continuous,
  Torus.memSobolev_of_memSobolev_of_lt}
\leanok
\uses{cor:torus-sobolev-besov-sandwich, thm:torus-embeddings, def:nikolskii-property,
  thm:approx-embedding}
Let $\#d \geq 1$, $1 < p < \infty$ and $m \in \N$.
\begin{enumerate}
  \item For every $r$ with $p \leq r \leq \infty$ and $m > \#d (1/p - 1/r)$, every
    $f \in \Sob{m}{p}(\T^d)$ is $\iota_{r,p} g$ for some
    $g \in \Besov{m - \#d(1/p - 1/r)}{r}{\max(p,2)}(\T^d)$.
  \item If $m > \#d/p$, then every $f \in \Sob{m}{p}(\T^d)$ has a continuous representative.
  \item For every $r$ with $p \leq r < \infty$ and every $k \in \N$ with
    $k < m - \#d (1/p - 1/r)$, every $f \in \Sob{m}{p}(\T^d)$ is $\iota_{r,p} g$ for some
    $g \in \Sob{k}{r}(\T^d)$.
\end{enumerate}
\end{theorem}

\begin{proof}
\leanok
In (1), $m \geq 1$ since $\#d (1/p - 1/r) \geq 0$ (\Cref{def:nikolskii-property}); in (2),
$m \geq 1$ since $\#d/p \geq 0$. So $f \in \Besov{m}{p}{\max(p,2)}(\T^d)$
(\Cref{cor:torus-sobolev-besov-sandwich}), and the first two items of \Cref{thm:torus-embeddings}
apply. In (3), $m > \#d (1/p - 1/r)$, so (1) gives $g$, and $g \in \Besov{k}{r}{\min(r,2)}(\T^d)$
since $0 \leq k < m - \#d(1/p - 1/r)$ (\Cref{thm:approx-embedding}); hence $g \in \Sob{k}{r}(\T^d)$
(\Cref{cor:torus-sobolev-besov-sandwich}, with $1 < p \leq r$).
\end{proof}

\chapter{Bessel potential spaces}
\label{ch:bessel-potential-spaces}

In this chapter $V$ is a finite-dimensional real inner product space of dimension $n \geq 1$
with Lebesgue measure, and $1 < p < \infty$. It identifies Mathlib's Bessel potential spaces
$H^{s,p}(V)$, the tempered distributions $f$ with $J_s f \in L_p(V)$, with the Triebel--Lizorkin
spaces $\TL{s}{p}{2}$. \Cref{sec:bessel-kernels} shows that the Bessel potentials of the
dilated Littlewood--Paley kernels are multiples of the dilates of Bessel kernels, which form
uniform Schwartz families whose dilates are Calder\'on--Zygmund families (\Cref{def:cz-family});
\Cref{sec:bessel-potential-lifting} proves $\TL{s}{p}{2} \subseteq H^{s,p}(V)$ for every real
$s$, $H^{s,p}(V) = \TL{s}{p}{2}$ for $s \geq 0$, and Calder\'on's theorem
$\Sob{m}{p}(V) = H^{m,p}(V)$ (\Cref{cor:sobolev-bessel}).
\Cref{sec:bessel-potential-distributions} proves $H^{s,p}(V) = \TL{s}{p}{2}$ for tempered
distributions and every real $s$, with the Fourier multipliers of \Cref{ch:fourier-multipliers}.
The constants come from uniform Schwartz bounds and are existential. The real interpolation spaces
of Bessel potential spaces are
Besov spaces (\Cref{sec:bessel-potential-interpolation}), and
\Cref{sec:bessel-potential-embedding} proves $H^{s,p}(V) \hookrightarrow L_r(V)$ for
$p \leq r < \infty$ and $s \geq n(1/p - 1/r)$.

\section{Uniform Schwartz families and the Bessel kernels}
\label{sec:bessel-kernels}

The Bessel potential $J_\sigma$ on Schwartz functions (\Cref{def:bessel-schwartz}) multiplies
the Fourier transform by the Bessel weight
$\langle \xi \rangle^\sigma = (1 + \|\xi\|^2)^{\sigma/2}$ (\Cref{def:bessel-weight}). Since
$1 + \|2^j \zeta\|^2 = 2^{2j} (2^{-2j} + \|\zeta\|^2)$, for every Schwartz function $\omega$ the
Bessel potential of the dilate $\omega_{2^j}$ is $2^{j\sigma} (K_j)_{2^j}$ with the Bessel kernel
$K_j = \mathcal F^{-1} \bigl( (2^{-2j} + \|\cdot\|^2)^{\sigma/2} \, \widehat\omega \bigr)$, also
for Mathlib's Bessel potential on tempered distributions. Let now $\widehat\omega$ vanish outside
an annulus $\{a \leq \|\zeta\| \leq b\}$ with $a > 0$, as for the block $\psi_1$ or the bases
$G^{\psi}_{j+1}$ and $\tilde G_{j+2}$ of \Cref{def:psi-base} and \Cref{def:psi-dual}. The
symbols $(\varepsilon + \|\zeta\|^2)^{\sigma/2}$, $0 < \varepsilon \leq 1$, have bounded
derivatives on the annulus uniformly in $\varepsilon$, so the kernels $K_j$ form a uniform
Schwartz family (\Cref{def:uniform-schwartz}), and their dilates form a Calder\'on--Zygmund family
(\Cref{def:cz-family}) with an explicit overlap. The seminorm bounds of the family come from
compactness and from the continuity of $\mathcal F^{-1}$ on Schwartz functions; they are
existential. Write $\|G\|_{k,l} = \sup_x \|x\|^k \|D^l G(x)\|$ for the seminorms of a Schwartz
function $G$.

\paragraph{Uniform families under linear maps.}

\begin{theorem}[Continuous linear maps preserve uniform families]
\label{thm:uniform-schwartz-linear-map}
\lean{Rd.exists_isUniformSchwartz_iff_isVonNBounded, Rd.IsUniformSchwartz.exists_map,
  Rd.exists_isUniformSchwartz_fourierInv}
\leanok
\uses{def:uniform-schwartz}
For every family $(G_j)_{j \in \N}$ of Schwartz functions $V \to \C$, there are bounds $A$ such
that $(G_j)$ is uniformly Schwartz with bounds $A$ if and only if the set
$\{G_j : j \in \N\}$ is bounded in $\mathcal S(V)$, that is, absorbed by every neighbourhood of
$0$. For every uniform family $(G_j)$ and every continuous real-linear map $T$ of
$\mathcal S(V)$ into itself, there are bounds $A'$ such that $(T G_j)$ is uniformly Schwartz with
bounds $A'$. In particular, $(\mathcal F^{-1} G_j)$ is uniformly Schwartz with some bounds.
\end{theorem}

\begin{proof}
\leanok
The topology of $\mathcal S(V)$ is generated by the seminorms $\|\cdot\|_{k,l}$, so a set is
bounded exactly when every seminorm is bounded on it (Mathlib's
\texttt{WithSeminorms.\allowbreak isVonNBounded\_\allowbreak iff\_\allowbreak
seminorm\_\allowbreak bddAbove} with \texttt{schwartz\_\allowbreak withSeminorms}); this is the
first claim. Continuous linear maps send bounded sets to bounded sets (Mathlib's
\texttt{Bornology.\allowbreak IsVonNBounded.\allowbreak image}), and $\{T G_j : j \in \N\}$ is
the image of $\{G_j : j \in \N\}$ under $T$. The inverse Fourier transform is a continuous linear
map of $\mathcal S(V)$.
\end{proof}

\paragraph{The shifted Bessel symbols.}

\begin{definition}[Shifted Bessel symbols]
\label{def:shifted-bessel-symbol}
\lean{Rd.besselSymbol, Rd.besselSymbol_eq_comp, Rd.besselSymbol_one}
\leanok
\uses{def:bessel-weight}
For all real $\varepsilon$ and $\sigma$ and every $\zeta \in V$ let
\[
  \beta_{\varepsilon,\sigma}(\zeta) = (\varepsilon + \|\zeta\|^2)^{\sigma/2} ,
\]
so that $\beta_{\varepsilon,\sigma} = g_\sigma \circ (\varepsilon + \|\cdot\|^2)$ with
$g_\sigma(t) = t^{\sigma/2}$, and $\beta_{1,\sigma} = \langle \cdot \rangle^\sigma$.
\end{definition}

\begin{lemma}[Scaling of the Bessel weight]
\label{lem:bessel-weight-scaling}
\lean{Rd.besselWeight_eq_mul_besselSymbol}
\leanok
\uses{def:shifted-bessel-symbol, def:bessel-weight}
For every real $c > 0$, every real $\sigma$ and every $\xi \in V$,
\[
  \langle \xi \rangle^\sigma = c^\sigma \, \beta_{c^{-2},\sigma}(c^{-1} \xi) .
\]
\end{lemma}

\begin{proof}
\leanok
$c^{-2} + \|c^{-1} \xi\|^2 = c^{-2} (1 + \|\xi\|^2)$, and
$(c^{-2})^{\sigma/2} = c^{-\sigma}$.
\end{proof}

\begin{lemma}[Smoothness and growth of the shifted symbols]
\label{lem:shifted-bessel-symbol-smooth}
\lean{Complex.contDiffOn_ofReal_rpow_const, Rd.contDiff_besselSymbol, Rd.besselSymbol_eq_mul,
  Rd.hasTemperateGrowth_besselSymbol}
\leanok
\uses{def:shifted-bessel-symbol, lem:bessel-weight-scaling}
For every real $\tau$, the function $t \mapsto t^\tau$ is smooth on $(0, \infty)$. For every
real $\varepsilon > 0$ and $\sigma$, $\beta_{\varepsilon,\sigma}$ is smooth, has temperate growth,
and $\beta_{\varepsilon,\sigma}(\zeta) = \varepsilon^{\sigma/2} \langle \varepsilon^{-1/2} \zeta
\rangle^\sigma$ for every $\zeta \in V$.
\end{lemma}

\begin{proof}
\leanok
The function $\zeta \mapsto \varepsilon + \|\zeta\|^2$ is smooth and positive, and real powers
are smooth on $(0, \infty)$. The identity is \Cref{lem:bessel-weight-scaling} with
$c = \varepsilon^{-1/2}$ at the point $c \zeta$. The Bessel weight
has temperate growth (Mathlib's
\texttt{Function.\allowbreak hasTemperateGrowth\_\allowbreak one\_\allowbreak add\_\allowbreak
norm\_\allowbreak sq\_\allowbreak rpow}), and so have its compositions with linear maps and its
multiples.
\end{proof}

\begin{lemma}[Bounded derivatives on compact sets]
\label{lem:bounded-derivatives-compact}
\lean{ContDiffOn.exists_norm_iteratedFDerivWithin_le}
\leanok
Let $E, F$ be real normed spaces, let $f \colon E \to F$ be smooth on a set $s \subseteq E$ of
unique differentiability, let $K \subseteq s$ be compact and $r \in \N$. Then there is $C$ with
$\|D^i_s f(x)\| \leq C$ for every $i \leq r$ and every $x \in K$, where $D^i_s$ is the $i$-th
derivative within $s$.
\end{lemma}

\begin{proof}
\leanok
The function $x \mapsto \sum_{i \leq r} \|D^i_s f(x)\|$ is continuous on $s$, hence bounded on
$K$.
\end{proof}

\begin{lemma}[Derivatives of the shifted square norm]
\label{lem:norm-sq-derivatives}
\lean{Rd.exists_norm_iteratedFDeriv_add_norm_sq_le}
\leanok
\uses{lem:bounded-derivatives-compact}
For every real $b$ and every $r \in \N$ there is $D$ with
$\|D^i (\varepsilon + \|\cdot\|^2)(\zeta)\| \leq D^i$ for every real $\varepsilon$, every
$1 \leq i \leq r$ and every $\zeta \in V$ with $\|\zeta\| \leq b$.
\end{lemma}

\begin{proof}
\leanok
The derivatives of positive order of $\varepsilon + \|\cdot\|^2$ are those of $\|\cdot\|^2$;
\Cref{lem:bounded-derivatives-compact} with $s = V$ and the closed ball of radius $b$ bounds them
by some $C$, and $D = \max(1, C)$ satisfies $C \leq D \leq D^i$ for $i \geq 1$.
\end{proof}

\begin{lemma}[Bounds of the shifted symbols on an annulus]
\label{lem:shifted-bessel-symbol-bounds}
\lean{Rd.exists_norm_iteratedFDeriv_besselSymbol_le, Rd.besselSymbolBound,
  Rd.besselSymbolBound_nonneg, Rd.norm_besselSymbol_le}
\leanok
\uses{lem:shifted-bessel-symbol-smooth, lem:bounded-derivatives-compact,
  lem:norm-sq-derivatives}
Let $a > 0$ and $b, \sigma$ be real. For every $r \in \N$ there is $C_r \geq 0$ with
$\|D^r \beta_{\varepsilon,\sigma}(\zeta)\| \leq C_r$ for every $\varepsilon \in (0, 1]$ and every
$\zeta \in V$ with $a \leq \|\zeta\| \leq b$. Moreover, with
\[
  M_{a,b,\sigma} = \max \bigl( (a^2)^{\sigma/2}, (1 + b^2)^{\sigma/2} \bigr) \geq 0 ,
\]
$|\beta_{\varepsilon,\sigma}(\zeta)| \leq M_{a,b,\sigma}$ for every $\varepsilon \in [0, 1]$ and
every $\zeta \in V$ with $a \leq \|\zeta\| \leq b$.
\end{lemma}

\begin{proof}
\leanok
Write $\beta_{\varepsilon,\sigma} = g_\sigma \circ f_\varepsilon$ with $f_\varepsilon(\zeta) =
\varepsilon + \|\zeta\|^2$. For $\varepsilon \in (0, 1]$, $f_\varepsilon > 0$ on all of $V$, so
its range lies in $(0, \infty)$ as Mathlib's bound for compositions requires. For
$\varepsilon \in [0, 1]$ and $a \leq \|\zeta\| \leq b$, $f_\varepsilon(\zeta)$ lies in the compact
interval $[a^2, 1 + b^2] \subseteq (0, \infty)$, where the derivatives of $g_\sigma$ of the orders
$i \leq r$ are bounded by some $C$ (\Cref{lem:shifted-bessel-symbol-smooth},
\Cref{lem:bounded-derivatives-compact}), and the derivatives of $f_\varepsilon$ of the orders
$1 \leq i \leq r$ are at most $D^i$ (\Cref{lem:norm-sq-derivatives}). The Fa\`a di Bruno bound
for compositions (Mathlib's
\texttt{norm\_\allowbreak iteratedFDeriv\_\allowbreak comp\_\allowbreak le'}) gives
$C_r = \max(0, r! \, C D^r)$. The second claim holds since $t \mapsto t^{\sigma/2}$ is monotone on
$[a^2, 1 + b^2]$, increasing for $\sigma \geq 0$ and decreasing for $\sigma < 0$.
\end{proof}

\paragraph{The Bessel kernels.}

\begin{definition}[Bessel symbol and kernel families]
\label{def:bessel-kernel-family}
\lean{Rd.besselSymbolFamily, Rd.coe_besselSymbolFamily, Rd.besselKernelFamily,
  Rd.fourier_besselKernelFamily}
\leanok
\uses{def:shifted-bessel-symbol, lem:shifted-bessel-symbol-smooth}
For every Schwartz function $\widehat\omega \colon V \to \C$, every real $\sigma$ and every
$j \in \N$ let $S_j = \beta_{2^{-2j},\sigma} \, \widehat\omega$ and
$K_j = \mathcal F^{-1} S_j$. Both are Schwartz functions, and
$\widehat{K_j} = \beta_{2^{-2j},\sigma} \, \widehat\omega$.
\end{definition}

The product $S_j$ is a Schwartz function because $\beta_{2^{-2j},\sigma}$ has temperate growth
(\Cref{lem:shifted-bessel-symbol-smooth}; Mathlib's
\texttt{SchwartzMap.\allowbreak smulLeftCLM}), and $\mathcal F$ inverts $\mathcal F^{-1}$ on
Schwartz functions.

\begin{lemma}[Leibniz bound against a Schwartz function]
\label{lem:leibniz-schwartz}
\lean{Rd.norm_iteratedFDeriv_mul_schwartz_le}
\leanok
For every smooth $g \colon V \to \C$, every Schwartz function $\omega \colon V \to \C$, every
$l \in \N$, every $x \in V$ and all reals $C_0, \dots, C_l$ with $\|D^i g(x)\| \leq C_i$ for
every $i \leq l$,
\[
  \|D^l (g \omega)(x)\| \leq \sum_{i \leq l} \binom{l}{i} C_i \, \|\omega\|_{0,l-i} .
\]
\end{lemma}

\begin{proof}
\leanok
The Leibniz rule (Mathlib's
\texttt{norm\_\allowbreak iteratedFDeriv\_\allowbreak mul\_\allowbreak le}) with
$\|D^{l-i} \omega(x)\| \leq \|\omega\|_{0,l-i}$.
\end{proof}

\begin{lemma}[Supports in an annulus]
\label{lem:tsupport-annulus}
\lean{Rd.mem_annulus_of_mem_tsupport}
\leanok
For all reals $a, b$ and every $f \colon V \to \C$ that vanishes outside the annulus
$\{a \leq \|\zeta\| \leq b\}$, the closure of the support of $f$ lies in that annulus.
\end{lemma}

\begin{proof}
\leanok
The annulus is closed.
\end{proof}

\begin{theorem}[The Bessel kernels are uniformly Schwartz]
\label{thm:bessel-kernels-uniform}
\lean{Rd.exists_isUniformSchwartz_besselSymbolFamily,
  Rd.exists_isUniformSchwartz_besselKernelFamily}
\leanok
\uses{def:bessel-kernel-family, lem:shifted-bessel-symbol-bounds, lem:leibniz-schwartz,
  lem:tsupport-annulus, thm:uniform-schwartz-linear-map, def:uniform-schwartz}
Let $a > 0$ and $b$ be real, and let $\widehat\omega \colon V \to \C$ be a Schwartz function
that vanishes outside the annulus $\{a \leq \|\zeta\| \leq b\}$. Then for every real $\sigma$
there are bounds $A$ such that $(S_j)$ is uniformly Schwartz with bounds $A$, and bounds $A'$
such that $(K_j)$ is uniformly Schwartz with bounds $A'$.
\end{theorem}

\begin{proof}
\leanok
Let $C_r \geq 0$ be the bounds of \Cref{lem:shifted-bessel-symbol-bounds}; note
$2^{-2j} \in (0, 1]$. At a point $x$ outside the closure of the support of $\widehat\omega$
every derivative of $S_j$ vanishes. At a point $x$ of that closure, $a \leq \|x\| \leq b$
(\Cref{lem:tsupport-annulus}), so $\|x\|^k \leq |b|^k$ and
\[
  \|x\|^k \|D^l S_j(x)\| \leq |b|^k \sum_{i \leq l} \binom{l}{i} C_i \,
  \|\widehat\omega\|_{0,l-i}
\]
by \Cref{lem:leibniz-schwartz}, for every $j$. The kernels follow by
\Cref{thm:uniform-schwartz-linear-map}.
\end{proof}

\begin{theorem}[Bessel potentials of dilates]
\label{thm:bessel-potential-dilate}
\lean{Rd.besselSchwartz_schwartzDilate, Rd.besselPotential_schwartzDilate,
  Rd.enorm_two_pow_rpow}
\leanok
\uses{def:bessel-kernel-family, lem:bessel-weight-scaling, def:schwartz-dilate,
  lem:fourier-dilate, lem:bessel-weight-basic, def:bessel-schwartz,
  lem:bessel-schwartz-distribution}
For every Schwartz function $\omega \colon V \to \C$, every real $\sigma$ and every $j \in \N$,
let $K_j$ be the Bessel kernels of $\widehat\omega$. Then
\[
  J_\sigma \omega_{2^j} = 2^{j\sigma} (K_j)_{2^j} ,
\]
and Mathlib's Bessel potential $J_\sigma$ maps the tempered distribution of $\omega_{2^j}$ to
that of $2^{j\sigma} (K_j)_{2^j}$. For every $j \in \N$ and every real $\sigma$,
$|(2^j)^\sigma| = 2^{j\sigma}$.
\end{theorem}

\begin{proof}
\leanok
Both sides are Schwartz functions, and the Fourier transform is injective on them. At $\xi$,
the Fourier transform of the left side is $\langle \xi \rangle^\sigma \widehat\omega(2^{-j}\xi)$
(\Cref{lem:fourier-dilate}), and that of the right side is
$2^{j\sigma} \, \beta_{2^{-2j},\sigma}(2^{-j} \xi) \, \widehat\omega(2^{-j} \xi)$
(\Cref{lem:fourier-dilate}, \Cref{def:bessel-kernel-family}). They agree by
\Cref{lem:bessel-weight-scaling} with $c = 2^j$. The statement on tempered distributions follows
with \Cref{lem:bessel-schwartz-distribution}, and $(2^j)^\sigma = 2^{j\sigma} > 0$.
\end{proof}

The identity holds for every Schwartz function $\omega$; the annulus condition of
\Cref{thm:bessel-kernels-uniform} makes the kernels $K_j$ uniform in $j$.

\paragraph{The dilated Bessel kernels.}

\begin{definition}[Overlap of the Bessel kernels]
\label{def:bessel-overlap}
\lean{Rd.besselOverlap, Rd.besselOverlap_ne_top}
\leanok
\uses{lem:shifted-bessel-symbol-bounds}
For every real $a$, every $N \in \N$, every real $\sigma$ and every Schwartz function
$\widehat\omega \colon V \to \C$ let
\[
  B_{a,N,\sigma,\widehat\omega} = (N + 1) \bigl( M_{a, 2^N a, \sigma} \,
  \|\widehat\omega\|_{0,0} \bigr)^2 < \infty .
\]
\end{definition}

\begin{lemma}[Frequencies of the dilated Bessel kernels]
\label{lem:bessel-kernel-frequency}
\lean{Rd.fourier_dilate_besselKernelFamily,
  Rd.mem_annulus_of_fourier_dilate_besselKernelFamily_ne_zero,
  Rd.enorm_fourier_dilate_besselKernelFamily_le,
  Rd.mem_Icc_of_fourier_dilate_besselKernelFamily_ne_zero}
\leanok
\uses{def:bessel-kernel-family, lem:fourier-dilate, lem:shifted-bessel-symbol-bounds,
  lem:dyadic-annulus-count}
Let $\widehat\omega \colon V \to \C$ be a Schwartz function, $\sigma$ real and $K_j$ the Bessel
kernels of $\widehat\omega$. For every $j \in \N$ and every $\xi \in V$, with
$\zeta = 2^{-j} \xi$,
\[
  \widehat{(K_j)_{2^j}}(\xi) = \beta_{2^{-2j},\sigma}(\zeta) \, \widehat\omega(\zeta) .
\]
Let $a > 0$ and $b$ be real, and let $\widehat\omega$ vanish outside the annulus
$\{a \leq \|\zeta\| \leq b\}$. Then for every $j \in \N$ and every $\xi \in V$,
$|\widehat{(K_j)_{2^j}}(\xi)| \leq M_{a,b,\sigma} \|\widehat\omega\|_{0,0}$, and if
$\widehat{(K_j)_{2^j}}(\xi) \neq 0$, then $a \leq 2^{-j} \|\xi\| \leq b$. If moreover
$b = 2^N a$ with $N \in \N$, and $i_0 \in \N$ is the least index with
$\|\xi\| \leq 2^{i_0} 2^N a$, then $\widehat{(K_j)_{2^j}}(\xi) \neq 0$ implies
$i_0 \leq j \leq i_0 + N$.
\end{lemma}

\begin{proof}
\leanok
The identity is \Cref{lem:fourier-dilate} with \Cref{def:bessel-kernel-family}. If
$\widehat\omega(\zeta) \neq 0$, then $a \leq \|\zeta\| = 2^{-j} \|\xi\| \leq b$, and
$|\beta_{2^{-2j},\sigma}(\zeta)| \leq M_{a,b,\sigma}$ (\Cref{lem:shifted-bessel-symbol-bounds}). For
$b = 2^N a$, $\|\xi\| \leq 2^j 2^N a$ gives $i_0 \leq j$, and \Cref{lem:dyadic-annulus-count}
with $t = \|\xi\|$ gives $j \leq i_0 + N$.
\end{proof}

\begin{theorem}[The dilated Bessel kernels form a Calder\'on--Zygmund family]
\label{thm:cz-family-bessel}
\lean{Rd.tsum_fourier_dilate_besselKernelFamily_sq_le, Rd.exists_isCZFamily_besselKernelFamily}
\leanok
\uses{lem:bessel-kernel-frequency, def:bessel-overlap, lem:dyadic-annulus-count,
  thm:bessel-kernels-uniform, thm:cz-family-dilation, def:hormander-const}
Let $a > 0$ be real, $N \in \N$, and let $\widehat\omega \colon V \to \C$ be a Schwartz function
that vanishes outside the annulus $\{a \leq \|\zeta\| \leq 2^N a\}$. Then for every real $\sigma$
and every $\xi \in V$,
\[
  \sum_j |\widehat{(K_j)_{2^j}}(\xi)|^2 \leq B_{a,N,\sigma,\widehat\omega} ,
\]
and there is $H < \infty$ such that $((K_j)_{2^j})_j$ is a Calder\'on--Zygmund family with
overlap $B_{a,N,\sigma,\widehat\omega}$ and H\"ormander constant $H$.
\end{theorem}

\begin{proof}
\leanok
By \Cref{lem:bessel-kernel-frequency} every term is at most
$(M_{a,2^N a,\sigma} \|\widehat\omega\|_{0,0})^2$, and a nonzero term needs
$a \leq 2^{-j} \|\xi\| \leq 2^N a$; \Cref{lem:dyadic-annulus-count} with $t = \|\xi\|$ bounds the
sum. The kernels $K_j$
are uniformly Schwartz with some bounds $A$ (\Cref{thm:bessel-kernels-uniform}), and
\Cref{thm:cz-family-dilation} gives the family with $H = 4 H_A$.
\end{proof}

\paragraph{Annuli of the Littlewood--Paley kernels.}

\begin{lemma}[Annuli of the Littlewood--Paley kernels]
\label{lem:littlewood-paley-annuli}
\lean{Rd.mem_annulus_of_fourier_psiOneS_ne_zero, Rd.mem_annulus_of_fourier_psiBase_ne_zero,
  Rd.mem_annulus_of_fourier_psiDualBase_ne_zero}
\leanok
\uses{lem:fourier-dilate, lem:psi-one, lem:psi-base, lem:psi-dual-basic,
  lem:psi-dual-frequency, lem:lp-kernels-frequency}
For every $j \in \N$ and every $\zeta \in V$:
if $\widehat{\psi_1}(\zeta) \neq 0$, then $1 \leq \|\zeta\| \leq 2^2 \cdot 1$; if
$\widehat{G^{\psi}_{j+1}}(\zeta) \neq 0$, then $1/2 \leq \|\zeta\| \leq 2^2 \cdot 1/2$; and if
$\widehat{\tilde G_{j+2}}(\zeta) \neq 0$, then $1/4 \leq \|\zeta\| \leq 2^4 \cdot 1/4$.
\end{lemma}

\begin{proof}
\leanok
The first claim is \Cref{lem:psi-one} with \Cref{lem:lp-kernels-frequency}. Further
$G^{\psi}_{j+1} = G^{\psi}_1$ and $\psi_1 = (G^{\psi}_1)_2$ (\Cref{lem:psi-base}), so
$\widehat{G^{\psi}_{j+1}}(\zeta) = \widehat{\psi_1}(2 \zeta)$ (\Cref{lem:fourier-dilate}), and the
first claim gives $1 \leq 2 \|\zeta\| \leq 4$. Likewise $\tilde G_{j+2} = \tilde G_2$ and
$\tilde\psi_2 = (\tilde G_2)_4$ (\Cref{lem:psi-dual-basic}), so
$\widehat{\tilde G_{j+2}}(\zeta) \neq 0$ gives $1 < 4 \|\zeta\| < 16$
(\Cref{lem:fourier-dilate}, \Cref{lem:psi-dual-frequency}).
\end{proof}

So \Cref{thm:cz-family-bessel} applies to $\widehat{\psi_1}$ with $a = 1$, $N = 2$, to
$\widehat{G^{\psi}_{j+1}}$ with $a = 1/2$, $N = 2$, and to $\widehat{\tilde G_{j+2}}$ with
$a = 1/4$, $N = 4$.

\section{Bessel potential spaces as Triebel--Lizorkin spaces}
\label{sec:bessel-potential-lifting}

Mathlib's Bessel potential space $H^{s,p}(V)$, $s \in \R$, consists of the tempered
distributions $f$ with $J_s f \in L_p(V)$, where $J_s$ multiplies the Fourier transform by
$\langle \xi \rangle^s$ (\texttt{TemperedDistribution.\allowbreak besselPotential} and
\texttt{TemperedDistribution.\allowbreak MemSobolev}). Its bundled form, Mathlib's
\texttt{BesselPotentialSpace}, has the norm $\|f\| = \|J_s f\|_p$. This section proves, for
$n \geq 1$ and $1 < p < \infty$, that $H^{s,p}(V) \cap L_p(V) = \TL{s}{p}{2}$ for every real $s$
and $H^{s,p}(V) = \TL{s}{p}{2}$ for $s \geq 0$ \cite{Triebel}, and hence
$\Sob{m}{p}(V) = H^{m,p}(V)$, Calder\'on's theorem~\cite{Calderon1961} (see also
\cite[Chapter V, \S 3, Theorem 3]{SteinSingular}). The constants come from the uniform Schwartz
bounds of
\Cref{sec:bessel-kernels} and are existential. The kernels are the Bessel potentials of the
blocks $\psi_j$ and of the synthesis kernels
$\tilde\psi_j$: up to the factors $2^{\pm js}$ they are dilates of Bessel kernels, and these form
Calder\'on--Zygmund families. So $J_s f$ is the limit of the synthesis sums
$\sum_{j < J} (J_s \tilde\psi_j) * (\psi_j * f)$, and the weighted blocks $2^{js} \psi_j * f$ of
$f = J_{-s} h$ are the analysis operator of a Calder\'on--Zygmund family applied to $h$. The
first two paragraphs hold for every even Fourier multiplier $m(D)$ of
\Cref{sec:sobolev-two-littlewood-paley}.

\paragraph{Even Fourier multipliers, translations and reflections.}

\begin{lemma}[Even Fourier multipliers commute with translations and reflections]
\label{lem:bessel-schwartz-translation}
\lean{Rd.fourier_besselSchwartz_apply, Rd.fourier_compSubConstCLM_apply,
  Rd.fourierMultiplierCLM_compSubConstCLM, Rd.besselSchwartz_compSubConstCLM,
  Rd.fourierMultiplierCLM_schwartzReflect, Rd.fourierMultiplierCLM_shiftReflect}
\leanok
\uses{def:bessel-schwartz, lem:bessel-weight-basic, lem:reflect-basic, lem:schwartz-shift-reflect,
  def:schwartz-dilate, lem:fourier-multiplier-fourier}
Let $m$ be a function on $V$, smooth with polynomially bounded derivatives. For all Schwartz
functions $\varphi, \eta$ and all $x, y, \xi \in V$,
\[
  \widehat{\varphi(\cdot - y)}(\xi) = e^{-2\pi i \langle y, \xi \rangle} \,
  \widehat\varphi(\xi), \qquad
  m(D) \bigl(\varphi(\cdot - y)\bigr) = (m(D) \varphi)(\cdot - y),
\]
in particular $\widehat{J_t \varphi} = \langle \cdot \rangle^t \widehat\varphi$ and
$J_t (\varphi(\cdot - y)) = (J_t \varphi)(\cdot - y)$ for every $t \in \R$. If $m$ is even, then
\[
  m(D) \bigl(\varphi(-\,\cdot)\bigr) = (m(D) \varphi)(-\,\cdot), \qquad
  m(D) \bigl(\eta(x - \cdot)\bigr) = (m(D) \eta)(x - \cdot) .
\]
\end{lemma}

\begin{proof}
\leanok
The first identity is the substitution $z \mapsto z + y$ (Mathlib's
\texttt{VectorFourier.\allowbreak fourierIntegral\_\allowbreak comp\_\allowbreak add\_\allowbreak
right}). The Fourier transform is injective on Schwartz functions, and by
\Cref{lem:fourier-multiplier-fourier} the Fourier transforms of the two sides of the second
identity are both $m(\xi) e^{-2\pi i \langle y, \xi \rangle} \widehat\varphi(\xi)$, those
of the third both $m(\xi) \widehat\varphi(-\xi) = m(-\xi) \widehat\varphi(-\xi)$
(\Cref{lem:reflect-basic}). The function $\eta(x - \cdot)$ is the translate by $x$ of
$\eta(-\,\cdot)$. For $J_t$, $m$ is the Bessel weight (\Cref{def:bessel-schwartz}).
\end{proof}

\paragraph{Even Fourier multipliers and convolutions.}

\begin{lemma}[Convolutions with Fourier multipliers, pointwise]
\label{lem:bessel-potential-convolution-pointwise}
\lean{Rd.conv_eq_toTemperedDistribution_shiftReflect, Rd.conv_fourierMultiplierCLM_eq,
  Rd.conv_eq_conv_fourierMultiplierCLM, Rd.conv_eq_conv_besselSchwartz}
\leanok
\uses{lem:schwartz-shift-reflect, lem:even-multiplier-transpose,
  lem:bessel-schwartz-translation, lem:bessel-weight-basic}
Let $1 \leq q, r \leq \infty$, let $m$ be an even function on $V$, smooth with polynomially bounded
derivatives, $\eta$ a Schwartz function and $x \in V$. For every $g \in L_q(V)$,
$(\eta * g)(x)$ is the tempered distribution of $g$ at $\eta(x - \cdot)$. For every
$h \in L_r(V)$,
\[
  \bigl((m(D) \eta) * h\bigr)(x) = (m(D) h)\bigl(\eta(x - \cdot)\bigr) .
\]
For all $g \in L_q(V)$ and $h \in L_r(V)$ with $g = m(D) h$ as tempered distributions,
$(\eta * g)(x) = ((m(D) \eta) * h)(x)$; in particular, for every $t \in \R$ and all
$g \in L_q(V)$, $h \in L_r(V)$ with $g = J_t h$, $(\eta * g)(x) = ((J_t \eta) * h)(x)$.
\end{lemma}

\begin{proof}
\leanok
The first claim is \Cref{lem:schwartz-shift-reflect}. For the second,
$(m(D) h)(\eta(x - \cdot)) = h(m(D)(\eta(x - \cdot))) = h((m(D) \eta)(x - \cdot))$ by
\Cref{lem:even-multiplier-transpose} and \Cref{lem:bessel-schwartz-translation}, and this is
$((m(D) \eta) * h)(x)$ by the first claim. The third combines the first two, and the Bessel
weight is even (\Cref{lem:bessel-weight-basic}).
\end{proof}

\begin{theorem}[The Bessel potential of a convolution]
\label{thm:bessel-potential-convolution}
\lean{Rd.integral_besselSchwartz_mul_conv, Rd.besselPotential_toLp_conv}
\leanok
\uses{thm:convolution-adjoint, lem:reflect-basic, lem:bessel-potential-transpose,
  lem:bessel-schwartz-translation, thm:young}
For every $1 \leq p < \infty$, every $F \in L_p(V)$, every $t \in \R$ and all Schwartz functions
$\eta, u$,
\[
  \int_V J_t u \, (\eta * F) = \int_V u \, \bigl((J_t \eta) * F\bigr) ,
\]
and $J_t(\eta * F) = (J_t \eta) * F$ as tempered distributions.
\end{theorem}

\begin{proof}
\leanok
Schwartz functions lie in $L_q(V)$ for the conjugate exponent $q$ of $p$, so by
\Cref{thm:convolution-adjoint} the two sides are $\int_V F(y) (\eta^\vee * J_t u)(y) \, dy$ and
$\int_V F(y) ((J_t \eta)^\vee * u)(y) \, dy$. By \Cref{lem:reflect-basic},
\Cref{lem:bessel-potential-transpose} and \Cref{lem:bessel-schwartz-translation},
\begin{align*}
  (\eta^\vee * J_t u)(y) &= \int_V \eta(z - y) \, J_t u(z) \, dz \\
  &= \int_V (J_t \eta)(z - y) \, u(z) \, dz = \bigl((J_t \eta)^\vee * u\bigr)(y) .
\end{align*}
Both convolutions lie in $L_p(V)$ by Young's inequality (\Cref{thm:young}), and the identity of
tempered distributions is the first identity with \Cref{lem:bessel-potential-transpose}.
\end{proof}

\begin{lemma}[The Bessel potential of a limit of synthesis sums]
\label{lem:bessel-potential-limits}
\lean{Rd.besselPotential_eq_of_tendsto}
\leanok
\uses{thm:bessel-potential-convolution}
Let $1 \leq p < \infty$, $t \in \R$, and for every $j \in \N$ let $\eta_j$ be a Schwartz
function and $F_j \in L_p(V)$. For all $f, g \in L_p(V)$ with
$\sum_{j < J} \eta_j * F_j \to f$ and $\sum_{j < J} (J_t \eta_j) * F_j \to g$ in $L_p(V)$ as
$J \to \infty$, $J_t f = g$ as tempered distributions.
\end{lemma}

\begin{proof}
\leanok
The partial sums satisfy the identity by \Cref{thm:bessel-potential-convolution} and the
linearity of $J_t$. The embedding of $L_p(V)$ into the tempered distributions and $J_t$ are
continuous, and limits of tempered distributions are unique.
\end{proof}

\paragraph{The Bessel kernels of the Littlewood--Paley kernels.}

\begin{definition}[Bessel kernels of a family]
\label{def:bessel-base}
\lean{Rd.besselBase}
\leanok
\uses{def:bessel-kernel-family}
For every real $\sigma$, every family $\omega = (\omega_j)_{j \in \N}$ of Schwartz functions and
every $j \in \N$, let $K^\sigma_j[\omega]$ be the $j$-th Bessel kernel of $\widehat{\omega_j}$
(\Cref{def:bessel-kernel-family}):
\[
  K^\sigma_j[\omega] = \mathcal F^{-1} \bigl( \beta_{2^{-2j},\sigma} \, \widehat{\omega_j}
  \bigr) .
\]
\end{definition}

By \Cref{thm:bessel-potential-dilate},
$J_\sigma (\omega_j)_{2^j} = 2^{j\sigma} (K^\sigma_j[\omega])_{2^j}$. For the bases $G^{\psi}$ and
$\tilde G$ of \Cref{def:psi-base} and \Cref{def:psi-dual}, $(G^{\psi}_j)_{2^j} = \psi_j$ and
$(\tilde G_j)_{2^j} = \tilde\psi_j$.

\begin{theorem}[The Bessel kernels of the Littlewood--Paley kernels]
\label{thm:cz-family-bessel-base}
\lean{Rd.exists_isCZFamily_besselBase, Rd.exists_isCZFamily_besselBase_psiDualBase,
  Rd.exists_isCZFamily_besselBase_psiBase}
\leanok
\uses{def:bessel-base, thm:bessel-kernels-uniform, thm:cz-family-bessel, def:bessel-overlap,
  lem:cz-family-exceptions, lem:littlewood-paley-annuli, def:psi-base, def:psi-dual}
Let $a > 0$ be real, $M, N \in \N$, and let $(\omega_j)$ be Schwartz functions with
$\omega_j = \omega_M$ for every $j \geq M$ such that $\widehat{\omega_M}$ vanishes outside the
annulus $\{a \leq \|\zeta\| \leq 2^N a\}$. Then for every real $\sigma$ there are
$B, H < \infty$ such that $((K^\sigma_j[\omega])_{2^j})_j$ is a Calder\'on--Zygmund family with
overlap $B$ and H\"ormander constant $H$. In particular, for every real $\sigma$ there are such
$B, H$ for $\omega = G^{\psi}$ and for $\omega = \tilde G$.
\end{theorem}

\begin{proof}
\leanok
For $j \geq M$, $K^\sigma_j[\omega]$ is the Bessel kernel $K_j$ of $\widehat{\omega_M}$. These
kernels are uniformly Schwartz, and their dilates have the overlap
$B_{a,N,\sigma,\widehat{\omega_M}}$ (\Cref{thm:bessel-kernels-uniform},
\Cref{thm:cz-family-bessel}, \Cref{def:bessel-overlap}); \Cref{lem:cz-family-exceptions} adds the
first $M$ kernels. The family $G^{\psi}$ is constant from $j = 1$ on, with $a = 1/2$ and
$N = 2$, and $\tilde G$ from $j = 2$ on, with $a = 1/4$ and $N = 4$ (\Cref{def:psi-base},
\Cref{def:psi-dual}, \Cref{lem:littlewood-paley-annuli}).
\end{proof}

\begin{lemma}[Convolutions with Bessel potentials of dilates]
\label{lem:bessel-dilate-convolution}
\lean{Rd.conv_besselSchwartz_schwartzDilate}
\leanok
\uses{def:bessel-base, thm:bessel-potential-dilate, def:schwartz-dilate}
For every $j \in \N$, every real $\sigma$, every family $(\omega_j)$ of Schwartz functions,
every $F \colon V \to \C$ and every $x \in V$,
\[
  \bigl( (J_\sigma (\omega_j)_{2^j}) * F \bigr)(x)
  = \bigl( (K^\sigma_j[\omega])_{2^j} * (2^{j\sigma} F) \bigr)(x) .
\]
\end{lemma}

\begin{proof}
\leanok
\Cref{thm:bessel-potential-dilate} and the linearity of the convolution.
\end{proof}

\begin{theorem}[The synthesis bound for Bessel potentials of dilates]
\label{thm:bessel-synthesis-bound}
\lean{Rd.eLpNorm_sum_conv_besselSchwartz_le}
\leanok
\uses{lem:bessel-dilate-convolution, thm:bessel-potential-dilate, def:cz-const,
  lem:scaled-synthesis-bound}
Let $n \geq 1$, $1 < p < \infty$, $\sigma$ real, and let $(\omega_j)$ be Schwartz functions such
that $((K^\sigma_j[\omega])_{2^j})_j$ is a Calder\'on--Zygmund family with overlap $B$ and
H\"ormander constant $H$. Then for every finite set $t \subseteq \N$ and all $g_j \in L_p(V)$,
\[
  \Bigl\| \sum_{j \in t} \bigl(J_\sigma (\omega_j)_{2^j}\bigr) * g_j \Bigr\|_p
  \leq S_{V,p}(B, H) \, \Bigl\| x \mapsto \bigl\| \bigl(\mathbf 1_t(j) \, 2^{j\sigma}
  |g_j(x)|\bigr)_j \bigr\|_{\ell_2} \Bigr\|_p .
\]
\end{theorem}

\begin{proof}
\leanok
By \Cref{lem:bessel-dilate-convolution} the sum is
$\sum_{j \in t} (K^\sigma_j[\omega])_{2^j} * ((2^j)^\sigma g_j)$, and $|(2^j)^\sigma| = 2^{j\sigma}$
(\Cref{thm:bessel-potential-dilate}); \Cref{lem:scaled-synthesis-bound} with
$c_j = (2^j)^\sigma$ applies.
\end{proof}

\paragraph{Triebel--Lizorkin spaces in Bessel potential spaces.}

\begin{theorem}[$\TL{s}{p}{2} \subseteq H^{s,p}$]
\label{thm:triebel-lizorkin-to-bessel}
\lean{Rd.exists_besselPotential_eq_of_memTriebelLizorkin,
  Rd.memSobolev_bessel_of_memTriebelLizorkin}
\leanok
\uses{thm:cz-family-bessel-base, thm:bessel-synthesis-bound, lem:synthesis-sums-convergent,
  lem:bessel-potential-limits, lem:square-function-reproducing, lem:psi-dual-derivatives,
  lem:tl-pointwise-blocks}
Let $n \geq 1$ and $1 < p < \infty$. For every real $s$ there is $C < \infty$ such that for every
$f \in \TL{s}{p}{2}$ there is $h \in L_p(V)$ with $J_s f = h$ as tempered distributions and
$\|h\|_p \leq C \|f\|_{F^s_{p,2}}$. In particular $f \in H^{s,p}(V)$ for every
$f \in \TL{s}{p}{2}$.
\end{theorem}

\begin{proof}
\leanok
Let $B, H$ be the constants of \Cref{thm:cz-family-bessel-base} for $\omega = \tilde G$ and
$\sigma = s$, and $C = S_{V,p}(B, H)$. Since $(\tilde G_j)_{2^j} = \tilde\psi_j$
(\Cref{lem:psi-dual-derivatives}), \Cref{thm:bessel-synthesis-bound} for the blocks
$g_j = \psi_j * f \in L_p(V)$ (\Cref{lem:tl-pointwise-blocks}) gives the hypothesis of
\Cref{lem:synthesis-sums-convergent} for $K_j = J_s \tilde\psi_j$, so by its item (3)
$\sum_{j < J} (J_s \tilde\psi_j) * (\psi_j * f)$ converges to some $h \in L_p(V)$ with
$\|h\|_p \leq C \|f\|_{F^s_{p,2}}$. The sums $\sum_{j < J} \tilde\psi_j * (\psi_j * f)$ tend to
$f$ (\Cref{lem:square-function-reproducing}), and \Cref{lem:bessel-potential-limits} gives
$J_s f = h$.
\end{proof}

\paragraph{Bessel potential spaces in Triebel--Lizorkin spaces.}

\begin{lemma}[The blocks of a Bessel potential]
\label{lem:bessel-potential-blocks}
\lean{Rd.dyadicWeight_mul_enorm_blockFun_eq}
\leanok
\uses{lem:bessel-potential-convolution-pointwise, lem:bessel-dilate-convolution, lem:psi-base,
  def:bessel-base, def:triebel-lizorkin}
Let $1 \leq q, r \leq \infty$ and $s$ real. For all $g \in L_q(V)$ and $h \in L_r(V)$ with
$g = J_{-s} h$ as tempered distributions, every $j \in \N$ and every $x \in V$,
\[
  2^{js} \, |(\psi_j * g)(x)| = \bigl| \bigl( (K^{-s}_j[G^{\psi}])_{2^j} * h \bigr)(x) \bigr| .
\]
\end{lemma}

\begin{proof}
\leanok
$\psi_j = (G^{\psi}_j)_{2^j}$ (\Cref{lem:psi-base}), so
$(\psi_j * g)(x) = ((J_{-s} \psi_j) * h)(x)$ (\Cref{lem:bessel-potential-convolution-pointwise})
$= 2^{-js} ((K^{-s}_j[G^{\psi}])_{2^j} * h)(x)$ (\Cref{lem:bessel-dilate-convolution}).
\end{proof}

\begin{theorem}[$H^{s,p} \subseteq \TL{s}{p}{2}$ on $L_p$]
\label{thm:bessel-to-triebel-lizorkin-lp}
\lean{Rd.exists_triebelLizorkinNorm_le_of_besselPotential_eq}
\leanok
\uses{lem:bessel-potential-blocks, thm:cz-family-bessel-base, thm:vector-valued-strong-type,
  def:cz-const}
Let $n \geq 1$ and $1 < p < \infty$. For every real $s$ there is $C < \infty$ such that
$\|g\|_{F^s_{p,2}} \leq C \|h\|_p$ for all $g, h \in L_p(V)$ with $J_s g = h$ as tempered
distributions.
\end{theorem}

\begin{proof}
\leanok
Let $B, H$ be the constants of \Cref{thm:cz-family-bessel-base} for $\omega = G^{\psi}$ and
$\sigma = -s$, and $C = S_{V,p}(B, H)$. Since $g = J_{-s} h$ (Mathlib's
\texttt{TemperedDistribution.\allowbreak besselPotential\_\allowbreak neg\_\allowbreak
apply\_\allowbreak eq\_\allowbreak iff}), \Cref{lem:bessel-potential-blocks} writes the weighted
blocks of $g$ as the analysis operator of the family $((K^{-s}_j[G^{\psi}])_{2^j})_j$ at $h$,
and \Cref{thm:vector-valued-strong-type} bounds it.
\end{proof}

\begin{theorem}[$H^{s,p} \subseteq \TL{s}{p}{2}$ for $s \geq 0$]
\label{thm:bessel-to-triebel-lizorkin}
\lean{Rd.exists_memTriebelLizorkin_of_besselPotential_eq}
\leanok
\uses{thm:triebel-lizorkin-to-bessel, thm:bessel-to-triebel-lizorkin-lp,
  thm:square-function-upper, prop:tl-elementary}
Let $n \geq 1$, $1 < p < \infty$ and $s \geq 0$ real. There is $C < \infty$ such that for every
tempered distribution $f$ and every $h \in L_p(V)$ with $J_s f = h$ there is $g \in L_p(V)$ with
$g = f$ as tempered distributions, $g \in \TL{s}{p}{2}$ and $\|g\|_{F^s_{p,2}} \leq C \|h\|_p$.
\end{theorem}

\begin{proof}
\leanok
Let $C$ be the constant of \Cref{thm:bessel-to-triebel-lizorkin-lp}. Since
$h \in L_p(V) = \TL{0}{p}{2}$ (\Cref{thm:square-function-upper}) and $-s \leq 0$,
$h \in \TL{-s}{p}{2}$ (\Cref{prop:tl-elementary}), so \Cref{thm:triebel-lizorkin-to-bessel} for
$-s$ gives $g \in L_p(V)$ with $J_{-s} h = g$. Then $g = J_{-s} J_s f = f$ (Mathlib's
\texttt{TemperedDistribution.\allowbreak besselPotential\_\allowbreak neg\_\allowbreak
apply\_\allowbreak eq\_\allowbreak iff}), $J_s g = h$, and
\Cref{thm:bessel-to-triebel-lizorkin-lp} bounds $\|g\|_{F^s_{p,2}}$.
\end{proof}

On $L_p(V)$ the two spaces agree for every real $s$ (\Cref{cor:bessel-triebel-lizorkin}). For
$s < 0$, $H^{s,p}(V)$ also contains tempered distributions outside $L_p(V)$, and
\Cref{thm:bessel-to-triebel-lizorkin} is the statement for $s \geq 0$, where every element of
$H^{s,p}(V)$ is an $L_p$ function.

\paragraph{The bundled spaces and Calder\'on's theorem.}

\begin{corollary}[The bundled spaces]
\label{cor:bessel-bundled}
\lean{Rd.exists_enorm_toBesselPotentialSpace_le, Rd.exists_memTriebelLizorkin_toDistr}
\leanok
\uses{thm:triebel-lizorkin-to-bessel, thm:bessel-to-triebel-lizorkin,
  lem:lp-distribution-determined, lem:bessel-space-enorm}
Let $n \geq 1$ and $1 < p < \infty$. For every real $s$ there is $C < \infty$ such that every
$f \in \TL{s}{p}{2}$, as an element of Mathlib's $H^{s,p}(V)$, has
$\|f\|_{H^{s,p}} \leq C \|f\|_{F^s_{p,2}}$. For every real $s \geq 0$ there is $C < \infty$ such
that for every $F \in H^{s,p}(V)$ there is $g \in \TL{s}{p}{2}$ with $g = F$ as tempered
distributions and $\|g\|_{F^s_{p,2}} \leq C \|F\|_{H^{s,p}}$.
\end{corollary}

\begin{proof}
\leanok
The norm of $F \in H^{s,p}(V)$ is $\|J_s F\|_p$ (\Cref{lem:bessel-space-enorm}), and the
$L_p$ function $J_s f$ of $f \in \TL{s}{p}{2}$ is the function $h$ of
\Cref{thm:triebel-lizorkin-to-bessel}, since both represent the same tempered distribution
(\Cref{lem:lp-distribution-determined}). The second claim is
\Cref{thm:bessel-to-triebel-lizorkin} for the distribution $F$ and its $L_p$ function $J_s F$.
\end{proof}

\begin{corollary}[$\TL{s}{p}{2} = H^{s,p} \cap L_p$]
\label{cor:bessel-triebel-lizorkin}
\lean{Rd.memTriebelLizorkin_iff_memSobolev_bessel}
\leanok
\uses{thm:triebel-lizorkin-to-bessel, thm:bessel-to-triebel-lizorkin-lp}
For every $n \geq 1$, every $1 < p < \infty$, every real $s$ and every $f \in L_p(V)$,
$f \in \TL{s}{p}{2}$ if and only if $f \in H^{s,p}(V)$.
\end{corollary}

\begin{proof}
\leanok
\Cref{thm:triebel-lizorkin-to-bessel}; conversely, if $J_s f = h \in L_p(V)$, then
$\|f\|_{F^s_{p,2}} \leq C \|h\|_p < \infty$ (\Cref{thm:bessel-to-triebel-lizorkin-lp}).
\end{proof}

\begin{corollary}[$\Sob{m}{p} = H^{m,p}$]
\label{cor:sobolev-bessel}
\lean{Rd.memSobolev_iff_memSobolev_bessel}
\leanok
\uses{cor:sobolev-triebel-lizorkin, cor:bessel-triebel-lizorkin}
For every $n \geq 1$, every $1 < p < \infty$, every $m \in \N$ and every $f \in L_p(V)$,
$f \in \Sob{m}{p}(V)$ if and only if $f \in H^{m,p}(V)$.
\end{corollary}

\begin{proof}
\leanok
Both are $\TL{m}{p}{2}$ (\Cref{cor:sobolev-triebel-lizorkin},
\Cref{cor:bessel-triebel-lizorkin}).
\end{proof}

For $p = 2$, \Cref{sec:sobolev-two-littlewood-paley} proves $H^s(V) = \TL{s}{2}{2}$ directly by
Plancherel's theorem, with explicit constants (\Cref{thm:sobolev-two-triebel-lizorkin}).

\section{Bessel potential spaces of tempered distributions}
\label{sec:bessel-potential-distributions}

For tempered distributions the quasi-norm of \Cref{def:tl-distr} describes $\TL{s}{p}{2}$ for
every real $s$, and $H^{s,p}(V) = \TL{s}{p}{2}$ holds without the restriction to $L_p$. One
inclusion is the multiplier theorem (\Cref{thm:tl-multiplier}) for the symbol
$\langle \xi \rangle^{-s}$ applied to $L_p = \TL{0}{p}{2}$; the other sums the blocks of $J_s u$,
which are convolutions of the blocks of $u$ with the Bessel potentials $J_s \tilde\psi_j$, by the
synthesis bound of their Calder\'on--Zygmund family. In this section $J_s$ is Mathlib's Bessel
potential, and the constants are existential.

\begin{lemma}[Synthesis sums of $L_p$ families]
\label{lem:synthesis-sums-lp-families}
\lean{Rd.exists_tendsto_sum_convLp}
\leanok
\uses{lem:conv-op, lem:mixed-norm-tails, lem:mixed-norm-tools, lem:conv-fourier-multiplier}
Let $1 \leq p < \infty$, $\sigma \in \R$, $C < \infty$, $g_j \in L_p(V)$ with
$\| x \mapsto \| (2^{j\sigma} |g_j(x)|)_j \|_{\ell_2} \|_p < \infty$, and $K_j$ Schwartz functions
with
\[
  \Bigl\| \sum_{j \in t} K_j * g_j \Bigr\|_p \leq C \, \Bigl\| x \mapsto \bigl\|
  \bigl(\mathbf 1_t(j) \, 2^{j\sigma} |g_j(x)|\bigr)_j \bigr\|_{\ell_2} \Bigr\|_p
\]
for every finite set $t \subseteq \N$. Then the partial sums $\sum_{j < J} K_j * g_j$ converge in
$L_p(V)$ to some $h$ with
$\|h\|_p \leq C \, \| x \mapsto \| (2^{j\sigma} |g_j(x)|)_j \|_{\ell_2} \|_p$.
\end{lemma}

\begin{proof}
\leanok
The convolutions lie in $L_p(V)$ (\Cref{lem:conv-op}). For $N \leq N'$ the difference of the
partial sums at $N'$ and at $N$ is the sum over $t = \{N, \dots, N' - 1\}$, so its norm is at most
$C$ times the mixed norm of the weighted family with the terms $j < N$ set to $0$. These tails
tend to $0$ (\Cref{lem:mixed-norm-tails}), so the partial sums form a Cauchy sequence
(\Cref{lem:mixed-norm-tools}) and converge in the complete space $L_p(V)$. Every partial sum has
norm at most $C \, \| x \mapsto \| (2^{j\sigma} |g_j(x)|)_j \|_{\ell_2} \|_p$, and so has the limit.
\end{proof}

For the blocks $g_j = \psi_j * f$ of $f \in L_p(V)$ this is \Cref{lem:synthesis-sums-convergent}.

\begin{lemma}[The blocks of a Bessel potential]
\label{lem:bessel-potential-distrib-blocks}
\lean{Rd.distribBlock_besselPotential}
\leanok
\uses{lem:distrib-block-multiplier, lem:bessel-weight-basic}
Let $1 \leq p < \infty$ and $s \in \R$. For every tempered distribution $u$ on $V$, every
$j \in \N$ and every $g \in L_p(V)$ with $\Delta_j u = g$, the block $\Delta_j(J_s u)$ is the
tempered distribution of $(J_s \tilde\psi_j) * g \in L_p(V)$.
\end{lemma}

\begin{proof}
\leanok
$J_s$ is the Fourier multiplier with the Bessel weight $\langle \xi \rangle^s$, which is smooth
with polynomially bounded derivatives (\Cref{lem:bessel-weight-basic}); apply
\Cref{lem:distrib-block-multiplier}.
\end{proof}

\begin{theorem}[$H^{s,p} \subseteq \TL{s}{p}{2}$ for tempered distributions]
\label{thm:bessel-to-triebel-lizorkin-distr}
\lean{Rd.exists_triebelLizorkinNormDistr_le_of_besselPotential_eq}
\leanok
\uses{thm:tl-multiplier, lem:symbol-examples, thm:tl-distr-lp, thm:square-function-upper}
Let $n \geq 1$ and $1 < p < \infty$. For every real $s$ there is $C < \infty$ such that for every
tempered distribution $u$ on $V$ and every $h \in L_p(V)$ with $J_s u = h$,
\[
  \|u\|_{F^s_{p,2}} \leq C \, \|h\|_p .
\]
\end{theorem}

\begin{proof}
\leanok
Since $J_s u = h$, $u = J_{-s} h$ (Mathlib's \texttt{TemperedDistribution.\allowbreak
besselPotential\_\allowbreak neg\_\allowbreak apply\_\allowbreak eq\_\allowbreak iff}), the
Fourier multiplier of $h$ with the symbol $\langle \xi \rangle^{-s}$ of order $-s$
(\Cref{lem:symbol-examples}). \Cref{thm:tl-multiplier} with the smoothness $0$ gives $C_1 < \infty$
with $\|u\|_{F^s_{p,2}} \leq C_1 \|h\|_{F^0_{p,2}}$, the quasi-norm of $h$ is that of
\Cref{def:triebel-lizorkin} (\Cref{thm:tl-distr-lp}), and
$\|h\|_{F^0_{p,2}} \leq Q_{V,p} \|h\|_p$ (\Cref{thm:square-function-upper}). Take
$C = C_1 Q_{V,p}$.
\end{proof}

\begin{theorem}[$\TL{s}{p}{2} \subseteq H^{s,p}$ for tempered distributions]
\label{thm:triebel-lizorkin-distr-to-bessel}
\lean{Rd.exists_besselPotential_eq_of_memTriebelLizorkinDistr,
  Rd.memSobolev_bessel_of_memTriebelLizorkinDistr}
\leanok
\uses{lem:tl-distr-eq, lem:bessel-potential-distrib-blocks, lem:synthesis-sums-lp-families,
  thm:cz-family-bessel-base, thm:bessel-synthesis-bound, lem:distrib-block-limit}
Let $n \geq 1$ and $1 < p < \infty$. For every real $s$ there is $C < \infty$ such that for every
tempered distribution $u \in \TL{s}{p}{2}$ there is $h \in L_p(V)$ with $J_s u = h$ as tempered
distributions and $\|h\|_p \leq C \|u\|_{F^s_{p,2}}$. In particular $u \in H^{s,p}(V)$ for every
tempered distribution $u \in \TL{s}{p}{2}$.
\end{theorem}

\begin{proof}
\leanok
Let $B, H < \infty$ be such that $((K^s_j[\tilde G])_{2^j})_j$ is a Calder\'on--Zygmund family
with overlap $B$ and H\"ormander constant $H$ (\Cref{thm:cz-family-bessel-base}), and
$C = S_{V,p}(B, H)$. For $u \in \TL{s}{p}{2}$ the blocks are $L_p$ functions $g_j$ and
$\|u\|_{F^s_{p,2}} = \| x \mapsto \| (2^{js} |g_j(x)|)_j \|_{\ell_2} \|_p < \infty$
(\Cref{lem:tl-distr-eq}). The blocks of $J_s u$ are $(J_s \tilde\psi_j) * g_j$, and their sums
over finite sets satisfy the synthesis bound with the constant $C$
(\Cref{lem:bessel-potential-distrib-blocks}, \Cref{thm:bessel-synthesis-bound} with
$\omega = \tilde G$). By \Cref{lem:synthesis-sums-lp-families} the partial sums
$\sum_{j < J} (J_s \tilde\psi_j) * g_j$ converge in $L_p(V)$ to some $h$ with
$\|h\|_p \leq C \|u\|_{F^s_{p,2}}$, and $J_s u = h$ by \Cref{lem:distrib-block-limit}.
\end{proof}

\begin{theorem}[$H^{s,p} = \TL{s}{p}{2}$ for every real $s$]
\label{thm:bessel-triebel-lizorkin-distr}
\lean{Rd.memSobolev_bessel_iff_memTriebelLizorkinDistr,
  Rd.exists_triebelLizorkinNormDistr_le_besselPotentialSpace,
  Rd.exists_enorm_besselPotentialSpace_le_triebelLizorkinNormDistr}
\leanok
\uses{thm:bessel-to-triebel-lizorkin-distr, thm:triebel-lizorkin-distr-to-bessel,
  lem:bessel-space-enorm, lem:lp-distribution-determined}
Let $n \geq 1$ and $1 < p < \infty$. For every real $s$ and every tempered distribution $u$ on
$V$, $u \in H^{s,p}(V)$ if and only if $u \in \TL{s}{p}{2}$. For every real $s$ there are
$C, C' < \infty$ such that $\|F\|_{F^s_{p,2}} \leq C \|F\|_{H^{s,p}}$ for every element $F$ of
Mathlib's bundled space $H^{s,p}(V)$, as a tempered distribution, and
$\|u\|_{H^{s,p}} \leq C' \|u\|_{F^s_{p,2}}$ for every tempered distribution
$u \in \TL{s}{p}{2}$, as an element of the bundled space.
\end{theorem}

\begin{proof}
\leanok
The equivalence is \Cref{thm:bessel-to-triebel-lizorkin-distr} and
\Cref{thm:triebel-lizorkin-distr-to-bessel}. The norm of the bundled space is $\|J_s F\|_p$
(\Cref{lem:bessel-space-enorm}), and for $u \in \TL{s}{p}{2}$ the $L_p$ function of the bundled
element is the function $h$ of \Cref{thm:triebel-lizorkin-distr-to-bessel}, since both represent
$J_s u$ (\Cref{lem:lp-distribution-determined}).
\end{proof}

For $f \in L_p(V)$ the theorem contains \Cref{cor:bessel-triebel-lizorkin}, since the two
quasi-norms of $f$ agree (\Cref{thm:tl-distr-lp}).

\section{Bessel potential spaces as interpolation spaces}
\label{sec:bessel-potential-interpolation}

The Bessel potential norm is equivalent to the quasi-norm of $\TL{s}{p}{2}$
(\Cref{lem:bessel-potential-norm-tl}), so the interpolation theorem for Triebel--Lizorkin spaces
(\Cref{thm:tl-interpolate}) identifies the real interpolation spaces of two Bessel potential
spaces, and of $L_p$ and one of them, with Besov spaces, in the norms $\|J_s f\|_p$
themselves; this is~\cite[Theorem~6.4.5(4)]{BerghLofstrom}. In this section $n \geq 1$,
$1 < p < \infty$, $J_s$ is Mathlib's Bessel potential, and $\Besov{s}{p}{q}$ is the
approximation space of the band-limited scheme in $L_p = L_p(V)$.

\begin{definition}[The Bessel potential norm]
\label{def:bessel-potential-norm}
\lean{Rd.besselPotentialNorm, Rd.besselPotentialNorm_le, Rd.besselPotentialNorm_eq,
  Rd.besselPotentialNorm_eq_enorm, Rd.besselPotentialNorm_eq_enorm_toBesselPotentialSpace,
  Rd.besselPotentialNorm_lt_top_iff}
\leanok
\uses{lem:lp-distribution-determined, lem:bessel-space-enorm}
For every $1 \leq p < \infty$, every $s \in \R$ and every $f \in L_p$, $\|f\|_{H^{s,p}}$ is the
infimum of $\|h\|_p$ over all $h \in L_p$ with $J_s f = h$ as tempered distributions. If
$J_s f = h \in L_p$, then $\|f\|_{H^{s,p}} = \|h\|_p$; if $f$ is the tempered distribution of an
element $U$ of Mathlib's $H^{s,p}(V)$, then $\|f\|_{H^{s,p}} = \|U\|$; and
$\|f\|_{H^{s,p}} < \infty$ if and only if $f \in H^{s,p}(V)$.
\end{definition}

The identity $\|f\|_{H^{s,p}} = \|h\|_p$ holds since an $L_p$ function is determined by its
tempered distribution (\Cref{lem:lp-distribution-determined}). So for $f$ in Mathlib's
$H^{s,p}(V)$ the Bessel potential norm is the norm of the bundled space
(\Cref{lem:bessel-space-enorm}).

\begin{lemma}[The Bessel potential norm and $\TL{s}{p}{2}$]
\label{lem:bessel-potential-norm-tl}
\lean{Rd.exists_triebelLizorkinNorm_le_besselPotentialNorm,
  Rd.exists_besselPotentialNorm_le_triebelLizorkinNorm}
\leanok
\uses{def:bessel-potential-norm, thm:triebel-lizorkin-to-bessel,
  thm:bessel-to-triebel-lizorkin-lp}
Let $n \geq 1$ and $1 < p < \infty$. For every real $s$ there are $C, C' < \infty$ such that for
every $f \in L_p$,
\[
  \|f\|_{F^s_{p,2}} \leq C \, \|f\|_{H^{s,p}}
  \quad \text{and} \quad
  \|f\|_{H^{s,p}} \leq C' \, \|f\|_{F^s_{p,2}} .
\]
\end{lemma}

\begin{proof}
\leanok
For each inequality take the larger of $1$ and the constant of
\Cref{thm:bessel-to-triebel-lizorkin-lp}, respectively of
\Cref{thm:triebel-lizorkin-to-bessel}. The first inequality is
\Cref{thm:bessel-to-triebel-lizorkin-lp} for every $h$ with $J_s f = h$, and it holds when there
is none, since then $\|f\|_{H^{s,p}} = \infty$ and $C \geq 1$. The second is
\Cref{thm:triebel-lizorkin-to-bessel} for $f \in \TL{s}{p}{2}$, and the right side is $\infty$
otherwise.
\end{proof}

\begin{definition}[Bessel potential couples]
\label{def:bessel-potential-couples}
\lean{Rd.besselPotentialCouple, Rd.besselPotentialCouple_norm0, Rd.besselPotentialCouple_norm1,
  Rd.lpBesselPotentialCouple, Rd.lpBesselPotentialCouple_norm0,
  Rd.lpBesselPotentialCouple_norm1}
\leanok
\uses{def:couple, def:bessel-potential-norm}
For all $s_0, s_1 \in \R$, the couple $(H^{s_0,p}, H^{s_1,p})$ on $L_p$ is
$(\|\cdot\|_{H^{s_0,p}}, \|\cdot\|_{H^{s_1,p}})$, and the couple $(L_p, H^{s_1,p})$ is
$(\|\cdot\|_p, \|\cdot\|_{H^{s_1,p}})$.
\end{definition}

\begin{theorem}[Bessel potential spaces interpolate to Besov spaces]
\label{thm:bessel-potential-interpolate}
\lean{Rd.memKMethod_besselPotentialCouple_iff, Rd.memKMethod_lpBesselPotentialCouple_iff}
\leanok
\uses{def:bessel-potential-couples, lem:bessel-potential-norm-tl, thm:tl-interpolate}
Let $n \geq 1$, $1 < p < \infty$, $0 < \eta < 1$, $q \in [1, \infty]$ and $f \in L_p$.
\begin{enumerate}
  \item For every $0 < s_0 < s_1$, $f \in (H^{s_0,p}, H^{s_1,p})_{\eta, q}$ if and only if
    $f \in \Besov{(1 - \eta) s_0 + \eta s_1}{p}{q}$.
  \item For every $s_1 > 0$, $f \in (L_p, H^{s_1,p})_{\eta, q}$ if and only if
    $f \in \Besov{\eta s_1}{p}{q}$.
\end{enumerate}
\end{theorem}

\begin{proof}
\leanok
The Bessel potential norms are equivalent to the quasi-norms of $\TL{s_i}{p}{2}$
(\Cref{lem:bessel-potential-norm-tl}); apply parts (3) and (4) of \Cref{thm:tl-interpolate} with
$q_0 = q_1 = 2$.
\end{proof}

\paragraph{The couple of $H^1(V) \hookrightarrow L_2(V)$.}
For $p = 2$ and $s_1 = 1$ the second space of the couple $(L_2, H^{1,2})$ is the image of the
embedding $\iota \colon H^1(V) \to L_2(V)$ of \Cref{def:sobolev-one-to-l2}, with the quotient norm
of the couple of $\iota$ (\Cref{def:map-couple}).

\begin{lemma}[The quotient norm of $\iota$]
\label{lem:sobolev-one-map-couple}
\lean{Rd.norm1_ofMap_toL2, Rd.ofMap_toL2_eq_lpBesselPotentialCouple}
\leanok
\uses{def:map-couple, def:sobolev-one-to-l2, lem:sobolev-one-to-l2, def:bessel-potential-norm,
  def:bessel-potential-couples}
For every $f \in L_2(V)$,
\[
  \inf \{\|U\| : U \in H^1(V), \ \iota U = f\} = \|f\|_{H^{1,2}} ,
\]
and the couple of $\iota$ on $L_2(V)$ is the couple $(L_2, H^{1,2})$ of
\Cref{def:bessel-potential-couples}.
\end{lemma}

\begin{proof}
\leanok
If $\|f\|_{H^{1,2}} < \infty$, then $f$ is the tempered distribution of some $U \in H^1(V)$ with
$\|U\| = \|f\|_{H^{1,2}}$ (\Cref{def:bessel-potential-norm}), and $\iota U = f$
(\Cref{lem:sobolev-one-to-l2}); so the infimum is at most $\|f\|_{H^{1,2}}$, also when
$\|f\|_{H^{1,2}} = \infty$. Conversely, for every $U$ with $\iota U = f$ the tempered distribution
of $f$ is that of $U$ (\Cref{lem:sobolev-one-to-l2}), so $\|U\| = \|f\|_{H^{1,2}}$
(\Cref{def:bessel-potential-norm}). Both couples have the first norm $\|\cdot\|_2$.
\end{proof}

\begin{lemma}[$(L_2(V), H^1(V))_{\theta, q} = B^\theta_{2,q}(V)$]
\label{lem:sobolev-whole-space-interpolation}
\lean{Rd.memKMethod_ofMap_toL2_iff}
\leanok
\uses{lem:sobolev-one-map-couple, thm:bessel-potential-interpolate, thm:rd-besov}
Let $n \geq 1$. For every $0 < \theta < 1$, every $q \in [1, \infty]$, every order $r \in \N$
with $r > \theta$ and every $F \in L_2(V)$, $F$ lies in the real interpolation space
$(L_2(V), \iota(H^1(V)))_{\theta, q}$ of the couple of $\iota$ if and only if
\[
  \|F\|_2 + \bigl\| \bigl(2^{k\theta} \omega_r(F, 2^{-k})_2\bigr)_k \bigr\|_{\ell_q} < \infty .
\]
\end{lemma}

\begin{proof}
\leanok
By \Cref{lem:sobolev-one-map-couple} the couple is $(L_2, H^{1,2})$, whose interpolation space is
$\Besov{\theta}{2}{q}$ (\Cref{thm:bessel-potential-interpolate} with $p = 2$, $s_1 = 1$ and
$\eta = \theta$), and part (1) of \Cref{thm:rd-besov} describes $\Besov{\theta}{2}{q}$ by the moduli
of order $r$.
\end{proof}

\section{Embeddings of Bessel potential spaces}
\label{sec:bessel-potential-embedding}

For $s \geq 0$ every element of $H^{s,p}(V)$ is an $L_p$ function in $\TL{s}{p}{2}$
(\Cref{thm:bessel-to-triebel-lizorkin}), so the critical embedding of
\Cref{thm:triebel-lizorkin-into-lp} applies to it. This gives the Sobolev embedding
$H^{s,p}(V) \hookrightarrow L_r(V)$ for $s \geq n(1/p - 1/r)$ and real $s$. Since
$\Sob{m}{p}(V) = H^{m,p}(V)$ (\Cref{cor:sobolev-bessel}), its case $s = m$ contains the Sobolev
embedding of \Cref{thm:critical-sobolev-embedding}. Stein proves it with the Bessel kernel and
fractional integration~\cite[Chapter~V]{SteinSingular}.

\begin{theorem}[$H^{s,p} \hookrightarrow L_r$]
\label{thm:bessel-potential-embedding}
\lean{Rd.exists_memLp_eq_of_memSobolev_bessel, Rd.exists_enorm_le_of_besselPotentialSpace}
\leanok
\uses{thm:bessel-to-triebel-lizorkin, cor:bessel-bundled, thm:triebel-lizorkin-into-lp,
  cor:triebel-lizorkin-into-lp-le, def:critical-embedding-const, lem:ae-equal-classes-distribution,
  thm:square-function-lower,
  prop:tl-elementary, def:nikolskii-property}
Let $n \geq 1$, $1 < p < \infty$, $p \leq r < \infty$ and $s \geq n(1/p - 1/r)$ real. Every
tempered distribution $f \in H^{s,p}(V)$ is a function $g \in L_r(V)$. There is $C < \infty$ such
that every $F$ in Mathlib's $H^{s,p}(V)$ is a function $g \in L_r(V)$ with
$\|g\|_r \leq C \|F\|_{H^{s,p}}$.
\end{theorem}

\begin{proof}
\leanok
Since $s \geq n(1/p - 1/r) \geq 0$, \Cref{thm:bessel-to-triebel-lizorkin} and
\Cref{cor:bessel-bundled} give an $L_p$ function $g_0 \in \TL{s}{p}{2}$ with $g_0 = f$ as
tempered distributions and $\|g_0\|_{F^s_{p,2}} \leq C_0 \|F\|_{H^{s,p}}$. For the first
claim, \Cref{cor:triebel-lizorkin-into-lp-le} gives $g \in L_r(V)$ with $g = g_0$ almost
everywhere, so $g = f$ as tempered distributions (\Cref{lem:ae-equal-classes-distribution}).
For the bound and $r = p$ take
$g = g_0$, with $\|g_0\|_p \leq \tilde Q_{V,p} \|g_0\|_{F^0_{p,2}}
\leq \tilde Q_{V,p} \|g_0\|_{F^s_{p,2}}$ (\Cref{thm:square-function-lower},
\Cref{prop:tl-elementary}). For $p < r$ and $\alpha = n(1/p - 1/r)$,
$\|g_0\|_{F^\alpha_{p,\infty}} \leq \|g_0\|_{F^s_{p,2}}$ (\Cref{prop:tl-elementary}), and
\Cref{thm:triebel-lizorkin-into-lp} gives $g \in L_r(V)$ with $g = g_0$ almost everywhere and
$\|g\|_r \leq K_{V,p,r} \|g_0\|_{F^\alpha_{p,\infty}}$; by
\Cref{lem:ae-equal-classes-distribution} $g = g_0 = f$ as tempered distributions.
\end{proof}

\chapter{Fourier multipliers}
\label{ch:fourier-multipliers}

In this chapter $V$ is a finite-dimensional real inner product space of dimension $n$ with
Lebesgue measure. A symbol of order $\sigma$ is a smooth function $m$ on $V$ whose derivatives
satisfy $\|D^k m(\xi)\| \leq C_k \langle \xi \rangle^{\sigma - k}$, with the Japanese bracket
$\langle \xi \rangle = (1 + \|\xi\|^2)^{1/2}$; this is H\"ormander's class
$S^\sigma_{1,0}$. \Cref{sec:symbol-classes} sets up the class, shows that
Schwartz functions, the Bessel weights, linear functions and $\|\xi\|^2$ are symbols, and proves
that a symbol rescaled to the dyadic annuli gives uniformly Schwartz kernels:
$m(D) \omega_{2^j} = 2^{j\sigma} (K_j)_{2^j}$ with $(K_j)$ uniformly Schwartz. Convolutions with
such kernels lie below the Hardy--Littlewood maximal function, uniformly in $j$.
\Cref{sec:tl-multipliers} proves with them that a symbol of order $\sigma$ maps $\TL{s}{p}{q}$ of
tempered distributions into $\TL{s-\sigma}{p}{q}$, for every real $s$. With
$L_p = \TL{0}{p}{2}$ and $H^{s,p} = \TL{s}{p}{2}$ (\Cref{thm:bessel-triebel-lizorkin-distr}) this
gives Mihlin's multiplier theorem on $L_p$ for symbols of order $0$
(\Cref{sec:mihlin-multipliers}) and the mapping properties of the Bessel potential spaces
$H^{s,p}$ for all $1 < p < \infty$ (\Cref{sec:bessel-potential-multipliers}). The constants come
from uniform Schwartz bounds and are existential.

\section{Symbols of order \texorpdfstring{$\sigma$}{sigma} and their kernel families}
\label{sec:symbol-classes}

For a function $m$ on $V$ that is smooth with polynomially bounded derivatives, $m(D)$ is Mathlib's
Fourier multiplier on Schwartz functions and on tempered distributions, as in the paragraph before
\Cref{lem:fourier-multiplier-fourier}. The derivative bounds for the powers of the Japanese bracket
below are mirrored from TauCeti (\texttt{Analysis/\allowbreak SpecialFunctions/\allowbreak
JapaneseBracket/\allowbreak Deriv.lean}, commit \texttt{b61a90d}, Apache-2.0), where they serve the
same purpose. TauCeti also proves Mihlin's multiplier theorem on $L_p$ for symbols satisfying
Mihlin's condition with finitely many derivatives, through its own Calder\'on--Zygmund theory;
the class $S^0_{1,0}$ lies in Mihlin's class. Write $\|G\|_{k,l} = \sup_x \|x\|^k \|D^l G(x)\|$
for the seminorms of a Schwartz function $G$.

\paragraph{Powers of the Japanese bracket.}

\begin{lemma}[Derivatives of powers of the Japanese bracket]
\label{lem:bracket-power-derivatives}
\lean{ContinuousLinearMap.norm_iteratedFDeriv_le_one_add_norm_sq_rpow,
  ContinuousLinearMap.exists_norm_iteratedFDeriv_smul_le_one_add_norm_sq_rpow,
  Rd.hasFDerivAt_one_add_norm_sq_rpow, Rd.exists_norm_iteratedFDeriv_one_add_norm_sq_rpow_le}
\leanok
Let $E$, $G$ be real normed spaces.
\begin{enumerate}
  \item For every continuous linear map $L \colon E \to G$, every $j \in \N$ and every
    $x \in E$, $\|D^j L(x)\| \leq \|L\| (1 + \|x\|^2)^{(1-j)/2}$.
  \item For every continuous linear map $L \colon E \to G$, every $k \in \N$, every real $\rho$
    and every $C^k$ function $g \colon E \to \R$ such that for every $i \leq k$ there is $C$
    with $\|D^i g(x)\| \leq C (1 + \|x\|^2)^{\rho - i/2}$ for every $x$, there is $C'$ with
    $\|D^k (g L)(x)\| \leq C' (1 + \|x\|^2)^{\rho + 1/2 - k/2}$ for every $x \in E$.
  \item If $E$ is an inner product space, then for every real $r$ and every $x \in E$ the
    derivative of $y \mapsto (1 + \|y\|^2)^r$ at $x$ is
    $(1 + \|x\|^2)^{r-1} \, 2r \langle x, \cdot \rangle$.
  \item If $E$ is an inner product space, then for every real $r$ and every $k \in \N$ there is
    $C$ with $\|D^k (1 + \|\cdot\|^2)^r(x)\| \leq C (1 + \|x\|^2)^{r - k/2}$ for every $x \in E$.
\end{enumerate}
\end{lemma}

\begin{proof}
\leanok
(1) $\|L x\| \leq \|L\| \, \|x\| \leq \|L\| (1 + \|x\|^2)^{1/2}$, $D L = L$ is constant, and the
higher derivatives vanish. (2) The Leibniz rule (Mathlib's
\texttt{norm\_\allowbreak iteratedFDeriv\_\allowbreak smul\_\allowbreak le}) and (1): the
exponents of $D^i g$ and $D^{k-i} L$ add up to $\rho + 1/2 - k/2$. (3) The chain rule for
$t \mapsto t^r$ and $\|\cdot\|^2$. (4) Induction on $k$: by (3) the derivative is
$(1 + \|\cdot\|^2)^{r-1}$ times the linear map $2r \langle \cdot, \cdot \rangle$, so (2) with
$\rho = r - 1$ and the induction hypothesis for $r - 1$ give the bound for $D^{k+1}$.
\end{proof}

\paragraph{The class.}

\begin{definition}[Symbols of order $\sigma$]
\label{def:symbol-growth}
\lean{Rd.HasSymbolGrowth, Rd.HasSymbolGrowth.contDiff, Rd.HasSymbolGrowth.exists_bound}
\leanok
Let $E$, $F$ be real normed spaces and $\sigma$ real. A function $m \colon E \to F$ is a
\emph{symbol of order $\sigma$} if it is smooth and for every $k \in \N$ there is $C_k$ with
\[
  \|D^k m(\xi)\| \leq C_k (1 + \|\xi\|^2)^{(\sigma - k)/2} \qquad \text{for every } \xi \in E .
\]
The constants $C_k$ can be taken nonnegative.
\end{definition}

\begin{lemma}[Symbols, Schwartz functions and temperate growth]
\label{lem:symbol-temperate-schwartz}
\lean{Rd.one_add_norm_sq_rpow_le_one_add_norm_pow, Rd.HasSymbolGrowth.hasTemperateGrowth,
  Rd.hasSymbolGrowth_schwartz}
\leanok
\uses{def:symbol-growth}
Let $E$, $F$ be real normed spaces.
\begin{enumerate}
  \item For every real $a$, every $N \in \N$ with $a \leq N$ and every $\xi \in E$,
    $(1 + \|\xi\|^2)^{a/2} \leq (1 + \|\xi\|)^N$.
  \item For every real $\sigma$, every symbol $m \colon E \to F$ of order $\sigma$ is smooth
    with polynomially bounded derivatives (Mathlib's
    \texttt{Function.\allowbreak HasTemperateGrowth}).
  \item For every Schwartz function $\varphi \colon E \to F$ and every real $\sigma$, $\varphi$
    is a symbol of order $\sigma$.
\end{enumerate}
\end{lemma}

\begin{proof}
\leanok
(1) $1 + \|\xi\|^2 \geq 1$, so $(1 + \|\xi\|^2)^{a/2} \leq (1 + \|\xi\|^2)^{N/2}$, and
$(1 + \|\xi\|^2)^{1/2} \leq 1 + \|\xi\|$. (2) By (1) with $a = \sigma - k \leq N$.
(3) With $N \geq k - \sigma$, $(1 + \|\xi\|)^N \|D^k \varphi(\xi)\|$ is at most $2^N$ times a
finite maximum of seminorms of $\varphi$ (Mathlib's
\texttt{SchwartzMap.\allowbreak one\_\allowbreak add\_\allowbreak le\_\allowbreak
sup\_\allowbreak seminorm\_\allowbreak apply}), and
$(1 + \|\xi\|^2)^{(k - \sigma)/2} \leq (1 + \|\xi\|)^N$ by (1).
\end{proof}

\begin{lemma}[Algebra of symbols]
\label{lem:symbol-algebra}
\lean{Rd.HasSymbolGrowth.mono, Rd.HasSymbolGrowth.add, Rd.HasSymbolGrowth.neg,
  Rd.HasSymbolGrowth.sub, Rd.hasSymbolGrowth_const, Rd.hasSymbolGrowth_continuousLinearMap,
  Rd.HasSymbolGrowth.mul, Rd.HasSymbolGrowth.const_mul, Rd.HasSymbolGrowth.ofReal}
\leanok
\uses{def:symbol-growth, lem:bracket-power-derivatives}
Let $E$, $F$ be real normed spaces and $\sigma, \tau$ real.
\begin{enumerate}
  \item If $\sigma \leq \tau$, every symbol of order $\sigma$ is a symbol of order $\tau$.
  \item For all symbols $m, m' \colon E \to F$ of order $\sigma$, $m + m'$, $-m$ and $m - m'$
    are symbols of order $\sigma$.
  \item For every $c \in F$, the constant function $c$ is a symbol of order $0$; every
    continuous linear map $L \colon E \to F$ is a symbol of order $1$.
  \item For every real normed algebra $A$, every symbol $m \colon E \to A$ of order $\sigma$,
    every symbol $m' \colon E \to A$ of order $\tau$ and every $c \in A$, $m m'$ is a symbol of
    order $\sigma + \tau$ and $c m$ is a symbol of order $\sigma$.
  \item For every symbol $m \colon E \to \R$ of order $\sigma$, $\xi \mapsto m(\xi) \in \C$ is
    a symbol of order $\sigma$.
\end{enumerate}
\end{lemma}

\begin{proof}
\leanok
(1) $1 + \|\xi\|^2 \geq 1$. (2) Derivatives are linear. (3) The derivatives of positive order of
a constant vanish; continuous linear maps by \Cref{lem:bracket-power-derivatives}~(1).
(4) By the Leibniz rule (Mathlib's
\texttt{norm\_\allowbreak iteratedFDeriv\_\allowbreak mul\_\allowbreak le}),
$\|D^k (m m')(\xi)\| \leq \sum_{i \leq k} \binom{k}{i} C_i C'_{k-i}
(1 + \|\xi\|^2)^{(\sigma - i)/2} (1 + \|\xi\|^2)^{(\tau - k + i)/2}$, and the exponents add up to
$(\sigma + \tau - k)/2$; $c m$ is the product with the constant $c$ of order $0$. (5) The
embedding $\R \to \C$ is a linear isometry, so it keeps the norms of the derivatives (Mathlib's
\texttt{LinearIsometry.\allowbreak norm\_\allowbreak iteratedFDeriv\_\allowbreak
comp\_\allowbreak left}).
\end{proof}

\begin{lemma}[Examples of symbols]
\label{lem:symbol-examples}
\lean{Rd.hasSymbolGrowth_one_add_norm_sq_rpow, Rd.hasSymbolGrowth_besselWeight,
  Rd.hasSymbolGrowth_inner, Rd.hasSymbolGrowth_norm_sq}
\leanok
\uses{def:symbol-growth, lem:bracket-power-derivatives, lem:symbol-algebra, def:bessel-weight}
For every real $r$, $\xi \mapsto (1 + \|\xi\|^2)^r$ is a symbol of order $2r$ on $V$. For every
real $s$, the Bessel weight $\xi \mapsto \langle \xi \rangle^s \in \C$ is a symbol of order $s$.
For every $v \in V$, $\xi \mapsto \langle \xi, v \rangle \in \C$ is a symbol of order $1$, and
$\xi \mapsto \|\xi\|^2 \in \C$ is a symbol of order $2$.
\end{lemma}

\begin{proof}
\leanok
The first claim is \Cref{lem:bracket-power-derivatives}~(4), since $r - k/2 = (2r - k)/2$; the
Bessel weight is the case $r = s/2$, taken into $\C$ by \Cref{lem:symbol-algebra}~(5). The
function $\langle \cdot, v \rangle$ is continuous and linear, and
$\|\xi\|^2 = \langle \xi \rangle^2 - 1$ with the constant $1$ of order $0 \leq 2$
(\Cref{lem:symbol-algebra}).
\end{proof}

The symbols $2\pi i \langle \xi, v \rangle$ and $-4\pi^2 \|\xi\|^2$ give the derivative
$\partial_v$ and the Laplacian of tempered distributions (Mathlib's
\texttt{TemperedDistribution.\allowbreak lineDeriv\_\allowbreak eq\_\allowbreak
fourier\allowbreak MultiplierCLM} and \texttt{TemperedDistribution.\allowbreak
laplacian\_\allowbreak eq\_\allowbreak fourier\allowbreak MultiplierCLM}).

\paragraph{Rescaled symbols.}

\begin{definition}[Rescaled symbols]
\label{def:scaled-symbol}
\lean{Rd.scaledSymbol, Rd.hasTemperateGrowth_scaledSymbol}
\leanok
For every $m \colon V \to \C$, every real $\sigma$ and every $j \in \N$ let
\[
  m^\sigma_j(\zeta) = 2^{-j\sigma} \, m(2^j \zeta) \qquad (\zeta \in V) .
\]
If $m$ is smooth with polynomially bounded derivatives, so is every $m^\sigma_j$.
\end{definition}

\begin{lemma}[Derivatives of rescaled symbols]
\label{lem:scaled-symbol-derivatives}
\lean{Rd.iteratedFDeriv_scaledSymbol, Rd.rpow_neg_mul_pow_mul_one_add_sq_rpow,
  Rd.norm_iteratedFDeriv_scaledSymbol_le}
\leanok
\uses{def:scaled-symbol}
For every smooth $m \colon V \to \C$, every real $\sigma$, all $j, k \in \N$ and every
$\zeta \in V$,
$D^k m^\sigma_j(\zeta) = 2^{-j\sigma} 2^{jk} \, D^k m(2^j \zeta)$. For every real $c > 0$,
every real $\sigma$, every $k \in \N$ and every real $x$,
\[
  c^{-\sigma} c^k (1 + (c x)^2)^{(\sigma - k)/2} = (c^{-2} + x^2)^{(\sigma - k)/2} .
\]
Hence, for every real $C$ with $\|D^k m(\xi)\| \leq C (1 + \|\xi\|^2)^{(\sigma - k)/2}$ for every
$\xi \in V$, $\|D^k m^\sigma_j(\zeta)\| \leq C (2^{-2j} + \|\zeta\|^2)^{(\sigma - k)/2}$ for
every $j \in \N$ and every $\zeta \in V$.
\end{lemma}

\begin{proof}
\leanok
The chain rule for $\zeta \mapsto m(2^j \zeta)$ (Mathlib's
\texttt{iteratedFDeriv\_\allowbreak comp\_\allowbreak const\_\allowbreak smul}). For the
identity, $c^{-2} + x^2 = c^{-2} (1 + (cx)^2)$ and $(c^{-2})^{(\sigma - k)/2} = c^{k - \sigma}$.
The bound combines both with $c = 2^j$ and $x = \|\zeta\|$.
\end{proof}

\begin{lemma}[Rescaled symbols on an annulus]
\label{lem:scaled-symbol-annulus}
\lean{Rd.HasSymbolGrowth.exists_norm_iteratedFDeriv_scaledSymbol_le}
\leanok
\uses{lem:scaled-symbol-derivatives, def:symbol-growth, lem:shifted-bessel-symbol-bounds}
Let $\sigma$ be real, $m \colon V \to \C$ a symbol of order $\sigma$, $a > 0$ and $b$ real. For
every $k \in \N$ there is $C \geq 0$ with $\|D^k m^\sigma_j(\zeta)\| \leq C$ for every $j \in \N$
and every $\zeta \in V$ with $a \leq \|\zeta\| \leq b$.
\end{lemma}

\begin{proof}
\leanok
By \Cref{lem:scaled-symbol-derivatives},
$\|D^k m^\sigma_j(\zeta)\| \leq C_k |\beta_{2^{-2j}, \sigma - k}(\zeta)|$, and
$2^{-2j} \in [0, 1]$, so \Cref{lem:shifted-bessel-symbol-bounds} gives the bound
$C = C_k M_{a,b,\sigma-k}$.
\end{proof}

\paragraph{The symbol and kernel families.}

\begin{definition}[Symbol families and symbol kernels]
\label{def:symbol-family}
\lean{Rd.symbolFamily, Rd.coe_symbolFamily, Rd.symbolKernelFamily,
  Rd.fourier_symbolKernelFamily, Rd.symbolBase}
\leanok
\uses{def:scaled-symbol}
For every $m \colon V \to \C$ that is smooth with polynomially bounded derivatives, every real
$\sigma$ and every Schwartz function $\widehat\omega \colon V \to \C$ let, for $j \in \N$,
\[
  S^m_j = m^\sigma_j \, \widehat\omega , \qquad K^m_j = \mathcal F^{-1} S^m_j ,
\]
Schwartz functions with $\widehat{K^m_j} = m^\sigma_j \, \widehat\omega$. For every family
$\omega = (\omega_j)_{j \in \N}$ of Schwartz functions and every $j \in \N$, let
$K^m_j[\omega]$ be the $j$-th symbol kernel $K^m_j$ of $\widehat{\omega_j}$.
\end{definition}

\begin{theorem}[The symbol kernels are uniformly Schwartz]
\label{thm:symbol-kernels-uniform}
\lean{Rd.HasSymbolGrowth.exists_isUniformSchwartz_symbolFamily,
  Rd.HasSymbolGrowth.exists_isUniformSchwartz_symbolKernelFamily}
\leanok
\uses{def:symbol-family, lem:scaled-symbol-annulus, lem:leibniz-schwartz, lem:tsupport-annulus,
  thm:uniform-schwartz-linear-map, def:uniform-schwartz, lem:symbol-temperate-schwartz}
Let $\sigma$ be real, $m \colon V \to \C$ a symbol of order $\sigma$, $a > 0$ and $b$ real, and
let $\widehat\omega \colon V \to \C$ be a Schwartz function that vanishes outside the annulus
$\{a \leq \|\zeta\| \leq b\}$. Then there are bounds $A$ such that $(S^m_j)$ is uniformly
Schwartz with bounds $A$, and bounds $A'$ such that $(K^m_j)$ is uniformly Schwartz with
bounds $A'$.
\end{theorem}

\begin{proof}
\leanok
As for the Bessel kernels (\Cref{thm:bessel-kernels-uniform}): let $C_r \geq 0$ be the bounds of
\Cref{lem:scaled-symbol-annulus}. Outside the closure of the support of $\widehat\omega$ every
derivative of $S^m_j$ vanishes; on it $a \leq \|x\| \leq b$ (\Cref{lem:tsupport-annulus}), so
$\|x\|^k \|D^l S^m_j(x)\| \leq |b|^k \sum_{i \leq l} \binom{l}{i} C_i
\|\widehat\omega\|_{0,l-i}$ by \Cref{lem:leibniz-schwartz}. The kernels follow by
\Cref{thm:uniform-schwartz-linear-map}.
\end{proof}

\begin{theorem}[Fourier multipliers of dilates]
\label{thm:symbol-dilate}
\lean{Rd.fourierMultiplierCLM_schwartzDilate}
\leanok
\uses{def:symbol-family, def:schwartz-dilate, lem:fourier-dilate, lem:fourier-multiplier-fourier}
For every $m \colon V \to \C$ that is smooth with polynomially bounded derivatives, every real
$\sigma$, every Schwartz function $\omega \colon V \to \C$ and every $j \in \N$, with the symbol
kernels $K^m_j$ of $\widehat\omega$,
\[
  m(D) \, \omega_{2^j} = 2^{j\sigma} (K^m_j)_{2^j} .
\]
\end{theorem}

\begin{proof}
\leanok
Both sides are Schwartz functions, and the Fourier transform is injective on them. At $\xi$ the
Fourier transform of the left side is $m(\xi) \, \widehat\omega(2^{-j} \xi)$, and that of the
right side is $2^{j\sigma} 2^{-j\sigma} m(2^j 2^{-j} \xi) \, \widehat\omega(2^{-j} \xi)$
(\Cref{lem:fourier-dilate}, \Cref{lem:fourier-multiplier-fourier}).
\end{proof}

\begin{corollary}[Symbols applied to the synthesis kernels]
\label{cor:symbol-psi-dual}
\lean{Rd.HasSymbolGrowth.exists_isUniformSchwartz_symbolBase,
  Rd.HasSymbolGrowth.exists_isUniformSchwartz_psiDualS}
\leanok
\uses{thm:symbol-kernels-uniform, thm:symbol-dilate, lem:uniform-schwartz-exceptions,
  lem:littlewood-paley-annuli, def:psi-dual, lem:psi-dual-derivatives}
Let $\sigma$ be real and $m \colon V \to \C$ a symbol of order $\sigma$.
\begin{enumerate}
  \item For every family $(\omega_j)$ of Schwartz functions, every $M \in \N$ with
    $\omega_j = \omega_M$ for $j \geq M$, and all reals $a > 0$, $b$ such that
    $\widehat{\omega_M}$ vanishes outside $\{a \leq \|\zeta\| \leq b\}$, the symbol kernels
    $j \mapsto K^m_j[\omega]$ are uniformly Schwartz with some bounds.
  \item There are a uniformly Schwartz family $(G_j)$ and bounds $A$ such that
    $m(D) \tilde\psi_j = 2^{j\sigma} (G_j)_{2^j}$ for every $j \in \N$, with the synthesis
    kernels $\tilde\psi_j$ of \Cref{def:psi-dual}.
\end{enumerate}
\end{corollary}

\begin{proof}
\leanok
(1) From $M$ on the kernels are those of $\widehat{\omega_M}$
(\Cref{thm:symbol-kernels-uniform}), and the first $M$ are exceptions
(\Cref{lem:uniform-schwartz-exceptions}). (2) $\tilde\psi_j = (\tilde G_j)_{2^j}$ with the bases
$\tilde G_j$ of \Cref{def:psi-dual}, constant from $j = 2$ on, with Fourier transform in
$\{1/4 \leq \|\zeta\| \leq 4\}$ (\Cref{lem:littlewood-paley-annuli}); take
$G_j = K^m_j[\tilde G]$ and apply (1) and \Cref{thm:symbol-dilate}.
\end{proof}

\paragraph{The maximal bound.}

\begin{theorem}[Uniform families below the maximal function]
\label{thm:uniform-family-maximal}
\lean{Rd.IsUniformSchwartz.enorm_conv_dilate_le_maximal,
  Rd.decayMaximalConst_mul_familyDecayConst_ne_top}
\leanok
\uses{def:uniform-schwartz, lem:uniform-schwartz-constants, lem:schwartz-decay,
  thm:decay-kernels-maximal, def:maximal-function, lem:maximal-function-kernels}
Let $(G_j)$ be uniformly Schwartz with bounds $A$ and let $N > n$. For every $j \in \N$, every
function $g \colon V \to \C$ and every $x \in V$,
\[
  |((G_j)_{2^j} \ast g)(x)| \leq C_{n,N} \, C_{A,0,N} \, M|g|(x)
\]
with the constants $C_{n,N}$ of \Cref{lem:maximal-function-kernels} and $C_{A,0,N}$ of
\Cref{def:uniform-schwartz}, independent of $j$; the product $C_{n,N} \, C_{A,0,N}$ is
finite.
\end{theorem}

\begin{proof}
\leanok
$|(G_j)_{2^j}(y)| \leq C_{A,0,N} 2^{jn} (1 + 2^j \|y\|)^{-N}$ (\Cref{lem:schwartz-decay},
\Cref{lem:uniform-schwartz-constants}), and kernels with this bound lie below the maximal
function (\Cref{thm:decay-kernels-maximal}).
\end{proof}

\section{Fourier multipliers on Triebel--Lizorkin spaces}
\label{sec:tl-multipliers}

A symbol $m$ of order $\sigma$ maps $\TL{s}{p}{q}$ into $\TL{s-\sigma}{p}{q}$, for tempered
distributions (\Cref{def:tl-distr}), every real $s$, $1 < p < \infty$ and $1 < q \leq \infty$;
Triebel proves this for the classical spaces~\cite[pp.~73--77]{Triebel}. The proof is
pointwise: since $\widehat{\tilde\psi_j} \, \widehat{\psi_j} = \widehat{\psi_j}$, the block
$\Delta_j(m(D) u)$ is the convolution of $\Delta_j u$ with $m(D) \tilde\psi_j =
2^{j\sigma} (G_j)_{2^j}$ (\Cref{cor:symbol-psi-dual}), whose weighted size lies below the maximal
function of the weighted block (\Cref{thm:uniform-family-maximal}); the Fefferman--Stein
inequality finishes. The constants are existential.

\begin{lemma}[The blocks of a Fourier multiplier]
\label{lem:distrib-block-multiplier}
\lean{Rd.distribBlock_fourierMultiplierCLM, Rd.distribBlock_fourierMultiplierCLM_eq_convLp}
\leanok
\uses{def:distrib-block, lem:psi-dual-reproduces, lem:psi-dual-derivatives,
  lem:conv-fourier-multiplier, lem:fourier-multiplier-fourier}
Let $m \colon V \to \C$ be smooth with polynomially bounded derivatives. For every tempered
distribution $u$ on $V$ and every $j \in \N$,
\[
  \Delta_j (m(D) u) = \bigl(\widehat{m(D) \tilde\psi_j}\bigr)(D) \, \Delta_j u .
\]
For every $1 \leq p < \infty$, every tempered distribution $u$, every $j \in \N$ and every
$g \in L_p(V)$ with $\Delta_j u = g$, $\Delta_j(m(D) u)$ is the tempered distribution of
$(m(D) \tilde\psi_j) * g \in L_p(V)$.
\end{lemma}

\begin{proof}
\leanok
Both sides of the identity are Fourier multipliers of $u$, with the symbols
$m \, \widehat{\psi_j}$ and $\widehat{\psi_j} \, m \, \widehat{\tilde\psi_j}$
(\Cref{lem:fourier-multiplier-fourier}), and these agree since
$\widehat{\tilde\psi_j} \, \widehat{\psi_j} = \widehat{\psi_j}$ (\Cref{lem:psi-dual-reproduces}).
The second claim follows with \Cref{lem:conv-fourier-multiplier}.
\end{proof}

\begin{lemma}[Weighted convolutions below the maximal function]
\label{lem:symbol-kernel-maximal}
\lean{Rd.dyadicWeight_mul_enorm_conv_le_maximal}
\leanok
\uses{thm:uniform-family-maximal, lem:almost-diagonal-pointwise, def:schwartz-dilate,
  thm:bessel-potential-dilate}
Let $(G_j)$ be uniformly Schwartz with bounds $A$, $N > n$, $\sigma$ real, $j \in \N$ and
$K = 2^{j\sigma} (G_j)_{2^j}$. For every real $s$, every $g \colon V \to \C$ and every $x \in V$,
\[
  2^{j(s - \sigma)} |(K * g)(x)| \leq C_{n,N} \, C_{A,0,N} \, M\bigl(2^{js} |g|\bigr)(x) .
\]
\end{lemma}

\begin{proof}
\leanok
$2^{j(s - \sigma)} |(K * g)(x)| = 2^{js} |((G_j)_{2^j} * g)(x)|$, since
$|(2^j)^\sigma| = 2^{j\sigma}$ (\Cref{thm:bessel-potential-dilate}); apply
\Cref{thm:uniform-family-maximal} and move the constant $2^{js}$ into the maximal function
(\Cref{lem:almost-diagonal-pointwise}).
\end{proof}

\begin{theorem}[Fourier multipliers on $\TL{s}{p}{q}$]
\label{thm:tl-multiplier}
\lean{Rd.triebelLizorkinNormDistr_fourierMultiplier_le,
  Rd.exists_triebelLizorkinNormDistr_fourierMultiplier_le,
  Rd.MemTriebelLizorkinDistr.fourierMultiplierCLM}
\leanok
\uses{def:tl-distr, lem:tl-distr-eq, cor:symbol-psi-dual, lem:distrib-block-multiplier,
  lem:symbol-kernel-maximal, lem:maximal-bounds-mixed-norm, thm:fefferman-stein,
  lem:maximal-function-kernels, def:uniform-schwartz, def:symbol-growth}
Let $1 < p < \infty$, $1 < q \leq \infty$, $\sigma$ real and $m \colon V \to \C$ a symbol of order
$\sigma$. For every real $s$ there is $C < \infty$ such that for every tempered distribution $u$
on $V$,
\[
  \|m(D) u\|_{F^{s-\sigma}_{p,q}} \leq C \, \|u\|_{F^s_{p,q}} .
\]
In particular $m(D)$ maps $\TL{s}{p}{q}$ into $\TL{s-\sigma}{p}{q}$. More precisely, for every
$m \colon V \to \C$ that is smooth with polynomially bounded derivatives, every uniformly
Schwartz family $(G_j)$ with bounds $A$ and $m(D) \tilde\psi_j = 2^{j\sigma} (G_j)_{2^j}$ for
every $j$, every $N > n$, every tempered distribution $u$ and every family $(g_j)$ in $L_p(V)$
with $\Delta_j u = g_j$ for every $j$,
\[
  \|m(D) u\|_{F^{s-\sigma}_{p,q}} \leq C_{n,N} \, C_{A,0,N} \, C^{\mathrm{FS}}_{p,q}
  \Bigl\| x \mapsto \bigl\| (2^{js} |g_j(x)|)_j \bigr\|_{\ell_q} \Bigr\|_{L_p}
\]
with the Fefferman--Stein constant $C^{\mathrm{FS}}_{p,q}$ of \Cref{thm:fefferman-stein}.
\end{theorem}

\begin{proof}
\leanok
By \Cref{lem:distrib-block-multiplier} the blocks of $m(D) u$ are the $L_p$ functions
$h_j = (m(D) \tilde\psi_j) * g_j$, and by \Cref{lem:symbol-kernel-maximal}, at almost every $x$,
$2^{j(s - \sigma)} |h_j(x)| \leq C_{n,N} C_{A,0,N} \, M(2^{js} |g_j|)(x)$ for every $j$. Hence
\Cref{lem:tl-distr-eq} and \Cref{lem:maximal-bounds-mixed-norm} give the second bound, which
uses only that $m$ has polynomially bounded derivatives. For a symbol of order $\sigma$ the
family $(G_j)$ exists by \Cref{cor:symbol-psi-dual}; with $N = n + 1$ and
$C = \max(1, C_{n,N} C_{A,0,N} C^{\mathrm{FS}}_{p,q})$, \Cref{lem:tl-distr-eq} gives the first
bound, also when the blocks of $u$ are no $L_p$ functions.
\end{proof}

\begin{corollary}[Fourier multipliers on $\TL{s}{p}{q}$ of $L_p$ functions]
\label{cor:tl-multiplier-lp}
\lean{Rd.exists_triebelLizorkinNorm_le_of_fourierMultiplier_eq}
\leanok
\uses{thm:tl-multiplier, thm:tl-distr-lp}
Let $1 < p < \infty$, $1 < q \leq \infty$, $\sigma$ real and $m \colon V \to \C$ a symbol of order
$\sigma$. For every real $s$ there is $C < \infty$ with
$\|g\|_{F^{s-\sigma}_{p,q}} \leq C \|f\|_{F^s_{p,q}}$ for all $f, g \in L_p(V)$ with
$g = m(D) f$ as tempered distributions.
\end{corollary}

\begin{proof}
\leanok
\Cref{thm:tl-multiplier} with $u = f$ and \Cref{thm:tl-distr-lp} for $f$ and $g$.
\end{proof}

\section{Mihlin's multiplier theorem}
\label{sec:mihlin-multipliers}

Mihlin~\cite{Mihlin1956} proved that a function $m$ with
$|\partial^\alpha m(\xi)| \leq C |\xi|^{-|\alpha|}$ for $\xi \neq 0$ and enough multi-indices
$\alpha$ is a bounded Fourier multiplier on $L_p(\R^n)$, $1 < p < \infty$; H\"ormander weakened
the condition~\cite[Theorem~2.5]{Hormander1960}, see also~\cite[Theorem~6.2.7]{Grafakos}. Symbols
of order $0$ satisfy Mihlin's condition, so the theorem below is a case of Mihlin's theorem. Here
it follows from $L_p = \TL{0}{p}{2}$ and the multiplier theorem on $\TL{0}{p}{2}$. TauCeti proves
Mihlin's theorem for symbols satisfying Mihlin's condition with finitely many derivatives and
Hilbert-valued $L_p$
(\texttt{Analysis/\allowbreak Fourier/\allowbreak MihlinHormander.lean}, commit
\texttt{b61a90d}), through its own Calder\'on--Zygmund theory. In this section $n \geq 1$ and
$1 < p < \infty$.

\begin{theorem}[Mihlin's multiplier theorem for symbols of order $0$]
\label{thm:mihlin-symbols}
\lean{Rd.exists_fourierMultiplier_eq_of_hasSymbolGrowth}
\leanok
\uses{thm:square-function-upper, thm:tl-distr-lp, thm:tl-multiplier,
  thm:triebel-lizorkin-distr-to-bessel}
Let $n \geq 1$, $1 < p < \infty$ and $m \colon V \to \C$ a symbol of order $0$. There is
$C < \infty$ such that for every $f \in L_p(V)$ the tempered distribution $m(D) f$ is a function
$g \in L_p(V)$ with $\|g\|_p \leq C \|f\|_p$.
\end{theorem}

\begin{proof}
\leanok
By \Cref{thm:tl-distr-lp} and \Cref{thm:square-function-upper},
$\|f\|_{F^0_{p,2}} \leq Q_{V,p} \|f\|_p$. \Cref{thm:tl-multiplier} with $\sigma = 0$ and $s = 0$
gives $C_1 < \infty$ with $\|m(D) f\|_{F^0_{p,2}} \leq C_1 Q_{V,p} \|f\|_p$, and
\Cref{thm:triebel-lizorkin-distr-to-bessel} with $s = 0$, where $J_0$ is the identity, gives
$C_2 < \infty$ and $g \in L_p(V)$ with $m(D) f = g$ and
$\|g\|_p \leq C_2 \|m(D) f\|_{F^0_{p,2}}$. Take $C = C_2 C_1 Q_{V,p}$.
\end{proof}

\begin{definition}[The multiplier on $L_p$]
\label{def:symbol-lp}
\lean{Rd.symbolLpConst, Rd.symbolLpConst_ne_top, Rd.symbolLpFun, Rd.symbolLp}
\leanok
\uses{thm:mihlin-symbols, lem:lp-distribution-determined}
Let $n \geq 1$, $1 < p < \infty$ and $m \colon V \to \C$ a symbol of order $0$, and fix a constant
$C_m < \infty$ of \Cref{thm:mihlin-symbols}. For $f \in L_p(V)$ let $m(D) f \in L_p(V)$ be the
function $g$ of \Cref{thm:mihlin-symbols} with $m(D) f = g$ as tempered distributions and
$\|g\|_p \leq C_m \|f\|_p$.
\end{definition}

The function $g$ is unique, since $L_p$ functions are determined by their tempered distributions
(\Cref{lem:lp-distribution-determined}).

\begin{lemma}[The multiplier on $L_p$ is bounded]
\label{lem:symbol-lp-bounded}
\lean{Rd.exists_fourierMultiplier_eq_symbolLpConst, Rd.toTemperedDistributionCLM_symbolLpFun,
  Rd.norm_symbolLpFun_le, Rd.toTemperedDistribution_symbolLp, Rd.enorm_symbolLp_le}
\leanok
\uses{def:symbol-lp, thm:mihlin-symbols, lem:lp-distribution-determined}
Let $n \geq 1$, $1 < p < \infty$ and $m \colon V \to \C$ a symbol of order $0$. The map
$f \mapsto m(D) f$ is a linear operator on $L_p(V)$, and for every $f \in L_p(V)$ the function
$m(D) f$ represents the tempered distribution $m(D) f$ and satisfies
$\|m(D) f\|_p \leq C_m \|f\|_p$.
\end{lemma}

\begin{proof}
\leanok
The representation and the bound hold by the choice of $C_m$ and of $m(D) f$
(\Cref{thm:mihlin-symbols}). The map is linear, since $m(D)$ is linear on tempered distributions
and $L_p$ functions are determined by their tempered distributions
(\Cref{lem:lp-distribution-determined}).
\end{proof}

\section{Fourier multipliers on Bessel potential spaces}
\label{sec:bessel-potential-multipliers}

A symbol of order $\sigma$ maps $H^{s,p}$ into $H^{s-\sigma,p}$ for every real $s$ and
$1 < p < \infty$; this gives the inclusions of the Bessel potential spaces, the action of
derivatives and of the Laplacian, and the characterisation of $H^{s+1,p}$ by first derivatives.
Mathlib proves the case $p = 2$ (\texttt{TemperedDistribution.\allowbreak MemSobolev.\allowbreak
mono}, \texttt{.lineDerivOp}, \texttt{.laplacian}) by Plancherel's theorem. The statements of
\Cref{cor:bessel-potential-mapping} (1), (2), of \Cref{lem:bessel-potential-closure}, of
\Cref{lem:bessel-two-laplacian} and of
\Cref{thm:bessel-potential-derivatives}, and the proofs of the last two, are mirrored from TauCeti
(\texttt{Analysis/\allowbreak Distribution/\allowbreak Sobolev.lean}, commit \texttt{b61a90d},
Apache-2.0), where Mihlin's theorem supplies the multipliers of order $0$; here
\Cref{thm:bessel-potential-multiplier} supplies all orders. Mathlib's Bessel potential spaces are
spaces of tempered distributions $u$ with $J_s u \in L_p(V)$; here they are complex-valued.

\begin{lemma}[Dimension $0$]
\label{lem:bessel-potential-dimension-zero}
\lean{Rd.besselPotential_of_subsingleton, Rd.memSobolev_fourierMultiplierCLM_of_subsingleton}
\leanok
Let $V = \{0\}$. Then $J_s u = u$ for every real $s$ and every tempered distribution $u$, and for
every $1 \leq p \leq \infty$, all reals $s, t$, every $m \colon V \to \C$ and every
$u \in H^{s,p}(V)$, $m(D) u \in H^{t,p}(V)$.
\end{lemma}

\begin{proof}
\leanok
The Bessel weight is $\langle 0 \rangle^s = 1$, and $m$ is the constant $m(0)$, so
$m(D) u = m(0) \, u$ (Mathlib's \texttt{fourierMultiplierCLM\_\allowbreak const} for tempered
distributions).
\end{proof}

\begin{lemma}[Sums and scalar multiples in $H^{s,p}$]
\label{lem:bessel-potential-closure}
\lean{TemperedDistribution.MemSobolev.sum, TemperedDistribution.MemSobolev.smul_of_tower}
\leanok
For every $1 \leq p \leq \infty$, every real $s$, every finite family $(u_i)$ in $H^{s,p}(V)$,
every $u \in H^{s,p}(V)$ and every scalar $c$ acting on tempered distributions through $\C$, such
as a real number, $\sum_i u_i \in H^{s,p}(V)$ and $c \, u \in H^{s,p}(V)$.
\end{lemma}

\begin{proof}
\leanok
Mathlib's \texttt{MemSobolev.\allowbreak add}, \texttt{MemSobolev.\allowbreak smul} and induction
over the finite set.
\end{proof}

\begin{theorem}[Fourier multipliers on Bessel potential spaces]
\label{thm:bessel-potential-multiplier}
\lean{TemperedDistribution.MemSobolev.fourierMultiplierCLM_of_hasSymbolGrowth}
\leanok
\uses{thm:bessel-triebel-lizorkin-distr, thm:tl-multiplier, lem:bessel-potential-dimension-zero}
Let $1 < p < \infty$ and $\sigma$ real. For every symbol $m \colon V \to \C$ of order $\sigma$,
every real $s$ and every tempered distribution $u \in H^{s,p}(V)$,
$m(D) u \in H^{s-\sigma,p}(V)$.
\end{theorem}

\begin{proof}
\leanok
For $n \geq 1$, $H^{s,p}(V) = \TL{s}{p}{2}$ for every real $s$
(\Cref{thm:bessel-triebel-lizorkin-distr}), and $m(D)$ maps $\TL{s}{p}{2}$ into
$\TL{s-\sigma}{p}{2}$ (\Cref{thm:tl-multiplier}). For $n = 0$ apply
\Cref{lem:bessel-potential-dimension-zero}.
\end{proof}

\begin{corollary}[Mapping properties of Bessel potential spaces]
\label{cor:bessel-potential-mapping}
\lean{TemperedDistribution.MemSobolev.mono_of_one_lt, Rd.memSobolev_bessel_of_nonpos,
  TemperedDistribution.MemSobolev.lineDerivOp_of_one_lt,
  TemperedDistribution.MemSobolev.laplacian_of_one_lt}
\leanok
\uses{thm:bessel-potential-multiplier, lem:symbol-examples, lem:symbol-algebra}
Let $s$ be real.
\begin{enumerate}
  \item For every $1 < p < \infty$ and every real $s' \leq s$, $H^{s,p}(V) \subseteq H^{s',p}(V)$.
    In particular, for every $1 < p < \infty$, every $f \in L_p(V)$ lies in $H^{s,p}(V)$ if
    $s \leq 0$.
  \item For every $1 < p < \infty$, every $v \in V$ and every $u \in H^{s,p}(V)$,
    $\partial_v u \in H^{s-1,p}(V)$.
  \item For every $1 < p < \infty$ and every $u \in H^{s,p}(V)$, $\Delta u \in H^{s-2,p}(V)$.
\end{enumerate}
\end{corollary}

\begin{proof}
\leanok
(1) $u \in H^{s',p}$ if and only if $J_{s'-s} u \in H^{s,p}$ (Mathlib's
\texttt{memSobolev\_\allowbreak besselPotential\_\allowbreak iff}), and
$\langle \xi \rangle^{s'-s}$ is a symbol of order $s' - s \leq 0$, hence of order $0$
(\Cref{lem:symbol-examples}, \Cref{lem:symbol-algebra}); apply
\Cref{thm:bessel-potential-multiplier}. An $L_p$ function lies in $H^{0,p}$, since $J_0$ is the
identity. (2) $\partial_v = 2\pi i \, \langle \xi, v \rangle(D)$ (Mathlib's
\texttt{TemperedDistribution.\allowbreak lineDeriv\_\allowbreak eq\_\allowbreak
fourierMultiplierCLM}) with the symbol $\langle \xi, v \rangle$ of order $1$. (3)
$\Delta = -(2\pi)^2 \|\xi\|^2(D)$ (Mathlib's \texttt{TemperedDistribution.\allowbreak
laplacian\_\allowbreak eq\_\allowbreak fourierMultiplierCLM}) with the symbol $\|\xi\|^2$ of
order $2$ (\Cref{lem:symbol-examples}).
\end{proof}

\begin{lemma}[$J_2 = 1 - (2\pi)^{-2} \Delta$]
\label{lem:bessel-two-laplacian}
\lean{TemperedDistribution.besselPotential_two_eq_sub_laplacian}
\leanok
For every tempered distribution $u$ on $V$ with values in a complex normed space,
$J_2 u = u - (2\pi)^{-2} \Delta u$.
\end{lemma}

\begin{proof}
\leanok
The symbol of $J_2$ is $1 + \|\xi\|^2$, and $\Delta = -(2\pi)^2 \|\xi\|^2(D)$ (Mathlib's
\texttt{TemperedDistribution.\allowbreak laplacian\_\allowbreak eq\_\allowbreak
fourierMultiplierCLM}); the factor comes from Mathlib's normalisation of the Fourier transform.
\end{proof}

\begin{theorem}[The derivative characterisation of $H^{s+1,p}$]
\label{thm:bessel-potential-derivatives}
\lean{TemperedDistribution.memSobolev_add_one_of_lineDerivOp,
  TemperedDistribution.memSobolev_add_one_iff, TemperedDistribution.memSobolev_natCast_add_one_iff}
\leanok
\uses{cor:bessel-potential-mapping, lem:bessel-two-laplacian, lem:bessel-potential-closure}
Let $1 < p < \infty$ and $s$ real. For every tempered distribution $u$ on $V$,
\[
  u \in H^{s+1,p}(V) \iff u \in H^{s,p}(V) \text{ and } \partial_v u \in H^{s,p}(V)
  \text{ for every } v \in V ,
\]
and it suffices that $\partial_{b_i} u \in H^{s,p}(V)$ for the vectors $b_i$ of one orthonormal
basis. For every $k \in \N$, $u \in H^{k+1,p}(V)$ if and only if $u \in H^{0,p}(V)$ and
$\partial_v u \in H^{k,p}(V)$ for every $v \in V$.
\end{theorem}

\begin{proof}
\leanok
If $u \in H^{s+1,p}$, then $u \in H^{s,p}$ and $\partial_v u \in H^{s,p}$ by
\Cref{cor:bessel-potential-mapping}. Conversely, $u \in H^{s+1,p}$ if and only if
$J_1 u \in H^{s,p}$, and by \Cref{lem:bessel-two-laplacian} and
$\Delta = \sum_i \partial_{b_i} \partial_{b_i}$ (Mathlib's \texttt{TemperedDistribution.\allowbreak
laplacian\_\allowbreak eq\_\allowbreak sum}),
\[
  J_1 u = J_{-1} J_2 u = J_{-1} u - (2\pi)^{-2} \sum_i J_{-1} \partial_{b_i} \partial_{b_i} u .
\]
Here $J_{-1} u \in H^{s+1,p} \subseteq H^{s,p}$ and
$J_{-1} \partial_{b_i} \partial_{b_i} u \in H^{s,p}$, since $\partial_{b_i} \partial_{b_i} u \in
H^{s-1,p}$ (\Cref{cor:bessel-potential-mapping}), so $J_1 u \in H^{s,p}$
(\Cref{lem:bessel-potential-closure}). The statement for $k \in \N$ follows by induction on $k$.
\end{proof}

\begin{corollary}[The lifting property for every $q$]
\label{cor:tl-lifting}
\lean{Rd.exists_triebelLizorkinNorm_besselPotential_le,
  Rd.exists_triebelLizorkinNorm_le_besselPotential}
\leanok
\uses{cor:tl-multiplier-lp, lem:symbol-examples}
Let $1 < p < \infty$, $1 < q \leq \infty$ and $\sigma, s$ real. There are $C, C' < \infty$ such that
for all $g, h \in L_p(V)$ with $J_\sigma g = h$,
\[
  \|h\|_{F^{s-\sigma}_{p,q}} \leq C \, \|g\|_{F^s_{p,q}}
  \quad \text{and} \quad
  \|g\|_{F^s_{p,q}} \leq C' \, \|h\|_{F^{s-\sigma}_{p,q}} .
\]
\end{corollary}

\begin{proof}
\leanok
$J_\sigma$ is the Fourier multiplier with the symbol $\langle \xi \rangle^\sigma$ of order
$\sigma$, and $g = J_{-\sigma} h$ with the symbol $\langle \xi \rangle^{-\sigma}$ of order
$-\sigma$ (\Cref{lem:symbol-examples}); apply \Cref{cor:tl-multiplier-lp} twice.
\end{proof}

This is Triebel's lifting property~\cite[Theorem~2.3.8]{Triebel} for $1 < q \leq \infty$.
